\ifx\danslelivre\undefined
\RequirePackage{luatex85}
\documentclass[10pt]{../fichiers-configuration-latex/smfart}
\usepackage{polyglossia}
%\usepackage[francais]{babel}
\usepackage{fontspec}
\usepackage[math-style=ISO,bold-style=upright,sans-style=upright]{unicode-math}
%\usepackage[mathletters]{ucs}
%\usepackage[utf8]{inputenc}
%\usepackage[T1]{fontenc}
\usepackage[unicode,colorlinks=true,linkcolor=blue,citecolor=Sepia,urlcolor=magenta]{hyperref}
\setmainlanguage{french}
\setmainfont[Ligatures=TeX]{Linux Libertine O}
% variante possible : TeXGyrePagella ?
\setmathfont{XITS Math}

% maths
\usepackage{amsmath}
%\usepackage{amssymb}
%\usepackage{mathrsfs}
%\usepackage{stmaryrd}
%\usepackage{euscript}

% couleurs
\usepackage[usenames,dvipsnames]{xcolor}

% graphiques
\usepackage{graphicx}
\usepackage[all]{xy}
% Xy-pic 3.8.9's PDF driver assumes that \fontdimen8\textfont3 is a
% legacy TeX rule width.  Under LuaTeX with unicode-math/XITS Math it is
% instead 6.62pt, which turns every arrow shaft into a thick black bar.
% Override only xypdf's private cached width; do not alter the math font.
\makeatletter
\AtBeginDocument{%
  \@ifundefined{xP@preclw}{}{%
    \edef\xP@preclw{0.4pt}%
    \edef\xP@lw{\xP@dim\xP@preclw}%
    \PackageInfo{sga4-xypdf-compat}{Xy-pic PDF line width set to \xP@preclw}%
  }%
}
\makeatother
%\usepackage{tikz}
%\usetikzlibrary{matrix,arrows,calc}
%\tikzset{column sep=3em,row sep=2em,text height=1.5ex, text depth=0.25ex,>=angle 90}
%\tikzset{description/.style={fill=white,inner sep=1pt}}

% hyperliens, 'etiquettes
\usepackage{xr-hyper}
\usepackage{../fichiers-configuration-latex/smfthm}
%\ifx\cacheretiquettes\undefined
%\usepackage[inline]{showlabels} \else \fi % pour rendre visibles les 'etiquettes
%\usepackage{../fichiers-configuration-latex/makeidx} % pour faire un index
%\usepackage[pagebackref,unicode]{hyperref}  % pour voir o`u sont cit'ees les r'ef'erences
\hypersetup{unicode,colorlinks=true,linkcolor=blue,citecolor=Sepia,urlcolor=magenta,
  hypertexnames=false}
%\hypersetup{breaklinks=true} %permet le retour `a la ligne dans les liens trop longs. Incompatible avec pdflatex
%\renewcommand*{\backref}[1]{{\small [cit\'e page(s)~#1.]}}

\usepackage{xspace}
\usepackage{pifont} % pour expose XVII

\renewcommand{\thefootnote}{\roman{footnote}} % pour mettre les notes en bas de page sous la forme (ii) et non (2), qui ressemble trop (?) a un exposant.

\newcounter{pageoriginal}
\newcommand{\pageoriginale}{\refstepcounter{pageoriginal}\marginpar{\footnotesize\textbf{\thepageoriginal}}\xspace}
\let\pageoriginaled\pageoriginale

\usepackage{ifthen}
\newboolean{original}
\def\sisi{\ifthenelse{\boolean{original}}}
\setboolean{original}{false}

\newcounter{cnote}
\def\nde#1{\ifthenelse{\boolean{original}}{}{\stepcounter{cnote}\def\thefootnote{\arabic{cnote}}\protect\footnote[\value{cnote}]{N.D.E.: #1}}}
\def\ndemark{\ifthenelse{\boolean{original}}{}{\stepcounter{cnote}\def\thefootnote{\arabic{cnote}}\protect\footnotemark[\value{cnote}]}}
\def\ndetext#1{\ifthenelse{\boolean{original}}{}{\def\thefootnote{\arabic{cnote}}\protect\footnotetext[\value{cnote}]{N.D.E.: #1}}}

\setmainlanguage{french}
\setotherlanguage{english}

% unicode-math deliberately disables the legacy \Arrowvert glyph.  The
% source uses it once as a rotated equality marker; \Vert has the intended
% appearance and keeps the historical diagram buildable with LuaLaTeX.
\renewcommand{\Arrowvert}{\Vert}


%%
%% Pour afficher le numéro de version Git (exige que LuaTeX soit appelé avec --shell-escape)
%%

\catcode`\%=12
\directlua{
function command_output(cmd)
  local output_pipe = io.popen(cmd)
  local line_iter = output_pipe and output_pipe:lines()
  local line = line_iter and line_iter()
  if output_pipe then output_pipe:close() end
  return line
end
local git_command1 = [[git log --pretty=format:"%h %ai" -1]]
local git_command2 = [[if git diff --quiet HEAD ; then echo clean ; else echo dirty ; fi]]
function git_status()
  local status = command_output(git_command1)
  if status then
    tex.tprint({luatexbase.catcodetables.CatcodeTableLaTeX, string.char(0x5c).."texttt{"}, {-2,status}, {luatexbase.catcodetables.CatcodeTableLaTeX, "}"})
    if not (command_output(git_command2) == "clean") then tex.sprint(-2, " (modifications non validées)") end
  else
    tex.sprint(-2, "?")
  end
end
}
\catcode`\%=14
\def\showgitstatus{\directlua{git_status()}}


% macros pour LaTeX
\renewcommand{\mathsf}{}

% Police différente pour U+2113 SCRIPT SMALL L
\setmathfont[range="2113]{latinmodern-math.otf}

% Give every explicit SGA label its own destination and preserve the printed
% number supplied as the optional argument.  The former definition inherited
% whichever counter anchor happened to precede it and ignored #1 entirely.
\makeatletter
\def\etiquette[#1]#2{%
  \phantomsection
  \protected@edef\@currentlabel{#1}%
  \label{#2}%
}
\makeatother
\def\eqhyperref[#1]#2{(\hyperref[#1]{#2})}
\def\pagehyperref[#1]#2{(page \hyperref[#1]{#2})}
\def\hyperrefext[#1][#2]#3{\hyperref[#2]{#1 #3}}
\def\eqhyperrefext[#1][#2]#3{\hyperref[#2]{#1 #3}}     %cf. exposes 17
\def\neweqhyperrefext[#1][#2]#3{\hyperref[#2]{#1 #3}}  %cf. exposes 17

\def\shat#1{\protect\widehat{#1}}
\def\stilde#1{\protect\widetilde{#1}}
%%%%%%%%%%%%%%%%%%


\def\etcn{etc.\xspace}
\def\etco{etc...\xspace}
\def\etc{\sisi{\etco}{\etcn}}
\def\ldotsn{\sisi{...}{\etcn}}

\def\cf{{cf.}\kern.3em}
\def\Cf{{Cf.}\kern.3em}
\def\ie{{i.e.}\kern.3em}
\def\iev{{i.e.},\xspace}
\def\loccito{loc.\kern2pt cit.\xspace}
\def\loccitn{\textit{loc.\kern2pt cit.}\xspace}
\def\loccit{\sisi{\loccito}{\loccitn}}
\def\EGA{\textup{EGA}\kern.2em}
\def\SGA{\textup{SGA}\kern.3em}

\def\quoi{}

\def\quav{\quad,\quad}
\def\quaet{\quad \text{et}\quad}
\def\fibreproduct#1#2#3{#1 ×_{#3} #2}
\let\fprod\fibreproduct
\def\ccheck#1#2{\mbox{${\check{\check{\rm #1}}}^{\lower2pt\hbox{$\scriptstyle #2$}}$}}
\def\sheaf#1{\protect\mathcal{#1}}
\def\vieuxsheaf#1{{#1}} % cf. problème avec sheaf{a\Phi_i} par exemple...
\DeclareMathOperator{\Hom}{Hom}
\DeclareMathOperator{\can}{can}
\DeclareMathOperator{\Ind}{Ind}
\DeclareMathOperator{\ind}{ind}
%\DeclareMathOperator{\Hhomtop}{\mathcal{H}\!\mathit{omtop}}
\DeclareMathOperator{\Hhomtop}{\mathbf{Homtop}}
\DeclareMathOperator{\Point}{Point}
\DeclareMathOperator{\Pro}{Pro}
\DeclareMathOperator{\Top}{Top}
\DeclareMathOperator{\Mod}{Mod}
\DeclareMathOperator{\Cart}{Cart}
%\DeclareMathOperator{\Mmorsite}{\mathcal{M}\!\mathit{orsite}}
\DeclareMathOperator{\Mmorsite}{\mathbf{Morsite}}
\DeclareMathOperator{\loc}{loc}
\DeclareMathOperator{\Loc}{Loc}
\DeclareMathOperator{\Pic}{Pic}
\DeclareMathOperator{\Et}{Et}
\DeclareMathOperator{\card}{card}
\DeclareMathOperator{\Ens}{Ens}
\DeclareMathOperator{\h}{\underline{H}}
\DeclareMathOperator{\Fonct}{Fonct}
\DeclareMathOperator{\Ker}{Ker}
\DeclareMathOperator{\but}{but}
\DeclareMathOperator{\Ob}{Ob}
\DeclareMathOperator{\ob}{ob}
\DeclareMathOperator{\Lib}{Lib}
\DeclareMathOperator{\esj}{esj}
\DeclareMathOperator{\id}{id}
\DeclareMathOperator{\Id}{Id}
\DeclareMathOperator{\ad}{ad}
\DeclareMathOperator{\Iim}{Im}
\DeclareMathOperator{\pr}{pr}
\DeclareMathOperator{\Cat}{Cat}
\DeclareMathOperator{\car}{car}
\DeclareMathOperator{\SysInd}{SysInd}
\DeclareMathOperator{\Sup}{Sup}
\DeclareMathOperator{\Coker}{Coker}
\DeclareMathOperator{\Filt}{Filt}
\DeclareMathOperator{\acc}{acc}
\DeclareMathOperator{\sur}{sur}
\DeclareMathOperator{\cat}{cat}
\DeclareMathOperator{\Ord}{Ord}
\DeclareMathOperator{\Gr}{Gr}
\DeclareMathOperator{\Cov}{Cov}
\DeclareMathOperator{\Caf}{Caf}
\DeclareMathOperator{\coker}{coker}
\DeclareMathOperator{\coim}{coim}
\DeclareMathOperator{\im}{im}
\DeclareMathOperator{\Ouv}{Ouv}
\DeclareMathOperator{\diag}{diag}
\DeclareMathOperator{\Esp}{Esp}
\DeclareMathOperator{\TOP}{TOP}
\DeclareMathOperator{\Unif}{Unif}
\DeclareMathOperator{\Iind}{\mathbf{Ind}}
\DeclareMathOperator{\Hhomcart}{\mathbf{Homcart}}
\DeclareMathOperator{\Hhomtopess}{\mathbf{Homtopess}}
\DeclareMathOperator{\Ppoints}{\mathbf{Points}}
\DeclareMathOperator{\Ppoint}{\mathbf{Point}}
\DeclareMathOperator{\Ttors}{\mathbf{Tors}}
\DeclareMathOperator{\Ssex}{\mathbf{Sex}}
\DeclareMathOperator{\Ffib}{\mathbf{Fib}}
\DeclareMathOperator{\Vvois}{\mathbf{Vois}}
\DeclareMathOperator{\Ppointess}{\mathbf{Pointess}}
\DeclareMathOperator{\scrHom}{\mathbf{Hom}}
\DeclareMathOperator{\sheafPic}{\mathbf{Pic}}
\DeclareMathOperator{\sheafAut}{\mathbf{Aut}}
\DeclareMathOperator{\Hhomtoplocan}{\mathbf{Homtoplocan}}
\DeclareMathOperator{\Hhomtopan}{\mathbf{Homtopan}}
\DeclareMathOperator{\Pptgeom}{\mathbf{Ptgéom}}
\DeclareMathOperator{\Morsite}{Morsite}
\DeclareMathOperator{\Fib}{Fib}
\DeclareMathOperator{\Fais}{Fais}
\DeclareMathOperator{\kar}{kar}
\DeclareMathOperator{\Groupes}{Groupes}
\DeclareMathOperator{\Restr}{Restr}
\DeclareMathOperator{\Isom}{Isom}
\DeclareMathOperator{\Prol}{Prol}
\DeclareMathOperator{\Res}{Res}
\DeclareMathOperator{\Tors}{Tors}
\DeclareMathOperator{\esp}{esp}
\DeclareMathOperator{\es}{es}
\DeclareMathOperator{\od}{od}
\DeclareMathOperator{\gr}{gr}
\DeclareMathOperator{\sob}{sob}
\DeclareMathOperator{\dfn}{dfn}
\DeclareMathOperator{\Sch}{Sch}
\DeclareMathOperator{\ouv}{ouv}
\DeclareMathOperator{\cons}{cons}
\DeclareMathOperator{\Nerf}{Nerf}
\DeclareMathOperator{\Points}{Points}
\DeclareMathOperator{\Espsob}{Espsob}
\DeclareMathOperator{\Homtop}{Homtop}
\DeclareMathOperator{\Pt}{Pt}
\DeclareMathOperator{\Quot}{Quot}
\DeclareMathOperator{\glob}{glob}
\DeclareMathOperator{\Homfib}{Homfib}
\DeclareMathOperator{\Qouv}{Qouv}
\DeclareMathOperator{\Quotouv}{Quotouv}
\DeclareMathOperator{\ens}{ens}
\DeclareMathOperator{\ab}{ab}
\DeclareMathOperator{\an}{an}
\DeclareMathOperator{\fs}{fs}
\DeclareMathOperator{\red}{red}
\DeclareMathOperator{\Ab}{Ab}
\DeclareMathOperator{\cosq}{cosq}
\DeclareMathOperator{\Ann}{Ann}
\DeclareMathOperator{\Zar}{Zar}


%%%%%%%%%%%%
%%%%% SGA4-2
\DeclareMathOperator{\Sex}{Sex}
\DeclareMathOperator{\vois}{vois}
\DeclareMathOperator{\Vois}{Vois}
\DeclareMathOperator{\Ext}{Ext}
\DeclareMathOperator{\ext}{ext}
\DeclareMathOperator{\Fer}{Fer}
\DeclareMathOperator{\Add}{Add}
\DeclareMathOperator{\Cocart}{Cocart}
\DeclareMathOperator{\cosk}{cosk}
\DeclareMathOperator{\plat}{plat}

\DeclareMathOperator{\faisc}{faisc}
\DeclareMathOperator{\sk}{sk}
\DeclareMathOperator{\ch}{ch}

\def\HH{\mathbf{H}}
\def\qc{\mathrm{qc}}
\def\et{\mathrm{et}}
\def\coh{\mathrm{coh}}

\def\Hvr#1{%
\begin{array}[b]{l}
{\kern2pt\scriptstyle r}\\[-8pt]
{\kern1pt\scriptstyle \vee}\\[-5pt]
H^{#1}
\end{array}}

\def\Hvv#1{%
\begin{array}[b]{l}
\kern2pt\rotatebox{90}{${\scriptscriptstyle\ll}$}\\[-6pt]
H^{#1}
\end{array}}

\def\Hvi#1{%
\begin{array}[b]{l}
{\kern0.2pt\scriptscriptstyle \infty}\\[-8pt]
{\kern1pt\scriptstyle \vee}\\[-5pt]
H^{#1}
\end{array}}

\def\Hv#1{%
\begin{array}[b]{l}
{\kern1pt\scriptstyle \vee}\\[-5pt]
H^{#1}
\end{array}}

\DeclareMathOperator{\cart}{cart}
\DeclareMathOperator{\Ffaisc}{\mathbf{Faisc}}


\def\Ev#1{%
\begin{array}[b]{l}
{\kern1pt\scriptstyle \vee}\\[-5pt]
E'_{#1}
\end{array}}

\DeclareMathOperator{\Cont}{Cont}
\DeclareMathOperator{\Ccont}{\mathscr{C}\!\mathit{ont}}
\DeclareMathOperator{\Site}{Site}
\newcommand\MM{\mathbf{M}}
\def\limtop#1{{\displaystyle{\varprojlim_{#1}\!\mbox{\rm top }}}}
\def\Limtop{\mathop{\varprojlim\!\mbox{\rm top}}}
\DeclareMathOperator{\spec}{spec}
\def\ss{\scriptstyle{\mathscr{S}}}
\DeclareMathOperator{\Gal}{Gal}
\def\smallm{\scriptstyle{\mathscr{M}}}
\DeclareMathOperator{\End}{End}
\DeclareMathOperator{\ass}{ass}
\DeclareMathOperator{\zar}{zar}
\DeclareMathOperator{\Ogg}{Ogg}
\DeclareMathOperator{\res}{res}
\DeclareMathOperator{\tr}{tr}
\DeclareMathOperator{\Supp}{Supp}
\DeclareMathOperator{\codim}{codim}
\DeclareMathOperator{\rk}{rk}
\DeclareMathOperator{\cur}{cur}
\DeclareMathOperator{\cdqc}{cdqc}
\newcommand\DD{\mathbf{D}}
\DeclareMathOperator{\cl}{cl}
\def\JJ{\mathbf{J}}
\def\FF{\mathbf{F}}
\def\BB{\mathbf{B}}
\DeclareMathOperator{\Alget}{Alget}
\DeclareMathOperator{\Ass}{Ass}
\DeclareMathOperator{\supp}{supp}
\DeclareMathOperator{\rad}{rad}
\DeclareMathOperator{\prof}{prof}
\DeclareMathOperator{\cint}{int}
\DeclareMathOperator{\sic}{sic}
\DeclareMathOperator{\Tor}{Tor}
\DeclareMathOperator{\tors}{tors}
\DeclareMathOperator{\Sym}{Sym}
\DeclareMathOperator{\parf}{parf}
\DeclareMathOperator{\tor}{tor}
\DeclareMathOperator{\defn}{defn}
\DeclareMathOperator{\Div}{Div}
\DeclareMathOperator{\Pol}{Pol}
\DeclareMathOperator{\Fex}{Fex}
\DeclareMathOperator{\Lie}{Lie}
\DeclareMathOperator{\tot}{tot}
\DeclareMathOperator{\Aut}{Aut}
\DeclareMathOperator{\HOM}{HOM}
\DeclareMathOperator{\pch}{pch}
\DeclareMathOperator{\Ch}{Ch}
\DeclareMathOperator{\EXT}{EXT}
\DeclareMathOperator{\PIC}{PIC}
\DeclareMathOperator{\Rf}{Rf}
\DeclareMathOperator{\inj}{inj}
\newcommand\VV{\mathbf{V}}

\def\Xv{%
\begin{array}[b]{l}
{\kern1pt\scriptstyle \vee}\\[-5pt]
X
\end{array}}

%%%%%%%
\newcommand\TT{\mathbf{T}}
\newcommand\XX{\mathbf{X}}
\newcommand\EE{\mathbf{E}}
%%%%%%

\def\qedsymbol{\textup{C.Q.F.D.}}
\let\CQFD\qed
\let\cqfd\qed

%%%%%\bigdot
\font\bigsymb=cmsy10 at 4pt
\def\bigdot{{\kern1.2pt\raise 1.5pt\hbox{\bigsymb\char15}}}
%%%%%

\def\Lright{\mathop{\vcenter{\hbox{\hbox{$\longrightarrow$}\lower8pt \hbox{\kern-16pt\hbox{$\longrightarrow$}}}}}}

\def\Aa{{\mathcal A}}
\def\Bb{{\mathcal B}}
\def\Cc{{\mathcal C}}
\def\Dd{{\mathcal D}}
\def\Ee{{\mathcal E}}
\def\Ff{{\mathcal F}}
\def\Gg{{\mathcal G}}
\def\Hh{{\mathcal H}}
\def\Ii{{\mathcal I}}
\def\Jj{{\mathcal J}}
\def\Kk{{\mathcal K}}
\def\Ll{{\mathcal L}}
\def\Mm{{\mathcal M}}
\def\Nn{{\mathcal N}}
\def\Oo{{\mathcal O}}
\def\Pp{{\mathcal P}}
\def\Qq{{\mathcal Q}}
\def\Rr{{\mathcal R}}
\def\Ss{{\mathcal S}}
\def\Tt{{\mathcal T}}
\def\Uu{{\mathcal U}}
\def\Vv{{\mathcal V}}
\def\Ww{{\mathcal W}}
\def\Xx{{\mathcal X}}
\def\Yy{{\mathcal Y}}
\def\Zz{{\mathcal Z}}

\newcommand{\CC}{\mathbf{C}}
\newcommand{\GG}{\mathbf{G}}
\newcommand{\KK}{\mathbf{K}}
\newcommand{\LL}{\mathbf{L}}
\newcommand{\NN}{\mathbf{N}}
\newcommand{\PP}{\mathbf{P}}
\newcommand\QQ{\mathbf{Q}}
\newcommand\RR{\mathbf{R}}
\newcommand\ZZ{\mathbf{Z}}

\def\to{\mathchoice{\longrightarrow}{\rightarrow}{\rightarrow}{\rightarrow}}
\def\from{\mathchoice{\longleftarrow}{\leftarrow}{\leftarrow}{\leftarrow}}
\def\mto{\mathchoice{\longmapsto}{\mapsto}{\mapsto}{\mapsto}}
\def\mfrom{\mathrel{\from\joinrel\mapstochar\,}}
\def\hto{\mathrel{\lhook\joinrel\to}}
\def\hfrom{\mathrel{\from\joinrel\rhook}}

\def\implies{\mathchoice{\Longrightarrow}{\Rightarrow}{\Rightarrow}{\Rightarrow}}
\def\implied{\mathchoice{\Longleftarrow}{\Leftarrow}{\Leftarrow}{\Leftarrow}}
\def\iff{\mathchoice{\Longleftrightarrow}{\Leftrightarrow}{\Leftrightarrow}{\Leftrightarrow}}

\def\To#1{\mathchoice{\xrightarrow{\labelstyle~~#1~}}{\stackrel{#1}{\longrightarrow}}{}{}}
\def\Mto#1{\mathrel{\mapstochar\To{#1}}}
\def\From#1{\mathchoice{\xleftarrow{\labelstyle~~#1~}}{\stackrel{#1}{\longleftarrow}}{}{}}

\newcommand{\isomto}{\overset{\sim}{\longrightarrow}}
\newcommand{\isomfrom}{\overset{\sim}{\longleftarrow}}

\newcommand{\IFF}{\unskip~$\iff$~\ignorespaces}
\newcommand{\THEN}{\unskip~$\implies$~\ignorespaces}

\let\leq\leqslant
\let\geq\geqslant
\let\le\leqslant
\let\ge\geqslant

\def\cf{{cf.}\kern.3em}
\def\Cf{{Cf.}\kern.3em}
\def\ie{{i.e.}\kern.3em}
\def\iev{{i.e.},\xspace}
%\def\loccito{loc.\kern2pt cit.\xspace}
%\def\loccitn{\textit{loc.\kern2pt cit.}\xspace}
%\def\loccit{\sisi{\loccito}{\loccitn}}
\def\EGA{\textup{EGA}\kern.2em}
\def\SGA{\textup{SGA}\kern.3em}

\let\oldS\S
\def\S{\oldS\kern .15em}
\def\ptbl{.\kern .2em}
\def\resp{resp.\kern.3em}
\def\Exp{Exp.\kern.3em}

\let\bar\overline
\let\hat\widehat
\let\tilde\widetilde
\let\dpl\displaystyle
\let\ldots\dots
\let\emptyset\varnothing
\let\skpt\skippointrait
\let\epsilon\varepsilon

\def\sfrac#1#2{{#1}/{#2}}

\DeclareMathOperator{\cd}{cd}
%\DeclareMathOperator{\Hhom}{\mathcal{H}\!\mathit{om}}
\DeclareMathOperator{\Hhom}{\mathbf{Hom}}
%\DeclareMathOperator{\Eext}{\mathcal{E}\!\mathit{xt}}
\DeclareMathOperator{\Eext}{\mathbf{Ext}}
\DeclareMathOperator{\Spec}{Spec}
\DeclareMathOperator{\Tr}{Tr}

\let\emptyset=\varnothing

%%%%%%%%%%%%%%%
%%% expose XVII
%%%%%%%%%%%%%%%

\makeatletter
\let\old@optexpose\optexpose
\newcommand{\empty@optexpose}[1]{}
\newcommand{\paren@optexpose}[1]{\old@optexpose{#1}(}
\newcommand{\refnoexpose}[1]{\let\optexpose\empty@optexpose\ref{#1}\let\optexpose\old@optexpose}%\refnoexpose{X.cor5.2} engendre "5.2" au lieu de "X 5.2"
\newcommand{\neweqref}[1]{\let\optexpose\paren@optexpose\textup{\ref{#1})}\let\optexpose\old@optexpose}%\neweqref{XII.eq4.1} engendre "XII (4.1)" au lieu de "(XII 4.1)"
\makeatother
\newcommand{\refvb}[2][VI${}_{\text{B}}$ ]{\sisi{#1\refnoexpose{#2}}{\ref{#2}}}

\newcommand{\leftexp}[2]{{\vphantom{#2}}^{#1}{#2}}
\newcommand{\AAA}{\mathbf{A}}
\newcommand{\Sf}{\mathfrak{S}}
\DeclareMathOperator*{\quotprojlim}{\text{``$\varprojlim$''}}
\DeclareMathOperator*{\quotinjlim}{\text{``$\varinjlim$''}}
\DeclareMathOperator{\Fl}{F\sisi{\ell}{l}}
\DeclareMathOperator{\Pl}{P\sisi{\ell}{l}}
%\DeclareMathOperator{\cl}{c\sisi{\ell}{l}}
\DeclareMathOperator{\tdf}{tdf}
\DeclareMathOperator{\rg}{rg}
\newcommand{\iInd}{\sisi{I}{i}nd}
\newcommand{\iind}{ind}
\newcommand{\ppro}{pro}
\newcommand{\tTor}{\sisi{T}{t}or}
\newcommand{\ttor}{tor}
\newcommand{\chap}{Chap.}
\newcommand{\Ffex}{\Ff\mathit{ex}}
\newcommand{\sheafk}{\underline{k}}
\newcommand{\sheafn}{\underline{n}}
\newcommand{\sheafGamma}{\underline{\Gamma}}

\sisi{
\DeclareMathOperator{\simarrow}{\rule{.15mm}{1.8mm}\kern -3.4pt\rightsquigarrow}
\newcommand{\sheafell}{\underline{\ell}}
\newcommand{\sheafZZ}{\underline{\ZZ}}
}{
\newcommand{\nderef}[1]{la N.D.E.~(\ref{#1}) page~\pageref{#1}}
}

%%%%%%%%%%%%%%%%%
%%%%% expose XVII
%%%%%%%%%%%%%%%%%



\externaldocument{../01/01.tex}
\externaldocument{../02/02.tex}
\externaldocument{../03/03.tex}
%\externaldocument{../04/04.tex}

\setboolean{original}{false}

\begin{document}

\begin{center}
IV. Topos\\
A. Grothendieck et J.-L. Verdier
\end{center}

\selectlanguage{english}version : \showgitstatus\selectlanguage{french}
\tableofcontents

\else
\chapter{Topos}
%\addtocontents{toc}{A. Grothendieck et J.-L. Verdier}
\begin{center}
A. Grothendieck et J.-L. Verdier
\end{center}
\fi

\setcounter{pageoriginal}{298}
\setcounter{section}{-1}
\section{Introduction}\etiquette[0]{IV.sec0}\pageoriginale

\subsection{}\etiquette[0.1]{IV.subsec0.1}
Nous avons vu dans \ref{II} diverses propriétés d'exactitude de
catégories
de la forme $\widetilde{𝒞}=$ catégorie des faisceaux
d'ensembles sur $𝒞$, où $𝒞$ est un petit site,
propriétés qu'on peut exprimer en disant qu'à beaucoup
d'égards, ces
catégories (que nous appellerons des \emph{topos}) héritent des
propriétés familières de la catégorie $(\Ens)$ des (petits)
ensembles. D'un autre côté, l'expérience a enseigné qu'il y a lieu
de considérer diverses situations en Mathématique \emph{surtout
comme un moyen technique pour construire
  les catégories de faisceaux} (d'ensembles) \emph{correspondantes,} \ie
\emph{les  « topos »  correspondants.} Il apparaît que toutes
les notions vraiment importantes liées à un site (par exemple
ses invariants cohomologiques, étudiés dans \ref{V}, divers autres
invariants « topologiques », tels ses invariants d'homotopie
étudiés récemment par M. \textsc{Artin} et B. \textsc{Mazur}
[1] et les notions étudiées dans le livre de J. \textsc{Giraud}
sur la cohomologie non commutative) s'expriment en fait directement
en termes du topos associé. Dans cette optique, il convient de
regarder deux sites comme étant essentiellement équivalents
lorsque les topos associés sont des catégories équivalentes,
et de considérer que la donnée d'un site (du moins dans le cas,
surtout important en pratique, où sa topologie est moins fine que
sa topologie canonique) revient à celle d'un topos $ℰ$
(savoir le topos associé, formé des faisceaux d'ensembles sur le
site), et d'une famille génératrice d'éléments de
$ℰ$ (\cf\pageoriginale II 4.9, et \hyperref[IV.cor1.2.1]{1.2.1}
ci-dessous). Ce point de vue est analogue à celui qui consiste à
associer un groupe à tout système de générateurs et tout
système de relations entre ces générateurs, et à attacher
son intérêt plutôt à la structure de ce groupe qu'au
système de générateurs et relations qui ont servi à
l'engendrer (considérés comme des données accessoires de la
situation). D'ailleurs le « lemme de comparaison »  III 4.1.
fournit de nombreux exemples de couples de sites $𝒞$,
$𝒞'$ non isomorphes, et même non équivalents en tant que
catégories, et donnant naissance à des topos équivalents, de
sorte qu'il y a lieu de considérer $𝒞$ et $𝒞'$
comme essentiellement équivalents.

\subsection{}\etiquette[0.2]{IV.subsec0.2}
Dans le présent exposé, nous donnons une caractérisation des topos
par des propriétés d'exactitude simples (due à J. \textsc{Giraud}),
nous étudions la notion naturelle de morphisme de topos, inspirée
par la notion d'application continue d'un espace topologique dans un
autre, et nous développons dans le cadre des topos certaines
constructions familières en théorie des faisceaux habituelle
(faisceaux $\Hhom$, faisceaux produit tensoriel, supports). Enfin,
nous montrons (en suivant M. \textsc{Artin}) comment on peut
reconstituer un topos à partir d'un « ouvert »  de celui-ci, du
« fermé » complémentaire, et d'un certain foncteur exact à
gauche qui les relie, qui, a peu de choses près, peut être choisi
d'ailleurs arbitrairement.

\subsection{}\etiquette[0.3]{IV.subsec0.3}
On a là un procédé de recollement de topos qui, appliqué à des topos
provenant d'espaces topologiques ordinaires, donnera en général un
topos qui ne sera plus du même type. Cela est une première
indication de la stabilité remarquable de la notion de topos par
diverses constructions naturelles, qui manque à la notion d'espace
topologique (dont la notion de\pageoriginale topos est inspirée).
Pour un deuxième exemple remarquable, signalons aussi celle de topos
classifiant relatif à un groupe d'un topos (\cf [1] ou
\hyperref[IV.exer5.9]{5.9} ci-dessous), inspirée de la notion classique
d'espace classifiant d'un groupe topologique, et la notion de « topos modulaire »  associé à divers « problèmes de modules »
en Géométrie Algébrique ou en Géométrie Analytique [10] [13].

D'autres topos, tel le \emph{topos étale} d'un schéma (étudié
systématiquement dans le présent Séminaire, à partir de \Exp
VII) ou le \emph{topos cristallin} d'un schéma relatif [6]
s'introduisent de façon naturelle lorsqu'on veut développer
pour des variétés algébriques abstraites (et plus
généralement des schémas) des théories de cohomologie
utilisables, qui remplacent la cohomologie de Betti classique des
variétés algébriques sur le corps des complexes.

\subsection{}\etiquette[0.4]{IV.subsec0.4}
On peut donc dire que la notion de topos, dérivé naturel du
\emph{point de vue faisceautique} en Topologie, constitue à son
tour un élargissement substantiel de la notion  d'espace
topologique \footnote[1]{\Cf [9], ou \hyperref[IV.subsec4.1]{4.1} et
\hyperref[IV.subsec4.2]{4.2} plus bas, pour les relations
  précises entre la notion de topos et celle d'espace
  topologique.}, englobant un grand nombre de situations qui
autrefois n'étaient pas considérées comme relevant de
l'intuition topologique. Le trait caractéristique de telles
situations est qu'on y dispose d'une notion de « localisation »,
notion qui est formalisée précisément par la notion de site
et, en dernière analyse, par celle de topos (via le topos
associé au site). Comme le terme de « topos »  lui-même est
censé précisément le suggérer, il semble raisonnable et
légitime aux auteurs du présent Séminaire de considérer que
l'objet de la Topologie est l'étude des \emph{topos} (et non des
seuls espaces topologiques).

\subsection{}\etiquette[0.5]{IV.subsec0.5}
Il nous a semblé utile d'inclure dans cet exposé général sur
les topos un assez grand nombre d'exemples, dont beaucoup n'ont que
des rapports\pageoriginale lointains avec le but initial que se
proposait ce séminaire, (c'est-à-dire l'étude de la
cohomologie étale). Le lecteur pressé, intéressé
exclusivement par la cohomologie étale, pourra bien entendu
omettre sans inconvénients la lecture de ces exemples, ainsi
d'ailleurs que de la plus grande partie du présent exposé,
auquel il lui suffira de se reporter en cas de besoin.

\section{Définition et caractérisation des topos}\etiquette[1]{IV.sec1}

\begin{enonce*}{Définition 1.1}\etiquette[1.1]{IV.def1.1}
On appelle $𝒰$-topos, ou simplement topos si aucune
confusion n'est à craindre, une catégorie $E$ telle qu'il existe
un site $C\in 𝒰$ tel que $E$ soit équivalente à la
catégorie $C^{\widetilde{\phantom{,}}}$ des $𝒰$ -faisceaux
d'ensembles sur $C$.
\end{enonce*}

\setcounter{subsection}{1}
\subsubsection{}\etiquette[1.1.1]{IV.subsubsec1.1.1}
Soit $E$ un $𝒰$-topos. Nous considérons toujours $E$ comme
muni de sa topologie canonique (\hyperrefext[II][II.def2.5]{2.5}), qui en fait donc
un \emph{site}, et même, en vertu de \hyperref[IV.subsubsec1.1.2]{1.1.2} d)
ci-dessous, un $𝒰$-site (\hyperrefext[II][II.def3.0.2]{3.0.2}). Sauf mention
expresse du contraire, nous ne considérerons pas d'autre topologie
sur $E$ que celle qu'on vient d'expliciter.

\subsubsection{}\etiquette[1.1.2]{IV.subsubsec1.1.2}
On a vu dans \hyperrefext[II][II.prop4.8]{4.8}, \hyperrefext[II][II.cor4.11]{4.11} qu'un
$𝒰$-topos $E$ est une $𝒰$-catégorie
(\hyperrefext[I][I.def1.1]{1.1}) satisfaisant aux conditions suivantes:
\begin{enumerate}
\renewcommand{\theenumi}{\alph{enumi}}
\renewcommand{\labelenumi}{\theenumi)}
\item Les limites projectives finies sont représentables dans $E$.

\item Les sommes directes indexées par un élément de $𝒰$
  sont représentables dans $E$. Elles sont disjointes et
  universelles (\hyperrefext[II][II.def4.5]{4.5}).

\item Les relations d'équivalence dans $E$ sont effectives
  universelles \sisi{(\hyperrefext[I][I.sec10.9]{10.9}).}{(\hyperrefext[I][I.sec10.10]{10.10}).}

\item $E$ admet une famille génératrice (\hyperrefext[II][II.subsec4.9]{4.9}) indexée par un
  élément de $𝒰$.
\end{enumerate}

En\pageoriginale fait, nous allons voir que ces propriétés
intrinsèques \emph{caractérisent} les $𝒰$-topos:

\begin{enonce*}{Théorème 1.2}\etiquette[1.2]{IV.theo1.2}
{\rm (J. Giraud).} Soit $E$ une $𝒰$-catégorie. Les
propriétés suivantes sont équivalentes:
\begin{itemize}
\item[i)] $E$ est un $𝒰$-topos $(\hyperref[IV.def1.1]{1.1})$.

\item[ii)] $E$ satisfait aux conditions {\rm a), b), c)} et {\rm d)
  de \hyperref[IV.subsubsec1.1.2]{1.1.2}.}

\item[iii)] Les $𝒰$-faisceaux sur $E$ pour la topologie
  canonique sont représentables, et $E$ possède une petite famille
  génératrice (condition {\rm \hyperref[IV.subsubsec1.1.2]{1.1.2} d)}).

\item[iv)] Il existe une catégorie $C\in 𝒰$ et un foncteur
  pleinement fidèle $i:E\rightarrow \widehat{C}$ (où $\widehat{C}$
  désigne la catégorie des $𝒰$-préfaisceaux sur $C$)
  admettant un adjoint à gauche a qui est exact à gauche.

\item[${\rm i')}$] Il existe un site $C\in 𝒰$, tel que les
  limites projectives soient représentables dans $C$ et que la
  topologie de $C$ soit moins fine que sa topologie canonique {\rm (II 2.5),} %% étiquette !
  tel que $E$ soit équivalent à la catégorie
  $C^{\widetilde{\phantom{,}}}$ des $𝒰$-faisceaux d'ensembles
  sur $C$.
\end{itemize}
\end{enonce*}

De plus:

\begin{enonce*}{Corollaire 1.2.1}\etiquette[1.2.1]{IV.cor1.2.1}
Soit $E$ un $𝒰$-topos, $C$ une sous-catégorie pleine de
$E$, qu'on munit de la topologie induite {\rm (\hyperrefext[III][III.subsec3.1]{3.1})}
par celle de $E$ (\hyperref[IV.subsubsec1.1.1]{1.1.1}). Considérons le
foncteur
$$
E\longrightarrow C^{\widetilde{\phantom{,}}}
$$
qui associe à tout $X\in \ob E$ la restriction à $C$ du foncteur
représenté par $X$. Ce foncteur est une équivalence de
catégories si et seulement si $\ob C$ est une famille
génératrice de $E$.
\end{enonce*}

L'équivalence\pageoriginale i) $\Leftrightarrow$ iv) résulte
aussitôt de \hyperrefext[II][II.theo5.5]{5.5}, et on a déjà rappelé plus haut
qu'on a i) $\Rightarrow$ ii). Comme ${\rm i'}$) $\Rightarrow$ i)
trivialement, il reste à prouver ii) $\Rightarrow$ iii) et iii)
$\Rightarrow$ ${\rm i'}$), ce qui sera fait dans
\hyperref[IV.subsubsec1.2.4]{1.2.4} et \hyperref[IV.subsubsec1.2.3]{1.2.3} ci-dessous. Le
corollaire s'obtient alors en remarquant que le foncteur envisagé
se factorise en $E\rightarrow E^{\widetilde{\phantom{,}}}\rightarrow
C^{\widetilde{\phantom{,}}}$, de sorte que, $E\rightarrow
E^{\widetilde{\phantom{,}}}$ étant une équivalence en vertu de
(iii), la question revient à celle de déterminer quand
$E^{\widetilde{\phantom{,}}}\rightarrow C^{\widetilde{\phantom{,}}}$
est une équivalence. On conclut alors grâce au « lemme de
comparaison »  \hyperrefext[III][III.theo4.1]{4.1}, \cf démonstration
\hyperref[IV.subsubsec1.2.3]{1.2.3} ci-dessous.

\setcounter{subsection}{2} \setcounter{subsubsection}{2}
\subsubsection{}\etiquette[1.2.3]{IV.subsubsec1.2.3}
\emph{Démonstration de} iii) $\Rightarrow$ ${\rm i'}$). Soit
$\mathfrak{X}=(X_i)_{i\in I}$, $I\in 𝒰$, une petite famille
génératrice de $E$. Comme $E$ est une $𝒰$-catégorie,
l'ensemble des classes d'isomorphie de diagrammes finis dans $E$
dont les objets sont des éléments $X_i$ est $𝒰$-petit. Par
suite le plus petit ensemble $\mathfrak{X}'$ d'objets de $E$,
contenant les limites projectives finies d'objets de $\mathfrak{X}$,
est une réunion dénombrable de petits ensembles et est donc
petit~\footnote[1]{Nous raisonnons ici sur les classes d'objets de
$E$ à isomorphisme près.}. Posant
$\mathfrak{X}^{(n+1)}=\sisi{\mathfrak{X}^{(n)(1)}}{(\mathfrak{X}^{(n)})'}$, et
$\mathfrak{X}=\bigcup_n \mathfrak{X}^{(n)}$, on voit que, quitte à
augmenter la famille de générateurs, on peut supposer que
$\mathfrak{X}$ est stable par limites projectives finies. Soit
$𝒱$ un univers contenant $𝒰$ tel que $E$ soit
$𝒱$-petite. Pour tout objet $H$ de $E$, notons $I(H)$
l'ensemble $\amalg_{i\in I}\Hom(X_i, H)$. Soient $H$ un objet de $E$
et $H'\hookrightarrow H$ le sous-faisceau de $H$, pour la topologie
canonique, « réunion »  des images des morphismes
$u:X_i\rightarrow H$, $(u,i)\in I(H)$ (\hyperrefext[II][II.theo4.1]{4.1}). Comme $H'$
est un sous-faisceau d'un $𝒰$-faisceau, $H'$ est un
$𝒰$-faisceau. Il est donc représentable. Comme la famille
$\mathfrak{X}$ est génératrice et que pour tout $i\in I$,
l'application $\Hom(X_i, H')\rightarrow \Hom(X_i,H)$ est bijective,
le morphisme $H'\hookrightarrow H$ est un isomorphisme
(\hyperrefext[II][II.subsec4.9]{4.9}). Par suite (\sisi{\hyperrefext[II][II.def6.1]{6.1}}{\hyperrefext[II][II.def5.2]{5.2}}) la famille
$(u:X_i\rightarrow H)$, $(u,i)\in I(H)$, est couvrante pour la
topologie canonique de $E$. Donc\pageoriginale tout objet de $E$
peut être recouvert, pour la topologie canonique de $E$, par des
objets de $\mathfrak{X}$. Soient $C$ la sous-catégorie pleine de $E$
définie par les objets de $\mathfrak{X}$ et $u:C\hookrightarrow E$
le foncteur d'inclusion. La topologie $\Cc$ induite sur $C$ par la
topologie canonique de $E$ est moins fine que la topologie canonique
de $C$, et lorsque $𝒰$ possède un élément de cardinal
infini, les limites projectives finies sont représentables dans $C$.
Il résulte de \sisi{\hyperrefext[III][III.subsec5.1]{5.1}}{\ref{III.theo4.1}} que le foncteur $F\mapsto F\circ
u$ est une équivalence de $E$ sur la catégorie des
$𝒰$-faisceaux sur $C$ pour la topologie $\Cc$, \cqfd

\subsubsection{}\etiquette[1.2.4]{IV.subsubsec1.2.4}
\emph{Démonstration de} ii) $\Rightarrow$ iii). La démonstration
comporte quatre pas.

Soit $𝒱$ un univers contenant $𝒰$, tel que $E$ soit
un élément de $𝒱$. Soient $\widetilde{E}$ la catégorie
des $𝒱$-faisceaux sur $E$ pour la topologie canonique, et
$J_E:E\rightarrow \widetilde{E}$ le foncteur canonique.

\paragraph{}\etiquette[1.2.4.1]{IV.par1.2.4.1}
Soit $(g_i:G_i\rightarrow H)$, $i\in I\in 𝒰$, une famille
épimorphique de $\widetilde{E}$. Si les $G_i$ et les produits
fibrés $\fprod{G_i}{G_j}{H}$ sont représentables, alors le
faisceau $H$ est représentable.

En effet, la somme directe $\coprod_{i\in I}G_i$ est représentable
dans $\widetilde{E}$ (\hyperrefext[II][II.theo4.1]{4.1}) par un faisceau
représentable $G$ (propriété b) et \hyperrefext[II][II.prop4.6]{4.6}). Il en
est de même pour la somme directe $K=\coprod_{(i,j)\in I\times
I}\fprod{G_i}{G_j}{H}$. De plus, le diagramme $K\rightrightarrows
G\rightarrow H$ est exact et $K$ est le carré fibré de $G$
au-dessus de $H$. (Cette dernière propriété est vraie dans la
catégorie des préfaisceaux, donc vraie dans la catégorie des
faisceaux (\hyperrefext[II][II.theo4.1]{4.1}.) On en déduit, d'après c) et
\hyperrefext[II][II.prop4.7]{4.7}, que la faisceau $H$ est représentable.

\paragraph{}\etiquette[1.2.4.2]{IV.par1.2.4.2}
Soit $X_{\alpha}$, $\alpha\in A\in 𝒰$, une famille de
générateurs de $E$. Pour tout faisceau $H$, désignons par
$I(H)$ l'ensemble $\coprod_{\alpha \in A}\Hom_E (X_{\alpha}, H)$. La
famille $(u:X_{\alpha}\rightarrow H, (u,\alpha)\in
I(H))$\pageoriginale est épimorphique dans $\widetilde{E}$.
Lorsque $H$ est un $𝒰$-faisceau, $I(H)$ est un élément
de $𝒰$.

En effet, tout faisceau $H$ étant but d'une famille épimorphique de
morphismes dont les sources sont des faisceaux représentables (
\hyperrefext[I][I.prop3.4]{3.4} et \hyperrefext[II][II.theo4.1]{4.1}), il suffit de montrer la
première assertion lorsque le faisceau $H$ est représentable. Soit
alors $G$ l'image de la famille $(u:X_{\alpha}\rightarrow H,
(u,\alpha)\in I(H))$. Le morphisme $G\rightarrow H$ est un
monomorphisme. Par suite, on est dans la situation du premier pas,
car $\fprod{X_{\alpha}}{X_{\beta}}{G}$ est isomorphe à
$\fprod{X_{\alpha}}{X_{\beta}}{H}$ qui est représentable, et de plus
$I(H)$ est un élément de $𝒰$. On en déduit que $G$ est
représentable. Mais $X_{\alpha}$ étant une famille génératrice,
$G\rightarrow H$ est un isomorphisme. La dernière assertion est
triviale par définition des $𝒰$-faisceaux.

\paragraph{}\etiquette[1.2.4.3]{IV.par1.2.4.3}
Tout sous-faisceau d'un faisceau représentable est
représentable.

En effet, soit $G\rightarrow H$ un sous-faisceau d'un faisceau
représentable. $G$ est alors un $𝒰$-faisceau. La famille
$(u:X_{\alpha}\rightarrow G, (u,\alpha)\in I(G))$ est donc
épimorphique dans $\widetilde{E}$ et indexée par un élément
de $𝒰$. De plus, les produits fibrés
$\fprod{X_{\alpha}}{X_{\beta}}{G}$ sont isomorphes aux produits
fibrés $\fprod{X_{\alpha}}{X_{\beta}}{H}$ qui sont
représentables. Par suite, on est dans la situation du premier pas
et $G$ est représentable.

\paragraph{}\etiquette[1.2.4.4]{IV.par1.2.4.4}
Tout $𝒰$-faisceau est représentable.

En effet, en vertu du premier et du deuxième pas, il suffit de
montrer que les produits fibrés $\fprod{X_{\alpha}}{X_{\beta}}{H}$
sont représentables. Or, ces produits fibrés sont des
sous-objets des produits $\fprod{X_{\alpha}}{X_{\beta}}{}$. On
conclut donc par le troisième pas. Ceci achève la
démonstration du théorème \hyperref[IV.theo1.2]{1.2}.

\begin{enonce*}[remark]{Remarque 1.3}\etiquette[1.3]{IV.rem1.3}
Bien\pageoriginale entendu, pour un $𝒰$-topos donné $E$, il
n'y a pas en général de façon privilégiée de le représenter à
équivalence près sous la forme $C^{\widetilde{\phantom{,}}}$, avec
$C$ un \emph{petit} site; ou, ce qui revient essentiellement au même
en vertu de \hyperref[IV.cor1.2.1]{1.2.1} lorsque l'on se borne aux $C$ dont la
topologie est moins fine que la topologie canonique, il n'y a pas de
\emph{petite} famille génératrice privilégiée dans
$E$. Lorsqu'on
cesse d'imposer à $C$ la condition que $C$ soit petit, il y a par
contre (en vertu de \hyperref[IV.theo1.2]{1.2} iii)) un choix tout à fait
canonique d'un $𝒰$-site $C$ tel que $E$ soit équivalent à
$C^{\widetilde{\phantom{,}}}$, à savoir $E$ lui-même! Ceci est une
des raisons techniques pour lesquelles il n'est pas commode en
pratique de travailler seulement avec des petits sites: en fait, les
sites les plus importants de tous, savoir les $𝒰$-topos, ne
sont pas petits! De plus, les sites générateurs de topos qui
s'introduisent dans de nombreuses questions en géométrie algébrique
(voire en topologie, \cf \hyperref[IV.subsec2.5]{2.5}) ne sont pas non plus
petits; exemple: le site étale d'un schéma \sisi{(VIII.1).}{(VII.1).}
\end{enonce*}

\begin{enonce*}{Proposition 1.4}\etiquette[1.4]{IV.prop1.4}
Soient $E$ un $𝒰$-topos, $E'$ une catégorie (pas
nécessairement une $𝒰$-catégorie), $F$ un préfaisceau
sur $E$ à valeurs dans $E'$. Pour que $F$ soit un faisceau à
valeurs dans $E'$ {\rm(\hyperrefext[II][II.def6.1]{6.1}),} il faut et il suffit que
$F$ transforme $𝒰$-limites inductives dans $E$ en limites
projectives dans $E'$.
\end{enonce*}

Soit $𝒱$ un univers contenant $𝒰$ et tel que $E'$
soit une $𝒱$-catégorie. Composant $F$ avec les foncteurs
$\Hom(X',-):E'\rightarrow 𝒱\text{-}\Ens$ définis par les
objets $X'$ de $E'$, on est ramené au cas où
$E'=𝒱\text{-}\Ens$, où $𝒱$ est un univers.
Supposons que $F$ transforme $𝒰$-limites inductives en
limites projectives, alors il résulte aussitôt des définitions
que $F$ est un faisceau, puisque une famille\pageoriginale couvrante
$X_i\rightarrow X$ dans $ℰ$, par définition de la
topologie canonique de $ℰ$, permet de considérer $X$ comme
une limite inductive du diagramme $ \left(
\fprod{X_i}{X_j}{X}\raise6pt\hbox{$\mathop{\nearrow}^{\mbox{$X_i$}}_{\displaystyle{\kern-10pt\raise-8pt\hbox{$\searrow_{\mbox{$X_j$}}$}}}$}
\right) $. Supposons que $F$ soit un faisceau, et prouvons qu'il
transforme $𝒰$-limites inductives en limites projectives. Si
$𝒱\subset 𝒰$, on peut supposer $𝒱=𝒰$
et il suffit d'appliquer le critère \hyperref[IV.theo1.2]{1.2} iii). Pour
traiter le cas général, il faut travailler un peu plus. Soient
$C$ une sous-catégorie pleine de $E$ engendrée par une famille
génératrice indexée par un élément de $𝒰$, et
$u:C\rightarrow E$ le foncteur d'inclusion. Munissons $C$ de la
topologie induite par la topologie canonique de $E$ \sisi{(III 4.5).}{(III 3.5).}
Notons $C^{\widetilde{\phantom{,}}}_{𝒰}$,
$C^{\widetilde{\phantom{,}}}_{𝒱}$ les catégories de
faisceaux sur $C$, $E^{\widetilde{\phantom{,}}}_{𝒱}$ la
catégorie des $V$-faisceaux sur $E$ pour la topologie canonique,
$J_E:E\rightarrow E_{𝒱}$ le foncteur canonique,
$i_{𝒰,𝒱}:C^{\widetilde{\phantom{,}}}_{𝒰}\rightarrow
C^{\widetilde{\phantom{,}}}_{𝒱}$ le foncteur d'inclusion. Le
diagramme:
\[\vcenter{
\xymatrix@C=3.25cm@R=1.25cm{
C^{\widetilde{\phantom{a}}}_{\underline{U}} \ar[r]^{i_{\underline{U},\underline{V}}} & C^{\widetilde{\phantom{a}}}_{\underline{V}}\\
E  \ar[r]^{J_E} \ar[u]&
E^{\widetilde{\phantom{a}}}_{\underline{V}}\ar[u]\quad,
}}\leqno{(*)}
\]
où les flèches verticales sont induites par le foncteur
$F\mapsto F\circ u$, est commutatif. Il résulte de la construction
explicite du foncteur faisceau associé (\hyperrefext[II][II.sec3]{3}) et de
\hyperrefext[II][II.theo4.1]{4.1} que le foncteur $i_{𝒰,𝒱}$ commute
aux $𝒰$-limites inductives. De plus, il résulte de
\hyperrefext[III][III.subsec5.1]{5.1} et de \hyperref[IV.cor1.5]{1.5} que les flèches
verticales de $(\ast)$ sont des équivalences de catégories. Par
suite $J_E:E\rightarrow E^{\widetilde{\phantom{,}}}_{𝒱}$
commute aux $𝒰$-limites inductives. Pour tout objet $X$ de
$E$, on a:
$$
F(X)\simeq
\Hom_{E^{\widetilde{\phantom{,}}}_{𝒱}}(J_E(X),F){\quoi}.
$$
Donc\pageoriginale $F$ transforme les $𝒰$-limites inductives
de $E$ en limites projectives.

\begin{enonce*}{Corollaire 1.5}\etiquette[1.5]{IV.cor1.5}
Soient $E$ un $𝒰$-topos, $E'$ une $𝒰$-catégorie,
$f:E\rightarrow E'$ un foncteur. Pour que $f$ admette un adjoint à
droite, il faut et il suffit que $f$ commute aux $𝒰$-limites
inductives.
\end{enonce*}

C'est une conséquence immédiate de \hyperref[IV.prop1.4]{1.4} pour le cas
$E'=𝒰\text{-}\Ens$.

\begin{enonce*}{Corollaire 1.6}\etiquette[1.6]{IV.cor1.6}
Soient $E$, $E'$ deux $𝒰$-topos, $f:E\rightarrow E'$ un
foncteur. Les conditions suivantes sont équivalentes:
\begin{enumerate}
\renewcommand{\theenumi}{\roman{enumi}}
\renewcommand{\labelenumi}{\theenumi)}
\item $f$ commute aux $𝒰$-limites inductives.

\item $f$ admet un adjoint à droite.

\item $f$ est continu {\rm(\hyperrefext[III][III.def1.1]{1.1}).}
\end{enumerate}
\end{enonce*}

L'équivalence i) $\Leftrightarrow$ ii) a été vue dans
\hyperref[IV.cor1.5]{1.5}. Pour prouver i) $\Leftrightarrow$ iii), appliquons
la définition des foncteurs continus, en choisissant un univers
$𝒱$ tel que $𝒰\in 𝒱$ (donc $E$, $E'$ sont
des sites $𝒱$-petits). Il faut exprimer que pour tout
$𝒱$-faisceau $F$ sur $E'$, le composé $F\circ f$ est un
faisceau sur $E$, c'est-à-dire (\hyperref[IV.prop1.4]{1.4}) qu'il transforme
$𝒰$-limites inductives en limites projectives. Il suffit
pour ceci que l'on ait i), puisque en vertu de \hyperref[IV.prop1.4]{1.4} $F$
lui-même transforme $𝒰$-limites inductives en limites
projectives; c'est aussi nécessaire comme on voit en prenant pour
$F$ un foncteur représentable.

\begin{enonce*}{Corollaire 1.7}\etiquette[1.7]{IV.cor1.7}
Avec les notations de $\hyperref[IV.cor1.6]{1.6}$, pour que $f$ soit le
foncteur image inverse $u^{\ast}$ pour un morphisme de topos
$(\hyperref[IV.def3.1]{3.1})$ $u:E'\rightarrow E$, il faut et il suffit que
$f$ soit exact à gauche et transforme familles épimorphiques en
familles épimorphiques.
\end{enonce*}

La\pageoriginale nécessité est triviale par définition (NB
tout foncteur exact à droite transforme épimorphisme en
épimorphisme). La suffisance résulte de \sisi{\hyperrefext[III][III.prop2.6]{2.6}}{\hyperrefext[III][III.prop1.6]{1.6}}, qui
implique que $f$ est continu sous les conditions indiquées, donc
commute aux $𝒰$-limites inductives en vertu de
\hyperref[IV.cor1.6]{1.6}.

\begin{enonce*}[remark]{Remarque 1.8}\etiquette[1.8]{IV.rem1.8}
Avec les notations de \hyperref[IV.cor1.5]{1.5}, on peut montrer que $f$ admet
un adjoint à gauche si et seulement si il commute aux
$𝒰$-limites projectives. En d'autres termes, un foncteur
\emph{covariant} $E\rightarrow (𝒰\text{-}\Ens)$ est
représentable si et seulement si il commute aux $𝒰$-limites
projectives \footnote[1]{Un énoncé plus général se trouve dans
\hyperrefext[I][I.cor8.12.8]{8.12.8}, \hyperrefext[I][I.exer8.12.9]{8.12.9}, l'esquisse de démonstration
qui suit a) à c) correspondant à la démonstration donnée dans
\loccit}. Nous indiquons seulement le principe de la démonstration,
qui se fait en deux étapes:
\begin{enumerate}
\renewcommand{\theenumi}{\alph{enumi}}
\renewcommand{\labelenumi}{\theenumi)}
\item Des arguments standards [5, ${\rm n^{\circ}}$ 195, \S 3]
montrent que si $F$ commute aux $\varprojlim$, il est
proreprésentable par un système projectif strict $(T_i)_{i\in
I}$, où $I$ est un ensemble ordonné filtrant, \emph{pas
nécessairement petit.} On peut supposer que si $i>j$, alors
$T_i\rightarrow T_j$ n'est pas un isomorphisme, et sous cette
hypothèse, $F$ est représentable si et seulement si $I$ est
petit (ce qui implique en fait que le système projectif est
essentiellement constant).

\item Pour prouver que $I$ est petit, sachant que pour tout objet $X$
de $E$, l'ensemble $F(X)=\varinjlim_i \Hom(T_i,X)$ l'est, il suffit
de disposer d'une \emph{petite famille cogénératrice}
$(X_j)_{j\in J}$ (\ie qui est génératrice pour la catégorie
opposée $E^{\circ}$). Or on montre que dans un $𝒰$-topos
$E$ existe toujours une petite famille cogénératrice.

\item Pour prouver ce dernier point, on note par des arguments
standards [Toh] que tout objet $X$ de $E$ admet un monomorphisme
dans un objet « injectif »; puis que pour toute famille
génératrice $(L_{\alpha})$ de $E$, si\pageoriginale on plonge
ainsi chaque $L_{\alpha}$ dans un objet injectif $I_{\alpha}$, la
famille $(I_{\alpha})$ est cogénératrice.
\end{enumerate}
\end{enonce*}

\section{Exemples de topos}\etiquette[2]{IV.sec2}

\setcounter{subsection}{-1}
\subsection{}\etiquette[2.0]{IV.subsec2.0}
Nous avons réuni dans le présent numéro un assez grand nombre
d'exemples typiques de topos, que le lecteur aura déjà eu l'occasion
de rencontrer par ailleurs, et qui sont destinés à lui faciliter
l'accès au « yoga »  des topos. Pour d'autres exemples (tirés de
la géométrie algébrique) de topologies sur des sites, donnant lieu à
autant de topos, il pourra consulter \SGA 3 IV 6, et (pour le topos
étale) l'exposé VII du présent séminaire. Comme nous ne référerons
guère par la suite au présent numéro que pour des questions de
notations ou de terminologie, nous laissons au lecteur le soin de
faire à titre d'exercice la vérification des énoncés dont nous avons
assorti ces exemples pour son instruction générale. Tous les
exemples du présent numéro seront précisés dans 4, où on examinera
leur dépendance fonctorielle par rapport aux données, et dans les
numéros suivants à titre d'illustration des notions générales
relatives aux topos.

\subsection{Topos associé à un espace topologique}\etiquette[2.1]{IV.subsec2.1}

Soient $X$ un petit espace topologique, $\Ouv(X)$ la catégorie des
ouverts de $X$, munie de la topologie canonique \sisi{(III 2.9.1).}{(III 1.9.1).} Nous
désignerons par $\Top(X)$ le topos des $𝒰$-faisceaux sur
$\Ouv(X)$. Ce topos est équivalent à la catégorie des espaces
topologiques étalés au-dessus de $X$, en associant à tout tel espace
$X'$ sur $X$ le faisceau $U\mapsto \Gamma(X'/U)=\Hom_X(U,X')$ sur
$\Ouv(X)$ [TF]\nde{Ici [TF] indique la référence [8]}. On ne pourra pas s'empêcher de noter parfois, par
abus de\pageoriginale langage, par la même lettre $X$ le topos
$\Top(X)$ défini par l'espace topologique $X$.

C'est évidemment l'exemple précédent qui a servi
principalement de guide et de support intuitif pour le
développement de la théorie des topos. On fera attention
cependant que les topos déduits des espaces topologiques sont de
nature très particulière, dû au fait notamment qu'ils ont
été décrits par des sites $C=\Ouv(X)$ où \emph{tous les
morphismes sont des monomorphismes} (donc dont la catégorie
sous-jacente se réduit à un ensemble préordonné). De ceci
résulte en particulier que les faisceaux représentés par les
objets de $C$ sont des sous-faisceaux du faisceau final, et par
suite que \emph{les sous-faisceaux du faisceau final forment une
famille génératrice du topos envisagé}. Cette propriété
n'est pas partagée par la plupart des topos qui s'introduisent de
façon naturelle en géométrie algébrique ou en algèbre,
\cf exemples plus bas. Elle est à peu de choses près
caractéristique des topos de la forme $\Top(X)$ (\cf
\hyperref[IV.exer7.1.9]{7.1.9} plus bas).

On vérifie facilement que l'application $\Ouv(X)\rightarrow
\Top(X)$, qui associe à tout ouvert de $X$ le faisceau qu'il
représente, est une \emph{bijection} de $\Ouv(X)$ avec l'ensemble
des sous-objets de l'objet final de $\Top(X)$, cette bijection étant
même un isomorphisme pour les structures d'ordre naturelles, \ie
induisant un isomorphisme des catégories correspondantes. Cela
suggère qu'il doit être possible de reconstituer à homéomorphisme
près l'espace topologique $X$, lorsqu'on connaît $\Top(X)$ à
équivalence près. Nous verrons plus bas (\hyperref[IV.subsec4.2]{4.2}) qu'il
en est bien ainsi, moyennant une légère restriction sur $X$.

\subsection{Topos ponctuel ou final, et topos vide ou initial}\etiquette[2.2]{IV.subsec2.2}\pageoriginale

Lorsque $X$ est un espace topologique réduit à un seul point, le
foncteur
$$
F\longrightarrow F(X):\Top(X)\longrightarrow 𝒰\mbox{-}\Ens
$$
est une équivalence de catégories. Ceci montre en particulier
que la \emph{catégorie} $𝒰$-$\Ens$ est un
$𝒰$-\emph{topos.} Nous avons vu sur des exemples dans \Exp
II que cet $𝒰$-topos est typique du point de vue
propriétés d'exactitude, en ce que la vérification de beaucoup
de propriétés (notamment des propriétés d'exactitude) des
topos généraux se ramène à ce topos particulier.
L'interprétation que nous en donnons ici en termes de l'espace
topologique ponctuel justifie l'abus de langage consistant à
appeler \emph{topos ponctuel} un topos équivalent à la
catégorie $𝒰$-$\Ens$ (bien que, en tant que catégorie,
il ne soit pas du tout équivalent à la catégorie ponctuelle
!). C'est la terminologie qui correspond à l'intuition
géométrique correcte du rôle joué par ces topos. On appelle
parfois, par abus de langage également, \emph{topos final} un
topos ponctuel, \cf \hyperref[IV.subsec4.3]{4.3}; on dira « \emph{le} topos
final »  pour le topos ($𝒰$-$\Ens$).

Lorsque $X$ est réduit à l'espace topologique vide, donc
$\Ouv(X)$ à la catégorie ponctuelle, alors un préfaisceau $F$
sur $\Ouv(X)$ est un faisceau si et seulement si sa valeur en
l'unique objet $\phi$ de $\Ouv(X)$ est un ensemble réduit à un
point. Il s'ensuit que $\Top(∅)$ est isomorphe à la catégorie
des $𝒰$-ensembles réduits à un point, catégorie qui
est équivalente à la catégorie ponctuelle. De ceci on conclut
en particulier que la catégorie ponctuelle (ainsi que toute
$𝒰$-catégorie équivalente à celle-ci) est un
$𝒰$-topos. On l'appelle parfois, par abus de\pageoriginale
langage, le \emph{topos vide} ou \emph{topos initial} (\cf
\hyperref[IV.subsec4.4]{4.4}); on prendra garde qu'il n'est pas équivalent
à la catégorie vide.

\subsection{Topos associé à un espace à opérateurs}\etiquette[2.3]{IV.subsec2.3}

Soient $X$ un espace topologique, $G$ un groupe discret opérant
sur $X$ par homéomorphismes. On a défini alors dans \sisi{[Toh.
\hyperref[IV.subsec5.1]{5.1}]}{[Toh]\nde{Ici [Toh] indique la référence [4].} \hyperref[IV.subsec5.1]{5.1}} la catégorie des $G$-\emph{faisceaux sur} $X$,
ou comme nous dirons aussi, des faisceaux sur $(X,G)$; ce sont les
faisceaux (d'ensembles) sur $X$, \emph{munis} d'opérations de $G$
compatibles avec celles de $G$ sur $X$. On constate aussitôt, à
l'aide du critère de Giraud \hyperref[IV.theo1.2]{1.2} iii), que cette %% \XXX ii) est probablement plus facile à vérifier
catégorie est un topos (NB $𝒰$ est sous-entendu dans tout
ceci), qu'on notera simplement $\Top(X,G)$. Lorsque $G$ est réduit
au groupe unité, on retrouve l'exemple \hyperref[IV.subsec2.1]{2.1}; lorsque
$X$ est réduit à l'espace ponctuel, on trouve le topos des
ensembles sur lesquels $G$ opère à gauche (ou
$G$-\emph{ensembles}), appelé aussi \emph{topos classifiant du
groupe discret} $G$, et noté $B_G$. On vérifie facilement que le
seul sous-objet de l'objet final $e$ du topos $B_G$ est $e$ ou le
faisceau vide $\phi$, en particulier, si $G$ n'est pas réduit au
groupe unité, les sous-objets de l'objet final du topos
classifiant $B_G$ ne forment pas une famille génératrice de
$B_G$. Donc $B_G$ n'est pas équivalent alors à un topos du type
$\Top(X)$ envisagé dans \hyperref[IV.subsec2.1]{2.1}.

La notion de $G$-faisceau était introduite dans \loccit pour
développer la théorie cohomologique des $G$-faisceaux
abéliens. Interprétant ces derniers comme les faisceaux
abéliens du topos $(X,G)$, ladite théorie se trouve incluse dans
celle de \Exp V, développée dans le cadre des topos
généraux.

On\pageoriginale peut se proposer plus généralement d'attacher
un topos approprié à un espace topologique $X$, muni d'un groupe
topologique $G$ (pas nécessairement discret) d'automorphismes, qui
donnerait naissance à une théorie cohomologique adéquate. De
même dans le contexte des variétés différentiables, ou
analytiques réelles ou complexes, ou des schémas. C'est
effectivement possible, \cf \hyperref[IV.subsec2.5]{2.5} ci-dessous.

\subsection{Topos classifiant d'un Groupe}\etiquette[2.4]{IV.subsec2.4}

Soit $E$ un topos, et $G$ un Groupe de $E$. Soit $(E,G)$ la
catégorie des objets de $E$ sur lesquels $G$ opère. On voit
aussitôt, grâce au critère de Giraud, que c'est un topos. On
l'appelle \emph{topos classifiant} du Groupe $G$, et on le note
$B_G$. Lorsque $E$ est le topos ponctuel (\hyperref[IV.subsec2.2]{2.2}) \ie
lorsque $G$ est un groupe ordinaire, on retrouve le topos
classifiant de \hyperref[IV.subsec2.3]{2.3}.

La terminologie adoptée ici se justifie, du fait que le topos
$B_G$ joue un rôle universel pour la classification des « torseurs » (ou fibrés principaux homogènes) sous $G$, ou plus
généralement sous les $G_{E'}=f^{\ast}(G)$, où $E'$ est un
topos « au-dessus de $E$ »  \ie muni d'un morphisme
$f:E'\rightarrow E$ (\cf \hyperref[IV.def3.1]{3.1} ci-dessous). Ce rôle,
explicité dans [3 Chap V] ou dans \hyperref[IV.exer5.9]{5.9} plus bas,
montre que $B_G$ joue, dans le contexte des topos, le même rôle
que les classiques espaces classifiants des groupes topologiques en
théorie homotopique des espaces topologiques. Ces derniers peuvent
être regardés (\cf \hyperref[IV.subsec2.5]{2.5}) comme une version
affaiblie des premiers, obtenue en ne retenant du topos classifiant
que le seul « type d'homotopie »  dudit topos, en un sens
convenable qu'il n'y a pas lieu de préciser ici.

\subsection{« Gros site »  et « gros topos »  d'un espace
topologique. Topos classifiant d'un groupe
topologique}\footnote[1]{L'introduction de ces sites et topos est
due à M. \textsc{Giraud}, qui a mis également en évidence
leurs avantages sur le « petit »  site
traditionnel.}\etiquette[2.5]{IV.subsec2.5}\pageoriginale

Soit $U-\Esp$ ou simplement $(\Esp)$ la catégorie des espaces
topologiques $\in 𝒰$. On sait que dans $(\Esp)$ les limites
projectives finies sont représentables. Considérons sur $(\Esp)$
la prétopologie (\sisi{\hyperrefext[I][I.condef1.3]{1.3}}{\ref{II.def1.3}}) pour laquelle $\Cov(X)$ est
l'ensemble des familles surjectives d'immersions ouvertes
$u_i:X_i\rightarrow X$. Nous considérerons $(\Esp)$ comme un site
à l'aide de la topologie engendrée par la prétopologie
précédente. Pour tout objet $X$ de $(\Esp)$, considérons la
catégorie
$$
(\Esp)_{/X}
$$
des objets de $(\Esp)$ au-dessus de $X$, \ie des espaces
topologiques au-dessus de $X$, comme un site, grâce à la topologie
induite par celle de $(\Esp)$ via le foncteur d'oubli
$(\Esp)_{/X}\rightarrow (\Esp)$ (\hyperrefext[III][III.prop5.2]{5.2} 4)). Ce site est
appelé le \emph{gros site} associé à $X$. On fera attention qu'il
n'est pas $\in 𝒰$; ce n'est pas non plus un $𝒰$-site
au sens de \hyperrefext[II][II.def3.0.2]{3.0.2}, donc des précautions sont nécessaires
pour lui appliquer les résultats habituels. Pour pallier cet
inconvénient, on peut choisir un univers $𝒱$ tel que
$𝒰\in 𝒱$, de sorte que $(\Esp)_{/X}$ devient un
$𝒱$-site, et on peut travailler avec le $𝒱$-topos
associé $(\Esp)_{/X}$, qui pourra être noté $\TOP(X)$ et sera appelé
le \emph{gros topos de} $X$. Si on répugne à agrandir $𝒰$,
on peut choisir un cardinal $c$ majorant les cardinaux de $X$ et de
tous les espaces topologiques qu'on compte faire intervenir dans les
raisonnements (le plus souvent, 
\sisi{$\Sup(\card X, \card ℛ)$}{$\Sup(\card X, \card \mathbb{R})$}
sera suffisant !), et on remplace $(\Esp)_{/X}$ par la
sous-catégorie $(\Esp)'_{/X}$ formée des $X'$ sur $X$ tels que
$\card X'\leq c$, munie de la topologie induite, et on note
$\TOP(X)$ le topos des faisceaux sur ce\pageoriginale site. Pour
fixer les idées, supposons que ce soit la première définition qui
ait été adoptée.

L'avantage du gros topos de $X$ sur le petit, c'est que le site qui
le définit contient $(\Esp)_{/X}$ comme sous-catégorie pleine; comme
la topologie de ce site est manifestement moins fine que la
canonique, on voit que le foncteur canonique de $(\Esp)_{/X}$ dans
$\TOP(X)$, associant à tout espace $X'$ sur $X$ le faisceau qu'il
représente, est \emph{pleinement fidèle.} Par suite, \emph{un espace
$X'$ sur $X$ est connu à $X$-isomorphisme près quand on connaît le
faisceau $(\in \TOP(X))$ qu'il définit}; donc la notion de faisceau
sur (le gros site de) $X$ peut être considéré comme une
\emph{généralisation} de celle d'espace topologique au-dessus de
$X$, à l'aide de laquelle toutes les constructions de la théorie des
faisceaux prennent un sens pour les espaces topologiques sur $X$.

Ainsi, lorsque $G$ est un objet-groupe de la catégorie $(\Esp)_{/X}$
des espaces topologiques au-dessus de $X$, on peut lui associer le
topos classifiant $B_G$ (\hyperref[IV.subsec2.4]{2.4}), d'où des groupes de
cohomologie classifiante, des groupes d'homotopie classifiante \etc
(définis comme les invariants correspondants du $𝒱$-topos
$B_G$). En particulier, lorsque $X$ est un espace ponctuel, $G$
s'identifie à un groupe topologique ordinaire. On peut vérifier,
moyennant les conditions locales habituelles assurant que la
cohomologie singulière des produits cartésiens $G^n$ coïncide avec
la cohomologie au sens des faisceaux (pour des coefficients
constants, disons), par exemple si $G$ est localement contractible,
que la cohomologie du topos classifiant de $G$ est canoniquement
isomorphe à celle de l'espace classifiant de $G$ au sens des
topologues.

L'introduction\pageoriginale des topos classifiants (via les « gros sites ») a sur les espaces classifiants l'avantage de fournir
une théorie plus riche, puisqu'ils fournissent notamment des
invariants cohomologiques utiles pour des coefficients plus
généraux que les coefficients constants ou localement constants.
De plus, la définition envisagée ici s'adapte de façon
évidente aux autres contextes habituels: variétés
différentiables, variétés ou espaces analytiques (réelles ou
complexes, au choix), schémas. Ce point de vue permet notamment de
faire le lien entre l'étude des classes caractéristiques du
point de vue traditionnel et du point de vue « arithmétique »,
en considérant les « groupes classiques »  comme provenant de
schémas définis sur l'anneau des entiers; \cf [\hyperref[IV.sec7]{7}]
pour des indications dans ce sens. De même, les résultats
généraux de J. \textsc{Giraud} [\hyperref[IV.sec3]{3}] sur la
classification des extensions de Groupes, développés dans le
cadre très général et très souple des topos, peuvent grâce aux « gros topos » se spécialiser en des résultats sur la
classification d'extensions de groupes topologiques, ou de groupes
de Lie réels ou complexes, qui ne semblaient guère connus des
topologues que dans le cas des extensions à noyau abélien
[\hyperref[IV.sec11]{11}].

\subsection{Topos de la forme $\widehat{C}$}\etiquette[2.6]{IV.subsec2.6}

Soit $C$ une petite catégorie. Alors la catégorie $\widehat{C}$
des $U$-préfaisceaux sur $C$ est évidemment un
$𝒰$-topos, puisqu'elle est de la forme
$C^{\widetilde{\phantom{,}}}$, où on munit $C$ de la topologie
chaotique. Nous donnerons plus bas quelques détails sur les
relations entre $C$ et $\widehat{C}$. Notons seulement ici qu'un
topos $E$ équivalent à un topos de la forme $\widehat{C}$ est de
nature assez spéciale, du fait qu'\emph{il admet une petite
famille génératrice $(X_i)$ formée d'objets connexes
projectifs,} \ie d'objets $X$ tels que le foncteur $Y\mapsto
\Hom(X,Y)$\pageoriginale transforme épimorphismes en
épimorphismes et sommes en sommes: il suffit en effet de prendre
dans $\widehat{C}$ la famille génératrice formée des foncteurs
représentés par les $X\in \ob C$. Notons d'ailleurs que si dans
un topos $E$ on a une famille couvrante $X_i\rightarrow X$, avec des
$X_i$ qui sont projectifs et connexes, alors toute autre famille
couvrante de $X$ est majorée (\hyperrefext[I][I.sec4.3.2]{4.3.2}, \hyperrefext[I][I.sec4.3.3]{4.3.3})
par la précédente. Par suite, dans un topos $E$ de la forme
$\widehat{C}$ tout objet $X$ admet une famille couvrante majorant
toutes les autres. Un topos de la forme $\Top(X)$
(\hyperref[IV.subsec2.1]{2.1}), avec $X$ un espace topologique dont les points
sont fermés, n'a la propriété précédente que si $X$ est
discret.

Lorsque la catégorie $C$ a un seul objet, $C$ s'identifie à un
monoïde $G$. Un préfaisceau sur $C$ s'identifie alors à un ensemble
sur lequel $G$ opère \emph{à droite} (puisque c'est un foncteur
$G^{\circ}\rightarrow (\Ens)$), et $\widehat{C}$ est le topos des
ensembles à monoïde d'opérateurs à droite, qu'on pourra aussi noter
$B_{G^{\circ}}$, compte tenu de \hyperref[IV.subsec2.3]{2.3}: c'est le
\emph{topos des ensembles à monoïde d'opérateurs $G^{\circ}$ (le
monoïde opposé à $G$)}. Lorsque $G$ est un groupe, utilisant
l'isomorphisme $g\mapsto g^{-1}$ de $G$ sur $G^{\circ}$, on retrouve
le topos classifiant $B_G$ de \hyperref[IV.subsec2.3]{2.3}.

\subsection{Topos classifiant d'un pro-groupe}\etiquette[2.7]{IV.subsec2.7}

\subsubsection{}\etiquette[2.7.1]{IV.subsubsec2.7.1}
Soit $𝒢=(G_i)_{i\in I}$ un système projectif de groupes,
avec $G_i$, $I\in 𝒰$. On suppose le système projectif
\emph{strict}. \ie les morphismes de transition $G_j\rightarrow G_i$
surjectifs. Si $E$ est un ensemble, on appelle \emph{opération de
$𝒢$ sur $E$} (à gauche, disons) la structure suivante: a)
une famille $(E_i)_{i\in I}$ de parties de $E$, de réunion $E$; b)
pour tout $i\in I$, une opération du groupe $G_i$ sur l'ensemble
$E_i$; on suppose de plus ces données soumises à\pageoriginale la
condition suivante: pour $j\geq i$, $E_i$ est le sous-ensemble de
$E_j$ formé des éléments fixes sous le groupe noyau de
$G_j\rightarrow G_i$. On dit aussi que $𝒢$ \emph{opère sur}
$E$ (à gauche) si on s'est donnée une opération de $G$ sur $E$ (à
gauche). Les ensembles $\in 𝒰$ munis d'une opération de
$𝒢$ forment une catégorie de façon évidente. On constate
aussitôt, grâce au critère de Giraud, que cette catégorie est un
$𝒰$-topos. On le note $B_{𝒢}$ et on l'appelle
\emph{topos classifiant de} $𝒢$. Lorsque $I$ admet un objet
initial $i_{\circ}$, posant $G=G_{i_\circ}$, on retrouve le topos
classifiant de \hyperref[IV.subsec2.3]{2.3}.

\subsubsection{}\etiquette[2.7.2]{IV.subsubsec2.7.2}
Un autre exemple important est celui où les groupe $G_i$ sont
finis, de sorte que
$$
G=\varprojlim G_i
$$
est un groupe topologique compact totalement discontinu, ou
\emph{groupe profini.} Une opération de $𝒢$ sur $E$
revient alors à une opération de $G$ sur $E$ qui est continue,
ou ce qui revient au même, telle que le stabilisateur de tout
point de $E$ est un sous-groupe ouvert de $G$. Le topos classifiant
$B_{𝒢}$ sera aussi noté $B_{G}$, où, bien entendu, $G$
doit être considéré comme muni de sa topologie profinie.

\subsubsection{}\etiquette[2.7.3]{IV.subsubsec2.7.3}
Il est facile de vérifier, utilisant les remarques de
\hyperref[IV.subsec2.6]{2.6}, que le topos $B_{𝒢}$ défini par un
système projectif strict $𝒢=(G_i)_{i\in I}$ de groupes
n'est équivalent à un topos de la forme $\widehat{C}$ que si ce
système projectif est essentiellement constant; dans le cas d'un
groupe profini, cela signifie que ce groupe est en fait fini.

\subsubsection{}\etiquette[2.7.4]{IV.subsubsec2.7.4}
L'interprétation géométrique suivante du topos classifiant
$B_G$ d'un groupe discret $G$ est utile, pour donner une intuition
géométrique correcte de ces topos. (\Cf aussi, dans le même
sens, \hyperref[IV.subsec4.5]{4.5}, \hyperref[IV.subsec5.8]{5.8}, \hyperref[IV.exer5.9]{5.9} et
\hyperref[IV.subsec7.2]{7.2} ci-dessous.) Soit $X$ un espace topologique
connexe, localement connexe et\pageoriginale localement simplement
connexe, $x$ un point de $X$, $G$ son groupe fondamentale en $x$.
(NB il est connu qu'à isomorphisme près, tout groupe discret $G$
peut s'obtenir ainsi.) Alors la théorie de Galois des
revêtements de $X$ fournit une équivalence entre la catégorie
$B_G$ des $G$-ensembles, et la catégorie des \emph{revêtements
étales} de $X$, \ie des espaces $X'$ sur $X$ qui sont localement
$X$-isomorphes à des $X$-espaces de la forme $X\times I$, où $I$
est un espace discret. (comparer \SGA 1 V 4,5). Cette dernière
peut d'ailleurs s'interpréter comme la catégorie des faisceaux
localement constants sur $X$, \ie la catégorie des objets
localement constants (IX 2.0) du topos $\Top(X)$.

Lorsque $G$ est un groupe profini, on a une interprétation
géométrique analogue de $B_G$, comme la catégorie des
$X$-schémas qui sont sommes de revêtements finis étales d'un
schéma $X$ connexe, muni d'un point géométrique $x$ et d'un
isomorphisme $G\simeq \pi_1(X,x)$. Il est connu encore que tout
groupe profini peut s'obtenir comme groupe fondamental d'un schéma
connexe convenable (spectre d'un corps si on veut). Enfin, on
rencontre également des pro-groupes (pas nécessairement profinis
ni essentiellement constants) pour la classification des
revêtements des espaces connexes et localement connexes qui ne
sont pas localement simplement connexes, ou la classification des
revêtements étales pas nécessairement finis ou ind-finis de
schémas connexes non normaux. Pour ce dernier cas, \cf \SGA 3 X 6.

\begin{enonce*}[remark]{Exercice 2.7.5}\etiquette[2.7.5]{IV.exer2.7.5}
Définir pour un topos $E$ la notion de connexité, de locale
connexité~\footnote[1]{\cf \sisi{\hyperref[IV.exer8.7]{8.7} 1).}{\hyperref[IV.exer8.7]{8.7} l).}}, de simple
connexité et de simple connexité locale. Définir la notion
d'objet constant et localement constant de $E$ (\cf IX 2.0). Soit
$f:E'\rightarrow E$\pageoriginale un morphisme de topos
(\hyperref[IV.def3.1]{3.1}), avec $E'$ simplement connexe (par exemple $E'$ le
topos ponctuel (\hyperref[IV.subsec2.2]{2.2})), et $E$ connexe et localement
connexe. Définir un pro-groupe strict $\pi_1(E,f)=(G_i)_{i\in
I}=𝒢$ (appelé \emph{pro-groupe fondamental de $E$ en $f$})
et une équivalence de catégories de $B_{𝒢}$ avec la
catégorie des objets localement constants de $E$.\footnote[2]{On
pourra s'inspirer de \SGA 3 X 6.}. Montrer que lorsque $E$ est
localement simplement connexe, $\pi_1(E,f)$ est essentiellement
constant et s'identifie donc à un groupe discret ordinaire
$\pi_1(E,f)$, qui s'appelle le \emph{groupe fondamental de $E$ en
$f$}. Montrer que tout pro-groupe strict $𝒢$ est isomorphe
(comme pro-groupe) au pro-groupe fondamental d'un topos connexe et
localement connexe convenable en un $f$ convenable, avec $E'$ le
topos ponctuel (prendre $E=B_{𝒢}$, et $f:E'\rightarrow E$
défini par le foncteur oubli $f^{\ast}:E\rightarrow E'=(\Ens)$).
Lorsque $𝒢$ est essentiellement constant, \ie isomorphe
(comme pro-groupe) à un groupe discret ordinaire, prouver qu'on
peut prendre ci-dessus $E$ localement simplement connexe (prendre
encore $E=B_G$).
\end{enonce*}

\subsection{Exemple d'un faux topos}\etiquette[2.8]{IV.subsec2.8}

Soit $𝒢=(G_i)_{i\in I}$ un pro-groupe strict, où $I$ est
ordonné filtrant et où $i>j$ implique que $G_i\rightarrow G_j$
n'est pas un isomorphisme. Supposons que l'on ait $\card(I)\not\in
𝒰$. Considérons la catégorie des ensembles $E\in
𝒰$ sur lesquels $𝒢$ opère à gauche
(\hyperref[IV.subsubsec2.7.1]{2.7.1}). C'est une $𝒰$-catégorie, et on
voit comme dans \hyperref[IV.subsubsec2.7.1]{2.7.1} que cette catégorie
satisfait aux conditions a), b), c) de \hyperref[IV.subsubsec1.1.2]{1.1.2}.
Cependant ce n'est pas un $𝒰$-topos, car on voit qu'il
n'admet pas de famille génératrice qui soit $𝒰$-petite.
On voit de même que ce n'est un $𝒱$-topos pour aucun
univers $𝒱$.

\section{Morphismes de topos}\etiquette[3]{IV.sec3}\pageoriginale

\begin{enonce*}{Définition 3.1}\etiquette[3.1]{IV.def3.1}
Soient $E$ et $E'$ deux $𝒰$-topos. On appelle morphisme\nde{Certains auteurs 
parlent de \emph{morphisme géométrique}.} de
$E$ dans $E'$, ou parfois (par abus de langage) application continue
de $E$ dans $E'$, un triple $u=(u_{\ast}, u^{\ast},\varphi)$,
formé de foncteurs
$$
u_{\ast}:E\longrightarrow E'\quav u^{\ast}:E'\longrightarrow E
$$
et d'un isomorphisme $\varphi$ « d'adjonction »  de bifoncteurs
en $X'\in \Ob E'$, $Y\in \Ob E$:
$$
\varphi:\Hom_E (u^{\ast}(X'),Y)\xrightarrow{\sim} \Hom_{E'}(X',
u_{\ast}(Y)){\quoi},
$$
le foncteur $u^{\ast}$ étant de plus soumis à la condition
d'être exact à gauche, \ie de commuter aux limites projectives
finies. Le foncteur $u_{\ast}$ est appelé le foncteur image
directe pour le morphisme de topos $u$, le foncteur $u^{\ast}$ est
appelé le foncteur image inverse pour le morphisme de topos $u$,
l'isomorphisme $\varphi$ est appelé l'isomorphisme d'adjonction,
pour $u$.
\end{enonce*}

\setcounter{subsection}{1}
\subsubsection{}\etiquette[3.1.1]{IV.subsubsec3.1.1}
Sauf mention expresse du contraire, on désignera par la suite,
pour un morphisme de topos $u:E\rightarrow E'$, par $u_{\ast}$ et
$u^{\ast}$ les foncteurs image directe et image inverse
correspondants \footnote[1]{Il y a lieu parfois d'écrire aussi
$u^{-1}$ au lieu de $u^{\ast}$, \cf \hyperref[IV.subsubsec13.2.3]{13.2.3}.}. On notera que, $u_{\ast}$
étant adjoint à droite de $u^{\ast}$ et $u^{\ast}$ étant
adjoint à gauche de $u_{\ast}$ par l'isomorphisme d'adjonction
$\varphi$, chacun des deux foncteurs $u_{\ast}$, $u^{\ast}$
détermine l'autre à isomorphisme unique près, d'après le
sorite bien connu des foncteurs adjoints [\hyperref[IV.sec14]{14}]. En
pratique, suivant les cas, il peut être plus commode de définir
un morphisme de topos $u:E\rightarrow E'$ soit par la donnée de
$u^{\ast}:E'\rightarrow E$, soit par la donnée de
$u_{\ast}:E\rightarrow E'$; dans le premier cas, il faut simplement
vérifier que le foncteur donné $u^{\ast}$ admet\pageoriginale
un adjoint à droite, et qu'il est exact à gauche. Dans le
deuxième, que le foncteur donné $u_{\ast}$ admet un adjoint à
gauche qui est exact à gauche. Dans l'un ou l'autre cas, on
déduit de la donnée partielle, grâce au \emph{choix} d'un
foncteur adjoint et d'un isomorphisme d'adjonction, un morphisme de
topos $u:E\rightarrow E'$, et ce dernier sera « unique à
isomorphisme unique près »  en termes de la donnée $u^{\ast}$
\resp $u_{\ast}$, en un sens assez clair, et qui sera d'ailleurs
entièrement explicité plus bas (3.2.1).

\subsubsection{}\etiquette[3.1.2]{IV.subsubsec3.1.2}
Si $u:E\rightarrow E'$ est un morphisme de topos, il résulte des
propriétés des foncteurs adjoints (\hyperrefext[I][I.prop2.11]{2.11}) que le
foncteur image directe $u_{\ast}:E\rightarrow E'$ commute aux
limites projectives, et le foncteur $u^{\ast}:E'\rightarrow E$
commute aux limites inductives \footnote[2]{D'ailleurs
(\hyperref[IV.cor1.5]{1.5} et \hyperref[IV.rem1.8]{1.8}), pour un foncteur
$u_{\ast}:E\rightarrow E'$ \resp $u^{\ast}:E'\rightarrow E$ donné,
ce foncteur admet un adjoint à gauche (\resp à droite), si et
seulement si il commute aux $𝒰$-limites projectives (\resp
aux $𝒰$-limites inductives).}\label{sga4[1]}. Comme on
suppose de plus que ce dernier est exact à gauche \ie commute aux
limites projectives finies, on voit donc en particulier que
$u^{\ast}$ \emph{est exact}. On peut donc dire que \emph{c'est le
foncteur image inverse $u^{\ast}$}, dans le couple $(u_{\ast},
u^{\ast})$, \emph{qui possède les propriétés d'exactitude les plus
remarquables.} Ces propriétés assurent que pour toute espèce de
structure algébrique $\Sigma$ dont les données peuvent se décrire en
termes de données de flèches entre les ensembles de base et des
ensembles déduits de ceux-ci par application répétée d'opérations de
limites projectives finies et de limites inductives quelconques, et
pour tout « objet de $ℰ'$ muni d'une structure d'espèce
$\Sigma$ » (notion qui a un sens grâce aux propriétés d'exactitude
internes du topos $ℰ'$ (\hyperrefext[II][II.theo4.1]{4.1})), son image par
$u^{\ast}$ est muni des mêmes structures. Plutôt que d'entrer dans
la tâche peu engageante de donner un sens précis à cet énoncé et de
le justifier\pageoriginale de façon formelle, nous conseillons au
lecteur de l'expliciter et de se convaincre de sa validité pour des
espèces de structure telles que celle de groupe, d'anneau, de module
sur un anneau, de comodule sur un coanneau, de bigèbre sur un anneau,
de torseur sous un groupe. (Dans ces exemples, les trois premières
espèces de structure se définissent en termes de limites projectives
finies exclusivement, tandis que les autres notions impliquent
implicitement des constructions faisant appel également à des
limites inductives.) De plus, le foncteur $u^{\ast}$ « \emph{commute} »  à toutes les opérations fonctorielles
habituelles faites en termes de telles structures, plus précisément
à toutes les opérations qui peuvent s'exprimer en termes de
$\varinjlim$, et de $\varprojlim$ finies: constructions d'objets
libres (groupes ou modules libres, p.ex.) engendrés par un objet,
produits tensoriels (\cf \S 12 plus bas) \etcn

Quant au foncteur image directe $u_{\ast}:E\rightarrow E'$, qui
commute aux limites projectives, il « respecte »  par suite
toute structure algébrique sur un objet (ou une famille d'objets)
de $E$, définissable en termes de limites projectives
exclusivement (telles que les structures de groupe, d'anneau ou de
module sur un anneau, parmi les exemples précédents). Par
contre, le foncteur $u_{\ast}$ n'est en général pas exact à
droite \ie il ne commute pas en général aux limites inductives
finies, et même ne transforme pas en général épimorphisme en
épimorphisme (c'est d'ailleurs ce défaut d'exactitude du
foncteur $u_{\ast}$ qui est la source de ces propriétés
cohomologiques, qui seront étudiées (du point de vue de
l'algèbre homologique commutative) dans l'exposé suivant). Par
suite, il ne s'étend pas, en général, en un foncteur sur des
objets du type comodule,\pageoriginale ou bigèbre, ou torseur sous
un groupe, et ne commute pas en général à des opérations
telles que « module libre engendré », produit tensoriel de
modules, \etc

\subsubsection{}\etiquette[3.1.3]{IV.subsubsec3.1.3}
En pratique, lorsqu'on est en présence d'un foncteur
$f:E\rightarrow F$ d'un $𝒰$-topos dans un autre, il y a lieu
d'en expliciter toutes les propriétés d'exactitude, y compris
l'existence éventuelle de foncteurs adjoints à gauche ou à
droite (\cf la note de bas de page \pageref{sga4[1]}), pour parvenir
à une compréhension de la « nature géométrique »  de
$f$, compréhension qui sera généralement un guide
indispensable pour une intuition géométrique correcte de la
situation. Ainsi, s'il s'avère que $f$ commute aux limites
inductives quelconques et aux limites projectives finies, il y a
lieu d'écrire $f$ sous la forme
$$
f=u^{\ast}{\quoi},
$$
où
$$
u:F\longrightarrow E
$$
est un morphisme de topos, \ie d'interpréter $f$ comme un foncteur
« image inverse »  par une « application continue »  de
topos. Lorsque $f$ commute aux limites projectives quelconques, donc
qu'il admet un adjoint à gauche, et si ce dernier (qui a priori
commute aux limites inductives quelconques) commute \emph{de plus}
aux limites projectives finies, il y a lieu d'écrire $f$ sous la
forme
$$
f=v_{\ast}{\quoi},
$$
où
$$
v:E\longrightarrow F
$$
est un morphisme de topos. Dans certains cas, il peut arriver que
$f$ satisfasse\pageoriginale aussi bien à l'une qu'à l'autre des
deux propriétés envisagées (\cf \hyperref[IV.subsec4.10]{4.10} pour un
exemple). Dans ce cas, il y a lieu  d'introduire simultanément les
morphismes de topos
$$
u:F\longrightarrow E \quad\text{et}\quad v:E\longrightarrow F{\quoi},
$$
qui donnent lieu à une suite de trois foncteurs adjoints
(\hyperrefext[I][I.sec5.3]{5.3}):
$$
v^{\ast} \quav v_{\ast}=u^{\ast}=f\quav u_{\ast}{\quoi}.
$$
Il convient de distinguer alors soigneusement entre les morphismes
de topos $u$ et $v$, sous peine de perdre l'intuition
géométrique de la situation.

On notera à ce propos que lorsque entre deux topos $E$, $F$ on a
une suite de trois foncteurs adjoints
$$
e, f, g\; (e,g: F\rightrightarrows E, f:E\rightarrow F){\quoi},
$$
alors $f$ commute aux limites inductives et aux limites projectives
quelconques, donc il peut toujours se mettre sous la forme
$u^{\ast}$, où $u:F\rightarrow E$ est un morphisme de topos. Alors
$g$ s'écrit donc $g=u_{\ast}$. Par contre, bien sûr, $e$ ne peut
s'écrire sous la forme $v^{\ast}$ (et alors $f$ sous la forme
$v_{\ast}$) que s'il commute de plus aux limites projectives finies.
Cela sera évidemment le cas s'il est lui-même l'adjoint à droite
d'un quatrième foncteur $d$. Sans condition de cette nature, on
écrira souvent
$$
e=u_{!}
$$
pour l'adjoint à gauche d'un foncteur image inverse $f=u^{\ast}$ ,
quand un tel adjoint à gauche existe, cette notation étant
suggérée par l'exemple d'une immersion ouverte $u:X\rightarrow
Y$ d'espaces topologiques. De même, si $e$ est de la forme
$v^{\ast}$, \ie $f$ de la forme $v_{\ast}$, on note parfois
$g=v^{!}$ l'adjoint à droite d'un foncteur image directe
$v_{\ast}$, lorsque cet adjoint existe \footnote[1]{\Cf plus bas \hyperref[IV.subsec14.4]{14.4},
dans le cas de faisceaux abéliens.}.

\subsection{}\etiquette[3.2]{IV.subsec3.2}
Soient\pageoriginale $E$, $E'$ deux $𝒰$-topos, et
$$
u=(u_{\ast},u^{\ast},\varphi)\quav
v=(v_{\ast},v^{\ast},\psi):E\rightrightarrows E'
$$
deux morphismes de topos de $E$ dans $E'$. On appelle
\emph{morphisme de $u$ dans $v$} tout morphisme de $u_{\ast}$ dans
$v_{\ast}$ (au sens de la catégorie $\Hhom(E,E')$ des foncteurs de
$E$ dans $E'$). Les morphismes de morphismes de topos se composent
de façon évidente, et on définit de cette façon une catégorie,
qui est en fait une $\sheaf{\mathcal{U}}$-\emph{catégorie}
(\hyperrefext[I][I.prop7.8]{7.8}) notée
$$
\Hhomtop(E,E'){\quoi},
$$
appelée \emph{catégorie des morphismes} (ou \emph{des
applications continues}) \emph{de $E$ dans $E'$}. On définit alors
de façon évidente un foncteur
$$
u\longmapsto u_{\ast}:\Hhomtop(E,E')\longrightarrow
\Hhom(E,E'){\quoi},
$$
foncteur qui est pleinement fidèle par définition des flèches
dans le premier membre (mais qui n'est pas injectif sur les objets).

\subsubsection{}
On fera attention que, si $u$ et $v$ sont donnés comme dessus, la
théorie des foncteurs adjoints fournit une bijection canonique
$$
\Hom(u_{\ast},v_{\ast})\xrightarrow{\sim}\Hom(v^{\ast},
u^{\ast}){\quoi};
$$
en particulier, on obtient un \emph{contrafoncteur} sur
$\Hhomtop(E,E')$:
$$
u\longmapsto u^{\ast}:\Hhomtop(E,E')^{\circ}\longrightarrow
\Hhom(E,E'){\quoi}.
$$
Conformément à l'intuition géométrique, suivant laquelle le
foncteur image directe $u_{\ast}$ « va dans le même sens »
que l'application continue qui lui donne naissance, il y a donc lieu
de définir le sens des flèches pour les morphismes entre
morphismes de topos \emph{en termes des foncteurs images directes,}
et non en termes des foncteurs images inverses (bien que ce soient
ces derniers qui, nous l'avons vu, possèdent les propriétés
d'exactitude caractéristiques de la notion de morphisme de topos).

\subsubsection{}\etiquette[3.2.2]{IV.subsubsec3.2.2}
Ayant\pageoriginale  défini la catégorie $\Hhomtop(E,E')$ des
morphismes du topos $E$ dans le topos $E'$, la notion
d'\emph{isomorphie} entre deux morphismes $u$, $v$ de $E$ dans $E'$
est également définie. En pratique, il n'y a pas lieu de
distinguer essentiellement entre deux morphismes de topos
isomorphes, tout au moins lorsqu'on dispose d'un isomorphisme
\emph{canonique} entre les deux (tout comme il n'y a pas lieu
souvent de distinguer entre deux éléments d'une catégorie,
lorsqu'on se donne un isomorphisme canonique entre eux). Signalons
à ce propos que le plus souvent, lorsqu'on traite de diagrammes de
morphismes de topos et de questions de commutativité de tels
diagrammes (notion qui a un sens grâce à (3.3)), il ne s'agit
que de commutativité à isomorphisme (« canonique ») près;
par abus de langage, on traite alors ces diagrammes comme des
diagrammes effectivement commutatifs.

On peut songer à justifier cet abus de langage en introduisant
l'ensemble \\ $\Homtop(E,E')/\Isom$ des morphismes de $E$ dans $E'$ à
isomorphisme près, et en appelant morphisme une telle classe
d'isomorphie, au lieu de suivre la définition \hyperref[IV.subsec3.2]{3.2}.
Mais ceci se heurte aux inconvénients très graves qui se présentent,
chaque fois qu'on essaie d'identifier deux objets isomorphes d'une
catégorie, sans disposer d'un isomorphisme canonique entre eux.
L'expérience prouve qu'un tel point de vue est impraticable, et
qu'il faut garder la notion « fine » \hyperref[IV.def3.1]{3.1} de la
notion de morphisme entre morphismes de topos, quitte à être obligé,
parfois, de se battre avec des compatibilités entre isomorphismes
canoniques \footnote[1]{Pour des exemples de telles batailles
(victorieuses, semble-t-il) nous renvoyons le lecteur au livre de
Mme M. \textsc{Hakim} sur les schémas relatifs [9].}.

\setcounter{subsection}{3} \setcounter{subsubsection}{0}
\subsubsection{}\etiquette[3.3.1]{IV.subsubsec3.3.1}
Soient\pageoriginale $E$, $E'$, $E''$ trois $𝒰$-topos, et
considérons des morphismes de topos
$$
u:E\longrightarrow E'\quav v:E'\longrightarrow E''{\quoi}.
$$
La théorie des foncteurs adjoints nous donne alors un isomorphisme
d'adjonction entre les foncteurs composés $v_{\ast}u_{\ast}$ et
$u^{\ast}v^{\ast}$, en termes des isomorphismes d'adjonction pour
les couples $(u_{\ast}, u^{\ast})$ et $(v_{\ast}, v^{\ast})$.
D'autre part, le foncteur $u^{\ast}v^{\ast}$ est exact à gauche, comme
composé de deux foncteurs exacts à gauche. Par suite, on trouve
un morphisme de $E$ dans $E''$, qu'on appelle \emph{composé des
morphismes $u$ et} $v$, et qu'on note $vu$:
$$
vu:E\longrightarrow E''{\quoi}.
$$
On vérifie alors trivialement que la composition des morphismes
est associative, et que pour tout $𝒰$-topos $E$, il y a un
morphisme de $E$ dans lui-même qui est une unité bilatère pour
la composition: c'est le morphisme $(\id_E, \id_E, \varphi)$, où
$\varphi$ est l'isomorphisme d'adjonction évident de $\id_E$ avec
lui-même. Soit alors $𝒱$ un univers tel que $𝒰\in
𝒱$. On définit une catégorie
$$
(𝒱\text{-}𝒰\text{-}\Top){\quoi},
$$
dont les objets sont les $𝒰$-topos qui sont $\in 𝒱$,
les flèches sont les morphismes entre de tels $𝒰$-topos,
et la composition des flèches étant celle qu'on vient
d'expliciter.

\subsubsection{}\etiquette[3.3.2]{IV.subsubsec3.3.2}
En fait, l'application de composition
$$
\Homtop(E, E')\times \Homtop(E',E'')\longrightarrow \Homtop(E,E'')
$$
est l'application induite sur les objets par un « foncteur
composition des morphismes »:
$$
\Hhomtop(E,E')\times \Hhomtop(E',E'')\longrightarrow
\Hhomtop(E,E''){\quoi},
$$
dont\pageoriginale l'effet sur les flèches est l'opération « produit de convolution »  habituel pour des morphismes entre
foncteurs (ici des foncteurs image directe). Ces foncteurs
composition satisfont à une propriété d'associativité
(stricte), pour quatre topos $E$, $E'$, $E''$, $E'''$, précisant
l'associativité de la composition des morphismes de topos. On peut
dire aussi, dans un langage qui commence à devenir familier [2],
[9], que les $𝒰$-topos sont les objets (ou 0-flèches)
d'une 2-\emph{catégorie}, dont les 1-flèches sont les morphismes
de topos, et les 2-flèches sont les morphismes de morphismes de
topos.

C'est le fait que les $𝒰$-topos (éléments d'un univers
$𝒱$) forment une 2-catégorie, et non plus seulement une
catégorie ordinaire comme les espaces topologiques ordinaires, qui
constitue du point de vue technique la différence la plus
importante entre la théorie des topos et celle des espaces
topologiques. Ce fait est la source de certaines complications
techniques auxquelles on a déjà fait allusion, mais aussi de
faits essentiellement nouveaux par rapport à la topologie
traditionnelle.

\subsection{}\etiquette[3.4]{IV.subsec3.4}
Le fait que les $𝒰$-topos (éléments d'un univers
$𝒱$) forment une 2-catégorie (\hyperref[IV.subsubsec3.3.2]{3.3.2})
permet en particulier de définir la notion d'\emph{équivalence de
deux $𝒰$-topos $E$, $E'$}: on dira que $E$ et $E'$ sont
\emph{équivalents} s'il existe des morphismes de topos
$u:E\rightarrow E'$ et $v:E'\rightarrow E$, tels que les composés
$uv$ et vu soient isomorphes respectivement au morphisme identique
de $E$ et de $E'$; on dit alors que les morphismes $u$ et $v$ sont
des\pageoriginale \emph{équivalences quasi-inverses l'une de
l'autre.}

On constate aussitôt que pour que le morphisme de topos
$u:E\rightarrow E'$ soit une équivalence, il faut et il suffit que
$u_{\ast}$ soit une équivalence, ou ce qui revient au même, que
$u^{\ast}$ soit une équivalence. (Utiliser le fait qu'un foncteur
$f:E\rightarrow E'$ entre deux topos qui est une équivalence est
à la fois de la forme $u_{\ast}$ et de la forme $v^{\ast}$, avec
$u$ et $v$ des morphismes de topos); et pour que $E$ et $E'$ soient
équivalents au sens de l'alinéa précédent, il faut et il
suffit qu'ils soient équivalents en tant que catégories (\ie
comme objets de la 2-catégorie $(𝒱\text{-}\cat)$). Comme
de juste, cela montre que les notions d'équivalence introduites ne
dépendent pas du choix de l'univers $𝒱$ de
\hyperref[IV.subsubsec3.3.1]{3.3.1}.

\subsubsection{}\etiquette[3.4.1]{IV.subsubsec3.4.1}
Pratiquement, il n'y a pas lieu le plus souvent de distinguer
essentiellement entre $𝒰$-topos équivalents, tout comme il
n'y a pas lieu souvent de distinguer essentiellement entre deux
catégories équivalentes, --- à condition toutefois qu'on
dispose d'une équivalence explicite de l'un à l'autre, ou tout
au moins une équivalence définie à isomorphisme unique près.
C'est la notion d'équivalence de topos qui remplace ici la notion
traditionnelle d'homéomorphie entre deux espaces topologiques.
Voir l'exemple \hyperref[IV.subsec4.2]{4.2} plus bas pour les relations
précises entre ces deux notions.

\section{Exemples de morphismes de topos}\etiquette[4]{IV.sec4}

Nous reprenons ici les exemples de 2, en utilisant la notion de
morphisme de topos. Les commentaires de \hyperref[IV.subsec2.0]{2.0}
s'appliquent également au présent numéro. L'univers
$𝒰$ sera généralement sous-entendu.

\subsection{Le topos $\Top(X)$ pour un espace topologique $X$
variable}\etiquette[4.1]{IV.subsec4.1}\pageoriginale

\subsubsection{}\etiquette[4.1.1]{IV.subsubsec4.1.1}
Soit une application continue
$$
f:X\longrightarrow Y
$$
d'espaces topologiques, on va lui associer canoniquement un
morphisme de topos
$$
\Top(f)\mbox{ ou } f:\Top(X)\rightarrow \Top(Y){\quoi},
$$
avec les notations de \hyperref[IV.subsec2.1]{2.1}. Lorsqu'on définit
$\Top(X)$ comme $\Ouv(X)^{\widetilde{\phantom{,}}}$, la description
la plus commode de $\Top(f)$ est par le foncteur image directe de
faisceaux
$$
f_{\ast}:\Top(X)\longrightarrow \Top(Y){\quoi},
$$
défini par la formule
$$
f_{\ast}(F)=F\circ f^{-1}{\quoi},
$$
où
$$
f^{-1}:\Ouv(Y)\longrightarrow \Ouv(X)
$$
est le foncteur évident $U\rightarrow f^{-1}(U)$. On a déjà noté que
ce foncteur est continu et exact à gauche, donc définit bien un
foncteur $f_{\ast}$ ci-dessus, admettant un adjoint à droite
$f^{\ast}$ qui est exact à gauche (\hyperrefext[III][III.exem1.9.1]{1.9.1}). Bien
entendu, en toute rigueur, le morphisme de Topos $\Top(f)$ dépend du
choix de l'adjoint à droite $f^{\ast}$ de $f_{\ast}$, donc n'est
défini qu'à isomorphisme canonique près. On se dispensera par la
suite de signaler expressément des phénomènes de ce genre.

Lorsqu'on adopte le point de vue « espaces étalés »  pour
définir $\Top(X)$, c'est le foncteur image inverse
$$
f^{\ast}:\Top(Y)\longrightarrow \Top(X)
$$
qui\pageoriginale est le plus commode pour définir le morphisme de
Topos $\Top(f)$, en posant simplement
$$
f^{\ast}(Y')=\fprod{X}{Y'}{Y}
$$
pour tout espace étalé $Y'$ sur $Y$; il est évident que le
produit fibré est bien un espace étale sur $X$, et que le
foncteur $f^{\ast}$ ainsi obtenu est exact à gauche et commute aux
$\varinjlim$ quelconques, et définit par suite un morphisme de
topos $\Top(f)$. Pour la compatibilité des deux définitions
obtenues, nous renvoyons à [TF].

Lorsqu'on a deux applications continues composables
$$
X\xrightarrow{f}Y\xrightarrow{g}Z{\quoi},
$$
on trouve un isomorphisme canonique
$$
\Top(gf)\simeq \Top(g)\Top(f)
$$
de morphismes de topos. Ces \emph{isomorphismes de transitivité},
pour trois applications continues composables $f$, $g$, $h$
satisfont à une relation de compatibilité que nous nous
dispensons d'écrire ici, et qui n'est autre que celle envisagée
dans \SGA 1 VI 7.4 B) (pour $\epsilon = (\Esp)^{\circ}$). On peut
l'exprimer en disant que pour $X$ variable dans la catégorie
$(\Esp)$,
$$
X\longmapsto \Top(X)
$$
est un « pseudo-foncteur »
\begin{equation*}
(𝒰\text{-}\esp)\longrightarrow
(𝒱\text{-}𝒰\text{-}\top){\quoi},\tag{4.1.1.1}
\end{equation*}\etiquette[4.1.1.1]{IV.eq4.1.1.1}
ou aussi, dans la terminologie des 2-catégories, qu'on a un
foncteur \emph{non strict} de 2-catégories [9]. En pratique, on se
permettra le plus souvent l'abus de langage consistant à
identifier $\Top(gf)$ et $\Top(g)\Top(f)$, c'est-à-dire de
raisonner comme si \eqhyperref[IV.eq4.1.1.1]{4.1.1.1} était un vrai foncteur de
catégories\pageoriginale ordinaires. On se permettra les abus de
langage analogues dans les autres exemples qui seront traités
ci-dessous.

\subsubsection{}\etiquette[4.1.2]{IV.subsubsec4.1.2}
Les considérations précédentes s'étendent immédiatement au
cas des topos associés aux espaces topologiques à groupes
d'opérateurs (\hyperref[IV.subsec2.3]{2.3}). Si $f=(f^{\es},f^{\gr})$,
$$
f:(X,G)\longrightarrow (Y,H)
$$
est un morphisme d'espaces à opérateurs (où
$$
f^{\es}:X\longrightarrow Y\quav f^{\gr}:G\longrightarrow H
$$
sont respectivement des applications continues et des morphismes de
groupes, compatibles dans un sens évident), on lui associe un
morphisme de topos
$$
\Top(f)\mbox{ ou } f:\Top(X,G)\longrightarrow \Top(Y,H){\quoi},
$$
dont la définition est laissée au lecteur. Lorsque les groupes
$G$ et $H$ sont les groupes unité, on retrouve la définition de
\hyperref[IV.subsubsec4.1.1]{4.1.1}; lorsque par contre ce sont les espaces $X$
et $Y$ qui sont réduits à un point, on trouve comme foncteur
image inverse $f^{\ast}$ le foncteur « restriction du groupe
d'opérateurs »
$$
f^{\ast}:B_H\longrightarrow B_G{\quoi},
$$
associant à tout $H$-ensemble le $G$-ensemble qu'il définit
grâce à $f:G\rightarrow H$. On retrouvera cet exemple sous
d'autres formes encore dans \hyperref[IV.subsec4.5]{4.5} et
\hyperref[IV.subsubsec4.6.1]{4.6.1}.

\subsubsection{}\etiquette[4.1.3]{IV.subsubsec4.1.3}
Étant donné une application continue $f:X\rightarrow Y$ d'espaces
topologiques, on lui associe également un morphisme sur les « gros
topos »  correspondants (\hyperref[IV.subsec2.5]{2.5})
$$
\TOP(f)\mbox{ ou } f:\TOP(X)\longrightarrow \TOP(Y){\quoi},
$$
défini\pageoriginale le plus commodément par le foncteur image
inverse
$$
f^{\ast}:\TOP(Y)\longrightarrow \TOP(X){\quoi},
$$
qui n'est autre que le \emph{foncteur restriction}. Ce morphisme
$\TOP(f)$ est un cas particulier du morphisme dit
« d'inclusion » pour un topos induit, qui sera étudié dans 5. On voit ainsi
que, sous les mêmes réserves que dans \hyperref[IV.subsubsec4.1.1]{4.1.1},
le topos $\TOP(X)$ peut être considéré comme un foncteur en
$X$, pour $X$ variable dans $(\Esp)$.

\subsection{Propriétés de fidélité de $X\mapsto \Top(X)$}\etiquette[4.2]{IV.subsec4.2}

Nous nous proposons de préciser dans quelle mesure un espace
topologique $X$ peut se reconstituer en termes du topos $\Top(X)$,
et dans ce but il convient de décrire, pour deux espaces $X$ et $Y$,
la catégorie des morphismes de $\Top(X)$ dans $\Top(Y)$
(\hyperref[IV.subsec3.2]{3.2}), de façon à pouvoir préciser les propriétés
de fidélité du « foncteur »  $X\mapsto \Top(X)$. Nous nous
bornons à énoncer les résultats auxquels on parvient, en renvoyant
le lecteur à [9] pour des détails. Le lecteur qui voudra faire
l'exercice de vérification lui-même pourra consulter l'exer.
\hyperref[IV.exer7.8]{7.8}.

\subsubsection{}\etiquette[4.2.1]{IV.subsubsec4.2.1}
Rappelons qu'un espace topologique $X$ est dit \emph{sobre} si toute
partie fermée irréductible de $X$ a exactement un point
générique. Signalons que presque tous les espaces utilisés en
pratique sont sobres; il en est en particulier ainsi d'un espace
séparé, plus généralement d'un espace dont tous les points
sont fermés, ou de l'espace sous-jacent à un schéma. Si $X$
est un espace topologique, on lui associe (\loccit ou \EGA
$0_I$, réédition) un espace topologique sobre $X_{\sob}$
et une application continue\pageoriginale
\begin{equation*}
\varphi: X\longrightarrow X_{\sob}\tag{4.2.1.1}
\end{equation*}\etiquette[4.2.1.1]{IV.eq4.2.1.1}
qui soit \emph{universelle} pour les applications continues de $X$
dans des espaces sobres; en d'autres termes, on construit un
foncteur adjoint à gauche $X\mapsto X_{\sob}$ du foncteur
d'inclusion $(\Espsob)\rightarrow (\Esp)$ de la catégorie des
espaces sobres dans celle des espaces topologiques
« quelconques » (les guillemets rappelant qu'il y a un univers !). La
construction explicite se fait en prenant comme points de $X_{\sob}$
les parties fermées irréductibles de $X$, comme ouverts les
ensembles de la forme $U'$, où $U$ est un ouvert de $X$ et où
$U'\subset X_{\sob}$ désigne l'ensemble des parties fermées
irréductibles de $X$ qui rencontrent $U$. L'application
\eqhyperref[IV.eq4.2.1.1]{4.2.1.1} est obtenue en associant à tout $x\in X$
l'adhérence de $\{x\}$. L'espace $X$ est sobre si et seulement si
l'application précédente est bijective, donc un
homéomorphisme.

On constate que le foncteur
$$
\varphi^{-1}:\Ouv(X_{\sob})\longrightarrow \Ouv(X)
$$
induit par $\varphi$ est un isomorphisme, ce qui implique que le
morphisme de topos
$$
\Top(\varphi):\Top(X)\longrightarrow \Top(X_{\sob})
$$
défini par $\varphi$ est également un isomorphisme. Ceci explique à
priori pourquoi $X_{\sob}$ doit s'introduire nécessairement dans la
question de reconstituer $X$ à partir de $\Top(X)$: comme ce dernier
ne dépend que de $X_{\sob}$ à isomorphisme près, la question ne
pourra avoir une réponse affirmative que si $X$ est sobre. Nous
préciserons plus bas (\hyperref[IV.subsec7.1]{7.1}) comment $X_{\sob}$ peut
effectivement se\pageoriginale reconstituer en termes de $\Top(X)$,
en interprétant ses points comme des « points »  du topos
$\Top(X)$ (ou encore comme des foncteurs fibres).

\subsubsection{}\etiquette[4.2.2]{IV.subsubsec4.2.2}
Sur tout espace topologique, il y a lieu d'introduire la relation
d'ordre $\leq$ pour laquelle on a
$$
x\leq y\Longleftrightarrow \bar{\{x\}}\subset \bar{\{y\}}\mbox{ \ie
} x\in \bar{\{y\}}
$$
(qu'on exprime encore en disant que $x$ est une
\emph{spécialisation} de $y$, ou que $y$ est une
\emph{générisation} de $x$). Pour un espace de la forme
$X_{\sob}$, ce n'est autre que la relation d'inclusion entre parties
fermées irréductibles de $X$.

Ceci posé, il y a lieu d'introduire, sur l'ensemble des
applications d'un espace $X$ dans un autre $Y$, la relation d'ordre
dite de « spécialisation », déduite de celle de $Y$, savoir
$$
f\leq g\Longleftrightarrow f(x)\leq g(x)\mbox{ pour tout } x\in
X{\quoi}.
$$
Avec ces conventions, on a le résultat suivant:

\subsubsection{}\etiquette[4.2.3]{IV.subsubsec4.2.3}
\emph{Soient $X$, $Y$ deux espaces topologiques, avec $Y$ sobre, et
soient $f$ et $g$ deux applications continues de $X$ dans $Y$. Alors
il y a au plus un morphisme de $\Top(f)$ dans $\Top(g)$, et pour
qu'il y en ait un, il faut et il suffit que $g$ soit une
spécialisation de $f$. Enfin, tout morphisme de topos
$\Top(X)\rightarrow \Top(Y)$ est isomorphe à un morphisme de la
forme $\Top(f)$, où $f:X\rightarrow Y$ est une application
continue (uniquement déterminée grâce à la première assertion).}

Si\pageoriginale on définit la catégorie $\cat(I)$ associé à
un ensemble ordonné I en déclarant que pour $i\geq j$, il y a
exactement une flèche de $i$ dans $j$, on peut résumer le
résultat précédent en disant qu'on a une équivalence de
catégories canonique
$$
\cat(\Hom_{(\esp)}(X,Y))\xrightarrow{\approx}\Hhomtop(\Top(X),\Top(Y))
$$
(où le deuxième membre est défini dans \hyperref[IV.subsec3.2]{3.2}).

On conclut formellement de ces résultats:

\begin{enonce*}{Corollaire 4.2.4}\etiquette[4.2.4]{IV.cor4.2.4}
\begin{enumerate}
\renewcommand{\theenumi}{\alph{enumi}}
\renewcommand{\labelenumi}{\theenumi)}
\item Soit $f:X\rightarrow Y$ une application continue. Alors
$\Top(f):\Top(X)\rightarrow \Top(Y)$ est une équivalence de topos
si et seulement si $f_{\sob}:X_{\sob}\rightarrow Y_{\sob}$ est un
homéomorphisme (donc, lorsque $X$ et $Y$ sont sobres, si et
seulement si $f$ est un homéomorphisme).

\item Soient $X$ et $Y$ des espaces topologiques. Pour que $\Top(X)$
et $\Top(Y)$ soient équivalents, il faut et il suffit que
$X_{\sob}$ et $Y_{\sob}$ soient homéomorphes (donc, si $X$ et $Y$
sont sobres, il faut et suffit que $X$ et $Y$ soient
homéomorphes).
\end{enumerate}
\end{enonce*}

\subsection{Morphismes dans le topos final: objets constants d'un topos,
foncteurs sections}\etiquette[4.3]{IV.subsec4.3}

Désignons par $P$ (initiale de « point ») le topos final type,
\ie $P=(\Ens)$ (\hyperref[IV.subsec2.2]{2.2}). Soit $E$ un topos quelconque,
on va voir qu'\emph{à isomorphisme unique près, il existe un
unique morphisme de topos}
$$
f:E\rightarrow P
$$
plus précisément, que \emph{la catégorie $\Hhomtop(E,P)$ est
équivalente à la catégorie ponctuelle}: pour deux objets de
cette catégorie, il existe donc une\pageoriginale unique flèche
de l'un dans l'autre, et c'est un isomorphisme. (Cela justifie dans
une certaine mesure l'appellation « topos final »). Pour ceci,
rappelons (3.2.1) que $\Hhomtop(E,P)$ est équivalente à la
sous-catégorie pleine de $\Hhom(P,E)^{\circ}$ formée des
foncteurs
$$
f^{\ast}:P=(\Ens)\longrightarrow E
$$
qui commutent aux $\varinjlim$ et sont exacts à gauche. Soit $e$ un
ensemble ponctuel, alors tout ensemble $X$ s'écrit canoniquement
comme « somme de $X$ copies de $e$ », d'où résulte que la
catégorie des foncteurs $g:(\Ens)\rightarrow E$ qui commutent aux
$\varinjlim$ est équivalente à la catégorie $E$, en associant à tout
$g$ l'objet $g(e)$ de $E$. On reconstitue $g$ en termes de $T=g(e)$,
à isomorphisme unique près, par $g(I)=T\times I$, où on pose
$T\times I=\amalg_{i\in I}T_i$, avec $T_i=T$ pour tout $i\in I$. Que
le foncteur en $I$ ainsi défini par $T$ commute bien aux
$\varinjlim$ résulte du fait qu'il est manifestement adjoint à
\sisi{droite}{gauche} du foncteur $X\mapsto \Hom(T,X)$ de $E$ dans $(\Ens)$. Ceci
dit, pour que $g$ soit exact à gauche, il est évidemment nécessaire
que $g(e)=T$ soit un objet final de $E$ (puisque $e$ est un objet
final de $(\Ens)$), et il résulte facilement du fait que dans $E$
« les sommes sont universelles »  (\hyperref[IV.subsubsec1.1.2]{1.1.2} b))
que cette condition est aussi suffisante. On trouve donc que la
catégorie des foncteurs $f^{\ast}$ images inverses est équivalente à
la catégorie des objets finaux de $E$, qui est évidemment elle-même
équivalente à la catégorie finale.

D'après ce qui précède, on voit que le choix d'un morphisme
(4.3.1) équivaut essentiellement à celui d'un objet final de
$E$, soit $e_E$. En termes de celui-ci, on a alors des isomorphismes
canoniques de foncteurs
\begin{equation*}
 f^{\ast}(I)\simeq e_E\times
I=\mbox{ somme de } I \mbox{ copies de } e_E \quad \mbox{ pour } I\in \ob
(\Ens),\tag{4.3.2}
\end{equation*}\etiquette[4.3.2]{IV.eq4.3.2}
et\pageoriginale
\begin{equation*}
f_{\ast}(X)=\Hom(e_E,X) \quad \quad \mbox{ pour } X\in \Ob E.\tag{4.3.3}
\end{equation*}\etiquette[4.3.3]{IV.eq4.3.3}


\setcounter{subsubsection}{3}
\subsubsection{}\etiquette[4.3.4]{IV.subsubsec4.3.4}
Les deux foncteurs précédents joueront par la suite un rôle
important. Pour un ensemble $I$, on appelle \emph{objet constant de
valeur} $I$ dans $E$ (ou, lorsque $E$ est réalisé comme une
catégorie $C^{\widetilde{\phantom{,}}}$ en termes d'un site $C$,
\emph{faisceau constant de valeur $I$ sur $C$}), l'objet
$f^{\ast}(I)=e_E\times I$ de \eqhyperref[IV.eq4.3.2]{4.3.2}. On le notera aussi
souvent $I_E$, où $I_C$ lorsque $E$ est défini par le site $C$.
Le fait que $I\mapsto I_E$ soit le foncteur image inverse d'un
morphisme de topos en précise les propriétés d'exactitude, qui
impliquent en particulier que ce foncteur respecte toutes les
espèces de structure algébriques habituelles, transformant un
groupe en un objet groupe de $E$ etc (\hyperref[IV.subsubsec3.1.2]{3.1.2}).
Lorsqu'on a par exemple un Groupe $G$ de $E$, on dira que c'est un
\emph{Groupe constant} (ou, le cas échéant, un \emph{faisceau en
groupes constant}) s'il est isomorphe à un groupe de la forme
$G_{\circ E}$, où $G_{\circ}$ est un groupe ordinaire. Même
terminologie pour toute autre espèce de structure
« algébrique », au sens précisé (plus ou moins) dans
\hyperref[IV.subsubsec3.1.2]{3.1.2}.

\subsubsection{}\etiquette[4.3.5]{IV.subsubsec4.3.5}
On fera attention que le foncteur $I\mapsto I_E$ n'est pas
nécessairement pleinement fidèle (ni même fidèle: prendre
pour $E$ le « topos vide »  (\hyperref[IV.subsec2.2]{2.2})), donc un objet
constant de $E$ ne détermine pas en général à isomorphisme
unique près l'ensemble $I$ qui lui donne naissance. On dit que $E$
est $0$-\emph{acyclique}, ou \emph{connexe-non vide}, si le foncteur
$I\mapsto I_E$ est pleinement fidèle. Il revient au même,
d'après les propriétés générales des foncteurs adjoints, de
dire que le morphisme d'adjonction
$$
I\longrightarrow f_{\ast}(f^{\ast}(I))=f_{\ast}(I_E)=\Hom(e_E, I_E)
$$
est\pageoriginale un isomorphisme (de sorte que le foncteur
\eqhyperref[IV.eq4.3.3]{4.3.3} permet de récupérer la « valeur »  d'un objet
constant de E). On vérifie facilement qu'il faut et il suffit pour
cela que $e_E$ ne soit pas l'objet initial $\phi_E$ de $E$, \ie que
$E$ ne soit pas un « topos vide » (ce qui exprime la
\emph{fidélité} du foncteur $I\mapsto I_E$ \footnote[1]{ou encore le
fait que ce foncteur est conservatif (\hyperrefext[I][I.cor6.3]{6.3}).}), \emph{et}
que $e_E$ soit un objet \emph{connexe} de $E$, \ie ne soit pas somme
de deux objets de $E$ qui ne soient pas « vides » (\ie qui ne
soient pas des objets initiaux de $E$).

\subsubsection{}\etiquette[4.3.6]{IV.subsubsec4.3.6}
Le foncteur \eqhyperref[IV.eq4.3.3]{4.3.3} est aussi souvent appelé \emph{foncteur
sections} et noté $\Gamma_E$ ou $\Gamma(E,-)$ ou simplement
$\Gamma$:
\begin{equation*}
\Hom(e_E,X)=\Gamma_E(X)=\Gamma(E,X)=\Gamma(X){\quoi}.\tag{4.3.6.1}
\end{equation*}
C'est un foncteur commutant aux limites projectives quelconques,
mais pas exact à droite en général, dont les foncteurs
dérivés (sur les objets groupes abéliens) seront étudiés
dans le prochain exposé.

\subsection{Morphismes du « topos vide »}\etiquette[4.4]{IV.subsec4.4}

Soit $\phi_{\text{top}}$ un topos vide, qui correspond donc à une
catégorie de faisceaux $\Phi$ équivalente à la catégorie finale
(\hyperref[IV.subsec2.2]{2.2}). Soit $E$ un topos. La catégorie des foncteurs
de $E$ dans $\Phi$ est évidemment équivalente à la catégorie
ponctuelle, et tout tel foncteur commute aux limites inductives et
projectives (sans aucun mérite d'ailleurs), donc peut être considéré
comme un foncteur image inverse $f$ pour un morphisme de topos
$\phi_{\text{top}}\rightarrow E$. Il en résulte que la
\emph{catégorie $\Hhomtop(\phi_{\text{top}},E)$ est équivalente à la
catégorie ponctuelle,} et en particulier qu'\emph{il existe à
isomorphisme\pageoriginale unique près un et un seul morphisme de
topos}
\begin{equation*}
\phi_{\text{top}}\longrightarrow E{\quoi}.\tag{4.4.1}
\end{equation*}\etiquette[4.4.1]{IV.eq4.4.1}
Ceci justifie dans une certaine mesure la terminologie « topos
initial »  introduite dans \hyperref[IV.subsec2.2]{2.2}.

On peut aussi déterminer les morphismes de topos
\begin{equation*}
E\longrightarrow \phi_{\text{top}}{\quoi}\ ;\tag{4.4.2}
\end{equation*}\etiquette[4.4.2]{IV.eq4.4.2}
on vérifie aussitôt qu'il existe un tel morphisme si et
seulement si l'objet initial de $E$ est aussi un objet final, \ie si
et seulement si $E$ lui-même est un « topos vide », et que
dans ce cas la catégorie $\Hhomtop(E,\phi_{\text{top}})$ est
encore équivalente à la catégorie ponctuelle. L'unique
morphisme \eqhyperref[IV.eq4.4.2]{4.4.2} (modulo isomorphie) est alors une
équivalence de topos.

\subsection{Le topos classifiant $B_G$ pour $G$ Groupe variable}\etiquette[4.5]{IV.subsec4.5}

\subsubsection{}\etiquette[4.5.1]{IV.subsubsec4.5.1}
Soient $E$ un topos, et
$$
f:G\longrightarrow H
$$
un morphisme de Groupes dans $E$. On en déduit un foncteur « restriction du Groupe d'opérateurs »
$$
f^{\ast}:B_H\longrightarrow B_G{\quoi},
$$
où les notations sont celles de \hyperref[IV.subsec2.4]{2.4}. Il est trivial
que ce foncteur commute aux limites inductives et aux limites
projectives, a fortiori il peut être interprété comme un
foncteur image inverse associé à un morphisme de topos
$$
B_f\mbox{ ou } f:B_G\longrightarrow B_H{\quoi}.
$$
On\pageoriginale explicite aisément le foncteur image directe
correspondant
$$
f_{\ast}:B_G\longrightarrow B_H
$$
par la formule
$$
f_{\ast}(X)=\Hhom_G(H_s,X){\quoi},
$$
où $X$ est un objet de \sisi{$X$}{$E$} avec $G$ opérant à gauche, où
$H_s$ est $H$ regardé comme muni des opérations à gauche par
$G$ déduites de \sisi{$F$}{$f$}, et où $\Hhom_G$ désigne le sous-objet
qu'on devine de l'objet $\Hhom$ défini plus bas (\hyperref[IV.subsec10.2]{10.2}); on fait
opérer $H$ à gauche sur $\Hhom_G(H_s,X)$ grâce aux
opérations droites de $H$ sur $H_s$ par translations à
\emph{droite}.

Comme le foncteur image inverse $f^{\ast}$ commute aux \sisi{$\varinjlim$}{$\varprojlim$}
quelconques (et non seulement aux \sisi{$\varinjlim$}{$\varprojlim$} finies), il est
lui-même l'adjoint à droite d'un foncteur
$$
f_{!}:B_G\longrightarrow B_H{\quoi},
$$
de sorte qu'on a une suite de trois foncteurs adjoints comme dans
\hyperref[IV.subsubsec3.1.3]{3.1.3}:
$$
f_{!},f^{\ast}, f_{\ast}{\quoi}.
$$
On explicite aisément $f_{!}$ par la formule
$$
f_{!}(X)=H\times {}^GX{\quoi},
$$
où le deuxième membre désigne le « produit contracté »,
déduit des opérations de $G$ sur $X$ (à gauche) et sur $H$
(à droite via translations à droite et $f$), défini comme le
quotient de $H\times X$ par $G$ opérant par la formule
$$
g\circ (h,x)=(hg^{-1}, gx){\quoi}.
$$

Le foncteur $f_{!}$, étant un adjoint à gauche, commute
évidemment aux limites inductives, mais il n'est pas en
général exact à gauche (\ie il\pageoriginale ne peut être
considéré à son tour comme un foncteur image inverse par un
morphisme de topos $B_H\rightarrow B_G$). En fait, on vérifie
facilement qu'il ne peut être exact à gauche que si
$f:G\rightarrow H$ est un isomorphisme. De même, le foncteur
$f_{\ast}$, qui commute aux limites projectives, n'est pas en
général exact à droite, et a fortiori n'admet pas en
général d'adjoint à droite. Tout au moins lorsque $E$ est le
topos ponctuel \ie que $G$ et $H$ sont des groupes ordinaires,
$f_{\ast}$ n'est exact à droite que si $f$ est un isomorphisme.

Lorsqu'on a un deuxième morphisme de groupes $g:H\rightarrow K$,
on trouve comme dans \hyperref[IV.subsubsec4.1.1]{4.1.1} une transitivité à
isomorphisme canonique près, de sorte que, sous la réserve
habituelle, on peut considérer que le topos classifiant $B_G$
dépend \emph{fonctoriellement} du Groupe $G$. On laisse au lecteur
le soin de généraliser ce comportement fonctoriel pour le cas où
on fait varier simultanément $G$ \emph{et} le topos $E$.

\subsubsection{Le topos $B_{𝒢}$ pour un pro-groupe variable
  $𝒢$}\etiquette[4.5.2]{IV.subsubsec4.5.2}

On laisse au lecteur le soin de préciser le caractère covariant
du topos $B_{𝒢}$ (\hyperref[IV.subsec2.7]{2.7}) par rapport à
$𝒢$, en calquant l'exposé que nous en donnons dans
\hyperref[IV.subsubsec4.5.1]{4.5.1}. On fera attention cependant que dans le cas
d'un morphisme $f:𝒢\rightarrow ℋ$ de pro-groupes qui
ne sont pas essentiellement constants, le morphisme de topos
correspondant $B_{𝒢}\rightarrow B_{ℋ}$ ne permet pas
en général la définition d'un foncteur $f_{!}$ (dont l'adjoint
à droite soit le foncteur $f^{\ast}$ de restriction du pro-groupe
d'opérateurs). Supposant, pour simplifier l'énoncé, que
$𝒢$ et $ℋ$ soient définis par des groupes
\emph{profinis} $G$ et $H$, on voit facilement que $f_{!}$ existe si
et seulement si l'image du morphisme envisagé $f:G\rightarrow H$
est d'indice fini dans $H$, et  dans\pageoriginale ce cas $f_{!}$
est donné par la même formule que dans \hyperref[IV.subsec4.5]{4.5}.

\subsection{Le topos $\widehat{C}$ pour $C$ catégorie variable}\etiquette[4.6]{IV.subsec4.6}

\subsubsection{}\etiquette[4.6.1]{IV.subsubsec4.6.1}
Soit
$$
f:C\longrightarrow C'
$$
un foncteur d'une catégorie $C\in 𝒰$ dans une autre $C'$,
d'où un foncteur
$$
f^{\ast}:\widehat{C}'\longrightarrow \widehat{C}\quav
f^{\ast}(F)=F\circ f{\quoi}.
$$
Il est trivial que ce foncteur commute aux limites projectives et
aux limites inductives, a fortiori il peut être considéré
comme le foncteur image inverse pour un morphisme de topos
$$
\widehat{f}\mbox{ ou } f:\widehat{C}\rightarrow \widehat{C}'{\quoi}.
$$
Le foncteur image directe correspondant
$$
f_{\ast}:\widehat{C}\longrightarrow \widehat{C}'
$$
n'est autre que le foncteur également noté $f_{\ast}$ dans I 5.1. De
plus (comme il était à prévoir du fait que $f^{\ast}$ commute
également aux $\varprojlim$ quelconques) $f^{\ast}$ admet aussi un
adjoint à gauche
$$
f_{!}:C\longrightarrow C'{\quoi},
$$
(qui était noté aussi $f_{!}$ dans \hyperrefext[I][I.prop5.1]{5.1}). On obtient donc
une suite de trois foncteurs adjoints
$$
f_{!}, f^{\ast}, f_{\ast}{\quoi},
$$
le premier étant d'ailleurs un prolongement de $f:C\rightarrow C'$
aux catégories de préfaisceaux (pour le plongement habituel de $C$,
$C'$ dans les catégories $\widehat{C}$, $\widehat{C}'$). On fera
attention que le foncteur $f_{!}$ n'est pas en général exact à
gauche, ni $f_{\ast}$ exact à droite, ce qui lève donc toute
ambiguïté sur la direction\pageoriginale de la variance du topos
$\widehat{C}$ associé à la catégorie variable $C$: ce topos est un
foncteur \emph{covariant} en $C$, sous la réserve habituelle
provenant des isomorphismes de transitivité (\cf
\hyperref[IV.subsubsec4.1.1]{4.1.1}). Lorsque l'ensemble des objets de $C$ et de
$C'$ est réduit à un élément, de sorte que $C$ et $C'$ s'identifient
à des monoïdes $G$, $G'$, les topos correspondants sont les topos
classifiants $B_{G^{\circ}}$ et $B_{{G'}^{\circ}}$, et on retrouve
la variance de ceux-ci pour $G$ variable rencontrée déjà à d'autres
points de vue dans \hyperref[IV.subsubsec4.1.2]{4.1.2} et \hyperref[IV.subsec4.5]{4.5}
(où la restriction à des groupes au lieu de monoïdes n'avait rien
d'essentiel).

\subsubsection{}\etiquette[4.6.2]{IV.subsubsec4.6.2}
On peut préciser la dépendance 2-fonctorielle du topos
$\widehat{C}$ par rapport à $C$, en introduisant pour deux
catégories $C$, $C'\in 𝒰$ un foncteur canonique
\begin{equation*}
\Hhom(C,C')\longrightarrow
\Hhomtop(\widehat{C},\widehat{C}'){\quoi}.\tag{4.6.2.1}
\end{equation*}\etiquette[4.6.2.1]{IV.eq4.6.2.1}
Il reste à définir ce foncteur sur les flèches, et pour ceci
on note que $f\mapsto f^{\ast}$ permet d'identifier (à
équivalence de catégories près) le deuxième membre à une
sous-catégorie pleine de $\Hhom(\widehat{C}',
\widehat{C})^{\circ}$ (3.2.1). Or si $f$, $g:C\rightarrow C'$ sont
deux foncteurs, tout morphisme $u:f\rightarrow g$ de foncteurs
définit un morphisme $f\circ g\rightarrow F\circ f$ de foncteurs
en $F\in \ob \widehat{C}'$ ($F$ étant un \emph{contra}foncteur),
qui est donc un morphisme $g^{\ast}\rightarrow f^{\ast}$ et
définit par suite un morphisme $\widehat{f}\rightarrow
\widehat{g}$ comme annoncé.

Lorsque $C$ est la catégorie ponctuelle, $\widehat{C}$ est le
topos ponctuel (\hyperref[IV.subsec2.2]{2.2}) noté $P$, et \eqhyperref[IV.eq4.6.2.1]{4.6.2.1}
s'interprète comme un foncteur naturel
\begin{equation*}
C'\longrightarrow \Hhomtop(P,\widehat{C}') \mathop{=}^{\dfn}
\Ppoints(\widehat{C}')\tag{4.6.2.2}
\end{equation*}\etiquette[4.6.2.2]{IV.eq4.6.2.2}
de $C'$ dans la « \emph{catégorie des points} » de
$\widehat{C}'$, qui sera étudiée dans 6. Ce\pageoriginale
foncteur n'est pas nécessairement une équivalence de
catégories (7.3.3), a fortiori \eqhyperref[IV.eq4.6.2.1]{4.6.2.1} n'est pas
nécessairement une équivalence de catégories.

Par contre, le foncteur \eqhyperref[IV.eq4.6.2.1]{4.6.2.1} est toujours
\emph{pleinement fidèle.} Pour voir ceci, notons que si $f$,
$g:E\rightarrow E'$ sont deux morphismes de topos tels que $f_{!}$
et $g_{!}$ soient définis (\hyperref[IV.subsubsec3.1.3]{3.1.3}), il résulte
de la théorie des foncteurs adjoints qu'on a des isomorphismes
canoniques
$$
\Hom(f_{!},g_{!})\simeq \Hom(g^{\ast}, f^{\ast})\simeq
\Hom(f_{\ast}, g_{\ast}) \mathop{=}^{\dfn} \Hom(f,g){\quoi}.
$$
Appliquant ceci à des foncteurs de la forme $\widehat{f}$,
$\widehat{g}$ associés à $f$, $g:C\rightarrow C'$, on trouve le
résultat annoncé, compte tenu que l'application naturelle de
prolongement $\Hom(f,g)\rightarrow \Hom(f_{!},g_{!})$ est bijective
(I 7.8).

\subsubsection{}\etiquette[4.6.3]{IV.subsubsec4.6.3}
On peut se demander quand le foncteur \eqhyperref[IV.eq4.6.2.1]{4.6.2.1} est une
équivalence de catégories, \ie quand il est essentiellement
surjectif, ce qui est un cas particulier de la question de
déterminer tous les morphismes de topos $\widehat{C}\rightarrow
\widehat{C}'$. Plus généralement, si $C$ est une catégorie
$\in 𝒰$ et $E$ un topos, on peut se proposer de déterminer
les morphismes de topos
$$
f:\widehat{C}\longrightarrow E{\quoi}.
$$
La catégorie de ces morphismes est équivalente à la
catégorie opposée de celle des foncteurs $f^{\ast}:E\rightarrow
\widehat{C}$ commutant aux limites inductives quelconques et aux
limites projectives finies. Interprétant les foncteurs
$E\rightarrow \widehat{C}$ comme des foncteurs $E\times
C^{\circ}\rightarrow (𝒰\text{-}\Ens)=(\Ens)$, ou encore
comme des foncteurs $F:C^{\circ}\rightarrow \Hhom(E, (\Ens))$, la
propriété d'exactitude envisagée s'exprime par la condition
que pour tout objet $X$ de $C$, le foncteur $F(X):E\rightarrow
(\Ens)$ commute aux limites inductives quelconques et aux limites
projectives finies, (ou, comme nous dirons dans 6, $F(X)$ est un « \emph{foncteur\pageoriginale fibre} »  sur $E$). Il revient au
même de dire que $F(X)$ est le foncteur image inverse pour un
morphisme de topos $P\rightarrow E$, de sorte qu'on trouve
finalement une équivalence de catégories canonique
\begin{equation*}
\Hhomtop(\widehat{C}, E) \xrightarrow{\approx} \Hhom(C,
\Ppoints(E)){\quoi},\tag{4.6.3.1}
\end{equation*}\etiquette[4.6.3.1]{IV.eq4.6.3.1}
où on désigne pour abréger par
$$
\Ppoints(E)=\Hhomtop(P,E)
$$
la catégorie des points du topos $E$.

Lorsque $E$ est de la forme $\widehat{C}'$, on voit aussitôt à
partir des définitions que le composé de \eqhyperref[IV.eq4.6.2.1]{4.6.2.1} et de
l'équivalence précédente \eqhyperref[IV.eq4.6.3.1]{4.6.3.1} est le foncteur
\begin{equation*}
\Hhom(C,C')\longrightarrow
\Hhom(C,\Ppoints(\widehat{C}'))\tag{4.6.3.2}
\end{equation*}\etiquette[4.6.3.2]{IV.eq4.6.3.2}
défini par $F\mapsto i\circ F$, où $i:C'\rightarrow
\Ppoints(\widehat{C}')$ est le plongement canonique \eqhyperref[IV.eq4.6.2.2]{4.6.2.2}.
Il s'ensuit aussitôt que \emph{pour que} \eqhyperref[IV.eq4.6.2.1]{4.6.2.1}
         \emph{soit une équivalence, il faut et il suffit que $C$
         soit vide ou que} \eqhyperref[IV.eq4.6.2.2]{4.6.2.2} \emph{soit essentiellement
         surjectif} (donc une équivalence de catégories). Cette
         dernière condition sur $C'$ est satisfaite dans certains
         cas intéressants, et notamment lorsque $C'$ est la
         catégorie à un seul objet définie par un groupe $G$
         (\hyperref[IV.subsubsec7.2.5]{7.2.5}).


\begin{enonce*}[remark]{Remarque 4.6.4}\etiquette[4.6.4]{IV.rem4.6.4}\nde{Pour
un point de vue récent de la notion de type d'homotopie de topos, voir
[15] $\S$ 7.1.6.}
Le fait d'avoir associé un topos $\widehat{C}$ à une catégorie
arbitraire $C$ suggère qu'une catégorie $C$ admet des invariants de
nature topologique (groupes de cohomologie, d'homotopie \etc), tout
comme un topos. Les groupes de cohomologie de $\widehat{C}$ à
coefficients dans un objet groupe abélien $F$ (au sens général
étudié dans V) ne sont autres que les valeurs des foncteurs
dérivés $\varprojlim^{(n)}$ du foncteur $\varprojlim$, déjà
familiers\pageoriginale aux topologues. J. L. Verdier et
(indépendamment) D. G. Quillen ont vérifié que lorsqu'on se borne
aux coefficients constants, ou plus généralement localement
constants, ces groupes de cohomologie s'identifient aux groupes de
cohomologie de l'ensemble semi-simplicial $\Nerf(C)$ canoniquement
associé à $C$ [5, ${\rm n^{\circ}}$ 212, prop. 4.1] et que de plus,
à isomorphie près dans la « catégorie homotopique » de [2], tout
ensemble semi-simplicial peut être obtenu à l'aide d'une catégorie
$C$. Moyennant une notion convenable de type d'homotopie pour des
topos, que nous ne précisons pas ici, on peut dire que les types
d'homotopie semi-simpliciaux des topologues ne sont autres 
que
les types d'homotopie des topos de la forme spéciale $\widehat{C}$,
plus généralement des topos $E$ où tout objet admette un
recouvrement qui raffine tous les autres (\hyperref[IV.subsec2.6]{2.6})
\footnote[1]{et, plus généralement encore, des topos qui sont « localement $\infty$ - connexes »  en un sens évident que nous
laissons au lecteur le soin de préciser.}. En l'absence de cette
condition sur $E$, on peut tout au mieux exprimer son type
d'homotopie par un \emph{système projectif} convenable d'ensembles
semi-simpliciaux [1].
\end{enonce*}

\subsection{Le topos $C^{\sim}$ pour un site $C$
variable (foncteurs cocontinus)}\etiquette[4.7]{IV.subsec4.7}

Soit
$$
f:C\longrightarrow C'
$$
un foncteur cocontinu (\hyperrefext[III][III.def2.1]{2.1}) entre sites $\in
𝒰$, \ie un foncteur tel que le foncteur $\widehat{f}_{\ast}$
de $\widehat{C}$ dans $\widehat{C}'$
(\hyperref[IV.subsubsec4.6.1]{4.6.1}) applique $C^{\sim}$ dans
$C'^{\sim}$, c'est-à-dire induise un foncteur
\begin{equation*}
\widetilde{f}_{\ast}:\widetilde{C}\longrightarrow
{C'}^{\widetilde{\phantom{,}}}\tag{4.7.1}
\end{equation*}\etiquette[4.7.1]{IV.eq4.7.1}
rendant commutatif le diagramme de foncteurs
\[\vcenter{
\xymatrix@C=2.5cm@R=1.25cm{ C^{\widetilde{\phantom{a}}}
\ar[r]^{\widetilde{f}_{\ast}} \ar[d]_{\sisi{i}{i_{\ast}}}&
C'^{\widetilde{\phantom{a}}} \ar[d]^{\sisi{i'}{i_{\ast}'}}\\
\widehat{C}  \ar[r]^{\widehat{f}_{\ast}}       & \widehat{C}'\quad,
}}\leqno{(4.7.2)}
\]\etiquette[4.7.2]{IV.eq4.7.2}
où\pageoriginale \sisi{$i$}{$i_{\ast}$}, \sisi{$i'$}{$i_{\ast}'$} sont les foncteurs d'inclusion. On a vu
alors (\hyperrefext[III][III.prop2.3]{2.3}) que le foncteur $\widetilde{f}_{\ast}$
admet un adjoint à gauche $\widetilde{f}^{\ast}$, et que ce
dernier est exact à gauche. En d'autres termes,
$\widetilde{f}_{\ast}$ est le foncteur image directe associé à
un morphisme de topos
$$
\widetilde{f}\mbox{ ou } f:C^{\sim}\longrightarrow
{C'}^{\widetilde{\phantom{,}}}{\quoi},
$$
le foncteur image inverse correspondant étant bien entendu
$\widetilde{f}^{\ast}$. Prenant les adjoints à gauche des
foncteurs en jeu, le diagramme commutatif \eqhyperref[IV.eq4.7.2]{4.7.2} donne
d'ailleurs un diagramme commutatif à isomorphisme près
\[\vcenter{
\xymatrix@C=2.5cm@R=1.5cm{
C^{\widetilde{\phantom{a}}} &  C'^{\widetilde{\phantom{a}}} \ar[l]_{\widetilde{f}^{\ast}}\\
\widehat{C}  \ar[u]^{\underline{a}}  & \widehat{C}'
\ar[l]_{\widehat{f}^{\ast}}\ar[u]_{\underline{a}'}\quad,
}}\leqno{(4.7.3)}
\]\etiquette[4.7.3]{IV.eq4.7.3}
où $\underline{a}$ et $\underline{a}'$ sont les foncteurs « faisceaux
associés », diagramme qui redonne aussitôt la formule
(\hyperrefext[III][III.prop2.3]{2.3})
$$
\widetilde{f}^{\ast}=\underline{a}\widehat{f}^{\ast} i'{\quoi}.
$$

La propriété de transitivité pour les morphismes de topos
$\widehat{f}:\widehat{C}\rightarrow \widehat{C}'$ implique la même
propriété pour les morphismes de topos
$\widetilde{f}:C^{\widetilde{\phantom{,}}}\rightarrow
{C'}^{\widetilde{\phantom{,}}}$ associés à des foncteurs
cocontinus, de sorte qu'on peut dire que \emph{le topos
$C^{\widetilde{\phantom{,}}}$ varie fonctoriellement en $C$} de fa\c
con \emph{covariante}, quand on prend comme « morphismes »  de
sites les foncteurs cocontinus.

Dans\pageoriginale tous les cas rencontrés jusqu'à présent, le
foncteur cocontinu $f$ utilisé est également continu,
c'est-à-dire (\hyperrefext[III][III.def1.1]{1.1}) se « prolonge »  en un
foncteur
$$
\widetilde{f}_{!}:C^{\widetilde{\phantom{,}}}\longrightarrow
{C'}^{\widetilde{\phantom{,}}}
$$
commutant aux limites inductives, et qui est adjoint à gauche de
$\widetilde{f}^{\ast}$, de sorte qu'on a une suite de trois
foncteurs adjoints
$$
\widetilde{f}_{!}, \widetilde{f}^{\ast},
\widetilde{f}_{\ast}{\quoi}.
$$
On fera attention que le foncteur $\widetilde{f}_{!}$ n'est pas en
général exact à gauche, ni $\widetilde{f}_{\ast}$ exact à droite, ce
qui lève toute ambiguïté sur la direction de la variance du topos
$C^{\widetilde{\phantom{,}}}$, pour $C$ variable par des foncteurs
dont on suppose seulement qu'ils sont cocontinus, ou même continus
et cocontinus.

\begin{enonce*}[remark]{Remarque 4.7.4}\etiquette[4.7.4]{IV.rem4.7.4}
Étant donné un morphisme de topos
$$
F:E\longrightarrow E'{\quoi},
$$
pour qu'il existe un foncteur $F_{!}$ adjoint à gauche de
$F^{\ast}$, il faut et il suffit qu'on puisse « réaliser »
(à équivalence près) $E$ et $E'$ sous la forme
$C^{\widetilde{\phantom{,}}}$ et ${C'}^{\widetilde{\phantom{,}}}$,
pour deux sites $C$, $C'\in 𝒰$, et qu'on puisse trouver un
foncteur \emph{continu et cocontinu} $f:C\rightarrow C'$ tel que $F$
s'identifie à $\widetilde{f}$. C'est évidemment suffisant, et
pour la nécessité, il suffit de prendre pour $C$ et $C'$ des
petites sous-catégories pleines génératrices de $E$ et $E'$
respectivement, telles que
$$
F_{!}(\ob C)\subset \ob C'{\quoi},
$$
munies des topologies induites par celles de $E$ et de $E'$, et de
prendre pour $f$ le foncteur induit par $F_{!}$ (\cf
\hyperref[IV.cor1.2.1]{1.2.1}). Rappelons (\hyperref[IV.rem1.8]{1.8}) que l'existence de
$F_{!}$ signifie aussi que $F^{\ast}$ commute aux limites
projectives, ou,\pageoriginale ce qui revient ici au même puisque
ce foncteur est exact à gauche, qu'il commute aux produits.
\end{enonce*}

\setcounter{subsubsection}{4}
\subsubsection{}\etiquette[4.7.5]{IV.subsubsec4.7.5}
Si on se demande quels sont les morphismes de topos $F:E\rightarrow
E'$ qui peuvent se réaliser par un foncteur cocontinu (pas
nécessairement continu) de sites, on voit de même que la
réponse est la suivante: l'ensemble des objets $X$ de $E$ tels que
le foncteur $X'\mapsto \Hom(X,F^{\ast}(X'))$ de $E'$ dans $(\Ens)$
commute aux limites projectives (ou, ce qui revient au même, aux
produits) doit être une \emph{famille génératrice de $E$.}\nde{Cet énoncé ne paraît pas
correct: la condition exprimée est équivalente au fait que $F^{\ast}$ commute aux produits quelconques, 
ou encore que $F^{\ast}$ possède un adjoint à gauche $F_{!}$. Pour tout site $C$, le foncteur
identité $F=id : C \longrightarrow C$ est cocontinu lorsque l'on munit le membre de droite
de la topologie triviale. Le foncteur $F^{\ast} : \hat{C} \longrightarrow \widetilde{C}$ 
est alors le foncteur faisceau associé, et ne commute pas en général avec les produits quelconques.}

\subsection{Le morphisme de topos $\widetilde{C}\rightarrow
\widehat{C}$ pour un site $C$}\etiquette[4.8]{IV.subsec4.8}

Soit $C$ un petit site, auquel sont donc associés les deux topos
$C^{\widetilde{\phantom{,}}}$ et $\widehat{C}$ (le deuxième ne
dépendant pas de la topologie mise sur $C$). On a défini dans II
3.4 le foncteur « faisceau associé »
$$
\underline{a}:\widehat{C}\longrightarrow
C^{\widetilde{\phantom{,}}}{\quoi},
$$
et établi qu'il est exact à gauche et commute aux limites
inductives. C'est donc le foncteur image inverse associé à un
morphisme de topos
$$
p:\widetilde{C}\longrightarrow \widehat{C},\;
p^{\ast}=\underline{a}{\quoi},
$$
le foncteur image directe correspondant étant l'inclusion
canonique
$$
i=p_{\ast}:C^{\widetilde{\phantom{,}}}\longrightarrow
\widehat{C}{\quoi}.
$$
Il est bien connu que ce dernier foncteur n'est pas en général exact
à droite (ses foncteurs dérivés sur les objets abéliens donnent
naissance aux préfaisceaux de cohomologie $\Xx^n(F)$ de V 2), et que
$\underline{a}$ ne commute pas en général aux $\varprojlim$ quelconques,
ce qui lève toute ambiguïté sur la direction du morphisme « naturel »  de topos entre $\widehat{C}$ et
$C^{\widetilde{\phantom{,}}}$.

Lorsqu'on a un foncteur cocontinu
$$
f:C\longrightarrow C'
$$
de\pageoriginale sites, on en déduit un diagramme de morphismes de
topos
\[
\xymatrix@C=2.25cm@R=1.25cm{ C^{\widetilde{\phantom{a}}} \ar[d]_{\sisi{p}{p_{\ast}}}
\ar[r]^{\widetilde{f}} &
C'^{\widetilde{\phantom{a}}} \ar[d]^{\sisi{p'}{p'_{\ast}}}\\
\widehat{C} \ar[r]^{\widehat{f}}&  \widehat{C}' }
\]
qui est commutatif à isomorphisme canonique près: c'est en effet
ce qu'exprime la commutativité du diagramme de foncteurs
\eqhyperref[IV.eq4.7.2]{4.7.2}. On peut donc dire que le morphisme de topos
$p:C^{\widetilde{\phantom{,}}}\rightarrow \widehat{C}$ est
fonctoriel en $C$, quand on varie $C$ par des foncteurs
\emph{cocontinus} entre sites.

\subsection{Effet d'un foncteur continu de sites. Morphismes de sites}\etiquette[4.9]{IV.subsec4.9}

\subsubsection{}\etiquette[4.9.1]{IV.subsubsec4.9.1}
Si
$$
f:C\longrightarrow C'
$$
est un foncteur continu de $C$ dans $C'$. \ie tel que le foncteur
$\widehat{f}^{\ast}:\widehat{C}'\rightarrow \widehat{C}$ applique
${C'}^{\widetilde{\phantom{,}}}$ dans $C^{\widetilde{\phantom{,}}}$,
donc induise un foncteur
$$
f_s:{C'}^{\widetilde{\phantom{,}}}\longrightarrow
C^{\widetilde{\phantom{,}}}{\quoi},
$$
on a vu (\hyperrefext[III][III.prop1.2]{1.2}) que ce dernier admet un adjoint à
gauche
$$
f^s:C^{\widetilde{\phantom{,}}}\longrightarrow
{C'}^{\widetilde{\phantom{,}}}{\quoi},
$$
qui « prolonge »  d'ailleurs $f$ dans un sens évident. On fera
attention qu'en général $f^s$ n'est pas exact à \sisi{droite}{gauche} (même
si $f$ est de plus cocontinu), ni $f_s$ ne commute aux limites
inductives, de sorte que la donnée de $f$ ne permet pas, sans
autre hypothèse, de décrire un morphisme de topos dans un sens
ou dans l'autre entre $C^{\widetilde{\phantom{,}}}$ et
${C'}^{\widetilde{\phantom{,}}}$. Le cas où $f$ est cocontinu, \ie
où $f_s$ commute aux limites inductives et peut\pageoriginale donc
être regardé comme un foncteur image inverse pour un morphisme
de topos $\widetilde{f}:C^{\widetilde{\phantom{,}}}\rightarrow
{C'}^{\widetilde{\phantom{,}}}$, a été examiné dans
\hyperref[IV.subsec4.7]{4.7}. Nous allons examiner le cas où le foncteur
$f^s$ est exact à gauche, donc peut être considéré comme un
foncteur image inverse pour un morphisme de topos \emph{en sens
inverse}:
\begin{equation*}
\Top(f)=g:{C'}^{\sim}\longrightarrow C^{\sim}{\quoi}. \tag{4.9.1.1}
\end{equation*}\etiquette[4.9.1.1]{IV.eq4.9.1.1}
On fera attention qu'\emph{on a pris garde de ne pas noter ce
morphisme par la lettre $\widetilde{f}$ ou $f$, pour éviter des
confusions avec la situation de $\hyperref[IV.subsec4.7]{4.7}$, suivant en
cela les recommandations générales de
$\hyperref[IV.subsubsec3.1.3]{3.1.3}$.} On dit parfois que le foncteur
$f:C\rightarrow C'$ est un \emph{morphisme de sites de $C'$ dans
$C$} (attention, \emph{pas de $C$ dans $C'$}) s'il est continu et si
le foncteur $f^s$ est exact à gauche, en d'autres termes s'il
existe un morphisme de topos \eqhyperref[IV.eq4.9.1.1]{4.9.1.1} tel que, le foncteur
image inverse correspondant
$$
g^{\ast}:C^{\widetilde{\phantom{,}}}\longrightarrow
{C'}^{\widetilde{\phantom{,}}}
$$
\emph{prolonge} le foncteur $f$, \ie rende commutatif à
isomorphisme près le diagramme de foncteurs
\[
\xymatrix@C=2cm@R=1.5cm{
C \ar[d]_{\epsilon} \ar[r]^{f}& C' \ar[d]^{{\epsilon}'}\\
C^{\widetilde{\phantom{a}}} \ar[r]^{g^{\ast}} &
C'^{\widetilde{\phantom{a}}} &, }
\]
%% widetilde + \phantom : bidouille horrible \XXX
où $\epsilon$, $\epsilon'$ sont les foncteurs canoniques de II
4.4.0.

\subsubsection{}\etiquette[4.9.2]{IV.subsubsec4.9.2}
Pratiquement, on reconnaît qu'un foncteur continu $f:C\rightarrow
C'$ est un morphisme de sites de $C'$ dans $C$, par le fait que dans
$C$ les limites projectives finies sont représentables, et que $f$ y
commute (\hyperrefext[III][III.prop1.3]{1.3} 5)). Moyennant\pageoriginale la condition
indiquée sur $C$ (presque toujours vérifiée en pratique), et
supposant de plus que la topologie de $C'$ est moins fine que la
topologie canonique (presque toujours vérifiée également), la
condition suffisante précédente ($f$ exact à gauche) pour que $f$
soit un morphisme de sites de $C'$ dans $C$ est d'ailleurs aussi
nécessaire.

\subsubsection{}\etiquette[4.9.3]{IV.subsubsec4.9.3}
Dans l'esprit de ce qui précède, si $C$ et $C'$ sont deux
$𝒰$-sites, il y a lieu de définir la catégorie des
morphismes de sites $\Mmorsite(C',C)$ de $C$ dans $C'$ comme la
sous-catégorie pleine de la catégorie \emph{opposée
$\Hhom(C,C')^{\circ}$} de la catégorie des foncteurs de $C$ dans
$C'$, formée des foncteurs qui veulent bien être des morphismes
de sites (de $C'$ dans $C$). De cette façon, on obtient un
foncteur canonique (défini à isomorphisme unique près)
$$
\Mmorsite(C',C)\longrightarrow
\Hhomtop({C'}^{\widetilde{\phantom{,}}},
C^{\widetilde{\phantom{,}}}){\quoi}.
$$

\begin{enonce*}{Proposition 4.9.4}\etiquette[4.9.4]{IV.prop4.9.4}
Soient $E$ un $𝒰$-topos, $C$ un $𝒰$-site, alors le
foncteur $f\mapsto f^{\ast}| C=f^{\ast}\circ \epsilon_{C}$,
associant à tout morphisme de topos $f:E\rightarrow
C^{\widetilde{\phantom{,}}}$ la « restriction »  à $C$ du
foncteur image inverse associé
$f^{\ast}:C^{\widetilde{\phantom{,}}}\rightarrow E$, induit une
équivalence de catégories
$$
\Hhomtop(E,C^{\widetilde{\phantom{,}}})\xrightarrow{\approx}\Mmorsite(E,C)(\hookrightarrow
\Hhom(C,E)^{\circ}){\quoi}.
$$
Lorsque dans $C$ les $\varprojlim$ finies sont représentables, le
foncteur pleinement fidèle correspondant
$$
\Hhomtop(E,C^{\widetilde{\phantom{,}}})\longrightarrow
\Hhom(C,E)^{\circ}
$$
a comme image essentielle l'ensemble des foncteurs $g:C\rightarrow
E$ qui sont exacts à gauche et continus, ou encore qui sont exacts
à gauche et transforment\pageoriginale famille couvrante en
famille couvrante.
\end{enonce*}

La dernière assertion résulte de la première grâce à
\hyperref[IV.subsubsec4.9.2]{4.9.2}. D'autre part, on déduit de
\hyperrefext[III][III.prop1.2]{1.2} iv) et de IV \hyperref[IV.theo1.2]{1.2} iii) que le foncteur
$G\rightarrow G\circ \epsilon$ induit une équivalence de la
catégorie des foncteurs \emph{continus} de
$C^{\widetilde{\phantom{,}}}$ dans $E$, et la catégorie des
foncteurs \emph{continus} de $C$ dans $E$, un foncteur quasi-inverse
étant obtenu par $g\rightarrow g^s$ dans les notations de \loccit
D'autre part, par définition même, $g$ est un morphisme de sites
si et seulement si $g^s$ est le foncteur image inverse associé à
un morphisme de topos $E\rightarrow C^{\widetilde{\phantom{,}}}$,
d'où la conclusion grâce à 3.2.1, le « ou encore »
provenant de \hyperrefext[III][III.prop1.6]{1.6}.

On retiendra surtout de \hyperref[IV.prop4.9.4]{4.9.4} que (lorsque dans $C$ les
$\varprojlim$ finies sont représentables) « il revient au même » de se donner un morphisme de topos $f:E\rightarrow
C^{\widetilde{\phantom{,}}}$, ou un foncteur $g:C\rightarrow E$ qui
est exact à gauche et transforme familles couvrantes en familles
couvrantes.

\setcounter{subsubsection}{4}
\subsubsection{}\etiquette[4.9.5]{IV.subsubsec4.9.5}
Utilisant les développements \hyperref[IV.subsubsec4.9.1]{4.9.1} et
\hyperref[IV.subsubsec4.9.3]{4.9.3}, on voit comme d'habitude que pour un
$𝒰$-site $C$ variable, via la notion de morphisme de sites
et de morphisme de morphismes de sites qu'on vient d'expliciter, le
topos $C^{\widetilde{\phantom{,}}}$ dépend fonctoriellement (ou
plus exactement, 2-fonctoriellement) du site $C$. On notera que
grâce à la terminologie introduite,
$C^{\widetilde{\phantom{,}}}$ dépend de façon \emph{covariante}
du site $C$.

\subsubsection{}\etiquette[4.9.6]{IV.subsubsec4.9.6}
Il est immédiat que (contrairement à ce qui se passe pour la notion
de foncteur cocontinu, \cf \hyperref[IV.subsec4.7]{4.7} 4) \emph{tout}
morphisme de topos $F:E\rightarrow E'$ peut se réaliser à l'aide
d'un morphisme de sites $f:C\rightarrow C'$ (\ie d'un foncteur
continu $f:C'\rightarrow C$ tel que\ldots): il suffit de choisir
dans $E$ et $E'$ des petites sous-catégories pleines génératrices
$C$\pageoriginale et $C'$ respectivement, munies des topologies
induites par celles de $E$ et de $E'$, telles que l'on ait
$$
F^{\ast}(\ob C')\subset \ob C{\quoi},
$$
et de prendre pour $f$ le foncteur induit par $F^{\ast}$. C'est ce
qui explique que la plupart des morphismes de topos qu'on rencontre
en pratique sont effectivement décrits à l'aide de morphismes de
sites (plutôt qu'à l'aide de foncteurs cocontinus comme en
\hyperref[IV.subsec4.7]{4.7}).

\subsection{Relations entre le petit et le gros topos associés à
un espace topologique $X$}\etiquette[4.10]{IV.subsec4.10}

On reprend les notations de \hyperref[IV.subsec2.5]{2.5}. En particulier,
$𝒱$ est un univers tel que $𝒰\in 𝒱$. Nous
nous écartons de la convention de \hyperref[IV.subsec2.1]{2.1}, en
désignant par $\Top(X)$ le topos des $𝒱$-faisceaux (et non
pas des $𝒰$-faisceaux) sur $X$. Ainsi, nous raisonnerons
avec des $𝒱$-topos et non des $𝒰$-topos. (NB il
serait possible de garder les $𝒰$-topos, en adoptant la
convention appropriée pour la définition $\TOP(X)$, de sorte que
celui-ci soit un $𝒰$-topos, \cf \hyperref[IV.subsec2.5]{2.5}). Ceci
posé, nous allons définir DEUX morphismes de topos
\begin{equation*}
\begin{cases}
 f: \Top(X)\longrightarrow \TOP(X)\\
 g: \TOP(X)\longrightarrow \Top(X)\quav gf\simeq \id_{\Top(X)}{\quoi},
\end{cases}\tag{4.10.1}
\end{equation*}\etiquette[4.10.1]{IV.eq4.10.1}
donnant lieu à une suite de trois foncteurs adjoints
\begin{equation*}
g^{\ast}=f_{!}\quav g_{\ast}=f^{\ast}\quav
g^{!}=f_{\ast}{\quoi},\tag{4.10.1.1}
\end{equation*}\etiquette[4.10.1.1]{IV.eq4.10.1.1}
les notations étant celles de \hyperref[IV.subsubsec3.1.3]{3.1.3}. Nous allons
définir successivement les\pageoriginale deux premiers foncteurs
de cette suite, le troisième étant défini alors comme adjoint
à droite du deuxième.

\setcounter{subsubsection}{1}
\subsubsection{}\etiquette[4.10.2]{IV.subsubsec4.10.2}
Le foncteur
$$
g_{\ast}=f^{\ast}:\TOP(X)\longrightarrow \Top(X)
$$
est défini simplement comme le \emph{foncteur restriction} de
$(\Esp)_{/X}$ au site $\Ouv(X)$ des ouverts de $X$, qui transforme
bien faisceaux en faisceaux, comme il résulte immédiatement de
la définition. Désignant par des lettres soulignées des
faisceaux sur le gros site de $X$, on désigne, pour un tel
faisceau $\underline{F}$, par $\underline{F}_X$ sa restriction au site
$\Ouv(X)$, d'où un foncteur
\begin{equation*}
\Restr:\underline{F}\longmapsto \underline{F}_X:\TOP(X)\longrightarrow
\Top(X){\quoi}.\tag{4.10.2.1}
\end{equation*}\etiquette[4.10.2.1]{IV.eq4.10.2.1}
Il est évident que ce foncteur commute aux $𝒱$-limites
projectives, puis-que celles-ci se calculent argument par argument
(on pourrait aussi invoquer l'existence de l'adjoint à gauche,
construit dans \hyperref[IV.subsubsec4.10.4]{4.10.4} ci-dessous.) Je dis qu'il
commute également aux $𝒱$-limites inductives. Pour s'en
convaincre, on va donner une interprétation fort commode des « gros »  faisceaux sur $X$, \ie des objets de $\TOP(X)$, en termes
de faisceaux ordinaires (NB il s'agit de $𝒱$-faisceaux) sur
les espaces topologiques $X'$ au-dessus de~$X$.

\subsubsection{}\etiquette[4.10.3]{IV.subsubsec4.10.3}
Pour un gros faisceaux $\underline{F}$ sur $X$, et pour tout espace $X'$
sur $X$ (sous-entendu: $X'\in 𝒰$), on définit de façon
évidente, comme dans \hyperref[IV.subsubsec4.10.2]{4.10.2}, le « petit »
faisceau $\underline{F}_{X'}$, restriction de $\underline{F}$ à $X'$. Si
$$
u:X''\longrightarrow X'
$$
est\pageoriginale  un morphisme de $(\Esp)_{/X}$, on définit alors
de façon évidente un morphisme
$u_{\ast}(\underline{F}_{X''})\rightarrow \underline{F}_{X'}$, ou ce qui
revient au même, un « \emph{morphisme de transition} »
\begin{equation*}
\varphi_u :u^{\ast}(\underline{F}_{X'})\longrightarrow
\underline{F}_{X''}{\quoi}. \tag{4.10.3.1}
\end{equation*}\etiquette[4.10.3.1]{IV.eq4.10.3.1}
Ces morphismes, pour $u$ variable, satisfont à une condition de
transitivité évidente pour un composé
$$
X'''\xrightarrow{u}X''\xrightarrow{v}X'
$$
de morphismes dans $(\Esp)_{/X}$, qu'on laisse au lecteur le soin
d'expliciter. On obtient de cette façon un foncteur naturel, qui
va de la catégorie $\TOP(X)$ des gros faisceaux sur $X$, dans la
catégorie des systèmes
$$
(\underline{F}_{X'})\quad(X'\in \ob (\Esp)_{/X}), (\varphi_u)\quad(u\in
Fl(\Esp)_{/X}){\quoi},
$$
satisfaisant à la condition de transitivité envisagée, et tels de
plus que pour tout morphisme $u:X''\rightarrow X'$ qui est une
\emph{immersion ouverte} (ou, plus généralement, un étalement), le
morphisme de transition $\varphi_u$ soit un isomorphisme. Comme les
foncteurs $u^{\ast}:\Top(X')\rightarrow \Top(X'')$ utilisés pour la
description des $\varphi_u$ commutent aux $𝒱$-limites
inductives, il en résulte aussitôt que dans la description
précédente des objets de $\TOP(X)$ en termes de « petits »
faisceaux $F_{X'}$, les $𝒱$-limites inductives se calculent
argument par argument, \ie les foncteurs de la forme
$\underline{F}\rightarrow \underline{F}_{X'}$ commutent aux
$𝒱$-limites inductives.

En particulier, il en est ainsi du foncteur $\underline{F}\mapsto
\underline{F}_X$ envisagé dans \hyperref[IV.subsubsec4.10.2]{4.10.2}. Il admet donc
bien un adjoint à droite (\hyperref[IV.cor1.5]{1.5}). Il reste à en construire
un adjoint à gauche (dont l'existence résulte d'ailleurs a priori de
\hyperref[IV.rem1.8]{1.8}), et de vérifier que ce dernier est exact à gauche.
Cela\pageoriginale achèvera la définition des morphismes de topos
annoncés \eqhyperref[IV.eq4.10.1]{4.10.1} en termes des foncteurs \eqhyperref[IV.eq4.10.1.1]{4.10.1.1}

\subsubsection{}\etiquette[4.10.4]{IV.subsubsec4.10.4}
Le foncteur
$$
g^{\ast}=f_{!}:\Top(X)\longrightarrow \TOP(X)
$$
s'obtient en associant à tout petit faisceau $F$ sur $X$ le système
de ses images inverses $(F_{X'})$, $X'\in \ob (\Esp)_{/X}$, par les
morphismes structuraux $X'\rightarrow X$, le morphisme de transition
$\varphi_u$ pour un $u:X''\rightarrow X'$ étant l'isomorphisme de
transitivité pour les images inverses de faisceaux
(\hyperref[IV.subsubsec4.1.1]{4.1.1}). On voit aussitôt qu'on obtient ainsi un
foncteur
$$
\Prol:\Top(X)\longrightarrow \TOP(X){\quoi},
$$
\emph{pleinement fidèle}, dont l'image essentielle est formée
des gros faisceaux $ℱ$ sur $X$ pour lesquels tous les
morphismes de transitivité $\varphi_u$ de \sisi{4.11.3}{\eqhyperref[IV.eq4.10.3.1]{4.10.3.1}} sont des
isomorphismes. Nous laissons au lecteur le soin de définir un
isomorphisme d'adjonction entre ce foncteur de prolongement
canonique et le foncteur restriction de \hyperref[IV.subsubsec4.10.2]{4.10.2},
prouvant que ce dernier est adjoint à droite du premier. C'est
immédiat en termes de la description \hyperref[IV.subsubsec4.10.3]{4.10.3} de
la catégorie $\TOP(X)$.

\begin{enonce*}[remark]{Remarques 4.10.5}\etiquette[4.10.5]{IV.rem4.10.5}
\begin{enumerate}
\renewcommand{\theenumi}{\alph{enumi}}
\renewcommand{\labelenumi}{\theenumi)}
\item La construction \hyperref[IV.subsubsec4.10.4]{4.10.4} montre en même temps que
$g^{\ast}=f_{!}$ est pleinement fidèle (ou, ce qui revient au
même, que $g_{\ast}=f^{\ast}$ est un foncteur de passage à une
catégorie de fractions, ou enfin que $g^{!}=f_{\ast}$ est
pleinement fidèle). Les gros faisceaux sur $X$ qui appartiennent
à l'image essentielle du foncteur $g^{\ast}=f_{!}=\Prol$
méritent le nom de gros \emph{faisceaux étales} sur $X$,
puisqu'ils forment une catégorie équivalente à celle des
faisceaux ordinaires sur $X$, ou encore à celle des espaces
étalés sur\pageoriginale $X$. On voit d'ailleurs facilement que
si $X'$ est un espace topologique au-dessus de $X$, alors le gros
faisceau sur $X$ qu'il représente est étale au sens
précédent si et seulement si $X'$ est un espace étalé sur
$X$.

\item On peut encore exprimer la pleine fidélité de $f_{\ast}$ en
écrivant que le morphisme d'adjonction
$f^{\ast}f_{\ast}\rightarrow \id_{\Top(X)}$ est un isomorphisme, \ie
que $g_{\ast}f_{\ast}\simeq \id$, \ie que $gf\simeq \id$. Ainsi $g$
fait de $\TOP(X)$ un \emph{topos sur} $\Top(X)$, admettant une « section »  $f$ sur $\Top(X)$.

\item Le fait que le foncteur $g^{\ast}=f_{!}$ soit pleinement
fidèle, exact et qu'il commute aux $𝒱$-limites inductives
justifie partiellement le point de vue fort commode (dû à J.
\textsc{Giraud}) suivant lequel il est inoffensif, dans pratiquement
toutes les questions de théorie des faisceaux sur $X$, de
remplacer les faisceaux habituels ou « petits »  faisceaux par
les « gros » faisceaux associés. Il en est en particulier
ainsi des questions cohomologiques, puisque $g_{\ast}$ étant
exact, les foncteurs $R^i g_{\ast}$ (\sisi{V\quad}{\hyperrefext[V][V.sec5]{5}}) sont nuls pour $i>0$,
donc (\sisi{V\quad}{\hyperrefext[V][V.prop5.4]{5.4}}) que pour tout \emph{gros} faisceau $ℱ$ sur
$X$, on a des isomorphismes canoniques
$$
H^i(\TOP(X),ℱ)\simeq H^i(\Top(X),
ℱ_X)=H^i(X,F){\quoi}.
$$
Appliquant ceci à un faisceau de la forme $\Prol(F)$, où $F$ est
un petit faisceau sur $X$, on en conclut (puisque $\Prol(F)_X\simeq
F$ canoniquement) un isomorphisme canonique
$$
H^i(X,F)=H^i(\TOP(X), \Prol(F)){\quoi}.
$$
Ainsi, \emph{les invariants cohomologiques de $X$, calculés via le
petit ou le gros topos de $X$, sont essentiellement identiques.} Le
même résultat est d'ailleurs valable en cohomologie non
commutative.
\end{enumerate}
\end{enonce*}

\begin{enonce*}[remark]{Exercice 4.10.6}\etiquette[4.10.6]{IV.exer4.10.6}
Soient\pageoriginale $S$ un $𝒰$-site dont la topologie est
moins fine que la topologie canonique, $\underline{M}$ une partie de
$Fl$ $S$ satisfaisant les conditions suivantes:
\begin{enumerate}
\renewcommand{\theenumi}{\alph{enumi}}
\renewcommand{\labelenumi}{\theenumi)}
\item Les morphismes de $ℳ$ sont quarrables (\sisi{\hyperrefext[I][I.sec10.3]{10.3}}{\hyperrefext[I][I.sec10.7]{10.7}}), et
$\underline{M}$ est stable par changement de base.

\item $\underline{M}$ contient les flèches identiques et est stable par
composition.

\item Une flèche $u:X\rightarrow Y$ telle qu'il existe une famille
couvrante $Y_i\rightarrow Y$, avec les $\fprod{X}{Y_i}{Y}\rightarrow
Y_i$ dans $\underline{M}$, est elle-même élément de $\underline{M}$.

\item Pour tout $X\in \ob S$, toute famille couvrante de $X$ est
raffinée par une famille couvrante $(f_i:X_i\rightarrow X)$, avec
les $f_i\in \underline{M}$.
\end{enumerate}

Pour tout $X\in \ob S$, considérons le site $S(X)$ (« petit site
de $X$ ») dont la catégorie sous-jacente est la sous-catégorie
pleine de $S_{/X}$ formée des objets dont le morphisme structural
est dans $\underline{M}$, munie de la topologie induite
(\hyperrefext[III][III.subsec3.1]{3.1}) par celle de $S$.
\begin{itemize}
\item[$1^{\circ})$] Pour toute flèche $u:X\rightarrow Y$ de $S$,
montrer que le changement de base par $u$ de $S(Y)$ dans $S(X)$ est
un foncteur continu, d'où un foncteur commutant aux limites
inductives (\sisi{III}{\hyperrefext[III][III.prop1.2]{1.2} iv)})
$$
S(u)^x:S(Y)^{\sim}\longrightarrow S(X)^{\sim}{\quoi}.
$$

\item[$2^{\circ})$] Définir une équivalence entre le topos
$S^{\sim}$ et la catégorie des systèmes
$$
(F_X)\quad (X\in \ob S)\quav (\varphi_u)\quad (u\in Fl S)
$$
formés\pageoriginale d'objets $F_X\in \ob S(X)^{\sim}$, et pour
toute flèche $u:X\rightarrow Y$ dans $S$, d'un morphisme
$\varphi_u:S(u)^x(F_Y)\rightarrow F_X$, ces systèmes étant
soumis à une condition de transitivité pour un composé $v\circ
u$ de flèches de $S$, et à la condition que $u\in \underline{M}$ implique
que $\varphi_u$ soit un isomorphisme.

\item[$3^{\circ})$] Définir des foncteurs « restrictions »  et « prolongement »
$$
\Res_X:(S_{/X})^{\sim}\longrightarrow S(X)^{\sim},
\Prol_X:S(X)^{\sim}\longrightarrow (S_{/X})^{\sim}
$$

Montrer que $\Res_X$ commute aux petites limites inductives et
projectives et que $\Prol_X$ est \emph{pleinement fidèle}, son
image essentielle étant formée des faisceaux $F$ tels que
$\varphi_u$ soit un isomorphisme pour toute flèche $u$ de
$S_{/X}$.

\item[$4^{\circ})$] Définir un morphisme d'adjonction faisant de
$\Res_X$ l'adjoint à droite de $\Prol_X$. En conclure qu'il existe
un morphisme de topos
$$
f:S(X)^{\sim}\rightarrow (S_{/X})^{\sim}{\quoi}
$$
faisant de $S(X)^{\sim}$ un sous-topos de $(S_{/X})^{\sim}$ et tel
que
\begin{align*}
f_{!} &= \Prol_X\\
f^x &= \Res_X{\quoi}.
\end{align*}

Montrer que $\Prol_X$ transforme faisceaux abéliens en faisceaux
abéliens.

\item[$5^{\circ})$] Montrer que si $S(X)$ admet des produits
fibrés, $\Prol_X$ est exact. En déduire qu'il existe alors un
morphisme de topos $g:(S_{/X})^{\sim}\rightarrow S(X)^{\sim}$ qui
soit une rétraction à gauche de $f$, \ie tel que
\begin{align*}
g^x &= \Prol_X\\
g_x &= \Res_X{\quoi}.
\end{align*}
(Pour\pageoriginale un exemple où $S(X)$ n'admet pas de produits
fibrés, prendre pour $S$ la catégorie des schémas munie de la
topologie étale, pour $M$ les morphismes lisses, pour $X$ un
schéma noethérien de dimension $>0$).

\item[$6^{\circ})$] Montrer que pour tout faisceau abélien $F$ de
$(S_{/X})^{\sim}$ on a un isomorphisme
$$
H^q(X,F)\simeq H^q(X,\Res_X F)\qquad \forall q{\quoi}.
$$
Montrer, en utilisant par exemple les hyperrecouvrements, que pour
tout faisceau abélien $G$ de $S(X)^{\sim}$ on a un isomorphisme
$$
H^q(X,G)\simeq H^q(X,\Prol_X G)\qquad \forall q{\quoi}.
$$

\item[$7^{\circ})$] Acheter une médaille en chocolat pour le rédacteur.
\end{itemize}
\end{enonce*}

\section{Topos induit}\etiquette[5]{IV.sec5}

\subsection{}\etiquette[5.1]{IV.subsec5.1}
Soient $E$ un topos, $X$ un objet de $E$. Alors la catégorie
$E_{/X}$ des objets de $E$ au-dessous de $X$ est un topos, comme il
résulte par exemple immédiatement du critère de Giraud
\hyperref[IV.theo1.2]{1.2} ii). On peut aussi, grâce à \hyperref[IV.cor1.2.1]{1.2.1}
réaliser $E$ comme une catégorie de faisceaux $C^{\sim}$, où
$C$ est une sous-catégorie\pageoriginale génératrice de $E$
qu'on peut choisir telle que $X\in \ob C$; alors on sait que
$E_{/X}=C^{\sim}_{/X}$ est équivalente à $(C_{/X})^{\sim}$
(\hyperrefext[III][III.prop5.4]{5.4}), donc c'est un topos.

On appelle le topos $E_{/X}$ le \emph{topos induit} sur l'objet $X$
de $E$.

\subsection{}\etiquette[5.2]{IV.subsec5.2}
On va définir un morphisme de topos canonique
\begin{equation*}
j_X:E_{/X}\longrightarrow E{\quoi},\tag{5.2.1}
\end{equation*}\etiquette[5.2.1]{IV.eq5.2.1}
qui est appelé \emph{morphisme d'inclusion} du topos induit
$E_{/X}$ dans le topos ambiant $E$, ou mieux (\cf
\hyperref[IV.subsec5.7]{5.7}), \emph{morphisme de localisation de $E$ en $X$.}
Ce morphisme correspond à une suite de trois foncteurs adjoints
(\cf \hyperref[IV.subsubsec3.1.3]{3.1.3})
\begin{equation*}
j_{X!}\quav j^{\ast}_X\quav j_{X^{\ast}}{\quoi},\tag{5.2.2}
\end{equation*}\etiquette[5.2.2]{IV.eq5.2.2}
qui peuvent s'expliciter de la façon suivante:
\begin{enumerate}
\renewcommand{\theenumi}{\alph{enumi}}
\renewcommand{\labelenumi}{\theenumi)}
\item Le foncteur
$$
j_{X!}:E_{/X}\longrightarrow E
$$
est le foncteur « \emph{oubli de la flèche structurale} » de
$E_{/X}$ dans $E$.

\item Le foncteur
$$
j^{\ast}_X:E\longrightarrow E_{/X}
$$
est défini par
$$
j^{\ast}_X(Z)=(X\times Z, \pr_1){\quoi},
$$
où $\pr_1:X\times Z\rightarrow X$ est la première projection. On
peut aussi l'interpréter comme le \emph{foncteur changement de
base} relativement au morphisme
$$
X\longrightarrow e_E{\quoi},
$$
où $e_E$ est l'objet final de $E$. Il est trivial, par
définition du foncteur changement de base, que $j^{\ast}_X$ est
bien adjoint à droite de $j_{X!}$. Il commute en particulier aux
limites projectives. Il commute également aux limites
inductives,\pageoriginale en vertu des propriétés d'exactitude
spéciales des topos (\hyperrefext[II][II.sec4]{4}).

\item De ce dernier fait résulte (\hyperref[IV.cor1.6]{1.6}) que $j^{\ast}_X$ admet un
adjoint à droite
$$
j_{X^{\ast}}:E_{/X}\longrightarrow E{\quoi}.
$$
Pour un objet $X'$ au-dessus de $X$, l'objet $j_{X^{\ast}}(X')$ est
aussi parfois noté par l'un des symboles $\prod_{X/e_E}(X'/X)$, ou
$\Hhom_{X/e_E}(X,X')$, ou $\Res_{X/e_E}(X'/X)$ (« restriction de
Weil »), dont l'un ou l'autre est sans doute déjà familier au
savant lecteur de notre modeste ouvrage.
\end{enumerate}

\subsection{}\etiquette[5.3]{IV.subsec5.3}
On fera attention que le foncteur $j_{X!}$ commute aux produits
fibrés, et transforme monomorphismes en monomorphismes, mais il
n'est pas en général exact à gauche pour autant. Il ne l'est
que si $X$ est un objet final de $ℰ$ (puisque
$X=j_{X!}(X,\id_X)$, et $(X,\id_X)$ est un objet final de $E_{/X}$),
c'est-à-dire si et seulement si $j_X$ est en fait une
équivalence de topos (donc si tous les foncteurs \eqhyperref[IV.eq5.2.2]{5.2.2}
sont des équivalences). \footnote[1]{Pour une propriété de
commutation de $j_{X!}$ au changement de topos, \cf XVII 5.1.2.}

De même, le foncteur $j_{X^{\ast}}$ n'est pas en général exact
à droite, et ne transforme pas nécessairement épimorphismes en
épimorphismes.

On conclut de ces observations que la direction du morphisme de
topos reliant $E_{/X}$ à $E$, pour un objet $X$ du topos $E$, est
déterminée sans ambiguïté possible.

\subsection{}\etiquette[5.4]{IV.subsec5.4}
Le foncteur $j^{\ast}_X$ est souvent appelé \emph{foncteur de
localisation} ou \emph{foncteur restriction}. Cette dernière
terminologie se justifie en identifiant les objets de $E$ \resp sur
$E_{/X}$ aux faisceaux sur $E$ \resp sur $E_{/X}$ (\hyperref[IV.theo1.2]{1.2}
iii)), et en notant qu'avec cette identification, la formule
d'adjonction entre $j_{X!}$ (foncteur oubli) et $j^{\ast}_X$
s'interprète en disant que $j^{\ast}_X(F)$ \emph{est\pageoriginale
le faisceau restriction à $E_{/X}$ du faisceau $F$ sur $E$} (en
prenant « restriction »  au sens généralisé: composé avec le
foncteur « d'inclusion » $j_{X!}:E_{/X}\rightarrow E$).
Conformément à des notations familières en d'autres contextes, on
écrira aussi souvent, indifféremment:
\begin{equation*}
j^{\ast}_X(F)=F| X=F_X=\mbox{ restriction de } F\mbox{ à }
X{\quoi}.\tag{5.4.1}
\end{equation*}\etiquette[5.4.1]{IV.eq5.4.1}

De même, lorsque $E$ est de la forme $C^{\sim}$, où $C$ est un
(petit) site, et que $X$ est de la forme $\epsilon(S)$ avec $S\in
\ob C$, de sorte qu'on a une équivalence de catégories rappelée dans
\hyperref[IV.subsec5.1]{5.1}
$$
E_{/X}=C^{\sim}_{/\epsilon(S)}\xrightarrow{\approx}(C_{/S})^{\sim}
$$
($C_{/S}$ étant munie de la topologie induite par celle de $C$),
le foncteur $j^{\ast}_X$ s'identifie simplement au foncteur « restriction »  d'un faisceau variable $F$ sur $C$ à la
catégorie $C_{/S}$.

Ces réflexions permettent de prévoir d'ailleurs le rôle
important que joueront les topos induits et les morphismes de
localisation \eqhyperref[IV.eq5.2.1]{5.2.1} dans toutes les questions où on est
amené à raisonner par « localisation sur $E$ »  (\cf 8),
c'est-à-dire pratiquement dans toutes les questions faisant
intervenir des sites ou des topos.

\subsection{}\etiquette[5.5]{IV.subsec5.5}
Soit
$$
f:X\longrightarrow Y
$$
une flèche de $E$, qui permet donc d'interpréter $X$ (ou plus
correctement, $(X,f)$) comme un objet de $E_{/Y}$. Il est évident
qu'on a un isomorphisme canonique
\begin{equation*}
(E_{/Y})_{/X}\xrightarrow{\sim}E_{/X}\tag{5.5.1}
\end{equation*}\etiquette[5.5.1]{IV.eq5.5.1}
(\emph{transitivité des topos induits}). Appliquant la
construction de \hyperref[IV.subsec5.1]{5.1} à $E_{/Y}$ au lieu de $E$, on
trouve donc un morphisme de topos canonique
\begin{equation*}
\loc(f)\mbox{ ou } f:E_{/X}\longrightarrow E_{/Y}{\quoi},\tag{5.5.2}
\end{equation*}\etiquette[5.5.2]{IV.eq5.5.2}
appelé\pageoriginale également \emph{morphisme de localisation
(associé à $f$).} Lorsque $Y$ est l'objet final de $E$, on
retrouve essentiellement \eqhyperref[IV.eq5.2.1]{5.2.1}. Le morphisme de localisation
pour $f$ est associé à une suite de trois foncteurs adjoints
\begin{equation*}
f_{!}\quav f^{\ast}\quav f_{\ast}{\quoi},\tag{5.5.3}
\end{equation*}\etiquette[5.5.3]{IV.eq5.5.3}
qui peuvent s'interpréter respectivement comme foncteur d'oubli,
comme foncteur restriction, et comme un foncteur noté au choix ($X'$
étant l'argument) $\prod_{X/Y}(X'/X)$, $\Hhom_{X/Y}(X,X')$ ou
$\Res_{X/Y}(X'/X)$.

\subsection{}\etiquette[5.6]{IV.subsec5.6}
La transitivité des foncteurs oubli implique que pour deux
morphismes composables
$$
X\xrightarrow{f} Y \xrightarrow{g} Z
$$
de $E$, on a un isomorphisme canonique de morphismes de topos
$$
\loc(gf)\simeq \loc(g)\loc(f):E_{/X}\longrightarrow
E_{/Y}\longrightarrow E_{/Z}{\quoi}.
$$
De plus, pour trois morphismes composables, on a la relation de
compatibilité habituelle pour les isomorphismes de transitivité
précédents. Ceci permet donc de considérer, moyennant l'abus de
langage habituel, que
$$
X\longmapsto E_{/X}
$$
est un foncteur covariant de $E$ dans la catégorie
($𝒱$-$𝒰$-$\text{top}$) (\hyperref[IV.subsubsec3.3.1]{3.3.1}).
Plus précisément, on trouve un 2-foncteur (non strict en
général) de $E$ dans la 2-catégorie
($𝒱$-$𝒰$-$\text{top}$).

Avant de continuer les généralités sur les topos induits
(\hyperref[IV.subsec5.10]{5.10}), donnons quelques exemples instructifs.

\subsection{}\etiquette[5.7]{IV.subsec5.7}
Soit\pageoriginale $X$ un espace topologique, d'où un topos
$\Top(X)$ (\hyperref[IV.subsec2.1]{2.1}). Soit $X'$ un objet de $\Top(X)$,
qu'il sera commode d'interpréter comme un espace étalé sur
$X$, $p:X'\rightarrow X$. La catégorie $\Top(X)_{/X'}$ s'identifie
alors à la catégorie des espaces étalés $X''$ sur $X$, munis
d'un $X$-morphisme $X''\rightarrow X'$. On sait qu'un tel morphisme
fait de $X''$ un espace étalé sur $X'$, et on trouve de cette
façon une équivalence de topos canonique
$$
\Top(X)_{/X'}\xrightarrow{\approx} \Top(X'){\quoi}.
$$
Le morphisme de localisation s'identifie donc à un morphisme de
topos $\Top(X')\rightarrow \Top(X)$, et on constate aussitôt que
ce dernier n'est autre que le morphisme
$$
\Top(p):\Top(X')\longrightarrow \Top(X)
$$
associé à l'application continue structurale $p:X'\rightarrow
X$. Intuitivement, le morphisme de localisation n'est donc autre que
la traduction, en langage de topos, du morphisme d'étalement
$p:X'\rightarrow X$. Ceci explique à la fois le bien-fondé de la
terminologie « morphisme de localisation », et incite à la
prudence dans l'emploi du terme de « morphisme d'inclusion »,
celui-ci semblant surtout approprié dans le cas où $p$ est une
immersion ouverte. Il sera prudent par conséquent de réserver
dans \hyperref[IV.subsec5.1]{5.1} le terme de « morphisme d'inclusion »
au cas où le morphisme $X\rightarrow e_E$ est un monomorphisme,
\ie où $X$ s'identifie à un sous-objet de l'objet final de $E$.

\subsection{}\etiquette[5.8]{IV.subsec5.8}
Soient $E$ un topos, $G$ un Groupe de $E$, $H$ un sous-groupe de
$G$,
$$
X=G/H
$$
l'espace homogène quotient, qu'on regarde comme un objet de $E$
avec groupe d'opérateurs gauche $G$, \ie comme un objet du topos
classifiant $B_G$ (\hyperref[IV.subsec2.4]{2.4}). Nous\pageoriginale allons
déterminer le topos induit
$$
(B_G)_{/X}\quad (X=G/H){\quoi},
$$
en définissant une équivalence de topos
\begin{equation*}
c:B_H \xrightarrow{\approx} (B_G)_{/X}{\quoi},\tag{5.8.1}
\end{equation*}\etiquette[5.8.1]{IV.eq5.8.1}
telle que l'on ait commutativité à isomorphisme canonique près
dans le diagramme
\[\vcenter{
\xymatrix@C=2cm{ B_H \ar[r]^{c} \ar[r] \ar
@{^{(}->}@<-9pt>[rr]_{B_i}& {(B_{G})_{/X}} \ar[r]^{j_X}& B_G
}}\leqno{(5.8.2)}
\]\etiquette[5.8.2]{IV.eq5.8.2}
où $j_X$ est le morphisme de localisation dans $B_G$, et où
$$
B_i:B_H\longrightarrow B_G
$$
est le morphisme déduit de l'inclusion $i:H\rightarrow G$
(\hyperref[IV.subsec4.5]{4.5}).

Pour définir \eqhyperref[IV.eq5.8.1]{5.8.1}, nous allons simplement indiquer la
description des foncteurs image inverse $c^{\ast}$ et image directe
$c_{\ast}$, laissant au lecteur le soin de vérifier que ce sont
bien là des équivalences quasi-inverses l'une de l'autre,
donnant lieu au diagramme commutatif \eqhyperref[IV.eq5.8.2]{5.8.2} de morphismes de
topos. Soit $e$ le sous-objet de $E$ sous-jacent à $X$, image de
la section de $X$ (sur un objet final $e_E$ choisi de $E$) déduite
de la section unité de $G$. Si $X'$ est un objet de $B_G$
au-dessus de \sisi{$G$}{$X$}, alors on désigne par $c^{\ast}(X')$ l'image
inverse de $e$ dans $X'$,
$$
c^{\ast}(X')=\fprod{X'}{e}{X}{\quoi},
$$
qui est stable par l'opération gauche de $H$ sur $X'$, restriction
de l'opération donnée de $G$ sur $X'$. Cela définit bien un
foncteur
$$
c^{\ast}:(B_G)_{/X}\longrightarrow B_H{\quoi}.
$$
D'autre\pageoriginale part, si $X'$ est un objet de $B_H$, alors le
morphisme de $X'$ dans l'objet final $e_H$ de $B_H$ (\ie l'objet
final $e_E$ de $E$ avec opération triviale de $H$) donne, par
application du foncteur $i_{!}$ (\hyperref[IV.subsec4.5]{4.5}) un morphisme de
$B_G$
$$
i_{!}(X')\longrightarrow i_{!}(e_H)\simeq X{\quoi},
$$
qui permet d'interpréter $i_{!}(X')$ comme un objet de
$(B_G)_{/X}$,  d'où le foncteur cherché
$$
c_{\ast}:X'\longmapsto (i_{!}(X')\longrightarrow
X):B_H\longrightarrow \sisi{B_G{\quoi}.}{(B_G)_{/X}{\quoi}.}
% correction SGA4-CORR-0002 / SGA-SRC-000389: SGA4-EIV-04-L2917-L2918-88747C81142F.
$$

\setcounter{subsubsection}{2}
\subsubsection{}\etiquette[5.8.3]{IV.subsubsec5.8.3}
On voit donc, grâce à \eqhyperref[IV.eq5.8.2]{5.8.2} que si $G$ est un Groupe
d'un topos, $H$ un sous-Groupe, \emph{le foncteur « restriction
des opérateurs de $G$ à $H$ »  peut s'interpréter comme un
foncteur de localisation.} En particulier, prenant pour $H$ le
sous-groupe unité, on voit que le foncteur « \emph{oubli des
opérations de $G$} » s'interprète comme un foncteur de
localisation. On en conclut, plus généralement, que si $(X,G)$
est un objet de $E$ avec opération de $G$ (\ie un objet du topos
classifiant $B_G$), alors le foncteur « oubli des opérations de
$G$ »,
$$
E'=(B_G)_{/(X,G)}\longrightarrow E_{/X}
$$
peut s'interpréter comme un foncteur de localisation, relativement
à un objet $E_{G,(X,G)}$ de $E'$ convenable couvrant l'objet
final. Il suffit de prendre $E_{G,(X,G)}=E_G\times (X,G)$, d'où le
diagramme cartésien (où $E_G=G_s=G$ avec opération de $G$ par
translation à gauche)
\[\vcenter{
\xymatrix@C=2cm@R=1.5cm{
E_G \ar[d]& **[r]E_{G,(X,G)} = Z \ar[l] \ar[d]\\
e_G & (X,G) \ar[l]          \quad , }}\leqno{(5.8.3.1)}
\]\etiquette[5.8.3.1]{IV.eq5.8.3.1}
et\pageoriginale de noter que par transitivité des topos induits,
le topos induit $E'=((B_G)_{X/G)/Z}$ est équivalent au topos
$((B_G)_{/E_G})_{/Z}$; comme $(B_G)_{/E_G}$ est équivalent à $E$
par le foncteur $c^{\ast}$ de \eqhyperref[IV.eq5.8.1]{5.8.1} (avec $H=e$), il suffit
de vérifier que cette équivalence transforme $Z$ en l'objet $X$
de $E$ (ce qui est trivial) pour déduire l'équivalence
cherchée de catégories
$$
E'_{/Z}\xrightarrow{\approx}E_{/X}{\quoi}.
$$
On laisse au lecteur le soin de vérifier que le composé de
celle-ci avec le foncteur de localisation
$$
E'=(B_G)_{/(X,G)}\longrightarrow E'_{/Z}
$$
est le foncteur « oubli des opérations de $G$ ». On voit ainsi
que le diagramme cartésien ci-dessus donne, par passage aux topos
induits, un diagramme, commutatif à isomorphisme canonique près
(\hyperref[IV.subsec5.6]{5.6}), de morphismes de topos:
\[\vcenter{
  \xymatrix@C=1.25cm@R=1cm{
**[r](B_G)_{/E_G}\approx E \ar[d]_{j} & **[r](B_G)_{/Z} \approx
E'_{/Z} \ar[l]_(.3){i}
  \approx E_{/X} \ar[d]^{j'}\\
B_G & **[r](B_G)_{/(X,G)} =  E' \ar[l]_(.5){f}\quad,
}}\leqno{(5.8.3.2)}
\]\etiquette[5.8.3.2]{IV.eq5.8.3.2}
où les foncteurs images inverses associés aux flèches verticales
$j$, $j'$ sont les foncteurs « oubli des opérations de $G$ », et
où le foncteur image inverse $i^{\ast}$ s'identifie au foncteur de
localisation $E\rightarrow E_{/X}$. Ce diagramme\pageoriginale est
d'ailleurs « 2-cartésien »  au sens \hyperref[IV.prop5.11]{5.11}
ci-dessous.

\subsubsection{}\etiquette[5.8.4]{IV.subsubsec5.8.4}
On laisse au lecteur le soin d'étendre les réflexions qui
précèdent au cas d'un pro-groupe $𝒢=(G_i)$, dans le cas
particulier où $H$ est défini comme « noyau »  de
$𝒢\rightarrow G_{i_{\circ}}$. (Dans le cas des groupes
pro-finis, cela signifie qu'on se limite aux sous-groupes $H$ de $G$
\emph{ouverts}, \ie fermés et d'\emph{indice fini}.)

\begin{enonce*}[remark]{Exercice 5.9}\etiquette[5.9]{IV.exer5.9}
Soient $E$ un topos, $G$ un Groupe de $E$, $B_G$ le topos
classifiant de $G$ (\hyperref[IV.subsec2.4]{2.4}),
$$
\pi:B_G\longrightarrow E
$$
le morphisme de topos déduit de l'homomorphisme de $G$ dans le
Groupe unité $e$ de $E$ (compte tenu que $B_e\approx E$).
\begin{enumerate}
\renewcommand{\theenumi}{\alph{enumi}}
\renewcommand{\labelenumi}{\theenumi)}
\item Soit
$$
E_G=G_s
$$
l'objet de $B_G$ dont le $E$-objet sous-jacent est $G$, les
opérations de $G$ étant définies par translation à gauche.
Considérons $\pi^{\ast}(G)$ comme un Groupe de $B_G$ (c'est le
Groupe $G$ de $E$ avec les opérations triviales de $G$ dessus).
Montrer que le morphisme de $E$
$$
G_s\times G\longrightarrow G_s
$$
défini par les translations droites de $G$ sur $G_s$, est
compatible avec les opérations de $G$ sur $G_s$ et sur
$G=\pi^{\ast}(G)$, et définit un morphisme de $B_G$
$$
E_G\times \pi^{\ast}(G)\longrightarrow E_G{\quoi}.
$$
Montrer que ce morphisme fait de $E_G$ un objet de $B_G$ avec une
opération à droite du $B_G$-Groupe $\pi^{\ast}(G)$, et que cette
dernière fait de $E_G$ un \emph{torseur} sous\pageoriginale
$\pi^{\ast}(G)$ (\ie l'objet final de $B_G$ peut se recouvrir par
des objets $X_i$ tels que pour tout $i$, la restriction de $E_G$ à
$X_i$ soit isomorphe au fibré à opérateurs trivial
$(\pi^{\ast}(G)_{X_i})_d$).

\item Soit
$$
q:E'\longrightarrow E
$$
un topos sur $E$. Pour tout topos $B$ sur $E$, définir la
catégorie $\Hhomtop_E(E',B)$ des morphismes de topos de $E'$ dans
$B$ compatibles avec les morphismes structuraux $E'\rightarrow E$ et
$B\rightarrow E$ à isomorphisme donné près. Prenant $B=B_G$,
définir par la formule $f\mapsto f^{\ast}(E_G)$ une équivalence
de catégories
$$
\Hhomtop_E(E',B_G)\longrightarrow \sisi{\Ttors(E,q^{\ast}(G))}{\Ttors(E',q^{\ast}(G))}{\quoi},
$$
où le deuxième membre désigne la catégorie des
$q^{\ast}(G)$-torseurs (à droite) \sisi{sur $E$.}{sur $E'$.}
% correction SGA4-CORR-0003 / SGA-SRC-000390: SGA4-EIV-04-L3025-L3039-12E69F8B139B.

\item Soit $H$ un sous-Groupe de $G$, de sorte que $\pi^{\ast}(H)$ est
un sous-Groupe de $p^{\ast}(G)$, et opère donc sur $E_G$ par
restriction du Groupe d'opérateurs, d'où un objet quotient
$$
X=E_G/\pi^{\ast}(H){\quoi},
$$
qui n'est autre que l'objet $G/H$ muni de l'opération à gauche
habituelle de $G$. Soit $e_G$ l'objet final de $B_G$ (\ie l'objet
final $e_E$ de $E$ avec opération triviale - il n'y a pas le choix
- de $G$), et considérons dans $B_G$ le diagramme de morphismes
\[
\xymatrix@C=-1cm@R=.75cm{
E_G \ar[dd] \ar[rd]& \\
    & \qquad\qquad\quad X = E_{G}/\pi^{*}(H)\ar[ld]\\
e_G & }
\]
d'où,\pageoriginale en passant aux topos induits par $B_G$ sur ces
objets, un diagramme de morphismes de topos, commutatif à
isomorphisme canonique près (\hyperref[IV.subsec5.6]{5.6}). Montrer que ce
diagramme est équivalent au diagramme
\[
\xymatrix@C=.05cm@R=.85cm{
E & = & B_e \ar[dd] \ar[rrrd]& & &\\
  &   &                    &  & &B_H \ar[llld]\\
  &   & B_G                &  & & &,
}
\]
dont les trois flèches sont les flèches associées par
\hyperref[IV.subsec4.5]{4.5} aux morphismes de Groupes $e\rightarrow
H\rightarrow G$ de $E$. (Ainsi, \emph{le topos classifiant $B_H$
s'interprète intuitivement comme un espace homogène au-dessus de
$B_G$, de groupe $\pi^{\ast}(G)$, associé au torseur} (= fibré
principal homogène) \emph{universel $E_G$ sur $B_G$}).

\item Reprendre l'exemple \hyperref[IV.subsec5.7]{5.7} dans le cas où $X'$ est un
\emph{revêtement} étale (localement trivial, \cf
\hyperref[IV.subsubsec2.7.4]{2.7.4}) de $X$, en l'interprétant en termes des
considérations du présent exercice.
\end{enumerate}
\end{enonce*}

\setcounter{subsection}{9}
\subsection{Images inverses de topos induits}\etiquette[5.10]{IV.subsec5.10}

Soient
$$
f:E'\rightarrow E
$$
un morphisme de $𝒰$-topos, $X$ un objet de $E$, et
$X'=f^{\ast}(X)$ \sisi{son image inverse dans $E$.}{son image inverse dans $E'$.} On va alors définir
% correction SGA4-CORR-0004 / SGA-SRC-000391: SGA4-EIV-04-L3089-L3090-12922A6899C0.
un diagramme de morphismes de topos
\[\vcenter{
\xymatrix@C=2.5cm@R=1.5cm{
E_{/X}  \ar[d]_{j_X} & E'_{/X'} \ar[l]_{f_{/X}} \ar[d]^{j_{X'}} \\
E      &  E'\ar[l]_{f}\quad , }}\leqno{(5.10.1)}
\]\etiquette[5.10.1]{IV.eq5.10.1}
où\pageoriginale $j_X$ et $j_{X'}$ sont les morphismes de
localisation (\hyperref[IV.subsec5.2]{5.2}), et où $f_{/X}$ est défini
ci-dessous, diagramme qui sera commutatif à isomorphisme canonique
près. Nous définirons ici $f_{/X}$ à l'aide du foncteur image
inverse
\begin{equation*}
(f_{/X})^{\ast}:E_{/X}\longrightarrow E'_{/X'}{\quoi},\tag{5.10.2}
\end{equation*}\etiquette[5.10.2]{IV.eq5.10.2}
pour lequel nous prendrons simplement le foncteur « induit »
par $f^{\ast}$, en un sens évident. Comme les foncteurs
d'inclusion $j_{X!}:E_{/X}\rightarrow E$ et
\sisi{$j_{X'!}:E_{/X}\rightarrow E$}{$j_{X'!}:E'_{/X'}\rightarrow E'$} commutent aux $𝒰$-limites
% correction SGA4-CORR-0005 / SGA-SRC-000392: SGA4-EIV-04-L3107-L3109-0A24DC060A72.
inductives et aux produits fibrés, et qu'ils sont conservatifs, il
s'ensuit aussitôt que le foncteur $(f_{/X})^{\ast}$ commute aux
$𝒰$-limites inductives et aux produits fibrés tout comme
le foncteur $f^{\ast}$ qui l'induit; comme de plus $(f_{/X})^{\ast}$
transforme évidemment objet final $X$ en l'objet final $X'$, il
est exact à gauche, donc est associé à un morphisme de topos
\begin{equation*}
f_{/X}:E'_{/X'}\longrightarrow E_{/X}{\quoi},\tag{5.10.3}
\end{equation*}\etiquette[5.10.3]{IV.eq5.10.3}
défini à isomorphisme unique près (\hyperref[IV.subsubsec3.1.1]{3.1.1}).
D'autre part, le diagramme déduit de \eqhyperref[IV.eq5.10.1]{5.10.1} par passage
aux foncteurs images inverses est commutatif (à isomorphisme
canonique près) par construction et grâce au fait que $f^{\ast}$
commute aux produits $X\times Y$, donc \eqhyperref[IV.eq5.10.1]{5.10.1} lui-même est
commutatif, à un unique isomorphisme près induisant
l'isomorphisme canonique sur les foncteurs images réciproques des
deux composés $f\circ j_{X'}$ et $j_X\circ f_{/X}$ (3.2.1).

\setcounter{subsubsection}{3}
\subsubsection{}\etiquette[5.10.4]{IV.subsubsec5.10.4}
Lorsque $f:E'\rightarrow E$ est associé à une application
continue $Y'\rightarrow Y$ d'espaces topologiques
(\hyperref[IV.subsec4.1]{4.1}), de sorte que $X$ s'identifie à un espace
étalé au-dessus de $Y$, et $X'$ au produit fibré
$\fprod{X}{Y'}{Y}$, alors, identifiant les topos induits
$\Top(Y)_{/X}$ et $\Top(Y')_{/X'}$ à $\Top(X)$ et $\Top(X')$
respectivement (\hyperref[IV.subsec5.7]{5.7}), on constate que le morphisme
$f_{/X}$ qu'on vient de construire\pageoriginale est (à
isomorphisme unique près) le morphisme $\Top(X')\rightarrow
\sisi{\Top(Y')}{\Top(X)}$ déduit de l'application continue canonique
\sisi{$X'\rightarrow Y'$}{$X'\rightarrow X$}. On peut donc dire, à la lumière de cet
% correction SGA4-CORR-0006 / SGA-SRC-000393: SGA4-EIV-04-L3131-L3139-EFEC2484AE43.
exemple, que le topos $E'_{/X'}$ joue le rôle d'un produit fibré
de $E_{/X}$ et de $E'$ sur $E$. Cette intuition se précise à
l'aide du résultat suivant:

\begin{enonce*}{Proposition 5.11}\etiquette[5.11]{IV.prop5.11}
Les notations étant celles de $\hyperref[IV.subsec5.10]{5.10}$, le diagramme
$\eqhyperref[IV.eq5.10.1]{5.10.1}$, muni de l'isomorphisme de compatibilité
$\alpha:f\circ j_{X'}\rightarrow j_X\circ f_{/X}$, est « $2$-cartésien »; de façon précise, pour tout
$𝒰$-topos $F$, si on associe à tout morphisme de topos
$g:F\rightarrow E'_{/X'}$ le triplet $(f_{/X}\circ g, j_{X'}\circ g,
\alpha\ast g)=(g_1, g_2,\beta)$, où $g_1:F\rightarrow E_{/X}$ et
$g_2:F\rightarrow E'$ sont des morphismes de topos, et
$\beta:j_X\circ g_1\rightarrow f\circ g_2$ est un isomorphisme de
morphismes de topos de $F$ dans $E$, on trouve une équivalence de
catégories de $\sisi{\Hhomtop(F,E'_{/X})}{\Hhomtop(F,E'_{/X'})}$ avec la catégorie
% correction SGA4-CORR-0007 / SGA-SRC-000394: SGA4-EIV-04-L3144-L3156-68CB9AC918A9.
\sisi{$\fprod{\Hhomtop(F,E_{/X})^2}{\Hhomtop(F,E')}{\Hhomtop(F,E)}$}{$\Hhomtop(F,E_{/X})\times^{2}_{\Hhomtop(F,E)}\Hhomtop(F,E')$} de
tous les triples $(g_1,g_2,\beta)$ comme ci-dessus.
\end{enonce*}

(NB. On a mis un 2 au-dessus du signe du produit cartésien pour
rappeler qu'il ne s'agit pas d'un produit fibré ordinaire de
catégories, mais d'un « 2-produit fibré », dans le sens
explicité dans l'énoncé.)

Pour prouver \hyperref[IV.prop5.11]{5.11}, on explicite les morphismes de topos
à l'aide des foncteurs image inverse associés. Nous aurons
besoin pour cela de résultats auxiliaires, donnés dans
\hyperref[IV.subsubsec5.11.1]{5.11.1} à \hyperref[IV.prop5.12]{5.12} plus bas.

\setcounter{subsection}{11} \setcounter{subsubsection}{0}
\subsubsection{}\etiquette[5.11.1]{IV.subsubsec5.11.1}
Considérons de façon générale la situation où on se donne
deux catégories $E$, $E'$, un objet $X$ de $E$ qu'on suppose
quarrable \ie tel que le foncteur
$$
j:A\longmapsto A\times X:E\longrightarrow E_{/X}
$$
soit\pageoriginale défini, enfin un foncteur exact à gauche
\begin{equation*}
\varphi:E_{/X}\longrightarrow E'{\quoi}.\tag{$\ast$}
\end{equation*}
Comme $E_{/X}$ a un objet final, savoir $(X,\id_X)=X$, il en est
donc de même de $E'$, qui admet l'objet final
$$
e'\simeq \varphi(X){\quoi}.
$$
Supposant choisi l'objet final $e'$ de $E'$, on va à tout $\varphi$
comme ci-dessus associer un couple
\begin{equation*}
(\Psi,u),\mbox{ avec } \Psi=\varphi\circ j:E\rightarrow E',\mbox{ et
} u\in \Hom(e',\Psi(X)){\quoi}.\tag{${\ast}{\ast}$}
\end{equation*}
Pour définir $u$, on note que
$$
\sisi{\Psi(X)=X\times X}{j(X)=X\times X}
% correction SGA4-CORR-0008 / SGA-SRC-000395: SGA4-EIV-04-L3186-L3200-72E953777946.
$$
a une section canonique $\delta_X$ sur l'objet final $X$ de
$E_{/X}$, savoir la section diagonale, et on définit
$$
u=\varphi(\delta_X):\varphi(X)=e'\longrightarrow \varphi(X\times
X)=\Psi(X){\quoi}.
$$
Je dis que la connaissance du couple $({\ast}{\ast})$ permet de
reconstituer le foncteur $\varphi$ de $(\ast)$ à isomorphisme unique
près. Pour ceci, pour tout objet
$$
p:X'\longrightarrow X
$$
de $E_{/X}$, considérons le diagramme cartésien suivant de
$E_{/X}$:
\[
\xymatrix@C=1cm@R=1.5cm{
X' \ar[d] \ar[r] & **[r]j(X') = X'\times X \ar[d]^{j(p)}\\
X  \ar[r]^{\delta_X}& **[r]j(X) = X \times X   &, }
\]
où la première flèche horizontale est le morphisme graphe
$\sisi{\Gamma_p=(\id_X,p)}{\Gamma_p=(\id_{X'},p)}$. Appliquant le foncteur exact à gauche
% correction SGA4-CORR-0009 / SGA-SRC-000396: SGA4-EIV-04-L3204-L3217-524011AA658A.
$\varphi$ à ce diagramme, et utilisant la définition de $\Psi$
comme $\varphi\circ j$, on trouve un diagramme cartésien
\[
\xymatrix@C=2.25cm@R=1.25cm{
\varphi(X',p) \ar[d] \ar[r]& \Psi (X') \ar[d]^{\Psi(p)}\\
e'  \ar[r]          & \Psi (X) \quad , }
\]
en\pageoriginale d'autres termes on trouve un isomorphisme canonique
\begin{equation*}
\varphi(X',p)\simeq
\fprod{\Psi(X')}{e'}{\Psi(X)}{\quoi}\tag{5.11.1.1}
\end{equation*}\etiquette[5.11.1.1]{IV.eq5.11.1.1}
On constate aussitôt qu'il est fonctoriel en l'objet $(X',p)$ de
$E_{/X}$, ce qui explicite comment on reconstitue à isomorphisme
près le foncteur $\varphi:E_{/X}\rightarrow E'$ à l'aide du couple
$(\Psi,u)$ de $({\ast}{\ast})$. Notons d'ailleurs que le foncteur
$\Psi$ est exact à gauche; plus généralement, $\Psi$ commute à tout
type de limites projectives auquel commute $\varphi$, car $j$
commute aux limites projectives.

Inversement, supposant maintenant que dans $E$ et $E'$ les limites
projectives finies sont représentables, et partant d'un couple
$$
(\Psi,u),\Psi:E\longrightarrow E', u\in \Hom(e',\Psi(X)){\quoi},
$$
où le foncteur $\Psi$ est exact à gauche, on définit par la
formule \eqhyperref[IV.eq5.11.1.1]{5.11.1.1} un foncteur $\varphi:E_{/X}\rightarrow E'$.
On vérifie aussitôt que ce foncteur est exact à gauche, plus
généralement, qu'il commute à tout type de limites projectives
représentables dans $E$ auquel commute $\Psi$. Cela résulte
aussitôt du diagramme cartésien
\[
\xymatrix@C=2.5cm@R=1cm{ \displaystyle{\varprojlim^{E}_{I} X'_i}
\ar[d] &
\displaystyle{\varprojlim^{E_{/X}}_{I} X'_i} \ar[d] \ar[l]\\
\displaystyle{\varprojlim^{E}_{I} X} & X \ar[l]_{\delta} }
\]
reliant\pageoriginale les limites projectives calculées dans $E$
ou dans $E_{/X}$ (où $\delta$ désigne le morphisme diagonal dans
la limite projective du foncteur constant de valeur $X$), en lui
appliquant le foncteur $\Psi$, d'où un diagramme cartésien dans
$E'$, et en composant ce dernier avec le carré cartésien
déduit de $\Psi(X)\leftarrow e'$, faisant apparaître un carré
cartésien composé qui exprime la compatibilité de $\varphi$
à la limite projective envisagée.

On reconstitue à isomorphisme près le couple $(\Psi,u)$ à l'aide de
$\varphi$, en notant que l'on a un isomorphisme canonique
\begin{equation*}
\Psi(X')=\varphi(j(X')){\quoi},\tag{5.11.1.2}
\end{equation*}\etiquette[5.11.1.2]{IV.eq5.11.1.2}
déduit de \eqhyperref[IV.eq5.11.1.1]{5.11.1.1} en y remplaçant $(X',p)$ par
$j(X')=(X'\times X, \pr_2)$, et notant que $\Psi(X'\times X)\simeq
\Psi(X')\times \Psi(X)$. L'isomorphisme précédent est manifestement
fonctoriel en $X'$, \ie il donne un isomorphisme
$$
\Psi\simeq \varphi\circ j{\quoi},
$$
et on vérifie de même que moyennant cet isomorphisme,
$u:e'\rightarrow \Psi(X)$ s'identifie à $\varphi(\delta_X)$. On a
ainsi obtenu l'essentiel du

\begin{enonce*}{Lemme 5.11.2}\etiquette[5.11.2]{IV.lem5.11.2}
Soient $E$ et $E'$ deux catégories où les limites projectives
finies sont représentables, $X$ un objet de $E$, $\Ssex(E,E')$ la
catégorie des foncteurs exacts à gauche $\Psi$ de $E$ dans $E'$, et
$\Ssex(E_{/X},E')$ la catégorie analogue des foncteurs exacts à
gauche $\varphi$ de $E_{/X}$ dans $E'$. On a alors une équivalence
entre la catégorie $\Ssex(E_{/X},E')$ et la catégorie
$\Ssex(E,E')_{/\Gamma}$ des couples $(\Psi,u)$, avec $\Psi\in \ob
\Ssex(E,E')$ et $u:e'\rightarrow \Psi(X)$ (où $e'$ est l'objet
final de $E'$), dont la définition est explicitée dans
$\hyperref[IV.subsubsec5.11.1]{5.11.1}$, ainsi que celle d'un foncteur
quasi-inverse. Si $\varphi$ et $(\Psi,u)$ se correspondent, alors
$\varphi$ commute à un type déterminé de limites projectives
si\pageoriginale et seulement si il en est ainsi de $\Psi$. Si dans
$E$ et $E'$ les foncteurs changement de base commutent à un type
déterminé de limites inductives, alors pour que $\varphi$ commute
aux limites inductives dudit type, il faut et il suffit qu'il en
soit ainsi de $\Psi$.
\end{enonce*}

Nous laissons au lecteur le soin d'expliciter la structure de
catégorie $\Ssex(E,E')_{/\Gamma}$ correspondant aux couples
$(\Psi,u)$. Définissant le foncteur $\Gamma:\Ssex(E,E')\rightarrow
(\Ens)$ par $\Gamma(\Psi)=\Hom(e',\Psi(X))$, on peut interpréter
$u$ dans le couple $(\psi,u)$ comme un élément de $\Gamma(\Psi)$
\ie un homomorphisme $\Psi\rightarrow \Gamma$ dans la catégorie
des préfaisceaux sur $\Ssex(E,E')^{\circ}$, ce qui justifie la
notation ${}_{/\Gamma}$ utilisée dans l'énoncé de
\hyperref[IV.subsubsec5.11.1]{5.11.1}. Il reste, pour prouver \hyperref[IV.lem5.11.2]{5.11.2},
à expliciter dans \hyperref[IV.subsubsec5.11.1]{5.11.1} le caractère
fonctoriel des constructions envisagées pour $\varphi$ \resp
$(\psi,u)$ variable, ce qui est essentiellement trivial et laissé
au lecteur, et enfin à vérifier l'assertion concernant les
propriétés de commutation aux limites inductives, ce qui est
également trivial à l'aide des formules explicites
\eqhyperref[IV.eq5.11.1.1]{5.11.1.1} et \eqhyperref[IV.eq5.11.1.2]{5.11.1.2} reliant ces foncteurs.

On trouve en particulier, lorsque $E$ et $E'$ sont des
$𝒰$-topos, et en interprétant les morphismes de topos à
l'aide des foncteurs images inverses associés:

\begin{enonce*}{Proposition 5.12}\etiquette[5.12]{IV.prop5.12}
Soient $E$ et $E'$ deux $𝒰$-topos, $X$ un objet de $E$, et
considérons sur $\Hhomtop(E',E)$ le foncteur contravariant
$$
\gamma:f\longmapsto
\Gamma(E',f^{\ast}(X))=\Hom(e',f^{\ast}(X)):\Hhomtop(E',E)^{\circ}\longrightarrow
(\Ens){\quoi}.
$$
On a alors une équivalence de catégories
\begin{equation*}
\Hhomtop(E',E_{/X})\xrightarrow{\approx}\Hhomtop(E',E)_{/\gamma}{\quoi},\tag{5.12.1}
\end{equation*}\etiquette[5.12.1]{IV.eq5.12.1}
où \pageoriginale le deuxième membre est la sous-catégorie
pleine de $\Hhomtop(E',E)^{\widehat{\phantom{,}}}_{/\gamma}$
formée des couples $(f,u)$, avec $f\in \Homtop(E',E)$ et $u\in
\gamma(f)$ \ie $u\in \Gamma(E',f^{\ast}(X))$. Ce foncteur s'obtient
en associant à tout morphisme de topos $h:E'\rightarrow E_{/X}$ le
couple $(f,u)$ formé du composé
$$
f=j_X\circ h:E'\rightarrow E_{/X}\rightarrow E{\quoi},
$$
et du morphisme
$$
u:e'=h^{\ast}(X,\id_X)\longrightarrow
f^{\ast}(X)=h^{\ast}(j^{\ast}_X(X))=h^{\ast}(X\times X)
$$
déduit du morphisme diagonal $\delta_X:X\rightarrow X\times X$ en
lui appliquant $h^{\ast}$.
\end{enonce*}

\setcounter{subsection}{12}
\subsubsection{}\etiquette[5.12.2]{IV.subsubsec5.12.2}
Utilisant \hyperref[IV.prop5.12]{5.12}, la démonstration de \hyperref[IV.prop5.11]{5.11}
devient à peu près évidente, et se réduit essentiellement
à ceci (qui tiendra lieu d'une démonstration « en forme »
qui ne serait pas plus instructive): se donner un morphisme de topos
$g:F\rightarrow E'_{/X'}$ « revient au même »  que de se
donner un morphisme de topos $g_2:F\rightarrow E'$ et une section
$u$ de $g^{\ast}_2(X')$; or
$g^{\ast}_2(X')=g^{\ast}_2(f^{\ast}(X))=(fg_2)^{\ast}(X)$, donc la
donnée de $u$ équivaut aussi à la donnée d'un relèvement
du morphisme de topos $fg_2:F\rightarrow E$ en un morphisme de topos
$g_1:F\rightarrow E_{/X}$. Cela exprime bien que la donnée d'un
morphisme de topos $g:F\rightarrow E'_{/X'}$ équivaut
essentiellement à la donnée d'un triple $(g_1,g_2,\alpha)$ comme
dans \hyperref[IV.prop5.11]{5.11}.

\begin{enonce*}[remark]{Remarque 5.13}\etiquette[5.13]{IV.rem5.13}
Nous verrons (\S 15\nde{Ce paragraphe n'existe pas. Voir cependant [16] proposition 4.47 pour 
un énoncé général d'existence de produits fibrés de topos (en notant que tous les morphismes
géométriques entre topos de Grothendieck sont bornés).}) que pour tout diagramme
\[
\xymatrix@C=2cm@R=1.25cm{
E_1 \ar[d]&\\
E   & E_2 \ar[l] }
\]
de\pageoriginale topos, il existe un 2-produit fibré au sens de la
2-catégorie ($𝒱$-$𝒰$-$\Top$) des $𝒰$-topos
$\in 𝒱$ (ne dépendant pas, à équivalence près, du
choix d'un univers $𝒱$ tel que $𝒰\in𝒱$).
Pour d'autres exemples naturels que \hyperref[IV.prop5.11]{5.11} de « produits fibrés de topos », voir le dernier chapitre du livre de
\textsc{Giraud} [3].
\end{enonce*}

\begin{enonce*}[remark]{Exercice 5.14}\etiquette[5.14]{IV.exer5.14}
\begin{enumerate}
\renewcommand{\theenumi}{\alph{enumi}}
\renewcommand{\labelenumi}{\theenumi)}
\item Soient $F$, $G$ deux topos au-dessus d'un topos $E$. Définir
la catégorie $\Hhomtop_E(F,G)$ des couples $(f,\alpha)$, où
$f:F\rightarrow G$ est un morphisme de topos et $\alpha:p\rightarrow
q\circ f$ est un isomorphisme de morphismes de topos \sisi{de $E$ dans $G$}{de $F$ dans $E$}
% correction SGA4-CORR-0010 / SGA-SRC-000397: SGA4-EIV-04-L3388-L3393-FCC0F20683E8.
($p:F\rightarrow E$ et $q:G\rightarrow E$ étant les morphismes de
topos structuraux). Définir des foncteurs d'accouplement
$\Hhomtop_E(F,G)\times \Hhomtop_E(G,H)\rightarrow \Hhomtop_E(F,H)$,
et des isomorphismes d'associativité.

\item Soient $X$ et $Y$ deux objets d'un topos $E$. Définir un
foncteur naturel de la catégorie discrète définie par
l'ensemble $\Hom(X,Y)$, dans la catégorie
$\Hhomtop_E(E_{/X},E_{/Y})$. Compatibilité avec les compositions
de morphismes dans $E$ et les accouplements envisagés dans a).

\item Prouver que le foncteur envisagé dans b) est une équivalence
de catégories. En particulier, un objet $X$ d'un topos $E$ se
reconstitue à isomorphisme unique près, quand on connaît « à
$E$-équivalence près »  le topos induit $E_{/X}$ comme topos
\emph{au-dessus de $E$.}~\footnote[1]{Résultat dû à P.
\textsc{Deligne}. Le cas où $E$ est le topos ponctuel avait été
traité auparavant par Mme \textsc{Hakim}.}
\end{enumerate}
\end{enonce*}

\section{Points d'un topos et foncteurs fibres}\etiquette[6]{IV.sec6}

\setcounter{subsection}{-1}
\subsection{}\etiquette[6.0]{IV.subsec6.0}
Soit $P$ le $𝒰$-\emph{topos ponctuel} type
(\hyperref[IV.subsec2.2]{2.2})
$$
P=(𝒰\text{-}\Ens){\quoi}.
$$
Étant donné l'intuition assez différente qui s'attache d'une part au
symbole\pageoriginale $P$, figurant un objet géométrique dans la
nature d'un point, et d'autre part à $(𝒰\text{-}\Ens)$,
figurant une « grosse catégorie », qu'on interprétera comme « la catégorie des faisceaux sur $P$ », il y a lieu suivant le
contexte d'utiliser l'un ou l'autre symbole exclusivement, bien que
du strict point de vue logique ils désignent un seul et même objet.

\begin{enonce*}{Définition 6.1}\etiquette[6.1]{IV.def6.1}
Soit $E$ un $𝒰$-topos. On appelle point de $E$ tout
morphisme de topos $P\rightarrow E$ du topos ponctuel
$P(\hyperref[IV.subsec6.0]{6.0})$ dans $E$. On appelle catégorie des points
de $E$, et on note $\Ppoints(E)$ ou $\sheaf{\Pt}(E)$, la catégorie
$\Hhomtop(P,E)$ $(\hyperref[IV.subsec3.2]{3.2})$. Si $p:P\rightarrow E$ est un
point de $E$, associé au foncteur image inverse
$p^{\ast}:E\rightarrow (𝒰\text{-}\Ens)$, et si $F$ est un
objet de $E$, l'ensemble $p^{\ast}(F)$ est appelé fibre de $E$ en
$p$, et noté $F_p$.
\end{enonce*}

\setcounter{subsection}{1}
\subsubsection{}\etiquette[6.1.1]{IV.subsubsec6.1.1}
Bien entendu, lorsque $F$ est un objet groupe (\resp anneau,
\resp\ldots) de $E$, sa fibre $F_p$ en $p$ est un groupe (\resp un
anneau, \resp\ldots) (\hyperref[IV.subsubsec3.1.2]{3.1.2}).

\setcounter{subsubsection}{1}

\subsubsection{}\etiquette[6.1.2]{IV.subsubsec6.1.2}
En vertu de 3.2.1, la catégorie des points de $E$ est
équivalente, via le foncteur $p\mapsto p^{\ast}$, à la
catégorie \emph{opposée} à celle des foncteurs
$$
\varphi:E\longrightarrow (𝒰\text{-}\Ens)
$$
qui commutent aux $𝒰$-limites inductives et qui sont exacts
à gauche. Il revient donc essentiellement au même de se donner
un point de $E$, ou un tel foncteur $\varphi$.

\begin{enonce*}{Définition 6.2}\etiquette[6.2]{IV.def6.2}
On appelle foncteur-fibre du $𝒰$-topos $E$ tout foncteur
$\varphi:E\rightarrow (𝒰\text{-}\Ens)$ qui commute aux
$𝒰$-limites inductives et qui est exact à gauche. On
appelle catégorie des foncteurs fibres sur $E$, et on note
$\Ffib(E)$, la\pageoriginale sous-catégorie pleine de
$\Hhom(E,(𝒰\text{-}\Ens))$ formée des foncteurs fibres sur
$E$.
\end{enonce*}

\setcounter{subsection}{2}
\subsubsection{}\etiquette[6.2.1]{IV.subsubsec6.2.1}
En vertu de \hyperref[IV.subsubsec6.1.2]{6.1.2}, la catégorie $\Ffib(E)$ des
foncteurs fibres sur $E$ est donc équivalente à
l'\emph{opposée} de la catégorie $\Ppoint(E)$ des points de $E$,
via le foncteur $p\mapsto p^{\ast}$ de cette dernière dans la
première:
\begin{equation*}
\Ppoint(E)\xrightarrow{\approx}\Ffib(E)^{\circ}{\quoi}.
\tag{6.2.1.1}
\end{equation*}\etiquette[6.2.1.1]{IV.eq6.2.1.1}

Pour qu'un foncteur
$$
\varphi:E\longrightarrow (𝒰\text{-}\Ens)
$$
soit un foncteur fibre, il faut et il suffit qu'il soit le foncteur
fibre $F\mapsto F_p$ associé à un point $p$ de $E$
(\hyperref[IV.def6.1]{6.1}). En d'autres termes, le foncteur \eqhyperref[IV.eq6.2.1.1]{6.2.1.1}
est surjectif. Signalons aussi qu'un foncteur $\varphi$ est un
foncteur fibre si et seulement si il est exact à gauche et s'il
transforme familles couvrantes en familles surjectives
(\hyperref[IV.cor1.7]{1.7}). Cette dernière condition s'exprime aussi en
disant que $\varphi$ commute aux $𝒰$-sommes et transforme
épimorphismes en épimorphismes.

\subsection{}\etiquette[6.3]{IV.subsec6.3}
Soit $C\in 𝒰$ un site. On appelle \emph{point du site} $C$
tout point du topos $C^{\sim}$, \emph{catégorie des points du
site} $C$ la catégorie
$$
\Ppoint(C)=\Ppoint(C^{\sim}){\quoi}.
$$
Considérons d'autre part le foncteur
\begin{equation*}
\varphi\longrightarrow \varphi|
C=\varphi\circ\epsilon_C:\Ffib(C^{\sim})\longrightarrow
\Hhom(C,(𝒰\text{-}\Ens)){\quoi},\tag{6.3.1}
\end{equation*}\etiquette[6.3.1]{IV.eq6.3.1}
qui est pleinement fidèle et dont l'image essentielle est la
sous-catégorie pleine
$\Mmorsite((𝒰\text{-}\Ens),C)^{\circ}$ (\hyperref[IV.prop4.9.4]{4.9.4}),
que nous noterons aussi $\Ffib(C)$. Un foncteur
$$
\Psi:C\longrightarrow (𝒰\text{-}\Ens)
$$
est\pageoriginale appelé un \emph{foncteur fibre sur le site} $C$
s'il est dans l'image essentielle de \eqhyperref[IV.eq6.3.1]{6.3.1}, qu'on vient
d'expliciter. Un tel foncteur est donc \emph{continu} (a fortiori
(III 1.6), il transforme familles couvrantes en familles couvrantes,
\ie en familles surjectives), et il est exact à gauche. Lorsque dans
$C$ les limites projectives finies sont représentables, (condition
vérifiée dans la quasi-totalité des cas qu'on rencontre en pratique)
les propriétés précédentes \emph{caractérisent} les foncteurs fibres
sur $C$, qui sont alors les foncteurs
$\Psi:C\rightarrow(𝒰\text{-}\Ens)$ qui sont exacts à gauche
et qui transforment familles couvrantes en familles surjectives
(\hyperref[IV.prop4.9.4]{4.9.4}).

On retiendra que le foncteur $\varphi\rightarrow \varphi| C$ est une
équivalence de catégories
$$
\Ffib(C^{\sim})\xrightarrow{\approx}\Ffib(C){\quoi},
$$
de sorte qu'\emph{il revient au même essentiellement de se donner un
foncteur fibre sur le topos $C^{\sim}$} (ou encore un point de ce
topos), \emph{ou un foncteur fibre sur le site $C$.} Étant donné un
foncteur fibre $\Psi$ sur $C$, provenant donc à isomorphisme près
d'un foncteur fibre $\varphi$ sur $C^{\sim}$, on reconstitue ce
dernier à l'aide de $\Psi$, à isomorphisme canonique près, par la
formule
\begin{equation*}
\varphi(F)\simeq \varinjlim_{C_{/F}}\Psi (X){\quoi},\tag{6.3.2}
\end{equation*}\etiquette[6.3.2]{IV.eq6.3.2}
où $C_{/F}$ est la catégorie des objets $X$ de $C$ munis d'un
morphisme $X\rightarrow F$ (dans $\widehat{C}$), \ie d'un élément de
$F(X)$. La formule \eqhyperref[IV.eq6.3.2]{6.3.2} est une conséquence immédiate du
fait que $\varphi$ commute aux limites inductives et qu'on a un
isomorphisme canonique fonctoriel en $F$ (\hyperrefext[II][II.cor4.1.1]{4.1.1}):
$$
F\simeq \varinjlim_{C_{/F}}\epsilon_C(X){\quoi}.
$$
Plus\pageoriginale généralement, lorsque $G$ est un préfaisceau sur
$C$, on a un isomorphisme canonique fonctoriel en $G$:
\begin{equation*}
\varphi(\underline{a}G)\simeq
\varinjlim_{C_{/G}}\Psi(X){\quoi},\tag{6.3.3}
\end{equation*}\etiquette[6.3.3]{IV.eq6.3.3}
où $\underline{a}$ $G$ est le faisceau associé à $G$
(\hyperrefext[II][II.cor4.1.1]{4.1.1}). En effet, on sait qu'on a la formule
$$
\underline{a}G\simeq \varinjlim_{C_{/G}}\epsilon_C(X){\quoi}.
$$

\setcounter{subsection}{4} \setcounter{subsubsection}{-1}
\subsubsection{}\etiquette[6.4.0]{IV.subsubsec6.4.0}
Nous renvoyons à \hyperrefext[I][I.def6.1]{6.1} pour la notion de \emph{famille
conservative de foncteurs}
$$
\varphi_i:E\longrightarrow E_i\qquad i\in I{\quoi}.
$$
On notera que lorsque $E$ et les $E_i$ sont des $𝒰$-topos,
et les foncteurs $\varphi_i$ sont des foncteurs images inverses
$f^{\ast}_i$ associés à des morphismes de topos
$$
f_i:E_i\longrightarrow E{\quoi},
$$
alors toutes les propriétés d'exactitude postulées dans
\hyperrefext[I][I.prop6.2]{6.2}, \hyperrefext[I][I.cor6.3]{6.3} et \hyperrefext[I][I.prop6.4]{6.4} sont vérifiées, à
condition dans \hyperrefext[I][I.prop6.2]{6.2} (v) non respé de supposer que $D$ est
fini. Il revient alors au même que $(f^{\ast}_i)_{i\in I}$ soit
conservative, ou conservative pour les \emph{monomorphismes}
(\hyperrefext[I][I.def6.1]{6.1}), ou conservative pour les \emph{épimorphismes}, ou
enfin qu'elle soit \emph{fidèle}.

Lorsque\pageoriginale la famille des foncteurs $f^{\ast}_i$ est
conservative, on dira aussi parfois que \emph{la famille des
morphismes de topos $(f_i)_{i\in I}$ est conservative}. En
particulier:

\begin{enonce*}{Définition 6.4.1}\etiquette[6.4.1]{IV.def6.4.1}
Soient $E$ un $𝒰$-topos, $(p_i:P\rightarrow E)_{i\in I}$ une
famille de points de $E$. On dit que cette famille de points est
conservative si la famille des foncteurs fibres associés $F\mapsto
F_{p_i}$ de $E$ dans $(𝒰\text{-}\Ens)$ est conservative
$(\hyperref[IV.subsubsec6.4.0]{6.4.0})$. On dit que le topos $E$ a suffisamment
de points lorsqu'il admet une famille conservative de points (\ie
lorsque la famille de tous les points du topos est conservative).
\end{enonce*}

\setcounter{subsection}{4} \setcounter{subsubsection}{1}
\subsubsection{}\etiquette[6.4.2]{IV.subsubsec6.4.2}
Signalons que tous les $𝒰$-topos utilisés jusqu'à présent
ont suffisamment de points. On peut cependant « en faisant
exprès » construire des topos qui n'ont pas suffisamment de points
(\hyperref[IV.exer7.2.6]{7.2.6} e) et \hyperref[IV.subsec7.4]{7.4}); on notera qu'un tel
topos est nécessairement non « vide » au sens géométrique de
\hyperref[IV.subsec2.2]{2.2}. Un topos admettant suffisamment de points admet
une famille conservative de points indexée par un $I\in𝒰$
(\hyperref[IV.prop6.5]{6.5}). Noter cependant à ce sujet que si $E$ est un
$𝒰$-topos, l'ensemble des classes d'isomorphie de points de
$E$ n'est pas nécessairement $𝒰$-petit (\hyperref[IV.subsec7.3]{7.3}).
Il l'est cependant dans de nombreux cas rencontrés en pratique.
Enfin, pour un intéressant théorème d'existence (dû à P. Deligne) de
suffisamment de points, couvrant tous les cas rencontrés en
géométrie algébrique\sisi{.}{, voir \hyperrefext[VI][VI.sec9]{9}.}

\subsubsection{}\etiquette[6.4.3]{IV.subsubsec6.4.3}
Lorsque $(p_i)_{i\in I}$ est une famille conservative de points du
topos $E$, on peut appliquer les remarques de \hyperrefext[I][I.prop6.2]{6.2}, qui
impliquent notamment que deux flèches $u$, $v:F\rightrightarrows G$
de $E$ sont égales si et seulement si pour tout $i\in I$, les
flèches induites sur les fibres $u_{p_i}$,
$v_{p_i}:F_{p_i}\rightrightarrows G_{p_i}$ sont\pageoriginale
égales; qu'une flèche $u:F\rightarrow G$ de $E$ est un monomorphisme
(\resp un épimorphisme) si et seulement si pour tout $i\in I$, il en
est ainsi de l'application $u_{p_i}:F_{p_i}\rightarrow G_{p_i}$;
qu'un objet $F$ de $E$ est initial (\resp final) si et seulement si
pour tout $i\in I$, $F_{p_i}$ est vide (\resp réduit à un point);
qu'un objet $F$ est un sous-objet de l'objet final $e_E$ si et
seulement si pour tout $i\in I$, $F_{p_i}$ a au plus un point.

\begin{enonce*}{Proposition 6.5}\etiquette[6.5]{IV.prop6.5}
\begin{enumerate}
\renewcommand{\theenumi}{\alph{enumi}}
\renewcommand{\labelenumi}{\theenumi)}
\item Soient $C$ un $𝒰$-site, $(\varphi_i)_{i\in I}$ une
famille de foncteurs fibres sur $C$ $(\hyperref[IV.subsec6.3]{6.3})$. Pour que
la famille $(p_i)_{i\in I}$ des points correspondants du topos
$E=C^{\sim}$ soit conservative, il faut et il suffit que pour toute
famille $(X_j\rightarrow X)_{j\in J}$ dans $C$, telle que pour tout
$i\in I$, la famille correspondante $(\varphi_i(X_j)\rightarrow
\varphi_i(X))_{j\in J}$ soit surjective, la famille donnée
$(X_j\rightarrow X)_{j\in J}$ soit couvrante.

\item Soit $E$ un $𝒰$-topos. Si $E$ admet suffisamment de
points $(\hyperref[IV.def6.4.1]{6.4.1})$, alors $E$ admet une famille
conservative de points qui est $𝒰$-petite.
\end{enumerate}
\end{enonce*}



L'assertion b) est un cas particulier de \hyperrefext[I][I.prop7.7]{7.7}. Pour
prouver a), appliquons \hyperrefext[I][I.prop7.7]{7.7} à la famille génératrice dans
$E$ formée des $\epsilon(X)(X\in \ob C)$. Rappelons
(\hyperrefext[II][II.theo4.4]{4.4}) que la famille $X_j\rightarrow X$ est couvrante
si et seulement si la famille $\epsilon(X_j)\rightarrow \epsilon(X)$
est épimorphique. Si la famille des points $(p_i)_{i\in I}$ est
conservative, il revient au même (\hyperref[IV.subsubsec6.4.3]{6.4.3}) de dire
que pour tout $i\in I$, la famille des $(X_j)_{p_i}\rightarrow
(X)_{p_i}$ soit surjective, \ie que la famille des
$\varphi_i(X_j)\rightarrow \varphi_i(X)$ soit surjective. Cela
établit le « il faut »  dans a). Pour le « il suffit », on
applique le critère \hyperrefext[I][I.prop7.7]{7.7}, qui nous ramène à vérifier que
tout monomorphisme $F\rightarrow \epsilon(X)$,\pageoriginale tel que
$F_{p_i}\rightarrow \epsilon(X)_{p_i}$ soit un isomorphisme pour
tout $i\in I$, est un isomorphisme, ou ce qui revient au même, un
épimorphisme. Or, on peut trouver une famille couvrante de
morphismes $\epsilon(X_j)\rightarrow F$, et quitte à raffiner
encore, on peut supposer que les morphismes
$\epsilon(X_j)\rightarrow \epsilon(X)$ sont induits par des
morphismes $X_j\rightarrow X$. Par hypothèse, pour tout $i\in I$,
les morphismes composés $(X_j)_{p_i}\rightarrow F_{p_i}\rightarrow
(X)_{p_i}$ forment une famille surjective (comme composée de deux
familles surjectives), \ie la famille des $\varphi_i(X_j)\rightarrow
\varphi_{i}(X)$ est surjective. Il en résulte par hypothèse que la
famille $X_{j}\to X$ est couvrante, donc que la famille des
$\epsilon(X_j)\rightarrow \epsilon(X)$ est épimorphisme, et a
fortiori que $F\rightarrow \epsilon(X)$ est épimorphique.

\begin{enonce*}{Corollaire 6.5.1}\etiquette[6.5.1]{IV.cor6.5.1}
Soit $C$ une catégorie. Une topologie sur $C$ faisant de $C$ un
$𝒰$-site admettant suffisamment de foncteurs fibres (\ie tel
que le topos associé admette suffisamment de points) est entièrement
connue quand on connaît la sous-catégorie pleine $\Phi$ de
$\Hhom(C,\Ens)$ formée des foncteurs fibres sur $C$.
\end{enonce*}

En effet, en vertu de \hyperref[IV.prop6.5]{6.5} a) on sait alors décrire les
familles couvrantes en termes de la famille des foncteurs fibres.

En fait, on va voir plus bas \sisi{\eqhyperref[IV.eq6.8.3]{6.8.3}}{\eqref{IV.eq6.8.5.1}} que la sous-catégorie
$\Phi$ est contenue dans la catégorie des foncteurs
proreprésentables sur $C$, donc $\Phi^{\circ}\simeq
\Ppoint(C^{\sim})$ s'identifie à équivalence près à une
sous-catégorie pleine de $\Pro-C$.

On comparera \hyperref[IV.cor6.5.1]{6.5.1} au théorème de \textsc{Giraud}
\hyperrefext[II][II.theo5.5]{5.5}, qui établit une correspondance biunivoque entre
l'ensemble des topologies sur $C$ et un certain ensemble de
sous-catégories pleines de $\Hhom(C^{\circ},\Ens)$.

\begin{enonce*}[remark]{Exercice 6.5.2}\etiquette[6.5.2]{IV.exer6.5.2}
Soient\pageoriginale $C$ une catégorie équivalente à une catégorie
$\in 𝒰$, $P$ une sous-catégorie de $\Pro(C)$. Pour tout
$p\in P$, on désigne par $X\mapsto X_p$ le foncteur proreprésenté
par $p$. Une famille $(X_i\rightarrow X)_{i\in I}$ dans $C$ est dite
$P$-\emph{couvrante} si pour tout $p\in P$, la famille
$(X_{ip}\rightarrow X_p)_{i\in I}$ est surjective. Montrer qu'il
existe une topologie \sisi{$T_p$}{$T_P$} sur $C$ pour laquelle les familles
couvrantes soient exactement les familles $P$-couvrantes, et que
\sisi{$T_p$}{$T_P$} est la topologie la plus fine sur $C$ pour laquelle les
foncteurs $X\mapsto X_p$ $(p\in P)$ soient des foncteurs fibres.
Montrer que la topologie \sisi{$T_p$}{$T_P$} fait de $C$ un site ayant
\emph{suffisamment} de foncteurs fibres. Montrer que pour toute
topologie $T$ sur $C$, parmi les topologies $T'$ plus fines que $T$
qui ont assez de foncteurs fibres, il en est une moins fine: c'est
la topologie \sisi{$T_p$}{$T_P$}, où $P$ est la sous-catégorie strictement pleine
% correction SGA4-CORR-0011 / SGA-SRC-000398: SGA4-EIV-04-L3703-L3721-426DABA42C96.
de $\Pro(C)$, équivalente à $\Ppoint(C)$, décrite dans
\sisi{\eqhyperref[IV.eq6.8.3]{6.8.3}}{\eqref{IV.eq6.8.5.1}}
ci-dessous.
\end{enonce*}

\begin{enonce*}[remark]{Problème 6.5.3}\etiquette[6.5.3]{IV.prop6.5.3}
Caractériser les sous-catégories (strictement pleines, stables par
$𝒰$-limites projectives filtrantes\ldots) $P$ de $\Pro(C)$
qui peuvent se déduire d'une topologie sur $C$ comme image
essentielle de $\Ppoint(C^{\sim})$ (\hyperref[IV.subsubsec6.8.5]{6.8.5}).
\footnote[1]{\cf \hyperref[IV.exer9.1.8]{9.1.8} a) et e) pour des conditions
nécessaires, et \hyperref[IV.exer9.1.12]{9.1.12} b) pour des conditions
nécessaires et suffisantes dans le cas de $\Top(X)$, $X$ espace
sobre localement noethérien.}
\end{enonce*}

\setcounter{subsection}{5}
\subsection{}\etiquette[6.6]{IV.subsec6.6}
Soit
$$
f:E\longrightarrow E'
$$
un morphisme de $𝒰$-topos, on en déduit un foncteur
canonique
$$
\Ppoint(f)=(p\longmapsto f\circ p):\Ppoint(E)\longrightarrow
\Point(E'){\quoi}.
$$
Lorsqu'on a deux morphismes de topos composables
$$
E\xrightarrow{f} E'\sisi{\xrightarrow{g} E'{\quoi}}{\xrightarrow{g} E''{\quoi}},
% correction SGA4-CORR-0012 / SGA-SRC-000399: SGA4-EIV-04-L3735-L3755-02F1B01CEFC6.
$$
on\pageoriginale vérifie trivialement que l'on a
$$
\Ppoint(gf)=\Ppoint(g)\circ\Ppoint(f):\Ppoint(E)\longrightarrow
\Ppoint(E')\longrightarrow \Ppoint(E''){\quoi}.
$$
Cela précise donc la dépendance fonctorielle (sans abus de langage
pour une fois) de $\Ppoint(E)$ par rapport au topos $E$: si
$𝒱$ est un univers tel que $𝒰\in𝒱$, on
trouve un (véritable !) foncteur
$$
\Ppoint:(𝒱\text{-}𝒰\text{-}\text{top})\longrightarrow
(𝒱\text{-}\cat)
$$
de la catégorie des $𝒰$-topos qui sont éléments de
$𝒱$ dans la catégorie des catégories éléments de
$𝒱$.

\subsection{}\etiquette[6.7]{IV.subsec6.7}
Soit $E$ un $𝒰$-topos. Alors tout objet $F$ de $E$ définit
un contrafoncteur
$$
p\longmapsto F_p=p^{*}(F):\Ppoint(E)^{\circ}\longrightarrow
(𝒰\text{-}\Ens){\quoi},
$$
d'où pour $F$ variable un foncteur canonique
\begin{equation*}
E\longrightarrow
\Ppoint(E)^{\widehat{\phantom{,}}}=\Hhom(\Ppoint(E)^{\circ},
(𝒰\text{-}\Ens)){\quoi},\tag{6.7.1}
\end{equation*}\etiquette[6.7.1]{IV.eq6.7.1}
qui par définition (\hyperref[IV.prop6.5]{6.5}) est fidèle si et seulement si
$E$ possède suffisamment de points. (Le foncteur \eqhyperref[IV.eq6.7.1]{6.7.1} a
une certaine analogie formelle avec une transformation de
Fourier\ldots)

Soit $X$ un objet de $E$, qui définit donc un préfaisceau
$\widehat{X}$ sur
$$
C=\Ppoint(E){\quoi}.
$$
Nous pouvons donc définir la catégorie
$$
C_{/\widehat{X}}=\Ppoint(E)_{/\widehat{X}}
$$
des morphismes $p\rightarrow \widehat{X}$ dans la catégorie
$\widehat{C}$ des préfaisceaux sur $C$, \ie la catégorie des couples
$(p,\xi)$, où $p$ est un point de $E$ et $\xi$ un élément de
$\widehat{X}(p)=X_p$, fibre de $X$ en $p$. Ceci posé, je dis
qu'\emph{on a une équivalence de catégories canonique}
\begin{equation*}
C'=\Ppoint(E_{/X})\xrightarrow{\approx}C_{/\widehat{X}}\quav \mbox{
où }\quad C=\Ppoint(E){\quoi},\tag{6.7.2}
\end{equation*}\etiquette[6.7.2]{IV.eq6.7.2}
dont\pageoriginale la composée avec le foncteur d'inclusion
$C_{/\widehat{X}}\rightarrow C$ n'est autre que le foncteur
$$
\Ppoint(j_X):C'=\Ppoint(E_{/X})\longrightarrow C=\Ppoint(E)
$$
déduit par fonctorialité (\hyperref[IV.subsec6.6]{6.6}) du morphisme de
localisation \eqhyperref[IV.eq5.2.1]{5.2.1}
$$
j_X:E_{/X}\longrightarrow E{\quoi}.
$$
C'est simplement le cas particulier de \hyperref[IV.prop5.12]{5.12} obtenu en
faisant $E'=P$.

\begin{enonce*}{Corollaire 6.7.3}\etiquette[6.7.3]{IV.cor6.7.3}
Soit $E$ un topos. Si $E$ a suffisamment de points, il en est de
même de tout topos induit $E_{/X}$ $(X\in \ob E)$. Inversement. si
$(X_i)_{i\in I}$ est une famille d'objets de $E$ qui couvre l'objet
final, telle que pour tout $i\in I$, $E_{/X_i}$ ait suffisamment de
points, alors $E$ a suffisamment de points.
\end{enonce*}

Pour la première assertion, soit $f:X'\rightarrow X''$ un morphisme
dans $E_{/X}$ qui n'est pas un isomorphisme, prouvons qu'il existe
un point $q$ de $E_{/X}$ tel que $f_q$ ne soit pas un isomorphisme.
Comme $E$ a assez de points, il existe un point $p$ de $E$ tel que
$f_p:X'_p\rightarrow X''_p$ ne soit pas un isomorphisme. Cela
implique qu'il existe un $q\in X_p$ tel que le morphisme induit par
$f_p$ pour les fibres $X'_q$, $X''_q$ de $X'_p$, $X''_p$ sur $X_p$
en $q$ ne soit pas un isomorphisme. Mais en vertu de l'équivalence
\eqhyperref[IV.eq6.7.2]{6.7.2}, $q$ peut s'identifier à un point de $E_{/X}$, et
l'application $X'_q\rightarrow X''_q$ qu'on vient d'envisager n'est
autre que l'application sur les fibres en $q$ induite par
$f:X'\rightarrow X''$, d'où la conclusion.

Inversement, supposons que les $X_i$ couvrent l'objet final de $E$,
et que les $E_{/X_i}$ aient assez de points, prouvons qu'il en est
de même de $E$. Soit\pageoriginale $f:X'\rightarrow X''$ un
morphisme dans $E$ qui n'est pas un isomorphisme. Il existe donc un
$i\in I$ tel que $f_i:X'_i\rightarrow X''_i$ n'est pas un
isomorphisme, où l'indice $i$ indique la restriction à $X_i$. Il
existe donc un point $p_i$ de $E_{/X_i}$ tel que
$f_{ip_i}:X'_{ip_i}\rightarrow X''_{ip_i}$ ne soit pas un
isomorphisme. Désignant par $p$ l'image de $p_i$ dans $E$ par le
morphisme de localisation $E_{/X_i}\rightarrow E$, cela signifie que
$X'_p\rightarrow X''_p$ n'est pas un isomorphisme, ce qui prouve que
$E$ a assez de points.

\begin{enonce*}[remark]{Remarque 6.7.4}\etiquette[6.7.4]{IV.rem6.7.4}
L'argument précédent montre que si $(p_{\alpha})$ est une famille
conservative de points de $E$, alors la famille des points de
$E_{/X}$ qui sont au-dessus d'un des $p_{\alpha}$ (famille indexée
par l'ensemble somme des \sisi{$E_{p_{\alpha}}$}{$X_{p_{\alpha}}$}) est conservative. De
% correction SGA4-CORR-0013 / SGA-SRC-000400: SGA4-EIV-04-L3851-L3860-79416B604DA1.
même, si les $X_i$ couvrent l'objet final de $E$, et si pour tout
$i\in I$, $P_i$ est un ensemble conservatif de points de $E_{/X_i}$,
alors l'ensemble des points de $E$ \emph{images} des $p\in P_i$ pour
les morphismes de localisation $E_{/X_i}\rightarrow E$ ($i\in I$ et
$p$ variables) est conservatif.
\end{enonce*}

\setcounter{subsection}{7}
\subsection{}\etiquette[6.8]{IV.subsec6.8}
Soient $E$ un topos, $p$ un point de $E$. On appelle \emph{voisinage
du point $p$ du topos $E$} un couple $(X,u)$, où $X\in \ob E$ et
$u\in X_p$. En vertu de \eqhyperref[IV.eq6.7.2]{6.7.2}, on peut aussi interpréter $u$
comme un relèvement de $p$ en un point du topos induit $E_{/X}$.
Introduisant le foncteur fibre $\varphi_p$ associé au point $p$, on
peut aussi interpréter un voisinage du point $p$ comme un objet de
la catégorie $E_{/\varphi_p}$ (sous-catégorie pleine de
$(E^\circ)^{\widehat{\phantom{,}}}_{/\varphi_p}$ formée des objets
dont la source est dans $E\subset \widehat{E}$). Cette dernière
catégorie, \ie la catégorie des couples $(X,u)$ comme ci-dessus,
s'appellera comme de juste la \emph{catégorie des voisinages} du
point $p$ du topos $E$; on la note aussi $\Vvois(p)$. Comme
$\varphi_p$\pageoriginale est exact à gauche et que les limites
projectives finies sont représentables dans $E$, les limites
projectives finies sont représentables dans la catégorie
$\Vvois(p)$, et le foncteur canonique $\Vvois(p)\rightarrow E$ y
commute (\sisi{I\quad }{\ref{I.sec8.10.10}}). A fortiori, la catégorie opposée
$\Vvois(p)^{\circ}$ est filtrante. De plus, il est clair
(\hyperrefext[I][I.prop3.4]{3.4}) qu'on a un isomorphisme canonique, fonctoriel en
$F\in \ob E$
\begin{equation*}
\varphi_p(F)=F_p\simeq
\varinjlim_{\Vvois(p)^{\circ}}F(X){\quoi}.\tag{6.8.1}
\end{equation*}\etiquette[6.8.1]{IV.eq6.8.1}
En d'autres termes \emph{le foncteur $\varphi_p$ « fibre en le
point $p$ »  est isomorphe à la limite inductive filtrante des
foncteurs représentés dans le topos $E$ par les voisinages du point
$p$ de $E$.} On voit même que \emph{le foncteur fibre est
ind-représentable} (\hyperrefext[I][I.sec8]{8}), ce qui signifie aussi qu'il est
représentable par un pro-objet $(X_i)_{i\in I}$ de $E$, où $I$ est
un ensemble ordonné filtrant tel que $I\in𝒰$. En effet, en
vertu de \loccit, il revient au même de dire que la catégorie
$\Vvois(p)$ admet une petite sous-catégorie pleine cofinale. Or soit
$C$ une petite sous-catégorie pleine génératrice de $E$, je dis que
la sous-catégorie pleine $C_{/\varphi_{p}}$ de
$\Vvois(p)=E_{/\varphi_p}$, formée des voisinages $(X,u)$ de $p$
tels que $X\in \ob C$ (catégorie qui est évidemment petite) est
cofinale dans $\Vvois(p)$. Soit en effet $(X,u)$ un voisinage de
$p$, il faut trouver un voisinage $(X',u')$ de $p$ au-dessus de
$(X,u)$, avec $X'\in \ob C$. Or il existe une famille couvrante
$X'_i\rightarrow X$, avec les $X'_i$ dans $C$, d'où résulte que les
$X'_{ip}$ couvrent $X_p$, donc il existe un $X'=X'_{i_{\circ}}$ et
un $u'\in X'_p$ tels que $(X',u')$ soit \sisi{au-dessus de $(X,p)$}{au-dessus de $(X,u)$}, ce
% correction SGA4-CORR-0014 / SGA-SRC-000401: SGA4-EIV-04-L3897-L3908-71889DC02146.
qu'on avait affirmé.

\setcounter{subsubsection}{1}
\subsubsection{}\etiquette[6.8.2]{IV.subsubsec6.8.2}
Soient\pageoriginale $C$ un $𝒰$-site, et $p$ un point de
$C$, \ie un point du topos $C^{\sim}$. Désignons encore, par abus de
notations, par $\varphi_p$ la « restriction »  du foncteur fibre
$\varphi_p$ à $C$, \ie le composé $\varphi_p\circ \epsilon$; nous
écrirons aussi souvent, pour un objet $X$ de $C$
$$
X_p=\varphi_p(X)=\epsilon(X)_p{\quoi}.
$$
On appelle \emph{voisinage du point $p$ dans le site $C$} un couple
$(X,u)$, où $X\in \ob C$ et $u\in X_p=\varphi_p(X)$. Ces voisinages
forment encore une catégorie $C_{/\varphi_p}$, qu'on pourra noter
$\Vvois_C(p)$. Montrons que \emph{la catégorie $\Vvois_C(p)^{\circ}$
est encore filtrante, qu'elle admet une sous-catégorie pleine
cofinale $𝒰$-petite, et qu'on a un isomorphisme fonctoriel
en le faisceau $F$ sur $C$:}
\begin{equation*}
F_p\simeq \varinjlim_{\Vvois_C(p)^{\circ}} F(X){\quoi}.\tag{6.8.3}
\end{equation*}\etiquette[6.8.3]{IV.eq6.8.3}
On note d'abord, en reprenant l'argument de \eqhyperref[IV.eq6.8.1]{6.8.1}, que si
$C'$ est une sous-catégorie pleine de $C$ qui est topologiquement
génératrice (II 3.0.1), alors $\Vvois_{C'}(p)^{\circ}$ est cofinale
dans $\Vvois_C(p)^{\circ}$. Prenant $C'$ $𝒰$-petite, on est
donc ramené, pour établir l'assertion, au cas où $C\in𝒰$.
Dans ce cas $\widehat{C}$ est un $𝒰$-topos, et on a un
foncteur faisceau associé $\underline{a}:\widehat{C}\rightarrow
C^{\sim}$, qui permet de construire sur $\widehat{C}$ le foncteur
fibre composé $\varphi_p\underline{a}$. Appliquant \eqhyperref[IV.eq6.8.1]{6.8.1} au
topos $\widehat{C}$ et à la sous-catégorie génératrice
pleine $C$ de
celui-ci, on trouve que la catégorie
$C_{/\varphi_p\underline{a}}^{\circ},$ qui est manifestement isomorphe à
$\Vvois_C(p)^{\circ}$, est filtrante, et qu'on a un isomorphisme
fonctoriel en le préfaisceau $P$ sur $C$:
$$
\varphi_p\underline{a}(P)\simeq \varinjlim_X P(X){\quoi}.
$$
Si $F$ est un faisceau sur $C$, on a $F=\underline{a}F$, et appliquant
l'isomorphisme  précédent\pageoriginale à $F$ considéré comme
préfaisceau, on obtient \eqhyperref[IV.eq6.8.3]{6.8.3}. En même temps,
cet argument
établit la formule plus générale
\begin{equation*}
(\underline{a}P)_p\simeq \varinjlim_{\Vvois_C(p)^{\circ}}
P(X){\quoi},\tag{6.8.4}
\end{equation*}\etiquette[6.8.4]{IV.eq6.8.4}
isomorphisme fonctoriel en le préfaisceau $P$ arbitraire sur $C$, du
moins lorsque $C$ est petit. Le cas général s'en déduit, en
introduisant un univers $𝒱$ tel que $C\in𝒱$, en
considérant $\varphi_p$ sur $C$ comme un foncteur fibre à valeurs
dans $(𝒱\text{-}\Ens)$, et utilisant la compatibilité de la
formation du préfaisceau associé avec l'agrandissement des univers
(\hyperrefext[II][II.prop3.6]{3.6}).

\setcounter{subsubsection}{4}
\subsubsection{}\etiquette[6.8.5]{IV.subsubsec6.8.5}
On a associé à tout point $p$ du topos $C^{\sim}$ un pro-objet du
$𝒰$-site $C$, défini par le foncteur canonique
$\Vvois_C(p)\rightarrow C$, compte tenu du fait que la catégorie des
voisinages de $p$ dans $C$ est cofiltrante et admet une petite
sous-catégorie pleine cofinale. On vérifie aussitôt que pour $p$
variable, on obtient ainsi un foncteur
\begin{equation*}
\Ppoint(C^{\sim})\longrightarrow \Pro(C){\quoi}.\tag{6.8.5.1}
\end{equation*}\etiquette[6.8.5.1]{IV.eq6.8.5.1}
Ce foncteur est \emph{pleinement fidèle.} En effet, via les
foncteurs fibres associés, il s'interprète comme un foncteur
$\Fib(C)\rightarrow \Fib(C^{\sim})\rightarrow \Pro(C)^{\circ}$. Le
premier foncteur est une équivalence en vertu de \hyperref[IV.prop4.9.4]{4.9.4}
(rappelé pour les foncteurs fibres en \hyperref[IV.subsec6.3]{6.3}), et le
deuxième est pleinement fidèle d'après le sorite général
(\hyperrefext[I][I.sec8]{8}) des pro-objets.

\subsubsection{}\etiquette[6.8.6]{IV.subsubsec6.8.6}
Prenons le cas particulier où $C$ est $𝒰$-petite et munie de
la topologie chaotique, de sorte que \sisi{$C=\widehat{C}$}{$C^{\sim}=\widehat{C}$}. Je dis que
% correction SGA4-CORR-0015 / SGA-SRC-000402: SGA4-EIV-04-L3984-L3991-1A1626E9E64F.
dans ce cas le foncteur précédent est même une équivalence de
catégories
\begin{equation*}
\Ppoint(\widehat{C})\xrightarrow{\approx}\Pro(C){\quoi}.\tag{6.8.6.1}
\end{equation*}\etiquette[6.8.6.1]{IV.eq6.8.6.1}
En\pageoriginale effet, il reste à prouver que ce foncteur est
essentiellement surjectif, ce qui résulte du fait que les foncteurs
de la forme $P\mapsto P(X)$ sur $\widehat{C}$ sont des foncteurs
fibres, et du fait évident que toute limite inductive filtrante de
foncteurs fibres sur un topos est encore un foncteur fibre.

Identifiant $C$ à une sous-catégorie pleine de $\Pro(C)$
(\hyperrefext[I][I.sec8]{8}), le foncteur
$$
C\longrightarrow \Ppoint(\widehat{C})
$$
induit par un foncteur quasi-inverse de \sisi{(6.8.5.2)}{\eqhyperref[IV.eq6.8.6.1]{6.8.6.1}} est manifestement
isomorphe au foncteur canonique déjà envisagé dans
\eqhyperref[IV.eq4.6.2.2]{4.6.2.2},
dont le quasi-inverse envisagé peut être considéré comme le
prolongement canonique (\hyperrefext[I][I.sec8]{8}), compte tenu du fait que
$\Ppoint(\widehat{C})$ est stable par petites limites projectives.

\subsubsection{}\etiquette[6.8.7]{IV.subsubsec6.8.7}
Soit de façon générale $(X_i)_{i\in I}$ un pro-objet du
$𝒰$-site $C$, d'où sur $C^{\sim}$ un foncteur
\begin{equation*}
\varphi(F)=\varinjlim_I F(X_i){\quoi},\tag{6.8.7.1}
\end{equation*}\etiquette[6.8.7.1]{IV.eq6.8.7.1}
et il est naturel de se demander, vu (\hyperref[IV.subsubsec6.8.5]{6.8.5}),
quand ce foncteur est un foncteur fibre. On trouve que les
conditions suivantes sont équivalentes:
\begin{enumerate}
\renewcommand{\theenumi}{\roman{enumi}}
\renewcommand{\labelenumi}{(\theenumi)}
\item Le foncteur $\varphi$ précédent est un foncteur fibre sur
$C^{\sim}$.

\item La « restriction »  de $\varphi$ à $C$ est un foncteur fibre
sur le site $C$ (\hyperref[IV.subsec6.3]{6.3}).

\item Pour toute famille couvrante $Y_{\alpha}\rightarrow Y$ d'un
objet $Y$ de $C$, tout $i_{\circ}\in I$ et tout morphisme
$X_{i_{\circ}}\rightarrow Y$, il existe un $i\geq i_{\circ}$, un
$\alpha$ et un morphisme $X_i\rightarrow Y_{\alpha}$ rendant
commutatif le diagramme
\[
\xymatrix@C=2cm@R=1.25cm{
X_i \ar[d] \ar[r]& Y_{\alpha} \ar[d]\\
X_{{i}_{o}} \ar[r]     &  Y \quad . }
\]

En\pageoriginale fait, la condition (iii) exprime simplement le fait
que $\varphi| C$ transforme familles couvrantes en familles
couvrantes. Les implications (i) $\Rightarrow$ (ii) $\Rightarrow$
(iii) sont triviales, et il reste à prouver (iii) $\Rightarrow$ (i).
Comme le foncteur $\varphi$ est manifestement exact à gauche, il
reste à prouver (\hyperref[IV.prop4.9.4]{4.9.4}) qu'il transforme familles
couvrantes $F_{\alpha}\rightarrow F$ dans $C^{\sim}$ en familles
couvrantes, \ie que pour tout $i_{\circ}$ et tout $x_{\circ}\in
F(X_{i_{\circ}})$, il existe un $i\geq i_{\circ}$, un $\alpha$ et un
$y\in F_{\alpha}(X_i)$, tels que $x_{\circ}$ et $y$ aient même image
dans $F(X_i)$. Or comme $F_{\alpha}\rightarrow F$ est couvrante, il
existe une famille couvrante $Z_{\lambda}\rightarrow X_{i_{\circ}}$
de $X_{i_{\circ}}$, tel que pour tout $\lambda$ l'image inverse de
$x_{\circ}$ dans $F(Z_{\lambda})$ se remonte en un élément d'un
$F_{\alpha}(Z_{\lambda})$. Par hypothèse (iii), il existe un $i\geq
i_{\circ}$ et un $X_{i_{\circ}}$-morphisme $X_i\rightarrow
Z_{\lambda}$. Alors l'image de $x_{\circ}$ dans $F(X_i)$ se remonte
donc en un élément $y$ d'un $F_{\alpha}(X_i)$, \cqfd
\end{enumerate}

\begin{enonce*}[remark]{Exercice 6.9}\etiquette[6.9]{IV.exer6.9}
Soient $C$ un site $\in 𝒰$, $\varphi$ un foncteur fibre sur
$C$, proreprésenté par un objet $(X_i)_{i\in I}$ de $\Pro(C)$. Pour
tout indice $i\in I$, tout objet $Y$ de $C$, tout morphisme
$u:X_i\rightarrow Y$ et tout crible couvrant $R$ de $Y$, choisissons
un indice $j\geq i$ tel que le composé $X_j\rightarrow
X_i\xrightarrow{u}Y$ se factorise par $R$. Pour $i$ fixé, soit
$J(i)$ l'ensemble formé de $i$, et des $j$ obtenus pour $Y$, $u$,
$R$ variables; pour une partie $I'$ de $I$, soit de même $J(I')$ la
réunion des $J(i)$ pour $i\in I'$. D'autre part, si $I'$ est une
partie finie de $I$, soit $m(I')\in I$ un majorant de $I'$, et pour
une partie quelconque $I'$ de $I$, soit $M(I')$
\sisi{la réunion des $M(I'')$}{l'ensemble des $m(I'')$}, où $I''$ parcourt
l'ensemble des parties finies
% correction SGA4-CORR-0016 / SGA-SRC-000403: SGA4-EIV-04-L4069-L4072-C1D342741BAA.
% correction SGA4-CORR-0143 / SGA-SRC-000403: full-phrase scope completion.
de\pageoriginale $I'$. On suppose qu'on choisit la fonction
$I'\rightarrow m(I')$ de façon que pour $I'$ réduit à un élément
$i$, on ait $m(I')=i$, ce qui implique que pour toute partie $I'$ de
$I$, on a $I'\subset M(I')$. Considérons les itérés $(MJ)^n$ de
l'application $MJ$ de l'ensemble des parties de $I$ dans lui-même,
et soit, pour toute partie $I'$ de $I$, $P(I')$ la réunion des
$(MJ)^n(I')$.
\begin{enumerate}
\renewcommand{\theenumi}{\alph{enumi}}
\renewcommand{\labelenumi}{\theenumi)}
\item Montrer que pour toute partie $I'$ de $I$, $P(I')$ est une
partie filtrante de $I$ pour l'ordre induit, et que $I$ est réunion
filtrante croissante des $P(I')$, lorsque $I'$ parcourt l'ensemble
des parties \emph{finies} de $I$.

\item Montrer qu'il existe un cardinal $c\in 𝒰$, ne dépendant
que du cardinal $\card \text{fl}(C)=a$, tel que pour toute partie
finie $I'$ de $I$, on ait $\card(P(I'))\leq c$. (Si $a$ est infini,
on peut prendre $c=2^{a}$).

\item Montrer, en utilisant \hyperref[IV.subsubsec6.8.7]{6.8.7}, que pour toute partie $I'$ de $I$
le foncteur $\varphi_{I'}$ sur $C$ proreprésenté par le pro-objet
\sisi{$(X_i)_{i\in I'}$}{$(X_i)_{i\in P(I')}$} est un foncteur fibre sur $C$, et que le foncteur
% correction SGA4-CORR-0017 / SGA-SRC-000404: SGA4-EIV-04-L4092-L4097-3DFA2E263D3A.
fibre $\varphi$ est limite inductive filtrante des foncteurs
$\varphi_{I'}$, lorsque $I'$ parcourt l'ensemble des parties finies
de $I$.

\item En conclure qu'il existe une petite sous-catégorie pleine $\Phi$
de $\sisi{\Ffib(C)\approxeq \Ffib(C)}{\Ffib(C)\approxeq \Ffib(C^{\sim})}$, telle que tout objet de $\Ffib(C)$
% correction SGA4-CORR-0018 / SGA-SRC-000405: SGA4-EIV-04-L4099-L4103-62091FBF3762.
soit limite inductive filtrante d'objets de $\Phi$. (Si $a$ est
infini, on peut prendre $\Phi$ de cardinal $\leq 2^{2^{a}}$; si $a$
est fini, on peut prendre $\Phi=C$ bien sûr).
\end{enumerate}
\end{enonce*}

\begin{enonce*}[remark]{Exercice 6.10}\etiquette[6.10]{IV.exer6.10}
Soient $C$ un $\mathcal{U}$-site, $D$ une
$\sheaf{\mathcal{U}}$-catégorie où les petites $\varinjlim$
filtrantes sont représentables, $\Fais(C,D)$ la catégorie des
faisceaux sur $C$ \sisi{à valeurs dans $F$}{à valeurs dans $D$} (\hyperrefext[II][II.def6.1]{6.1}). Définir, en
% correction SGA4-CORR-0019 / SGA-SRC-000406: SGA4-EIV-04-L4107-L4115-0501749AFE31.
étendant \eqhyperref[IV.eq6.8.3]{6.8.3}, un « foncteur fibre »
\begin{equation*}
F\longrightarrow F_p:\Fais(C,D)\longrightarrow D{\quoi}.\tag{6.10.1}
\end{equation*}\etiquette[6.10.1]{IV.eq6.10.1}
Si dans $D$ les $\varprojlim$ finies sont représentables, et
commutent aux $\varinjlim$ filtrantes, alors le foncteur précédent
est exact à gauche.
\end{enonce*}

\section{Exemples de foncteurs fibres et de points de topos}\etiquette[7]{IV.sec7}\pageoriginale

\subsection{Cas de $\Top(X)$ pour un espace topologique
  $X$}\etiquette[7.1]{IV.subsec7.1}

Soit $x$ un point de $X$. Regardons l'ensemble ponctuel $\{x\}$
comme un espace topologique, et considérons l'inclusion
$$
i_x:\{x\}\longrightarrow X{\quoi}.
$$
En vertu de \hyperref[IV.subsubsec4.1.1]{4.1.1} il définit un morphisme de topos
(défini à isomorphisme unique près)
\begin{equation*}
\Top(i_x):\Top(\{x\})\longrightarrow \Top(X){\quoi}.\tag{7.1.1}
\end{equation*}\etiquette[7.1.1]{IV.eq7.1.1}
D'autre part, on peut identifier $P$ à $\Top(\{x\})$
(\hyperref[IV.subsec2.2]{2.2}), et on trouve ainsi un point $p_x$ de
$\Top(X)$:
\begin{equation*}
p_x:P\longrightarrow \Top(X){\quoi},\tag{7.1.2}
\end{equation*}\etiquette[7.1.2]{IV.eq7.1.2}
défini par $x$ à isomorphisme unique près. Le foncteur fibre associé
est donné par la formule bien connue $[TF]$, résultant d'ailleurs
immédiatement de \hyperrefext[I][I.prop5.1]{5.1}:
\begin{equation*}
p^{\ast}_x(F)=F_x=F_{p_x}=\varinjlim_{U\ni x}F(U){\quoi},\tag{7.1.3}
\end{equation*}\etiquette[7.1.3]{IV.eq7.1.3}
la limite étant prise suivant les voisinages ouverts $U$ de $x$ dans
$X$.

Comme on sait que $\Top(X)$ est isomorphe à $\Top(X_{\sob})$
(\hyperref[IV.subsubsec4.2.1]{4.2.1}), on voit plus généralement que tout point
de $X_{\sob}$, \ie toute partie fermée irréductible $Z$ de $X$,
définit un point de $\Top(X)$, ou encore un foncteur fibre sur
$\Top(X)$, qu'on notera encore $F\mapsto F_Z$. Revenant à sa
définition par la formule \eqhyperref[IV.eq7.1.3]{7.1.3} sur $X_{\sob}$, on trouve
\begin{equation*}
F_Z=\varinjlim_{U\cap Z\neq \phi} F(U){\quoi},\tag{7.1.4}
\end{equation*}\etiquette[7.1.4]{IV.eq7.1.4}
la\pageoriginale limite inductive étant prise suivant les ouverts
$U$ de $X$ qui rencontrent $Z$.

\setcounter{subsubsection}{4}
\subsubsection{}\etiquette[7.1.5]{IV.subsubsec7.1.5}
Il est bien connu que les foncteurs fibres $F\mapsto F_x(x\in X)$
forment déjà une famille conservative. C'est particulièrement
évident sur l'interprétation des faisceaux sur $X$ en termes
d'espaces étalés, la fibre de $F$ en $x$ n'étant alors autre que sa
fibre au sens des espaces fibrés sur $X$. On notera que l'ensemble
d'indices de la famille conservative de foncteurs fibres envisagée
est $X$, donc est $𝒰$-petite.

\subsubsection{}\etiquette[7.1.6]{IV.subsubsec7.1.6}
En vertu de \hyperref[IV.subsubsec4.2.3]{4.2.3}, la catégorie $\Ppoint(\Top(X))$
est équivalente à la catégorie associée à l'ensemble $X_{\sob}$
ordonné par la relation de spécialisation. En d'autres termes: tout
foncteur fibre sur $\Top(X)$ est isomorphe à un foncteur fibre
$F\mapsto F_Z$, où $Z$ est un élément uniquement déterminé de
$X_{\sob}$, \ie une partie fermée irréductible uniquement
déterminée
de $X$; d'autre part, si $Z$, $Z'\in X_{\sob}$, alors l'ensemble
$\Hom(F_Z,F_{Z'})$ est vide ou réduit à un point, ce dernier cas se
présentant si et seulement si $Z$ est une spécialisation de $Z'$
dans $X_{\sob}$, \ie si et seulement si $Z\subset Z'$ comme partie
de $X$. Lorsque $X$ lui-même est sobre, on peut dans ces énoncés
remplacer $X_{\sob}$ par $X$ lui-même.

On notera que le groupe des automorphismes d'un foncteur fibre de
$\Top(X)$ (opposé au groupe des automorphismes du point
correspondant de $\Top(X)$) est toujours réduit au groupe unité
(plus généralement, \hyperref[IV.subsubsec4.2.3]{4.2.3} nous apprend la même
chose pour le groupe des automorphismes de tout morphisme de topos
associés à des espaces topologiques). C'est là un phénomène très
spécial au cas particulier envisagé: voir l'exemple
\hyperref[IV.subsec7.2]{7.2}, ainsi\pageoriginale que VIII 7. Dans le cas des
topos associés à des « problèmes de modules »  (\cf par exemple
[10]), les groupes d'automorphismes des foncteurs fibres ont
d'ailleurs une interprétation remarquable, comme les groupes
d'automorphismes des structures algébriques (sur des corps
algébriquement clos) qu'on se propose de classifier.

\subsubsection{}\etiquette[7.1.7]{IV.subsubsec7.1.7}
On vient de voir comment on peut reconstituer un espace sobre $X$
(ou l'espace sobre $X_{\sob}$ associé à un espace topologique
quelconque), du moins en tant qu'ensemble ordonné par la relation de
spécialisation, à l'aide du topos $\Top(X)$ (qu'il suffit même de
connaître à équivalence près), comme l'ensemble des classes
d'isomorphie de « points »  de $\Top(X)$. D'après ce qui a été
dit dans \hyperref[IV.subsec2.1]{2.1}, on reconstitue également la topologie
de $X$, \ie la famille de ses ouverts, de la façon suivante: pour
tout sous-objet $U$ de l'objet final $e_E$ de $E$, soit $\Pt(U)$
l'ensemble des $x\in X$ tels que $U_x\neq\phi$ (qui s'identifie
d'ailleurs, en vertu de \eqhyperref[IV.eq6.7.2]{6.7.2}, à l'ensemble des classes
d'isomorphie de points du topos induit $E_{/U}$). Alors $U\mapsto
\Pt(U)$ est un isomorphisme d'ensembles ordonnés de l'ensemble des
sous-objets de $e_E$ sur l'ensemble des ouverts de $X$.

\subsubsection{}\etiquette[7.1.8]{IV.subsubsec7.1.8}
La détermination \hyperref[IV.subsubsec7.1.6]{7.1.6} de la catégorie des points
du topos $\Top(X)$ défini par un espace topologique $X$ conduit à
adopter la terminologie suivante pour les points $p$, $p'$ d'un
topos quelconque $E$: on dit que $p$ est une \emph{spécialisation}
de $p'$, ou que $p'$ est une \emph{générisation} de $p$, lorsqu'il
existe un morphisme de $p'$ dans $p$ (au sens de la catégorie
$\Ppoint(E)$ de \hyperref[IV.def6.1]{6.1}). Ces relations sont encore
transitives, mais on trouve facilement grâce à
\hyperref[IV.subsubsec6.8.6]{6.8.6} des exemples où $p$ et $p'$ sont chacun
spécialisation de\pageoriginale l'autre, sans que $p$ et $p'$ soient
isomorphes (ni a fortiori égaux). La catégorie des
générisations
d'un point $p$ du topos $E$, par quoi on entend la catégorie
$\Ppoint(E)_{/p}$, joue à certains égards un rôle analogue à celui
du passage au localisé d'un schéma en un point. Ce rôle sera
précisé
au \sisi{Chap. VI}{\hyperrefext[VI][VI.subsec8.4]{8.4}} avec la construction du topos localisé de $E$ en le
point $p$, dont la catégorie des points est canoniquement
équivalente à la catégorie des générisations de $p$.

\begin{enonce*}[remark]{Exercice 7.1.9}\etiquette[7.1.9]{IV.exer7.1.9}
Soit $E$ un $𝒰$-topos. Prouver que $E$ est équivalent à un
topos de la forme $\Top(X)$, où $X$ est un espace topologique $\in
𝒰$, si et seulement si il satisfait aux deux conditions
suivantes:
\begin{enumerate}
\renewcommand{\theenumi}{\alph{enumi}}
\renewcommand{\labelenumi}{\theenumi)}
\item La famille des sous-objets de l'objet final $e_E$ est
génératrice pour le topos $E$.

\item $E$ a suffisamment de points (\hyperref[IV.prop6.5]{6.5}). (Cette condition n'est pas
superflue: \cf \hyperref[IV.subsec7.4]{7.4}.)
\end{enumerate}

Comparer avec \hyperref[IV.exer7.8]{7.8} c).
\end{enonce*}

\begin{enonce*}[remark]{Exercice 7.1.10}\etiquette[7.1.10]{IV.exer7.1.10}
Soient $X$ un espace topologique muni d'un groupe d'opérateurs $G$
$(X,G\in𝒰)$ et soit $E=\Top(X,G)$ (\hyperref[IV.subsec2.3]{2.3}).
\begin{enumerate}
\renewcommand{\theenumi}{\alph{enumi}}
\renewcommand{\labelenumi}{\theenumi)}
\item Montrer que pour tout objet $(X',G)$ de $E$, le topos induit
$E_{/(X',G)}$ est canoniquement équivalent au topos $\Top(X',G)$
(comparer \hyperref[IV.subsec5.7]{5.7}).

\item Lorsque $G$ opère proprement et librement sur $X'$ (Bourbaki,
Top. Gen. Chap. \hyperrefext[III][III.sec4]{4}), montrer que le morphisme d'espaces à
opérateurs $(X',G)\rightarrow (X'/G,e)$ induit
(\hyperref[IV.subsubsec4.1.2]{4.1.2}) une équivalence de topos
$$
\Top(X',G)\longrightarrow \Top(X'/G){\quoi}.
$$

\item Conclure\pageoriginale de b) qu'il existe un objet $Z$ de $E$
tel que le topos induit $E_{/Z}$ soit équivalent à $\Top(X)$, le
foncteur de localisation $E\rightarrow E_{/Z}$ étant isomorphe au
foncteur « oubli des opérations de $G$ ». (Calquer le
raisonnement de \hyperref[IV.subsubsec5.8.3]{5.8.3})

\item Conclure de c) que tout foncteur fibre sur $E$ est induit, via
le foncteur « oubli des opérations de $G$ », par un foncteur
fibre sur $\Top(X)$, donc définissable par un point de $X_{\sob}$.

\item Déterminer la structure de la catégorie $\Ppoint(E)$ (à
équivalence près) en termes de l'espace à opérateurs $(X_{\sob},G)$.
En conclure en particulier que deux points de $X_{\sob}$ définissent
des foncteurs fibres sur $E$ isomorphes si et seulement si ils sont
conjugués sous l'action de $G$.
\end{enumerate}
\end{enonce*}

\begin{enonce*}[remark]{Exercice 7.1.11}\etiquette[7.1.11]{IV.exer7.1.11}
Soit $X\in 𝒰$ un espace topologique. Prouver que le
$V$-topos $\TOP(X)$ (\hyperref[IV.subsec2.5]{2.5}) a suffisamment de points.
Plus précisément, pour tout objet $X'$ de $\TOP(X)$, et tout $x'\in
X'$, définir un foncteur fibre $F\mapsto F_{X',x'}$ sur $\TOP(X)$,
ne dépendant que du « germe » d'espace $(X',x')$ au-dessus de
$X$, et prouver que cette famille de foncteurs fibres est
conservative (et indexée par un ensemble d'indices
$𝒱$-petit). Montrer, en utilisant un système projectif
filtrant convenable $(X'_i,x'_i)_{i\in I}$, avec $I$ non
$𝒰$-petit, que l'on n'obtient pas de cette façon tous les
foncteurs fibres du $𝒱$-topos $\TOP(X)$. Montrer que
l'ensemble des classes d'isomorphie de tels foncteurs fibres n'est
pas de cardinal $\in 𝒱$.
\end{enonce*}

\subsection{Points d'un topos classifiant
  $B_G$}\etiquette[7.2]{IV.subsec7.2}\pageoriginale

Soient $E$ un $𝒰$-topos, $G$ un Groupe de $E$, d'où un topos
classifiant $B_G$ (\hyperref[IV.subsec2.4]{2.4}). Si $e$ est le Groupe
ponctuel, les morphismes $e\rightarrow G\rightarrow e$ définissent
(\hyperref[IV.subsec4.5]{4.5}) des morphismes de topos (compte tenu que
$B_e\approx E$)
\begin{equation*}
E\longrightarrow B_G \longrightarrow E{\quoi},\tag{7.2.1}
\end{equation*}\etiquette[7.2.1]{IV.eq7.2.1}
dont le composé est isomorphe à $\id_E$, d'où par passage aux
catégories de points des foncteurs
\begin{equation*}
\Ppoint(E)\longrightarrow \Ppoint(B_G)\longrightarrow
\Ppoint(E){\quoi},\tag{7.2.2}
\end{equation*}\etiquette[7.2.2]{IV.eq7.2.2}
dont le composé est isomorphe à l'identité. Utilisons le fait que le
premier morphisme de topos \eqhyperref[IV.eq7.2.1]{7.2.1} s'identifie à un morphisme
de localisation relativement à l'objet $E_G=G$ de $B_G$
(\hyperref[IV.subsec5.8]{5.8}), d'où résulte en vertu de \eqhyperref[IV.eq6.7.2]{6.7.2} que le
premier foncteur \eqhyperref[IV.eq7.2.2]{7.2.2} s'identifie au foncteur naturel « d'inclusion »
\begin{equation*}
\Ppoint(B_G)_{/\widehat{E}_G}\longrightarrow
\Ppoint(B_G){\quoi},\tag{7.2.3}
\end{equation*}\etiquette[7.2.3]{IV.eq7.2.3}
où $\widehat{E}_G$ désigne le préfaisceau $p'\mapsto (E_G)_{p'}$ sur
$\Ppoint(B_G)$. Comme $E_G$ couvre évidemment l'objet final de
$B_G$, ses fibres $(E_G)_p$ sont non vides, d'où résulte que
\eqhyperref[IV.eq7.2.3]{7.2.3} est essentiellement surjectif, ce qui signifie que
\emph{tout foncteur fibre sur $B_G$ est induit (à isomorphisme près)
par un foncteur fibre de $E$, via le foncteur « oubli des
opérations de $G$ » $B_G\rightarrow E$.} On en conclut en
particulier que \emph{si $E$ a suffisamment de points (\resp une
petite famille conservative de points)} il en est de même de $B_G$.

\begin{enonce*}[remark]{Remarque 7.2.4}\etiquette[7.2.4]{IV.rem7.2.4}
On peut aller plus loin et déterminer (à équivalence près) la
structure de la catégorie $C'=\Ppoint(B_G)$ en termes de la
catégorie $C=\Ppoint(E)$ et du préfaisceau en groupes
$$
\widehat{G}:p\longmapsto G_p: C^{\circ}\longrightarrow
(\text{Groupes})
$$
sur\pageoriginale celle-ci. Définissons en effet une nouvelle
catégorie $C_1$, ayant mêmes objets que $C$, et telle que pour deux
objets $p$, $q$ de $C$, on ait
$$
\Hom_{C_1}(p,q)=\Hom_C(p,q)\times \widehat{G}(p){\quoi},
$$
la composition des flèches dans $C_1$ étant induite par celle de
$C$, et par la loi de groupe du préfaisceau $\widehat{G}$. La donnée
d'un foncteur $\varphi_1:C_1\rightarrow D$ dans une catégorie
quelconque $D$ équivaut alors à la donnée d'un foncteur
$\varphi:C\rightarrow D$, muni d'une opération du préfaisceau
$\widehat{G}$ sur ce dernier (\ie la donnée, pour tout $p\in \ob C$,
d'une opération de $\widehat{G}(p)$ sur $\varphi(p)$, satisfaisant à
une condition de fonctorialité évidente pour $p$ variable).
Utilisant cette observation, on définit un foncteur canonique
\begin{equation*}
\varphi_1:C_1\longrightarrow C'{\quoi},\tag{7.2.4.1}
\end{equation*}\etiquette[7.2.4.1]{IV.eq7.2.4.1}
correspondant au foncteur $\varphi:C\rightarrow C'$ de
\eqhyperref[IV.eq7.2.2]{7.2.2}, associant à tout point $p$ de $E$ le point
$\varphi(p)$ induit sur $B_G$, avec les opérations naturelles de
$\widehat{G}(p)=G_p$ sur ce dernier, provenant des opérations de
$G_p$ sur les $X_p$ lorsque $X$ parcourt $B_G$. On laisse au lecteur
le soin de prouver que \eqhyperref[IV.eq7.2.4.1]{7.2.4.1} est une équivalence de
catégories, en utilisant la structure \eqhyperref[IV.eq7.2.3]{7.2.3} du premier
foncteur de \eqhyperref[IV.eq7.2.2]{7.2.2}, et en notant que le foncteur naturel
d'inclusion $C\rightarrow C_1$ admet une structure analogue,
relativement à l'image inverse $P_1$ du préfaisceau
$P'=\widehat{E}_G$ sur $C'$.
\end{enonce*}

\setcounter{subsubsection}{4}
\subsubsection{}\etiquette[7.2.5]{IV.subsubsec7.2.5}
On\pageoriginale conclut en particulier que les \emph{classes
d'isomorphie de points de $B_G$ correspondent exactement aux classes
d'isomorphie de points de $E$: tout point du topos classifiant $E_G$
est induit, à isomorphisme} non unique \emph{près, par un point de
$E$}. En particulier, lorsque $E$ est le topos ponctuel $P$, donc
que $G$ est un groupe ordinaire, alors la catégorie $\Ppoint(B_G)$
est un groupoïde connexe à groupe fondamental $G$: tout foncteur
fibre sur $B_G$ est isomorphe (de façon non canonique) au
foncteur $\omega$ « oubli des opérations de $G$ », et le monoïde
des endomorphismes de ce dernier foncteur est le groupe $G$ (donc
tout endomorphisme de $\omega$ est un automorphisme).

\begin{enonce*}[remark]{Exercice 7.2.6}\etiquette[7.2.6]{IV.exer7.2.6}
\begin{enumerate}
\renewcommand{\theenumi}{\alph{enumi}}
\renewcommand{\labelenumi}{\theenumi)}
\item Soit (Tors) la catégorie des couples $(\Gamma,P)\in 𝒰$,
avec $\Gamma$ un groupe et $P$ un torseur à droite sous $\Gamma$. Le
foncteur $(\Gamma,P)\mapsto \Gamma$
\begin{equation*}
(\Tors)\longrightarrow (\Groupes)\tag{7.2.6.1}
\end{equation*}\etiquette[7.2.6.1]{IV.eq7.2.6.1}
est un foncteur cofibrant. Avec les notations de \hyperref[IV.rem7.2.4]{7.2.4},
considérons le foncteur
$$
\widehat{G}:C^{\circ}\longrightarrow (\Groupes){\quoi},
$$
et soit $D'$ la catégorie cofibrée sur $C^{\circ}$ image inverse de
la catégorie cofibrée \eqhyperref[IV.eq7.2.6.1]{7.2.6.1}, $D$ la
catégorie fibrée
correspondante sur $C$, ayant même catégories fibres que $D'$, de
sorte que pour $p\in \ob C$, on ait
$$
D_p=\mbox{ catégorie des $G_p$-torseurs à droite}{\quoi}.
$$
Construire des foncteurs canoniques
\begin{equation*}
D\longrightarrow C'=\Ppoint(B_G)\quav C'\longrightarrow
D{\quoi},\tag{7.2.6.2}
\end{equation*}\etiquette[7.2.6.2]{IV.eq7.2.6.2}
\emph{quasi-inverses l'un de l'autre.} En particulier, en conclure
que la catégorie des points de $B_G$ au-dessus d'un point donné $p$
de $E$ est canoniquement équivalente à la catégorie des torseurs à
droite sous le groupe $G_p$.

\item Soit\pageoriginale $𝒢=(G_i)_{i\in I}$ un système
projectif strict de Groupes du $𝒰$-topos $E$ ($I\in
𝒰$ un ensemble préordonné filtrant), d'où un topos
classifiant \sisi{$B_{𝒰}$}{$B_{𝒢}$}, en calquant la définition de
% correction SGA4-CORR-0020 / SGA-SRC-000407: SGA4-EIV-04-L4431-L4436-FC403EE16D11.
\hyperref[IV.subsubsec2.7.1]{2.7.1}. Déterminer la catégorie des points de
$B_{𝒢}$, en calquant a). (NB. On remplacera les catégories
$(\Tors)$ et $(\Groupes)$ par les catégories de systèmes projectifs
indexés par $I$ de ces catégories.) En conclure que si $I$ admet un
ensemble cofinal dénombrable, alors tout foncteur fibre sur
$B_{𝒢}$ est isomorphe à un foncteur induit par un foncteur
fibre de $E$ (via le foncteur « oubli des opérations de $G$ »),
ce dernier étant déterminé à isomorphisme non unique près.

\item Prenant pour $E$ le topos ponctuel, de sorte que $𝒢$ est
un pro-groupe ordinaire, montrer que la conclusion de b) est
également valable si $I$ est quelconque, mais en revanche
$𝒢$ est profini (\ie les $G_i$ sont finis): tout foncteur
fibre est isomorphe au foncteur « oubli des opérations de $G$ ».

\item Prenant toujours le cas où $E$ est le topos ponctuel, donner un
exemple d'un foncteur fibre sur $B_{𝒢}$, pour un système
projectif convenable $𝒢=(G_i)_{i\in I}$, qui n'est pas
isomorphe au foncteur « oubli des opérations de $G$ ».

\item Considérons $𝒢=(G_i)$ comme un Groupe du topos
$\widehat{I}$, et soit $𝒯$ une \emph{gerbe} sur $I$ de
\emph{lien} $G$ [\hyperref[IV.sec3]{3}]. Montrer comment on
peut « tordre » le topos classifiant $B_G$ à l'aide de la gerbe $T$, pour
obtenir un topos $B^T_{𝒢}$, limite inductive de
sous-catégories $B^T(i)$ équivalentes (non canoniquement) aux topos
classifiants $B_{G_i}$. Montrer que le topos $B^T_{𝒢}$ admet
un foncteur fibre si et seulement si la gerbe $T$ est
« neutre », en établissant une\pageoriginale équivalence de catégories
entre la catégorie des foncteurs fibres de $B^T_{𝒢}$ et la
catégorie des sections de $T$ sur $I$. En conclure, si les $G_i$
sont commutatifs, de sorte que la classification des gerbes sur $I$
de lien $𝒢$ se fait par
$\displaystyle{\varprojlim_I}^{(2)}G_i=H^{2}(I,𝒢)$
(\loccit), un exemple d'un topos (non « vide »
(\hyperref[IV.subsec2.2]{2.2})) de la forme $B^T_{𝒢}$ qui n'a pas de
points. (Prendre un exemple où
$\displaystyle{\varprojlim_I}^{(2)}G_i\neq 0$.)
\end{enumerate}
\end{enonce*}

\subsection{Points des topos $\widehat{C}$; exemples de
$𝒰$-topos $\widehat{C}$ dont la catégorie
des points ne soit pas équivalente à une petite
catégorie}\etiquette[7.3]{IV.subsec7.3}

Soient $E$ un $𝒰$-topos, $(\varphi_i)_{i\in I}$ une famille
filtrante de foncteurs fibres sur $E$. Il résulte des propriétés
d'exactitude des foncteurs $\varinjlim$ (\sisi{I\quad et}{\hyperrefext[I][I.prop2.8]{2.8}}) que le foncteur
$\varphi=\varinjlim\varphi_i$ est également un foncteur fibre:
\emph{toute limite inductive filtrante de foncteurs fibres est un
foncteur fibre.} Par suite, la catégorie $\Ffib(E)$ admet des
$𝒰$-limites inductives filtrantes (et le foncteur
d'inclusion $\Ffib(E)\rightarrow \Hhom(E,(𝒰\text{-}\Ens))$ y
commute); en d'autres termes, la catégorie $\Ppoint(E)$ admet des
$𝒰$-limites projectives filtrantes. Ce fait a déjà été
utilisé dans exercice \sisi{\hyperref[IV.exer7.1.10]{7.1.10}}{\hyperref[IV.exer7.1.11]{7.1.11}} pour donner un exemple de
« grosses »  catégories de foncteurs fibres.

On peut construire des exemples nettement plus simples, avec des
topos de la forme $E=\widehat{C}$, $C\in𝒰$, en utilisant
\eqhyperref[IV.eq6.8.6.1]{6.8.6.1}. Ceci nous donne aussitôt des exemples où
$\Ffib(\widehat{C})$ n'est pas équivalente à une catégorie
$\in𝒰$ \ie où le cardinal de l'ensemble des classes
d'isomorphie d'objets de cette $𝒰$-catégorie n'est pas
$\in𝒰$. Ceci signifie que la catégorie $\Pro-(C)$ n'est pas
équivalente à une catégorie $\in 𝒰$. Il suffit par exemple
de prendre pour $D=C^{\circ}$ la catégorie des ensembles finis de la
forme $[0,n]$, où $n\geq 0$ est un\pageoriginale entier, auquel cas
$\Ind(D)\simeq \Pro(C)^{\circ}$ est équivalente à la catégorie des
ensembles non vides. Le topos $E$ est dans ce cas la catégorie bien
connue des \emph{ensembles cosimpliciaux}. On pourrait aussi prendre
pour $C$ la catégorie des ensembles finis non vides, on trouve que
la catégorie des points sur le topos $\widehat{C}$ des ensembles
simpliciaux est équivalente à la catégorie des ensembles profinis,
ou encore à la catégorie des espaces compacts totalement
discontinus.

\subsection{Topos non vides sans points}\etiquette[7.4]{IV.subsec7.4}

L'exemple suivant est dû à $P$. Deligne. (Pour un autre exemple, \cf
\hyperref[IV.exer7.2.6]{7.2.6} d).) On prend un espace compact $K$ muni d'une
mesure $\mu$, et l'ensemble ordonné $U$ des parties mesurables de
$K$ à ensemble de mesure nulle près. On fait de $U$ une catégorie
$\underline{U}$ telle que $\ob \underline{U}=U$, les morphismes de
$\underline{U}$ étant les « morphismes d'inclusion »  entre éléments
de $U$. On fait de $\underline{U}$ un site en prenant la prétopologie
pour laquelle $\Cov(E)$ (pour $E\in U$) est formé des familles
\emph{dénombrables} d'éléments $E_i$ de $E$ majorés par $E$, telles
que $E$ soit la réunion des $E_i$ à ensemble de mesure nulle près.
On en déduit un Topos $\Top(\mu)=\underline{U}^{\sim}$, admettant
l'ensemble des sous-objets de l'objet final comme famille
génératrice (lequel topos semble avoir échappé à l'attention des
probabilistes). Ce topos est un « topos vide »
(\hyperref[IV.subsec2.2]{2.2}) si et seulement si $\mu=0$. D'autre part, la
catégorie des points de ce topos est équivalent à la catégorie
discrète définie par l'ensemble des points $x\in K$ tels que
$\mu(\{x\})\neq 0$ (démonstration au lecteur). Elle est donc vide si
$K$ n'admet pas de tels points, par exemple si $K$ est le segment
unité de la droite, avec la mesure induite par la mesure de
Lebesgue.

\begin{enonce*}[remark]{Exercice 7.5}\etiquette[7.5]{IV.exer7.5}
(\emph{catégories\pageoriginale \sisi{Karoubiennes}{karoubiennes} et morphismes de
topos})\footnote[1]{Comparer \hyperrefext[I][I.exer8.7.8]{8.7.8}.}
\begin{enumerate}
\renewcommand{\theenumi}{\alph{enumi}}
\renewcommand{\labelenumi}{\theenumi)}
\item Soient $C$ une catégorie, $X$ un objet de $C$, $p$ un
endomorphisme de $X$. On dit que $p$ est un \emph{projecteur} si
$p^2=p$. Prouver que si $p$ est un projecteur, pour que
$\Ker(\id_X,p)$ soit représentable, il faut et il suffit que
$\Coker(\id_X,p)$ le soit, et que les deux objets de $C$ ainsi
obtenus sont canoniquement isomorphes. On dira alors que le
projecteur $p$ \emph{admet une image,} et $\Ker(\id_X,p)\simeq
\Coker(\id_X,p)$ est appelé l'\emph{image du projecteur $p$}; on
l'identifie suivant le contexte à un sous-objet ou à un objet
quotient de $C$. Un objet isomorphe à un $\text{Im\,} p$, pour un
projecteur convenable $p$ dans $X$, est appelé un \emph{facteur
direct de l'objet $X$ de $C$.} On dit que $C$ est une
\emph{catégorie avec facteurs directs}, ou une \emph{catégorie
\sisi{Karoubienne}{karoubienne},} si tout projecteur dans un objet de $C$ admet une
image. Si $F:C\rightarrow C'$ est un foncteur qui commute aux noyaux
ou aux conoyaux, alors $F$ transforme un projecteur admettant une
image en un projecteur admettant une image\nde{Pas besoin de conditions sur
$F$ ici, tous les foncteurs préservent les projecteurs admettant des images.}.

\item Montrer que pour toute catégorie $C$, on peut trouver un
foncteur $\varphi:C\rightarrow \kar (C)$ de $C$ dans une catégorie
karoubienne (déterminée à équivalence près), tel que pour toute
catégorie karoubienne $C'$, le foncteur
$$
f\mapsto f\circ \varphi:\Hhom(\kar(C),C')\longrightarrow \Hhom(C,C')
$$
soit une équivalence de catégories; on appellera $\kar(C)$
l'\emph{enveloppe de Karoubi} de la catégorie $C$. (Hint: prendre
pour $\ob \kar (C)$ l'ensemble des couples $(X,p)$, avec $X\in \ob
C$ et $p$ un projecteur dans $X$, et pour $\Hom((X,p),(Y,q))$ la
partie de $\Hom(X,Y)$ formée des $f:X\rightarrow Y$ tels que
$f=qfp.$) Montrer que $\varphi$ est pleinement fidèle.

\item Soit $F:E\rightarrow E'$ un foncteur d'une catégorie dans une
autre. Montrer que si $F$ commute à un certain type de limites
inductives ou projectives,\pageoriginale il en est de même de tout
facteur direct de $F$. En particulier, si $E$ et $E'$ sont des
$𝒰$-topos et si $F$ est un foncteur image inverse pour un
morphisme de topos $E'\rightarrow E$, alors il en est de même de
tout facteur direct de $F$; par suite, la catégorie $\Hhomtop(E',E)$
est karoubienne, et en particulier la catégorie $\Ppoint(E)$ est
karoubienne.

\item Montrer que pour toute $𝒰$-catégorie $C$, la catégorie
$\Ind(C)$ est karoubienne (où $\Ind(C)$ est formée avec les systèmes
inductifs de $C$ indexés par un ensemble préordonné filtrant $I\in
𝒰$). (Si $C\in 𝒰$, utiliser par exemple le fait
(\hyperref[IV.subsec7.3]{7.3}) que $\Ind(C)$ est équivalente à la catégorie
$\Ffib((C^{\circ})^{\widehat{\phantom{,}}})=\Ppoint((C^{\circ})^{\widehat{\phantom{,}}})^{\circ}.$)
En conclure une autre construction de $\kar(C)$ comme sous-catégorie
pleine de $\Ind(C)$ formée des images dans $\Ind(C)$ des projecteurs
d'objets $X$ de $C$.
\end{enumerate}
\end{enonce*}

\begin{enonce*}[remark]{Exercice 7.6}\etiquette[7.6]{IV.exer7.6}
(\emph{Morphismes essentiels de topos, points essentiels}).
\begin{enumerate}
\renewcommand{\theenumi}{\alph{enumi}}
\renewcommand{\labelenumi}{\theenumi)}
\item Soit $f:E\rightarrow F$ un morphisme de topos. Montrer que les
conditions suivantes sont équivalentes:
\begin{itemize}
\item[i)] $f_{!}$ existe, \ie $f^{\ast}$ admet un adjoint à gauche.

\item[ii)] $f^{\ast}$ commute aux $𝒰$-limites projectives.

\item[${\rm ii')}$] $f^{\ast}$ commute aux $𝒰$-produits.
\end{itemize}
On dira alors que $f$ est un \emph{morphisme essentiel} du topos $E$
\sisi{dans le topos $E'$.}{dans le topos $F$.}
% correction SGA4-CORR-0021 / SGA-SRC-000408: SGA4-EIV-04-L4603-L4613-3A8D0250B4A0.

Montrer que si $f$ satisfait à ces conditions, il en est de même de
tout facteur direct de $f$. (Utiliser \hyperref[IV.rem1.8]{1.8} et
\hyperref[IV.exer7.5]{7.5} c).)

\item Soit $E$ un topos. On appelle \emph{point essentiel du topos}
$E$ tout point $p:P\rightarrow E$ de $E$ tel que $p_{!}$ existe.
Montrer que si $E=\Top(X)$, $X$ espace topologique, et si $p$ est le
point de $E$ défini par un $x\in X$, alors $p$ est
essentiel\pageoriginale si et seulement si $x$ admet un plus petit
voisinage ouvert $U$ (ce qui signifie, si les points de $X$ sont
fermés, que $x$ est un \emph{point isolé} de $X$). Pour qu'un
morphisme de topos $f:E\rightarrow F$ soit essentiel (a)), il faut
que $f$ transforme points essentiels en points essentiels, et cette
condition est également suffisante lorsque $E$ admet suffisamment de
points essentiels (\ie si la famille de ces points est conservative)
(\cf d)).

\item Pour qu'un point $p$ du topos $E$ soit essentiel, il faut et il
suffit que le foncteur fibre associé soit représentable par un objet
$X$ de $E$. Pour que le foncteur covariant $\varphi:E\rightarrow
(\Ens)$ représenté par un objet donné $X$ de $E$ soit un foncteur
fibre (\ie définisse un point, nécessairement essentiel, de $E$), il
faut et il suffit que $X$ soit connexe non vide
(\hyperref[IV.subsubsec4.3.5]{4.3.5}), et projectif (\ie que le foncteur
$\varphi$ transforme épimorphismes en épimorphismes). (Hint: montrer
d'abord que pour que $\varphi$ commute aux sommes, il faut et il
suffit que $X$ soit connexe non vide, puis utiliser le critère
\hyperref[IV.prop4.9.4]{4.9.4}.) En conclure une équivalence entre la catégorie
$\Ppointess(E)$ et la sous-catégorie pleine $P$ de $E$ formée des
objets connexes non vides projectifs.

\item Montrer que la topologie de $P$ induite par celle de $E$ est la
topologie chaotique. En conclure l'équivalence des conditions
suivantes sur $E$ (due à J.E ROOS [12 c), prop. 1]): (i) La famille
des points essentiels de $E$ est conservative; (ii) La
sous-catégorie pleine $P$ de $E$ formée des objets connexes-non
vides projectifs est génératrice; (iii)~$E$ est équivalent à un
topos de la forme $\widehat{C}$, où $C$ est une catégorie
équivalente à une catégorie $\in 𝒰$ (ou, si on préfère,
$C\in 𝒰$). (Utiliser $c$) et le fait que si $P$ est
génératrice, le foncteur naturel $E\rightarrow P^{\sim}$ est une
équivalence de\pageoriginale catégories.)

\item La catégorie des points essentiels d'un topos $E$ est
karoubienne (\cf (\hyperref[IV.exer7.5]{7.5} c)). Soit $C$ une catégorie
équivalente à une catégorie $\in 𝒰$. Prouver que le foncteur
pleinement fidèle canonique (\hyperref[IV.subsubsec4.6.2]{4.6.2})
$$
C\longrightarrow \Ppoint(\widehat{C})
$$
se factorise par un foncteur
$$
C\longrightarrow \Ppointess(\widehat{C}){\quoi},
$$
qui fait de $\Ppointess(\widehat{C})$ une \emph{enveloppe de
Karoubi} (\hyperref[IV.exer7.5]{7.5} b)) \emph{de} $C$. En particulier, ce
foncteur est une équivalence de catégories si et seulement si la
catégorie $C$ est karoubienne. (Hint: utilisant $c$), prouver que
\sisi{tout point isolé}{tout point essentiel} de $\widehat{C}$ est isomorphe à un facteur direct
% correction SGA4-CORR-0022 / SGA-SRC-000409: SGA4-EIV-04-L4668-L4674-C44E048C17C9.
d'un $X\in \ob C$.)

\item Soient $C$ et $C'$ deux catégories équivalentes à des catégories
$\in 𝒰$. Prouver que le foncteur pleinement fidèle canonique
(\hyperref[IV.subsubsec4.6.2]{4.6.2})
$$
\Hhom(C,C')\longrightarrow \Hhomtop(\widehat{C},\widehat{C}')
$$
prend ses valeurs dans la sous-catégorie pleine
$\Hhomtopess(\widehat{C},\widehat{C}')$ des morphismes essentiels
\emph{de topos} (a)), et que le foncteur induit
$$
\Hhom(C,C')\longrightarrow \Hhomtopess(\widehat{C},\widehat{C}')
$$
est une équivalence de catégories si et seulement si $C$ est vide ou
$C'$ est karoubienne. (Utiliser $g$) ci-dessous.)

\item Soient $C$ une catégorie équivalente à une catégorie $\in
𝒰$, et $E$ un topos. Définir une équivalence entre la
catégorie $\Hhomtopess(\widehat{C},E)$ des morphismes essentiels de
topos $\widehat{C}\rightarrow E$, et la catégorie $\Hhom(C,
\Ppointess(E))$. (Utiliser $b$)\pageoriginale et le fait que
$\widehat{C}$ a suffisamment de points essentiels).

\item Soit $C$ une catégorie finie. Prouver que tout point de
$\widehat{C}$ est essentiel, et que $\Ppoint(\widehat{C})$ est
équivalente à une catégorie finie. (Noter que l'enveloppe de Karoubi
de $C$ est finie, et utilisant \eqhyperref[IV.eq6.8.6.1]{6.8.6.1}, se ramener à prouver
que si $C$ est une catégorie karoubienne finie, alors le foncteur
canonique $C\rightarrow \Pro(C)$ est une équivalence de catégories).
Utilisant b), en conclure que \emph{tout} morphisme de topos
$\widehat{C}'\rightarrow \widehat{C}$ est essentiel, et
$\Hhom(C,C')\rightarrow \sisi{\Hom(\widehat{C},\widehat{C}')}{\Hhomtop(\widehat{C},\widehat{C}')}$ est une
% correction SGA4-CORR-0023 / SGA-SRC-000410: SGA4-EIV-04-L4704-L4707-87F94451A0FC.
équivalence si et seulement si $C$ est vide ou $C'$ karoubienne.

\item Soit $X$ un espace topologique, $f:\Top(X)\rightarrow P$ le
morphisme de topos déduit de l'application continue de $X$ dans
l'espace ponctuel. Montrer que $f$ est essentiel si et seulement si
l'espace $X$ est \emph{localement connexe}, \ie satisfait la
condition suivante: pour tout $x\in X$ et tout voisinage ouvert $U$
de $x$, la composante connexe de $x$ dans $U$ est un voisinage de
$x$, ou encore: pour tout ouvert $U$ de $X$, les composantes
connexes de $U$ sont ouvertes. (\cf \hyperref[IV.exer8.7]{8.7} b) pour
généralisation).
\end{enumerate}
\end{enonce*}

\begin{enonce*}[remark]{Exercice 7.7}\etiquette[7.7]{IV.exer7.7}
(\emph{Points inhabituels d'un topos classifiant})\footnote[1]{\cf
aussi \hyperref[IV.exer7.2.6]{7.2.6} d).}.
\begin{enumerate}
\renewcommand{\theenumi}{\alph{enumi}}
\renewcommand{\labelenumi}{\theenumi)}
\item Soit $G$ un monoïde. Pour que $G$ contienne un élément $g$ tel
que $g\neq e$, $g^2=g$, il faut et il suffit que le topos
classifiant $B_G$ admette un point essentiel qui n'est pas isomorphe
au point banal (correspondant au foncteur fibre « oubli des
opérations de $G$ »). (Utiliser \hyperref[IV.exer7.6]{7.6} e)).

\item Soit $G$ le monoïde additif des entiers $\geq 0$. Montrer que
tout point essentiel du topos classifiant $B_G$ est isomorphe au
point banal. Construire un point de $B_G$ qui n'est pas isomorphe au
point banal. Montrer que la catégorie $\Ppoint(B_G)$ n'est pas
équivalente à une catégorie $\in 𝒰$.
\end{enumerate}
\end{enonce*}

\begin{enonce*}[remark]{Exercice 7.8}\etiquette[7.8]{IV.exer7.8}
(\emph{Topologie sur $\Point(E)$, et topos associés aux ensembles
ordonnés}).
\begin{enumerate}
\renewcommand{\theenumi}{\alph{enumi}}
\renewcommand{\labelenumi}{\theenumi)}
\item Soit $E$ un topos. Désignons par $\Point(E)$ l'ensemble des
classes, à isomorphisme près, de points de $E$. Pour tout ouvert $U$
de $E$, soit $U'$ l'ensemble\pageoriginale des classes de points $p$
de $E$ tels que $U_p\neq \phi$. Montrer que $U\mapsto U'$ est une
application croissante de l'ensemble ordonné $\Ouv(E)$ des ouverts
de $E$ dans l'ensemble des parties de $X=\Point(E)$, et que cette
application commute aux Inf finis aux $\Sup$ quelconques. En
conclure que l'ensemble des parties de $X$ de la forme $U'$ définit
sur $X$ une topologie (dite \emph{topologie canonique sur l'ensemble
des classes de points du topos $E$}). Montrer que si $E$ a
suffisamment de points, l'application $U\mapsto U'$ est un
isomorphisme de l'ensemble ordonné $\Ouv(E)$ avec l'ensemble ordonné
$\Ouv(X)$.

\item Soit $\Oo\in 𝒰$ un ensemble ordonné admettant des Inf
finis et des $\Sup$ quelconques, qui définit donc une catégorie
$\cat(\Oo)$ ayant des produits finis et des sommes quelconques. Nous
dirons qu'une famille de morphismes $U_i\rightarrow U$ de même but
$U$ est couvrante si $U$ est le $\Sup$ des $U_i$. Supposant que dans
$\cat(\Oo)$ les sommes sont universelles \ie que dans $\Oo$ les
$\Sup$ quelconques commutent aux Inf, on définit ainsi sur
$\cat(\Oo)$ une topologie qui en fait un site, d'où un topos
$\cat(\Oo)^{\sim}$, noté aussi simplement $\Oo^{\sim}$. Montrer que
$\Oo$ est canoniquement isomorphe à $\Ouv(\Oo^{\sim})$.

Montrer que la catégorie $\Hhomtop(E,\Oo^{\sim})$ est équivalente à
la catégorie associée à l'ensemble ordonné des morphismes
d'ensembles ordonnés $\Oo\rightarrow \Ouv(E)$ qui commutent aux Inf
finis et aux $\Sup$ quelconques. (Utiliser le critère
\hyperref[IV.prop4.9.4]{4.9.4}.) Si $\Oo'$ est un ensemble ordonné satisfaisant
aux mêmes conditions que $\Oo$, conclure que $\Hhomtop(\Oo^{\sim},
{\Oo'}^{\sim})$ est équivalente à la catégorie associée à l'ensemble
ordonné de tous les morphismes d'ensembles ordonnés
\sisi{$\Oo\rightarrow{\Oo}'$}{${\Oo}'\rightarrow\Oo$} qui commutent aux Inf finis et aux $\Sup$
% correction SGA4-CORR-0024 / SGA-SRC-000411: SGA4-EIV-04-L4772-L4781-10E61872BA64.
quelconques.

\item Montrer\pageoriginale que pour que le topos $E$ soit équivalent
à un topos de la forme $\Oo^{\sim}$, avec $\Oo$ comme dans b), il faut
et il suffit que la famille des ouverts de $E$ soit génératrice dans
$E$. Pour qu'il soit équivalent à un $\Top(X)$ pour $X\in
𝒰$, il faut et il suffit qu'il satisfasse à la condition
précédente et qu'il ait suffisamment de points. Supposant réalisée
la première condition, exprimer la seconde en termes de la structure
de l'ensemble ordonné $\Oo=\Ouv(E)$, en utilisant la dernière
assertion de b).

\item Définir un morphisme de topos canonique
$$
E\longrightarrow \Ouv(E)^{\sim}{\quoi},
$$
qui soit universel (à équivalence près) pour les morphismes de $E$
dans des topos engendrés par leurs ouverts (\cf c)).

\item Soit $X'$ un \emph{petit} sous-ensemble de $X=\Point(E)$, qu'on
munit de la topologie induit par celle de $X$. Définir un morphisme
canonique de topos
$$
\Ouv(E)^{\sim}\longrightarrow \Top(X'){\quoi},
$$
qui est une équivalence lorsque la famille $X'$ de points de $E$ est
conservative. En conclure alors un morphisme canonique de topos
$$
E\longrightarrow \Top(X'){\quoi}.
$$
Donner une caractérisation universelle (à équivalence près) de ce
morphisme de topos pour les morphismes du topos $E$ dans des topos
de la forme $\Top(Y)$. (Noter à ce propos qu'à équivalence près,
$\Top(X')$ ne dépend pas de la petite famille conservative de
foncteurs fibres choisie, ou ce qui revient au même
(\hyperref[IV.subsec4.2]{4.2}), que $X'_{\sob}$ ne dépend pas à homéomorphisme
près du choix d'une telle famille).

\item Supposons\pageoriginale que la catégorie $\Ppoint(E)$ ne
soit pas équivalente à une petite catégorie (\cf
\hyperref[IV.subsec7.3]{7.3}), et que $E$ admette une petite famille
conservative de foncteurs fibres. Montrer que si une telle famille
est prise assez grande, la partie $X'$ de $\Ppoint(E)$ qu'elle
définit n'est pas un espace topologique sobre pour la topologie
induite.

\item Pour les relations avec l'exercice \hyperref[IV.exer7.6]{7.6}, \cf \hyperref[IV.sec8]{8}.
\end{enumerate}
\end{enonce*}

\section{Localisation. Ouverts d'un topos}\etiquette[8]{IV.sec8}

\subsection{}\etiquette[8.1]{IV.subsec8.1}
Soit $C$ une catégorie, et $T(X)$ une relation faisant intervenir un
objet $X$ de $C$. On dit que la relation $T(X)$ est \emph{stable par
changement de base} (en $X$), si pour tout morphisme $X'\rightarrow
X$ de $C$, la relation $T(X)$ implique $T(X')$. Il revient au même
de dire que la sous-catégorie pleine $R$ de $C$, formée des objets
$X$ de $C$ tels qu'on ait $T(X)$, est un \emph{crible}
(\hyperrefext[I][I.def4.1]{4.1}).

\begin{enonce*}[remark]{Remarque 8.1.1}\etiquette[8.1.1]{IV.rem8.1.1}
Supposons que $C$ soit une $𝒰$-catégorie (\hyperrefext[I][I.def1.1]{1.1}),
de sorte que $C$ peut être identifié à une sous-catégorie pleine de
$\widehat{C}$ (\hyperrefext[I][I.prop1.4]{1.4}). Il est parfois commode d'étendre la
définition de la relation $T$ dans $\ob C$ en une relation
$\widehat{T}$ portant sur un objet $F$ de $\widehat{C}$, en
désignant par $\widehat{T}(F)$ la relation: pour toute flèche
$X\rightarrow F$ dans $\widehat{C}$, avec $X\in \ob C$, on a $T(X)$.
La relation $\widehat{T}(F)$ est évidemment encore stable par
changement de base en $F$. Désignant encore par $R$ le sous-objet de
l'objet final $e$ de $\widehat{C}$ défini par le crible $R$ de $C$
(\hyperrefext[I][I.rem4.2.1]{4.2.1}), la relation $\widehat{T}(F)$ peut aussi
s'exprimer par $\Hom_C^{\widehat{\phantom{,}}}(F,R)\neq ∅$, \ie
elle signifie que l'unique morphisme $F\rightarrow e$ se factorise
par le sous-objet $R$ de $e$.
\end{enonce*}

\subsection{}\etiquette[8.2]{IV.subsec8.2}
Supposons maintenant que $C$ soit un \emph{site}. Une relation
$T(X)$ en un argument $X\in \ob C$ est dite \emph{stable par
descente} si pour toute famille couvrante\pageoriginale
$(X_i:X_i\rightarrow X)_{i\in I}$ d'un objet $X$ de $C$, la relation
« $T(X_i)$ pour tout $i\in I$ »  implique la relation $T(X)$. On
dit que $T$ est une \emph{relation de nature locale} si elle est
stable par changement de base (\hyperref[IV.subsec8.1]{8.1}) et stable par
descente. Lorsque $C$ est une $𝒰$-catégorie, et que $T$ est
déjà supposée stable par changement de base, \ie qu'elle est
définissable par un crible $R$ de $C$, qu'on identifiera à un
sous-préfaisceau du préfaisceau final sur $C$, on voit aussitôt que
\sisi{$C$ est de nature locale}{$T$ est de nature locale} si et seulement si $R$ est un
% correction SGA4-CORR-0025 / SGA-SRC-000412: SGA4-EIV-04-L4867-L4875-0C740D22CD39.
\emph{faisceau}. Ainsi, \emph{les relations de nature locale en un
argument $X\in \ob C$ correspondent exactement aux sous-faisceaux du
faisceau final de $C$, \ie aux sous-objets de l'objet final du topos
$C^{\sim}$.} Elles ne dépendent donc essentiellement que du topos
$C^{\sim}$ défini par le site $C$ (comme toutes les notions
importantes associées à un site !). En termes du sous-faisceau $R$
du faisceau final, la relation $T(X)$ s'exprime comme
$\Hom(\epsilon(X),R)\neq ∅$. Elle se prolonge canoniquement en la
relation $T^{\sim}(F):\Hom(F,R)\neq ∅$ sur le topos $C^{\sim}$.

\begin{enonce*}[remark]{Remarque 8.2.1}\etiquette[8.2.1]{IV.rem8.2.1}
Avec la notation introduite dans \hyperref[IV.rem8.1.1]{8.1.1}, pour un
préfaisceau $F$ sur $C$, comme $\Hom(F,R)\simeq
\Hom(\underline{a}(F),R)$, la relation $\widehat{T}(F)$ équivaut à la
relation $\widehat{T}(\underline{a}(F))$.
\end{enonce*}

\begin{enonce*}{Définition 8.3}\etiquette[8.3]{IV.def8.3}
On appelle ouvert d'un $𝒰$-topos $E$ tout sous-objet de
l'objet final $e_E$ de $E$; si $X$ est un objet de $E$, on appelle
parfois ouvert de $X$ tout ouvert du topos induit $E_{/X}$, \ie tout
sous-objet de $X$.

On appelle ouvert d'un $𝒰$-site $C$ un ouvert du
$𝒰$-topos associé $C^{\sim}$.
\end{enonce*}

On notera que, les objets finaux de $E$ étant canoniquement
isomorphes, l'ambiguïté introduite dans \hyperref[IV.def8.3]{8.3} par le choix
\sisi{de $e_E$ est inoffensive}{de $e_E$ est inoffensive}; on\pageoriginale peut aussi, de façon plus
% correction SGA4-CORR-0026 / SGA-SRC-000413: SGA4-EIV-04-L4899-L4901-76EEDF7946A3.
intrinsèque mais moins maniable du point de vue pratique, définir
les ouverts de $E$ comme étant les sous-catégories pleines $R$ de
$E$ qui sont telles que la relation $F\in \ob R$ pour un objet $F$
de $E$ soit une relation de nature locale. De même, les ouverts d'un
$𝒰$-site $C$ correspondent biunivoquement aux
sous-catégories pleines de $C$ telles que la relation $F\in \ob R$
pour un objet de $C$ soit de nature locale; cette notion est donc
essentiellement indépendante de l'univers choisi $𝒰$. Enfin,
en vertu de ce qui a été dit dans \hyperref[IV.subsec8.2]{8.2}, une relation
de nature locale sur un topos (ou sur un site) est essentiellement
la même chose qu'un ouvert dudit topos (ou du site).

\setcounter{subsection}{3}
\subsection{Exemples d'ouverts}\etiquette[8.4]{IV.subsec8.4}

\subsubsection{}\etiquette[8.4.1]{IV.subsubsec8.4.1}
Soit $X$ un espace topologique, et interprétons $\Top(X)$
(\hyperref[IV.subsec2.1]{2.1}) comme la catégorie des espaces étales sur $X$.
Comme il est clair que les monomorphismes dans $\Top(X)$ sont les
applications injectives, il s'ensuit que les ouverts du \sisi{topos $T(X)$}{topos $\Top(X)$}
% correction SGA4-CORR-0027 / SGA-SRC-000414: SGA4-EIV-04-L4918-L4922-7AD68EE312D2.
s'identifient aux ouverts de l'espace $X$. Plus généralement, les
ouverts d'un topos $\Top(X,G)$ (\hyperref[IV.subsec2.4]{2.4}) s'identifient
aux ouverts de $X$ \emph{invariants} par l'action de $G$.

\subsubsection{}\etiquette[8.4.2]{IV.subsubsec8.4.2}
Les ouverts d'un $𝒰$-topos $E$ forment un ensemble
$𝒰$-petit (\hyperrefext[I][I.prop7.4]{7.4}) ordonné par l'inclusion,
admettant des sup quelconques et des inf finis. Il admet l'objet
initial (ou « objet vide ») $\phi_E$ comme plus petit élément,
et l'objet final $e_E$ de $E$ comme plus grand élément. Ces deux
éléments ne sont identiques que si $E$ est un « topos vide »
(\hyperref[IV.subsec2.2]{2.2}).

\subsubsection{}\etiquette[8.4.3]{IV.subsubsec8.4.3}
Soit $G$ un Monoïde dans $E$, d'où un topos $B_G$
(\sisi{\hyperref[IV.subsec2.4]{2.4}}{\hyperref[IV.subsec2.6]{2.6}}), catégorie des objets de $E$ sur lesquels $G$
opère à gauche. L'objet final $e_G$ de $B_G$ est
l'objet\pageoriginale final $e_E$ de $E$ avec opération triviale de
$G$. On voit par suite que l'ensemble ordonné des ouverts de $B_G$
est canoniquement isomorphe à la catégorie des ouverts de $E$. En
particulier, si $E$ est le topos ponctuel donc $G$ est un monoïde
ordinaire, l'ensemble des ouverts de $B_G$ est exactement formé de
deux éléments, savoir $\phi_{B_G}$ et $e_G$.

\subsubsection{}\etiquette[8.4.4]{IV.subsubsec8.4.4}
Si $E$ est un topos de la forme $\widehat{C}$, où $C$ est une
$𝒰$-catégorie, alors (\hyperref[IV.subsec8.1]{8.1}) l'ensemble ordonné
des ouverts de $E$ est isomorphe à l'ensemble ordonné des cribles de
$C$.

\subsubsection{}\etiquette[8.4.5]{IV.subsubsec8.4.5}
Soit $X$ un espace topologique sobre non discret. Montrer que
$\TOP(X)$ (\hyperref[IV.subsec2.5]{2.5}) a des ouverts qui ne proviennent pas
d'ouverts de $X$, \ie d'ouverts de $\Top(X)$, par image inverse à
l'aide du morphisme canonique \eqhyperref[IV.eq4.10.1]{4.10.1} $\TOP(X)\rightarrow
\Top(X)$.

\subsection{Exemples de relations de nature locale}\etiquette[8.5]{IV.subsec8.5}

Pour tout objet $X$ du topos $E$, on désigne par $F\rightarrow F_X$
le \sisi{\emph{foncteur de localisation}}{\emph{foncteur de localisation}} $j^{\ast}_X:E\rightarrow
% correction SGA4-CORR-0028 / SGA-SRC-000415: SGA4-EIV-04-L4961-L4963-3D3FEC3BEF6A.
E_{/X}$, $Y\mapsto Y\times X$ (\hyperref[IV.subsec5.2]{5.2}).

\subsubsection{}\etiquette[8.5.1]{IV.subsubsec8.5.1}
Soit $u:F\rightarrow G$ un morphisme de $E$. La relation en
l'argument $X\in \ob E$,

$u_X:F_X\rightarrow G_X$ est un monomorphisme (\resp un
épimorphisme, \resp un isomorphisme)

\noindent est une relation de nature locale. En particulier, si $G$
est un Groupe de $E$, la relation « $G_X$ est le groupe unité sur
$X$ » est de nature locale; l'ouvert de $E$ qui lui correspond est
appelé le \emph{cosupport du Groupe $G$.}

\subsubsection{}\etiquette[8.5.2]{IV.subsubsec8.5.2}
Soient\pageoriginale $u$, $v:F\rightrightarrows G$ deux morphismes
de $E$. La relation en l'argument $X\in \ob E$
$$
u_X=v_X:F_X\longrightarrow G_X
$$
est une relation de nature locale. En particulier, si $G$ est un
Groupe de $E$, et $u$ une \emph{section} de $G$
(\hyperref[IV.subsubsec4.3.6]{4.3.6}), la relation en $X$
$$
\mbox{la section $u_X$ du Groupe $G_X$ de $E_{/X}$ est nulle}
$$
est de nature locale; l'ouvert de $E$ qui lui correspond est appelé
le \emph{cosupport de la section $u$ du Groupe $G$.}

\begin{enonce*}[remark]{Remarque 8.5.3}\etiquette[8.5.3]{IV.rem8.5.3}
Lorsque $E=\Top(X)$, le cosupport d'un faisceau en groupes $G$,
\resp d'une section d'un tel faisceau $G$, n'est autre que l'ouvert
de $X$ complémentaire du \emph{support de} $G$ \resp du
\emph{support de} $u$ dans la terminologie classique [TF].
\end{enonce*}

\setcounter{subsubsection}{3}
\subsubsection{}\etiquette[8.5.4]{IV.subsubsec8.5.4}
Soit $P(f)$ une relation en l'argument $f\in {\rm f\ell\,} C$, $C$
un site. Lorsque dans $C$ les produits fibrés sont représentables,
on dit que la relation $P(f)$ est \emph{stable par changement de
base (\resp stable par descente), \resp de nature locale) sur le
but} (ou \emph{sur la base}, ou « \emph{en bas} »), si pour tout
morphisme $f:X\rightarrow Y$ de $C$, la relation en l'argument
$Y'\in \ob C_{/Y}$:

« le morphisme $f':X'=\fprod{X}{Y'}{Y}\rightarrow Y'$ déduit de
$f$ par changement de base par le morphisme structural
$Y'\rightarrow Y$ satisfait $P(f')$ »

est stable par changement de base, \resp stable par descente, \resp
est de nature locale. Lorsque la \sisi{topologie de $E$}{topologie de $C$} est définie à
% correction SGA4-CORR-0029 / SGA-SRC-000416: SGA4-EIV-04-L5009-L5018-76B83D11DDE8.
l'aide d'une prétopologie, on dit que la relation $P(f)$ est
\emph{de nature locale sur la source} (ou \emph{« en haut »}) si
pour tout morphisme $f:X\rightarrow Y$ de $C$ et toute famille
$(g_i:X_i\rightarrow X)_{i\in I}$\pageoriginale de $\Cov(X)$, la
relation $P(f)$ équivaut à la relation « $P(fg_i)$ pour tout $i\in
I$ ». Lorsque $\Cov(X)$ est formée de toutes les familles
couvrantes de $C$, cette condition signifie donc aussi que la
relation en l'objet $X'$ du site induit $C_{/X}$
\begin{center}
$P(fg)$, où $g:X'\rightarrow X$ est le morphisme structural
\end{center}
est de nature locale. Noter que si la relation $P(f)$ est de nature
locale en haut, elle est a fortiori de nature locale en bas.

\subsubsection{}\etiquette[8.5.5]{IV.subsubsec8.5.5}
Prenons par exemple $C=(\Sch)$, catégorie des schémas $\in
𝒰$, avec la topologie fidèlement plate quasi-compacte (\SGA
3 IV 6.3) ou une topologie moins fine. Alors chacune des propriétés
suivantes d'un morphisme est de nature locale en bas:

surjectif, radiciel, universellement ouvert, \sisi{universellement finie}{universellement fini},
% correction SGA4-CORR-0030 / SGA-SRC-000417: SGA4-EIV-04-L5031-L5041-911577E8356A.
propre, quasi-compact, quasi-compact et dominant, homéomorphisme
universel, séparé, quasi-séparé, localement de type fini, localement
de présentation finie, de type fini, de présentation finie, un
isomorphisme, un monomorphisme, une immersion ouverte, une immersion
fermée, une immersion quasi-compacte, affine, quasi-affine, fini,
quasi-fini, entier, plat, fidèlement plat, net, lisse, étale

\noindent (\EGA IV 2.6.4, 2.7.1 et \EGA IV 17.7.1). D'autre part,
munissons $(\Sch)$ de la prétopologie pour laquelle, pour tout
schéma $X\in 𝒰$, $\Cov(X)$ est formé des familles de
\sisi{morphismes $(f_i}{morphismes $(f_i}:X_i\rightarrow X)_{i\in I}$ qui sont surjectives et
% correction SGA4-CORR-0031 / SGA-SRC-000418: SGA4-EIV-04-L5044-L5047-EF75F33BD855.
telles que les $f_i$ soient plats et localement de présentation
finie. Alors chacune des propriétés suivantes est \emph{de nature
locale en haut}:

localement de type fini, localement de présentation finie, de type
fini, plat, net, lisse, étale.

\noindent (\EGA IV 17.7.5, 17.7.7).

\begin{enonce*}[remark]{Exercice 8.6}\etiquette[8.6]{IV.exer8.6}
Soit\pageoriginale $f:E\rightarrow F$ un morphisme de
$𝒰$-topos. Considérons la relation en l'argument $X\in \ob
F$: le morphisme induit $E_{/f^{\ast}(X)}\rightarrow F_{/X}$ est une
équivalence de topos. Prouver que cette relation est de nature
locale.
\end{enonce*}

\begin{enonce*}[remark]{Exercice 8.7}\etiquette[8.7]{IV.exer8.7}
(\emph{Partitions d'un topos, somme de topos}). Soit $E$ un
$𝒰$-topos. Une famille $(e_i)_{i\in I}$ d'ouverts de $E$,
\ie de sous-objets de l'objet final $e$ de $E$, est appelée une
\emph{partition de} $e$ (ou encore, du topos $E$) si le morphisme
canonique $\amalg_{i\in I}e_i\rightarrow e$ est un isomorphisme, \ie
(\hyperrefext[II][II.cor4.6.2]{4.6.2})) si $e$ est le sup des $e_i$ et si $i\neq j$
implique $e_i\cap e_j=∅_E$.
\begin{enumerate}
\renewcommand{\theenumi}{\alph{enumi}}
\renewcommand{\labelenumi}{\theenumi)}
\item Montrer que pour que la famille $(e_i)_{i\in I}$ d'objets de $E$
soit une partition de $E$, il faut et il suffit que le foncteur
\begin{equation*}
E\longrightarrow \prod_i E_{/e_i}\tag{8.7.1}
\end{equation*}\etiquette[8.7.1]{IV.eq8.7.1}
défini par \sisi{les foncteurs de localisation}{les foncteurs de localisation} $E\rightarrow E_{/e_i}$ soit
% correction SGA4-CORR-0032 / SGA-SRC-000419: SGA4-EIV-04-L5076-L5081-4292B2A3710F.
une équivalence de catégories.

\item Soit réciproquement $(E_i)_{i\in I}$ une famille de
$𝒰$-topos, avec $I\in 𝒰$. Prouver que $E=\Pi_{i\in
I} E_i$ est un $𝒰$-topos (qui sera appelé le \emph{topos
somme} de la famille des topos $(E_i)_{i\in I}$). Définir une
partition $(e_i)_{i\in I}$ de $E$ et des équivalences de catégories
$E_{/e_i}\approx E_i$, de telle façon que les foncteurs de
projection $E\rightarrow E_i$ s'identifient aux foncteurs de
localisation $E\rightarrow E_{/e_i}$.

\item Soit $E$ le topos somme de la famille des topos $(E_i)_{i\in
I}$. Montrer que pour tout $i\in I$ le foncteur projection
$E\rightarrow E_i$ est de la forme $u^{\ast}_i$, où
\begin{equation*}
u_i:E_i\longrightarrow E\tag{8.7.2}
\end{equation*}\etiquette[8.7.2]{IV.eq8.7.2}
est\pageoriginale un morphisme de topos. Prouver que pour tout
$𝒰$-topos $F$, le foncteur $v\mapsto (v\circ u_i)_{i\in I}$
\begin{equation*}
\Hhomtop(E,F)\longrightarrow \prod_i\Hhomtop(E_i,F)\tag{8.7.3}
\end{equation*}\etiquette[8.7.3]{IV.eq8.7.3}
est une équivalence de catégories.

\item Soit $(e_i)_{i\in I}$ une partition du topos $E$. Pour tout
morphisme de topos $g:F\rightarrow E$, la famille $(f_i)_{i\in I}$,
avec $f_i=g^{\ast}(e_i)$, est une partition de $F$, et les
morphismes induits \sisi{$g_i:E_{/e_i}\rightarrow F_{/f_i}$}{$g_i:F_{/f_i}\rightarrow E_{/e_i}$} permettent de
% correction SGA4-CORR-0033 / SGA-SRC-000420: SGA4-EIV-04-L5105-L5111-9F3C95BD7504.
reconstituer $g$ (à isomorphisme unique près), le foncteur
$g^{\ast}$ s'identifiant au produit cartésien des foncteurs
$g^{\ast}_i$ (compte tenu des équivalences du type (\hyperref[IV.exer8.7]{8.7}
1)). En conclure une description complète de la catégorie
$\Hhomtop(F,E)$, en termes des catégories de la forme
$\Hhomtop(F',E_i)$, où $F'$ est un topos de la forme $F_{/f}$ ($f$
sommande directe de l'objet final de $F$) et $E_i=E_{/e_i}$.

\item En particulier, si $F$ est connexe non vide (\cf
  \hyperref[IV.subsubsec4.3.5]{4.3.5}) \ie si
pour toute partition $(f_i)_{i\in I}$ de $F$ il existe un $i\in I$
et un seul tel que $f_i\neq \phi_F$, prouver que le foncteur
canonique
\begin{equation*}
\coprod_{i\in I}\Hhomtop(F,E_i)\longrightarrow
\Hhomtop(F,E)\tag{8.7.4}
\end{equation*}\etiquette[8.7.4]{IV.eq8.7.4}
déduit des morphismes de topos \eqhyperref[IV.eq8.7.2]{8.7.2} est une équivalence
de catégories. Plus particulièrement, la famille des morphismes de
topos \eqhyperref[IV.eq8.7.2]{8.7.2} induit une équivalence de catégories
$$
\coprod_{i\in I}\Ppoint(E_i)\xrightarrow{\approx}\Ppoint(E){\quoi}.
$$

\item Une\pageoriginale partie $I$ de l'ensemble des ouverts de $E$
est appelée une \emph{partition réduite} de $E$ si la famille
identique $(e_i)_{i\in I}$ indexée par $I$ est une partition, et si
les $e_i$ sont $\neq \phi_E$. Montrer que la relation de raffinement
(\hyperrefext[I][I.sec4.3.2]{4.3.2}) entre partitions réduites fait de l'ensemble $P$
des partitions réduites un ensemble ordonné $𝒰$-petit, tel
que la borne supérieure de deux éléments de $P$ existe. On trouve
ainsi un système projectif strict $(I_p)_{p\in P}$, qui sera
considéré comme un pro-ensemble et noté $\pi_{\circ}(E)$. On
l'appelle le \emph{pro-$\pi_{\circ}$ du topos} $E$.

\item Prouver que $\pi_{\circ}(E)$ pro-représente le foncteur
$$
T\longmapsto f_{\ast}f^{\ast}(T)=\Gamma(E,
T_E):(𝒰\text{-}\Ens)\longrightarrow (𝒰\text{-}\Ens)
$$
associé au morphisme canonique $f:E\rightarrow P$ de $E$ dans le
topos ponctuel $P$ (\hyperref[IV.subsec4.3]{4.3}). En particulier, pour que ce
foncteur soit représentable, il faut et il suffit que
$\pi_{\circ}(E)$ soit essentiellement constant, \ie isomorphe (en
tant que pro-ensemble) à un ensemble « ordinaire », qui sera
noté encore $\pi_{\circ}(E)$. Pour que $\pi_{\circ}(E)$ soit réduit
à un point, il faut et il suffit que $E$ soit « connexe non
vide ». Pour que $\pi_{\circ}(E)$ soit vide, il faut et il suffit que
$E$ soit le « topos vide »  (\hyperref[IV.subsec2.2]{2.2}).

\item Supposons $E$ quasi-compact, \ie que tout recouvrement de son
objet final $e$ admette un sous-recouvrement fini. Montrer qu'alors
$\pi_{\circ}(E)$ est un ensemble profini, et peut s'identifier par
suite (moyennant l'équivalence de catégories bien connue entre
pro-objets de la catégorie des ensembles finis, et espaces compacts
totalement discontinus) à un espace compact totalement discontinu,
qu'on notera également $\pi_{\circ}(E)$.

\item Soit $X$ un espace topologique, $\pi_{\circ}(X)$ l'ensemble des
composantes connexes de $X$, considéré comme un pro-ensemble
constant. Définir un morphisme\pageoriginale canonique de
pro-ensembles
\begin{equation*}
\pi_{\circ}(X)\longrightarrow
\pi_{\circ}(\Top(X)){\quoi},\tag{8.7.5}
\end{equation*}\etiquette[8.7.5]{IV.eq8.7.5}
ou ce qui revient au même, une application canonique
\begin{equation*}
\pi_{\circ}(X)\longrightarrow \varprojlim
\pi_{\circ}(\Top(X)){\quoi}.\tag{8.7.6}
\end{equation*}\etiquette[8.7.6]{IV.eq8.7.6}
Montrer que le premier morphisme est un isomorphisme si et seulement
si $X$ est homéomorphe à un espace somme d'espaces connexes (ce qui
est le cas en particulier si $X$ est localement connexe).

\item Supposons que $X$ soit un schéma quasi-compact. Prouver qu'alors
\eqhyperref[IV.eq8.7.6]{8.7.6} est bijectif.

\item Établir le « caractère fonctoriel »  en $X$ des morphismes
\eqhyperref[IV.eq8.7.5]{8.7.5} et \eqhyperref[IV.eq8.7.6]{8.7.6}, en précisant pour
commencer le sens de cette expression.


\item (Comparer \hyperref[IV.exer7.6]{7.6} i)). Montrer que les deux conditions
suivantes sur le
   $𝒰$-topos $E$ sont équivalentes: (i) Le morphisme canonique de
   $E$ dans le topos ponctuel est « essentiel »  (exercice \hyperref[IV.exer7.6]{7.6} a))
   \ie le foncteur $I\mapsto I_E$ de $(𝒰\text{-}\Ens)$ dans $E$ commute
   aux produits indexés par des ensembles $\in 𝒰$. (ii) Pour tout
   objet $X$ de $E$, le topos induit $E_{/X}$ a un
   $\pi_{\circ}(E_{/X})$ qui est un ensemble ordinaire, \ie $X$ est
   isomorphe à la somme d'une famille d'objets connexes de $E$. - On
   dit alors que $E$ est un \emph{topos localement connexe} (comparer \sisi{IV }{} \hyperref[IV.exer2.7.5]{2.7.5}).
\end{enumerate}
\end{enonce*}

\begin{enonce*}[remark]{Exercice 8.8}\etiquette[8.8]{IV.exer8.8}
(\emph{Morphismes dominants de topos}).
\begin{enumerate}
\renewcommand{\theenumi}{\alph{enumi}}
\renewcommand{\labelenumi}{\theenumi)}
\item Soit $f:E'\rightarrow E$ un morphisme de topos. Montrer
l'équivalence des conditions suivantes:

\begin{itemize}
\item[(i)] $f_{\ast}(\phi_{E'})=\phi_E$ (où $\phi_E$, \sisi{$=\phi_{E'}$}{$\phi_{E'}$}
% correction SGA4-CORR-0034 / SGA-SRC-000421: SGA4-EIV-04-L5210-L5215-B7E7AEA7ADBB.
désignent les objets initiaux de $E$, $E'$).

\item[(ii)] Pour\pageoriginale tout objet $X$ de $E$ qui est « non
vide »  \ie non isomorphe à $\phi_E$, $f^{\ast}(X)$ est un objet
non vide de $E'$.

\item[$\rm(ii')$] Comme (ii), avec $X$ un sous-objet de l'objet final
$e_E$ de $E$, \ie $X$ un ouvert du topos $E$.
\end{itemize}

On dit alors que le morphisme $f$ de topos est \emph{dominant}.

\item Soit $f:X'\rightarrow X$ une application continue d'espaces
topologiques, d'où un morphisme de topos
$\Top(f):\Top(X')\rightarrow \Top(X)$ (\hyperref[IV.subsec4.1]{4.1}). Montrer
que $\Top(f)$ est dominant si et seulement si $f$ est dominant, \ie
$f(X')$ est dense dans $X$.

\item Montrer que si $f:E'\rightarrow E$ est un morphisme « conservatif »  de topos (\hyperref[IV.subsubsec6.4.0]{6.4.0}), \ie tel que
$f^{\ast}$ soit fidèle, alors $f$ est dominant. \sisi{Montrer que le}{Montrer que la}
% correction SGA4-CORR-0035 / SGA-SRC-000422: SGA4-EIV-04-L5233-L5236-48FAFC859A25.
réciproque n'est pas nécessairement vraie, même pour un
morphisme de
localisation $E_{/U}\rightarrow E$, où $U$ est un ouvert du topos
$E$, ((Montrer que dans ce cas $f$ n'est conservatif que si $U$ est
l'objet final de $E$ (donc si $f$ est une équivalence de topos), et
utiliser l'exemple b).)

\item Soit $f:X'\rightarrow X$ une flèche dans un topos $E$. On dit
que $f$ est un \emph{morphisme dominant} dans le topos $E$ si le
morphisme correspondant des topos induits $E_{/X'}\rightarrow
E_{/X}$ \eqhyperref[IV.eq5.5.2]{5.5.2} est dominant. Montrer que cette
propriété ne
dépend que de l'image $f(X')$ de $X'$ dans $X$, en tant qu'élément
de l'ensemble ordonné $\ouv(X)$ des sous-objets de $X$. Montrer que
le morphisme $E_{/X'}\rightarrow E_{/X}$ est conservatif si et
seulement si $f$ est couvrant.
\end{enumerate}
\end{enonce*}

\section{Sous-topos et recollement de
  topos}\etiquette[9]{IV.sec9}\pageoriginale

\subsection{Sous-topos et plongement de
  topos}\etiquette[9.1]{IV.subsec9.1}

\begin{enonce*}{Définition 9.1.1}\etiquette[9.1.1]{IV.def9.1.1}
\begin{enumerate}
\renewcommand{\theenumi}{\alph{enumi}}
\renewcommand{\labelenumi}{\theenumi)}
\item Un morphisme de topos $f:F\rightarrow E$ est appelé un
plongement si le foncteur $f_{\ast}:F\rightarrow E$ est pleinement
fidèle (\ie si le morphisme d'adjonction
$f^{\ast}f_{\ast}\rightarrow \id_F$ est un isomorphisme).

\item On appelle sous-topos d'un topos $E$ une sous-catégorie
strictement pleine $E'$ de $E$ qui est un topos \sisi{et tel que le}{et telle que le}
% correction SGA4-CORR-0036 / SGA-SRC-000423: SGA4-EIV-04-L5269-L5277-A6346A0190E6.
foncteur d'inclusion $\alpha:E'\rightarrow E$ soit de la forme
$i_{\ast}$, où $i:E'\rightarrow E$ est un morphisme de topos (\ie
$(\hyperref[IV.subsubsec3.1.1]{3.1.1})$ $\alpha$ admet un adjoint à gauche
$i^{\ast}$ qui est exact à gauche). Le morphisme $i$ (qui est
déterminé à isomorphisme unique près) est alors appelé
le morphisme
d'inclusion du sous-topos $E'$ dans $E$.
\end{enumerate}
\end{enonce*}

\begin{enonce*}[remark]{Remarque 9.1.2}\etiquette[9.1.2]{IV.rem9.1.2}
\begin{enumerate}
\renewcommand{\theenumi}{\alph{enumi}}
\renewcommand{\labelenumi}{\theenumi)}
\item Il est clair que le morphisme d'inclusion d'un sous-topos est un
plongement, et qu'une sous-catégorie $E'$ du topos $E$ est un
sous-topos si et seulement si c'est l'image essentielle d'un
foncteur image directe $f_{\ast}$ associé à un morphisme de topos
\sisi{$g:F\rightarrow E$}{$f:F\rightarrow E$} qui est un plongement.
% correction SGA4-CORR-0037 / SGA-SRC-000424: SGA4-EIV-04-L5285-L5289-0A0F21448241.

\item Il résulte aussitôt des définitions qu'un morphisme de topos
$f:F\rightarrow E$ est un plongement si et seulement s'il est
isomorphe à un morphisme composé
$$
F\xrightarrow{g}E'\xrightarrow{i}E{\quoi},
$$
où $g$ est une équivalence de topos et $i$ le morphisme d'inclusion
d'un sous-topos $E'$ de $E$. Ce dernier est déterminé de façon
unique par les conditions précédentes, comme l'image essentielle de
$f_{\ast}$, et la factorisation précédente est également unique à
isomorphisme unique près. Bien entendu, il y a lieu en pratique
d'identifier, au moyen de $g$, $F$ au sous-topos $E'$ de $E$
(comparer \hyperref[IV.subsubsec3.4.1]{3.4.1}).

\item Une équivalence de topos (\hyperref[IV.subsec3.4]{3.4})
  $u:E\xrightarrow{\approx}E_1$
définit de façon évidente une bijection entre l'ensemble des
sous-topos de $E$ et l'ensemble des sous-topos de $E_1$. En effet,
$u_{\ast}$ étant une équivalence de catégories,\pageoriginale
établit une bijection entre l'ensemble des sous-catégories
strictement pleines de $E$ et l'ensemble des sous-catégories
strictement pleines de $E_1$ (à $E'\subset E$ correspondant l'image
essentielle $E'_1$ de $u_{\ast}| E'$), et on constate aussitôt que
$E'$ est un sous-topos de $E$ si et seulement si $E'_1$ est un
sous-topos de $E_1$. On peut donc, pour l'étude des sous-topos, se
ramener au cas où $E$ est de la forme $\widetilde{C}$, où $C$ est un
petit site. Dans ce cas, il résulte de \hyperrefext[II][II.theo5.5]{5.5} qu'une
sous-catégorie strictement pleine $E'$ de $E$ est un sous-topos si
et seulement si le foncteur d'inclusion $\alpha:E'\rightarrow E$
admet un adjoint à gauche qui est exact à gauche (cela implique donc
déjà que $E'$ est un topos), et que l'ensemble de ces sous-topos est
en correspondance biunivoque avec l'ensemble des topologies $T'$ sur
$C$ qui sont plus fines que la topologie donnée $T$ de $C$, en
associant à toute telle $T'$ la sous-catégorie strictement pleine
$(C,T')^{\sim}$ de $C^{\sim}$, formée des faisceaux pour la
topologie $T'$. Ceci montre en même temps, pour \emph{tout}
$\sheaf{\mathcal{U}}$-topos $E$: une sous-catégorie strictement
pleine $E'$ de $E$ est un sous-topos si et seulement si le foncteur
d'inclusion $\alpha:E'\rightarrow E$ admet un adjoint à gauche qui
est exact à gauche; \emph{l'ensemble ordonné $\mathfrak{E}(E)$ des
sous-topos de $E$ est $\mathcal{U}$-petit, et toute partie de
$\mathfrak{E}(E)$ admet une borne supérieure et une borne
inférieure.}

\item Si $E'$ et $E''$ sont deux sous-topos du topos $E$, leur
intersection est un sous-topos de $E$. En effet, comme cette
intersection est strictement pleine, on est ramené au cas où $E$ est
de la forme $C^{\sim}$. Avec les notations de \sisi{b)}{c)}, $E'$ et $E''$
correspondent alors à deux topologies $T'$, $T''$ sur $C$ plus fines
que $T$, et il suffit de voir qu'alors $E'\cap E''$ \emph{est la
catégorie des faisceaux pour la topologie $T'''$ borne supérieure de
\sisi{$T'$ et de $T''$}{$T'$ et de $T''$}}. Indiquons la démonstration de ce fait (qui
% correction SGA4-CORR-0038 / SGA-SRC-000425: SGA4-EIV-04-L5336-L5342-89B9DB6258D8.
aurait dû figurer en corollaire après \hyperrefext[II][II.theo4.4]{4.4} ou
\hyperrefext[II][II.theo5.5]{5.5} !). Soit $E'''$ la catégorie des faisceaux pour
$T'''$, évidemment contenue dans $E'$ et $E''$ donc dans $E'\cap
E''$. Soit, pour toute sous-catégorie pleine $D$ de $E$, $t(D)$ la
topologie la plus fine parmi celles pour lesquelles les éléments de
$D$ sont des faisceaux\pageoriginale (\hyperrefext[II][II.prop2.2]{2.2}). Les
inclusions $E'''\subset E'\cap E''\subset E'$, $E''$ impliquent
évidemment $t(E')$, $t(E'')\leq t(E'\cap E'')\leq t(E''')$, d'autre
part on a (\sisi{\hyperrefext[II][II.theo4.4]{4.4} I}{\hyperrefext[II][II.cor4.4.4]{4.4.4}}) $t(E')=T'$, $t(E'')=T''$,
$t(E''')=T'''$, de sorte que les inégalités précédentes et la
définition de $T'''$ comme borne supérieure impliquent que
$T'''=t(E'\cap E'')$, donc $E'\cap E''\subset E'''$, \cqfd\  Plus
généralement, cet argument montre que si $T'$ est une topologie
borne supérieure d'une famille quelconque de topologies $(T_i)$, la
catégorie $E'$ des faisceaux pour $T'$ est l'intersection des
catégories $E_i$ de faisceaux pour les $T_i$. On en conclut en
particulier:
\end{enumerate}
\end{enonce*}

\begin{enonce*}[remark]{Proposition 9.1.3}\etiquette[9.1.3]{IV.prop9.1.3}
Soit $E$ un topos. Alors pour toute famille de sous-topos
$(E_i)_{i\in I}$ de $E$, l'intersection $\bigcap_{i\in I} E_i$ est
un sous-topos de $E$.
\end{enonce*}

\begin{enonce*}{Proposition 9.1.4}\etiquette[9.1.4]{IV.prop9.1.4}
Soient $i:E'\rightarrow E$ un plongement de topos, et $F$ un topos.
Le foncteur
\begin{equation*}
f\longmapsto i\circ f:\Hhomtop(F,E')\longrightarrow
\Hhomtop(F,E)\tag{9.1.4.1}
\end{equation*}\etiquette[9.1.4.1]{IV.eq9.1.4.1}
est pleinement fidèle, et son image essentielle est formée des
morphismes de topos $g:F\rightarrow E$ tels que l'image essentielle
de $g_{\ast}$ soit contenue dans celle de $i_{\ast}$.
\end{enonce*}

Considérons en effet le diagramme commutatif évident
\[
\xymatrix@C=1.75cm@R=1cm@M=6pt{
\Hhomtop (F,E') \ar @{^{(}->}[d] \ar[r]& \Hhomtop (F,E) \ar @{^{(}->}[d]\\
\Hhom (F,E') \ar[r] &\Hhom (F,E) , }
\]
où les flèches verticales sont les foncteurs $h\mapsto h_{\ast}$,
qui sont pleinement fidèles par définition de $\Hhomtop$ (\sisi{IV}{}
\hyperref[IV.subsec3.2]{3.2}). D'autre part la deuxième flèche horizontale
$G\mapsto i_{\ast} G$ est pleinement fidèle, puisque $i_{\ast}$
l'est, donc il en est de même de la première flèche horizontale.
Reste à prouver la caractérisation de l'image essentielle de
\eqhyperref[IV.eq9.1.4.1]{9.1.4.1}. La nécessité de la condition étant
évidente par la
formule $(if)_{\ast}=i_{\ast}f_{\ast}$, il reste à prouver que si
l'image essentielle de $g_{\ast}$ est contenue dans celle de
$i_{\ast}$, \ie si on peut écrire \sisi{$g_{\ast}=f_{\ast}i_{\ast}$}{$g_{\ast}=i_{\ast}f_{\ast}$}, avec
$f_{\ast}:F\rightarrow E'$ un foncteur, alors $g$ est dans l'image
essentielle de \eqhyperref[IV.eq9.1.4.1]{9.1.4.1}. Il revient au même de dire que
$f_{\ast}$ admet un adjoint à gauche qui est \sisi{exact à droite}{exact à gauche}.
% correction SGA4-CORR-0039 / SGA-SRC-000426: SGA4-EIV-04-L5395-L5403-931842135FA7.
Or\pageoriginale il est clair que $f_{\ast}$ admet
$g^{\ast}i_{\ast}$ comme adjoint à gauche, et ce foncteur est exact
à gauche, comme composé des foncteurs exacts à gauche $g^{\ast}$ et
$i_{\ast}$, \cqfd

Pour une réciproque de \hyperref[IV.prop9.1.4]{9.1.4}, \cf exercice
\hyperref[IV.exer9.1.6]{9.1.6}.

Prenant pour $F$ le topos ponctuel (\sisi{IV }{} \hyperref[IV.subsec2.2]{2.2}), on déduit
en particulier \sisi{de 9.1.5:}{de 9.1.4:}
% correction SGA4-CORR-0040 / SGA-SRC-000427: SGA4-EIV-04-L5409-L5412-E671F75D540C.

\begin{enonce*}{Corollaire 9.1.5}\etiquette[9.1.5]{IV.cor9.1.5}
Si $f:E'\rightarrow E$ est un morphisme de plongement de topos, le
foncteur correspondant
\begin{equation*}
\Point (f):\Ppoint(E')\longrightarrow \Ppoint(E)\tag{9.1.5.1}
\end{equation*}\etiquette[9.1.5.1]{IV.eq9.1.5.1}
est pleinement fidèle.
\end{enonce*}

\begin{enonce*}[remark]{Exercice 9.1.6}\etiquette[9.1.6]{IV.exer9.1.6}
(\emph{Plongements de topos et 2-monomorphismes}).

Soit $i:E'\rightarrow E$ un morphisme de topos.

\begin{enumerate}
\renewcommand{\theenumi}{\alph{enumi}}
\renewcommand{\labelenumi}{\theenumi)}
\item Si $F$ et $G$ sont deux topos \emph{au-dessus de} $E'$, définir
(avec les notations de \hyperref[IV.exer5.14]{5.14} a)) un foncteur canonique
\begin{equation*}
f\longmapsto f\ast i:\Hhomtop_{E'}(F,G)\longrightarrow
\Hhomtop_E(F,G){\quoi}.\tag{9.1.6.1}
\end{equation*}\etiquette[9.1.6.1]{IV.eq9.1.6.1}

\item Supposons que $i$, en tant que 1-flèche de la 2-catégorie des
$𝒱\text{-}𝒰$-topos (\hyperref[IV.subsubsec3.3.2]{3.3.2}), soit
un $2$-monomorphisme, \ie tel que pour tout $𝒰$-topos $F$ qui
soit élément de $𝒱$, le foncteur (9.1.4.1) soit pleinement % mettre un ref
fidèle. Montrer que pour tout couple de $𝒰$-topos $F$, $G$
(pas nécessairement $\in 𝒱$) le foncteur (9.1.6.1) est %% mettre un ref
pleinement fidèle.

\item Montrer que $i$ est un plongement si et seulement si $i$ est un
\sisi{$2$-monomorphisme}{$2$-monomorphisme}. (Pour la suffisance, appliquer b) à des topos
% correction SGA4-CORR-0041 / SGA-SRC-000428: SGA4-EIV-04-L5444-L5446-5A9A088A4C85.
induits $E'_{/X'}$ et $E'_{/Y'}$, et utiliser \hyperref[IV.exer5.14]{5.14} c)).

\item Soit $g:F\rightarrow E$ un morphisme de topos, et supposons que
le $2$-produit fibré (\hyperref[IV.prop5.11]{5.11}) 
\sisi{$F'=\fprod{F^2}{E'}{E}$}{$F'=F\times_{E}^{2}E'$}
existe (condition toujours vérifiée, comme l'a montré P. DELIGNE\nde{Voir aussi 
[16] proposition 4.47.}).
Soit $j:F'\rightarrow F$ le morphisme de projection. Prouver que si
$i$ est un plongement, il en est de même de $j$. (Utiliser c), et
noter que la stabilité de la notion de 2-monomorphisme par
2-changement de base est formelle). Dans le cas où $E'$ est un
sous-topos de $E$ et $i:E'\rightarrow E$ est le morphisme canonique
d'inclusion, l'image essentielle de $F'$ par $j_{\ast}$ (qui est un
sous-topos de $F$ qui ne dépend\pageoriginale pas de
l'indétermination dans la construction d'un produit 2-fibré)
s'appelle le \emph{sous-topos de $F$ image inverse par
$g:F\rightarrow E$ du sous-topos $E'$ de $E$.}

\item Avec les notations de la fin de d), choisissons comme produit
$2$-fibré le sous-topos de $F$ image inverse du sous-topos $E'$ de
$E$. Soit $X$ un objet de $F$. Montrer que pour que $X$ appartienne
à $F'$, il faut qu'il satisfasse la condition suivante: pour tout
objet $U$ de $F$, désignant par \sisi{$g_U:F_{/U}\rightarrow E$ la}{$g_U:F_{/U}\rightarrow E$ le}
morphisme de topos \sisi{induit par $f$}{induit par $g$}, on a $g_{U\ast}(X| U)\in \Ob E'$,
% correction SGA4-CORR-0042 / SGA-SRC-000429: SGA4-EIV-04-L5464-L5473-2DEB2F76F63F.
où $X| U=(X\times U\xrightarrow{\pr_2}U)$ est la « restriction de
$X$ à $U$ ». (NB: On peut noter que si $i_U:F_{/U}\rightarrow F$
est l'inclusion canonique, on a $g_U=g\circ i_U$, d'où $g_{U\ast}(X|
U)=g_{\ast}(i_{U\ast}(X| U))=g_{\ast}(\Hhom(U,X))$ où $\Hhom(U,X)$
est l'objet défini dans 10.1 plus bas).

\noindent \emph{Problème} Réciproquement, la condition précédente
sur $X$ est-elle suffisante pour que $X$ appartienne à $F'$ ? Il
revient au même de demander si la sous-catégorie strictement pleine
de $F$ formée par ces objets $X$ est un \emph{sous-topos} de $F$

\item Soient $X$ un objet de $E$, $X'=i^{\ast}(X)$, et
$i_{/X}:E'_{/X'}\rightarrow E_{/X}$ le morphisme de topos induit par
$i$ \eqhyperref[IV.eq5.10.1]{5.10.1}. Montrer que si $i$ est un plongement, il en est
de même de $i_{/X}$. (Utiliser d) et \hyperref[IV.prop5.11]{5.11}).
Réciproquement, si $(X_{\alpha})_{\alpha}$ est une famille d'objets
couvrant l'objet final, et si pour tout $\alpha$, \sisi{$i_{/X\alpha}$}{$i_{/X_\alpha}$} est
% correction SGA4-CORR-0043 / SGA-SRC-000430: SGA4-EIV-04-L5481-L5487-CC92B56E8270.
un plongement, il en est de même de $i$.
\end{enumerate}
\end{enonce*}

\begin{enonce*}[remark]{Exercice 9.1.7}\etiquette[9.1.7]{IV.exer9.1.7}
(\emph{Sous-topos et ensembles multiplicatifs de flèches}).

\begin{enumerate}
\renewcommand{\theenumi}{\alph{enumi}}
\renewcommand{\labelenumi}{\theenumi)}
\item Soient $i:E'\rightarrow E$ un morphisme de topos, $S$ l'ensemble des
flèches $u$ de $E$ telles que $i^{\ast}(u)$ soit un isomorphisme, et
considérons le foncteur canonique induit par $i^{\ast}$ (\cf [2,
Chap I] ou \sisi{VI}{la preuve de \hyperrefext[VI][VI.prop6.2]{6.2}}):
\begin{equation*}
\rho:E[S^{-1}]\longrightarrow E'{\quoi}.\tag{9.1.7.1}
\end{equation*}\etiquette[9.1.7.1]{IV.eq9.1.7.1}

Montrer que $i$ est un plongement si et seulement si le foncteur
précédent $\rho$ est une équivalence de catégories, et dans ce cas
$S$ admet un calcul des fractions à gauche et un calcul des
fractions à droite. (Utiliser [2, Chap 1 1.3]). Dans ce cas, \sisi{l'image}{l'image}
% correction SGA4-CORR-0044 / SGA-SRC-000431: SGA4-EIV-04-L5505-L5512-27FBADCE4CE2.
essentielle de $i_{\ast}$ se retrouve à partir de $S$ comme formée
des $X\in \Ob \sisi{\not{E}}{(E)}$ tels que pour tout\pageoriginale
$u:Y\rightarrow Y'$ dans $S$ $\Hom(u,X):\Hom(Y',X)\rightarrow
\Hom(Y,X)$ soit bijective.

\item Nous supposons désormais que i est un plongement. Montrer que
pour tout topos $F$, le foncteur
$$
f^{\ast}\longmapsto f^{\ast}i^{\ast}:\Hhom(E',F)\longrightarrow
\Hhom(E,F)
$$
induit une équivalence entre la catégorie des $f^{\ast}$ provenant
de morphismes de topos $f:F\rightarrow E'$ (\ie exacts à gauche et
admettant un adjoint à droite) et la catégorie des foncteurs
$g^{\ast}:E\rightarrow F$ provenant de morphismes de topos
$g:F\rightarrow E$, et tels que $g^{\ast}$ transforme objets de $S$
en isomorphismes.

\item Soit $S$ une partie de $Fl$ $E$, où $E$ est un
$𝒰$-topos. Montrer que $S$ est associé à un sous-topos $E'$
de $E$ par le procédé de a) si et seulement si $S$ satisfait aux
conditions suivantes:
\begin{itemize}
\item[ST 1)] $S$ contient les flèches inversibles, est stable par
composition et par changement de base.

\item[ST 2)] Si un composé vu est dans $S$, alors $u\in S$ si et
seulement si $v\in S$.

\item[ST 3)] Si $u:X\rightarrow Y$ est tel qu'il existe une famille
\sisi{couvrante $Y_i\rightarrow Y$, $i\in I$, tel que}{couvrante $Y_i\rightarrow Y$, $i\in I$, telle que}
% correction SGA4-CORR-0045 / SGA-SRC-000432: SGA4-EIV-04-L5538-L5541-2E4DACFD35C6.
$u_i:X_i=\fprod{X}{Y_i}{Y}\to Y_{i}$ est dans $S$ pour tout $i\in
I$, alors $u\in S$.
\end{itemize}

Montrer qu'on obtient ainsi une correspondance biunivoque entre
l'ensemble des sous-topos $E'$ de $E$, et l'ensemble des parties $S$
de $Fl$ $E$ satisfaisant aux conditions ST 1), ST 2), ST 3). ($S$
étant donné, lui associer une $𝒰$-topologie $T'$, plus fine
que la topologie canonique $T$ de $E$, dont les familles couvrantes
$(X_i\rightarrow X)_{i\in I}$ sont celles pour lesquelles
l'inclusion $X'\hookrightarrow X$ de l'image des $X_i$ est $\in S$.
Montrer que la catégorie des faisceaux pour $T'$ est équivalente à
un sous-topos $E'$ de $E$, et utiliser \hyperrefext[II][II.theo4.4]{4.4}. pour
prouver que celui-ci redonne $S$.

\item Montrer que $S$ est connu quand on connaît la partie $S_0$ de
$S$ formée des \emph{monomorphismes}. Montrer que l'application
$S\mapsto S_0$ établit une bijection entre l'ensemble des parties
$S$ de $Fl$ $E$ satisfaisant les conditions ST 1) à ST 3) de c), et
l'ensemble des parties de $Fl$ $E$ formées de
monomorphismes\pageoriginale et satisfaisant aux mêmes conditions ST
1) à ST 3).

\item Soit $(E'_i)_{i\in I}$ une famille de sous-topos de $E$, et
$(S_i)_{i\in I}$ la famille des parties correspondantes de $Fl$ $E$.
Montrer que $\bigcap_{i\in I} S_i$ satisfait aux conditions ST 1) à
ST 3) de c), donc correspond à un sous-topos $E'$ de $E$, et que ce
dernier est le \emph{sous-topos engendré par les} $E_i$, \ie
\emph{la borne supérieure} (\cf \hyperref[IV.rem9.1.2]{9.1.2} c)) des sous-topos
$E'_i$.

\item Soit $A$ une partie de $Fl$ $E$. Montrer qu'il existe une plus
petite partie $S$ de $Fl$ $E$ parmi celles contenant $A$ et
satisfaisant aux conditions ST 1) à ST 3), et que $S$ correspond au
plus grand sous-topos $E'$ de $E$ tel que la partie multiplicative
de $Fl$ $E$ \sisi{correspondante contienne $S$}{correspondante contienne $A$}. Montrer que $S$ est égale
% correction SGA4-CORR-0046 / SGA-SRC-000433: SGA4-EIV-04-L5571-L5575-EB91D3DEBFCD.
à la réunion des parties $A_n$ $(n\geq 0)$ de $Fl$ $E$, construites
par récurrence de la façon suivante: $A_0=A$, $A_{n+1}$ est
formée des flèches $u:X\rightarrow Y$ qui sont \sisi{d'un des trois types}{d'un des quatre types}
% correction SGA4-CORR-0047 / SGA-SRC-000434: SGA4-EIV-04-L5577-L5592-B6B4A5FEB546.
suivants (i) à (iv):
\begin{enumerate}
\renewcommand{\theenumii}{\roman{enumii}}
\renewcommand{\labelenumii}{(\theenumii)}
\item $u$ est inversible, ou le composé de deux flèches de $A_n$,
ou déduit par changement de base d'une flèche de $A_n$.

\item $u$ est l'un des facteurs d'\sisi{une flèche composées}{une flèche composée} dont
% correction SGA4-CORR-0048 / SGA-SRC-000435: SGA4-EIV-04-L5583-L5587-CCE44DE2FF07.
l'autre facteur ainsi que le composé sont dans $A_n$.

\item $u$ est une somme d'une petite famille de flèches $\in
A_n$.

\sisi{\item Il existe}{\item Il existe} un épimorphisme $Y'\rightarrow Y$, tel que
% correction SGA4-CORR-0049 / SGA-SRC-000436: SGA4-EIV-04-L5589-L5594-92C604343DB3.
$u':X'=\fprod{X}{Y'}{Y}\rightarrow Y'$ déduit de $u$ par changement
de base soit dans $A_n$.
\end{enumerate}

\item Avec les notations de e), soient $A$ la réunion des $S_i$, $S$
la partie de $Fl$ $E$ déduite de $A$ par le procédé de f), et $E'$
le sous-topos correspondant de $E$. Montrer que $E'$ est
l'intersection (= borne inférieure) des sous-topos $E'_i$.
\end{enumerate}
\end{enonce*}

\begin{enonce*}[remark]{Exercice 9.1.7.2}\etiquette[9.1.7.2]{IV.exer9.1.7.2}
(\emph{Factorisation canonique d'un morphisme de topos}).
\begin{enumerate}
\renewcommand{\theenumi}{\alph{enumi}}
\renewcommand{\labelenumi}{\theenumi)}
\item Soit $f:F\longrightarrow E$
un morphisme de topos. Montrer qu'on peut trouver une factorisation
de $f$, à isomorphisme près, en
\begin{equation*}
F\xrightarrow{f'}E'\xrightarrow{i} E,\tag{9.1.7.3}
\end{equation*}\etiquette[9.1.7.3]{IV.eq9.1.7.3}
où $f'$ est \emph{conservatif} (\ie (\hyperref[IV.subsubsec6.4.0]{6.4.0})
${f'}^{\ast}$ est conservatif, ou encore fidèle)\pageoriginale et où
$i$ est un plongement, et que cette factorisation est unique à
équivalence près (\cf \hyperref[IV.prop9.1.4]{9.1.4}). Pour ceci, introduire
l'ensemble $S_f\subset Fl$ $E$ des flèches $u$ de $E$ telles que
$f^{\ast}(u)$ soit un isomorphisme, prouver que $S_f$ satisfait aux
conditions de 9.1.7. c), prendre pour $E'$ le sous-topos
correspondant de $E$, et définir ${f'}^{\ast}$ par (9.1.7.1), qui
correspond à un morphisme de topos $f'$ grâce à 9.1.7 b).

\item Le sous-topos de $E$ défini par le plongement $i$ est appelé le
sous-topos de $E$ \emph{image du morphisme de topos} $f$. Montrer
que c'est le plus petit des sous-topos $E_1$ de $E$ tels que $f$ se
factorise, à isomorphisme près, à travers $E_1$, \ie
(\hyperref[IV.prop9.1.4]{9.1.4}) le plus petit des sous-topos de $E$ contenant
l'image de $f_{\ast}$ (ou encore, l'image essentielle de
$f_{\ast}$). Pour que $f$ soit un plongement, il faut et il suffit
que $f$ induise une équivalence $f'$ de $E$ avec son image; pour que
$f$ soit conservatif, il faut et il suffit que son image soit égale
à $E$ tout entier.

\item Supposons que $f=\Top(f_0)$ soit associé à une application
continue d'espaces topologiques sobres $f_0:Y\rightarrow X$
(\hyperref[IV.subsec4.1]{4.1}) Considérons la factorisation canonique de
$f_0$ en
$$
Y\xrightarrow{{f'}_0}X'=f_0(Y)\xrightarrow{i_0}X{\quoi},
$$
où $i_0$ est l'inclusion. Montrer que la factorisation canonique de
$f$ s'identifie alors à la factorisation correspondante
\[
\xymatrix@C=2cm@M=.2cm{ \Top(Y) \ar[r]^{\Top(f'_0)}& \Top(X')
\ar@{^{(}->}[r]^{\Top(i_0)}& \Top(X). }
\]
Soit $Z$ le plus petit sous-espace sobre de $X$ contenant \sisi{$X'=f(Y)$}{$X'=f_0(Y)$},
\ie l'ensemble des points $x$ de $X$ tels que $\{\bar{x}\}\cap X'$
soit dense dans $\{\bar{x}\}$. (On notera que $Z=X'$ si les points
de $X$ sont fermés, ou si $X'$ est une partie \emph{constructible}
de $X$ supposé localement noethérien). Alors $f$ est conservatif,
\ie le sous-topos de $E=\Top(X)$ image de $F=\Top(Y)$ par
$f=\Top(f_0)$ est égal à $E$ lui-même, si et seulement si $Z=X$, \ie
si et seulement si \sisi{$f(Y)$ est}{$f_0(Y)$ est} \emph{très dense} (EGA IV 10.13) dans
% correction SGA4-CORR-0050 / SGA-SRC-000437: SGA4-EIV-04-L5636-L5656-24C579146AF1.
$X$. Ceci précise dans quelle mesure il est licite de regarder la
notion de morphisme conservatif de topos comme la généralisation
naturelle de la notion d'application continue \emph{surjective}
d'espaces topologiques.

\item Soit\pageoriginale
\[
\xymatrix{
Y \ar[d]& Y'\ar[l] \ar[d]\\
X & X'\ar[l] }
\]
un diagramme cartésien d'espaces topologiques. Si $Y$ et $X'$ sont
des espaces discrets, donc $Y'$ discret, alors le diagramme
correspondant de topos est 2-cartésien (\hyperref[IV.prop5.11]{5.11}). (NB.
Pour un contre-exemple lorsque l'on ne suppose plus $Y$, $X'$
discrets, $Y$, $X'$ étant des sous-espaces de l'espace séparé $X$,
\cf \hyperref[IV.exer9.1.10]{9.1.10} c)). En déduire un exemple d'un diagramme
2-cartésien de topos
\[
\xymatrix{
F \ar[d]_{f}& F'\ar[l] \ar[d]^{f'}\\
E & E'\ar[l] }
\]
tel que $f$ soit conservatif, mais $f'$ non conservatif, plus
précisément $E'=$ topos ponctuel, $F'=$ topos vide
(\hyperref[IV.subsec2.2]{2.2}). (Prendre $Y$ tel que l'image de $Y$ dans $X$
soit très dense et $\neq X$, et prendre $X'$ réduit à un point qui
n'est pas dans l'image).

\item Soit $(E_i)_{i\in I}$ une famille de sous-topos d'un topos
$E$. Prouver que $E$ est le sup des $E_i$ si et seulement si la
famille des morphismes de plongement $E_i\rightarrow E$ est
conservative (\hyperref[IV.subsubsec6.4.0]{6.4.0}).
\end{enumerate}
\end{enonce*}

\begin{enonce*}[remark]{Exercice 9.1.8}\etiquette[9.1.8]{IV.exer9.1.8}
(\emph{Sous-topos et points de $E$. Cas des espaces topologiques}).

Soit $E$ un topos.
\begin{enumerate}
\renewcommand{\theenumi}{\alph{enumi}}
\renewcommand{\labelenumi}{\theenumi)}
\item Soit $Z$ une sous-catégorie pleine de $\Ppoint(E)$ (ou ce qui
revient au même: une partie de $\Ob \Ppoint(E)$). Montrer que
l'ensemble $S(Z)$ des flèches de $E$ qui sont transformées en
bijections par les foncteurs fibre correspondants aux $z\in \Ob Z$
satisfait aux conditions de 9.1.7. c), et définit donc un sous-topos
$E_Z$ de $E$. Montrer que celui-ci admet suffisamment de points, et
que $Z$ est équivalente à une sous-catégorie pleine de
$\Ppoint(E_Z)$. Montrer qu'on peut obtenir par le procédé précédent
tout sous-topos $E'$ de $E$ qui admet suffisamment de points.
(Prendre pour $Z$ l'image essentielle du foncteur (9.1.5.1.)).
Montrer qu'on obtient\pageoriginale une bijection entre l'ensemble
des sous-topos $E'$ de $E$ admettant suffisamment de points, et un
sous-ensemble de l'ensemble des sous-catégories strictement pleines
$Z$ de $\Ppoint(E)$ qui sont stables par facteurs directs
(\sisi{\hyperrefext[I][I.sec10]{10}}{\hyperrefext[I][I.sec10.6]{10.6}}) et par petites limites projectives filtrantes (\cf
\hyperref[IV.subsec7.3]{7.3}.). (Problème: Quelles sont les sous-catégories de
$\Ppoint(E)$ qu'on trouve ainsi ? C'est essentiellement, sous une
autre forme, le problème déjà soulevé dans 6.5.3., compte-tenu du
dictionnaire \hyperref[IV.rem9.1.2]{9.1.2} c). Noter qu'en plus des conditions
nécessaires signalées à l'instant, il y a aussi des conditions plus
subtiles sur la nature topologique de $\Point(E)$, genre, sobriété,
\cf e)).

\item Soit $(Z_i)_{i\in I}$ une famille de sous-catégories pleines de
$\Ppoint(E)$, et soit $Z$ la sous-catégorie pleine réunion des
$Z_i$. Montrer que le sous-topos correspondant $E_Z$ défini en a)
est le sous-topos de $E$ engendré par les $E_{Z_i}$. (Utiliser
9.1.7. e)). En conclure que le sous-topos de $E$ engendré par une
famille de sous-topos ayant suffisamment de points, a également
suffisamment de points.

\item Soit $f:E'\rightarrow E$ un plongement de topos. Montrer que
l'application $U\mapsto f^{\ast}(U)$
\begin{equation*}
\Ouv(E)\longrightarrow \Ouv(E')\tag{9.1.8.1}
\end{equation*}\etiquette[9.1.8.1]{IV.eq9.1.8.1}
est surjective, et en conclure que l'application
\begin{equation*}
\Point(E')\longrightarrow \Point(E)\tag{9.1.8.2}
\end{equation*}\etiquette[9.1.8.2]{IV.eq9.1.8.2}
induite par $f$ sur les classes d'isomorphie de \sisi{points induit un}{points est un}
% correction SGA4-CORR-0051 / SGA-SRC-000438: SGA4-EIV-04-L5732-L5742-9BC865D9157B.
homéomorphisme de $\Point(E')$ sur un sous-espace de $\Point(E)$,
pour les topologies définies dans \hyperref[IV.exer7.8]{7.8} a). (Pour
l'injectivité, utiliser 9.1.5.).

\item Soit $f:X'\rightarrow X$ une application continue, d'où un
morphisme de topos (\hyperref[IV.subsec4.1]{4.1})
$$
\Top(f):\Top(X')\longrightarrow \Top(X){\quoi}.
$$
Prouver que ce dernier est un plongement si et seulement si la
topologie de $X'$ est image inverse par $f$ de celle de $X$ (\ie,
lorsque $X'$ est un espace de Kolmogoroff, par exemple un espace
sobre, si et seulement si $f$ induit un homéomorphisme de $X'$ sur
un sous-espace de $X$). (Pour la nécessité,\pageoriginale utiliser
c), pour la suffisance, \hyperref[IV.exer9.1.7.2]{9.1.7.2} c)).

\item Soit $X$ un espace topologique sobre, et $E=\Top(X)$. Utilisant
a) et d), établir un isomorphisme d'ensembles ordonnés $X'\mapsto
E_{X'}$ entre l'ensemble des sous-espaces sobres $X'$ de $X$ (\ie,
lorsque $X$ est séparé, entre l'ensemble de \emph{toutes} les
\sisi{parties de $X$ et l'ensemble}{parties de $X$) et l'ensemble} des sous-topos de $E$ ayant
% correction SGA4-CORR-0052 / SGA-SRC-000439: SGA4-EIV-04-L5757-L5762-4490514F6A8F.
suffisamment de points. Montrer que $E_{X'}$ est équivalent à
$\Top(X')$. Si $(X'_i)_{i\in I}$ est une famille de sous-espaces
sobres de $X$, montrer qu'il existe un plus petit sous-espace sobre
$X'$ de $X$ contenant les $X'_i$, égal à la réunion des $X'_i$ si I
est fini ou $X$ est séparé, et que $E_{X'}$ est le sous-topos de $E$
engendré par les $E_{X'_i}$. (Utiliser b)).

\item Donner un exemple d'un topos $E$ de la forme $\Top(X)$, et d'un
sous-topos $E'$ qui n'a pas suffisamment de points, et même qui est
non vide (\hyperref[IV.subsec2.2]{2.2}) sans points. (Prendre $E$ de la forme
$\widehat{U}$, $E'=U^{\sim}$, où $U$ est l'ensemble ordonné défini
dans \hyperref[IV.subsec7.4]{7.4}. Ou mieux, prendre l'exemple de \sisi{9.1. 10. b)}{\hyperref[IV.exer9.1.10]{9.1.10} b)},
qui montre qu'on peut prendre pour $X$ n'importe quel espace compact
non vide sans point isolé).
\end{enumerate}
\end{enonce*}

\begin{enonce*}[remark]{Exercice 9.1.9}\etiquette[9.1.9]{IV.exer9.1.9}
(\emph{Cas des topos définis par un ensemble ordonné}).

\begin{enumerate}
\renewcommand{\theenumi}{\alph{enumi}}
\renewcommand{\labelenumi}{\theenumi)}
\item On utilise le yoga de \hyperref[IV.exer7.8]{7.8} b) et c), donnant un dictionnaire
entre, d'une part les topos engendrés par leurs ouverts et les
morphismes de tels topos, d'autre part \sisi{les ensembles ordonnées}{les ensembles ordonnés} ayant
% correction SGA4-CORR-0053 / SGA-SRC-000440: SGA4-EIV-04-L5785-L5791-FDDCC930003C.
des $\Sup$ quelconques, des Inf finis et où le Inf de deux éléments %% uniformiser notations sup/Sup/\Sup etc.
est distributif par rapport aux $\Sup$  quelconques, les morphismes
entre de tels ensembles ordonnés étant les applications croissantes
commutant aux Inf finis et aux $\Sup$ quelconques.

\item Soit $E$ un topos, $f:E'\rightarrow E$ un plongement de
topos. Montrer que pour toute famille génératrice $(X_i)_{i\in I}$
dans $E$, la famille $(f^{\ast}(X_i))_{i \in I}$ est génératrice. En
déduire que si $E$ est engendré par ses ouverts, il en est de même
de $E'$ (Utiliser \hyperref[IV.exer9.1.8]{9.1.8}. c)).

\item Soit $f:E'\rightarrow E$ un morphisme de topos engendrés par
leurs ouverts, associé à un morphisme $g:\Oo\rightarrow \Oo'$ sur
les ensembles ordonnés correspondants. Montrer que $f$ est un
plongement si et seulement si $g$ est surjectif, \emph{et}
$g$\pageoriginale est surjectif sur les graphes des relations
d'ordre, \ie l'ordre de $\Oo'$ est quotient de celui de $\Oo$.

\item Soit $E$ un topos de la forme $\Oo^{\sim}$ comme dans
a). Déduire de b) et c) une correspondance biunivoque entre
l'ensemble des sous-topos de $E$, et l'ensemble des relations
d'équivalence $R$ dans $\Oo$ ayant les propriétés suivantes,
analogues aux conditions ST 1) à ST 3) de 9.1.7 c):
\begin{itemize}
\item[Q0 1)] La relation $R$ est compatible avec les Inf finis, ou
encore: si $U$, $U'$ sont équivalents par $R$, il en est de même de
$U\cap V$ et $U'\cap V$ pour tout $V\in \Oo$.

\item[Q0 2)] La relation $R$ est compatible aux $\Sup$ quelconques,
ou encore: si $U_i$ est équivalent à $U'_i$ pour tout $i\in I$,
alors $\Sup_i U_i$ est équivalent à $\Sup_i U'_i$.
\end{itemize}

Montrer que si des sous-topos \sisi{$E^{!}_i$}{$E'_i$} de $E$ correspondent aux
relations d'équivalence $R_i$, alors le sous-topos engendré par les
\sisi{$E^{!}_i$}{$E'_i$} correspond à la relation intersection des~$R_i$.
% correction SGA4-CORR-0054 / SGA-SRC-000441: SGA4-EIV-04-L5821-L5823-4FF74F776D91.
\end{enumerate}
\end{enonce*}

\begin{enonce*}[remark]{Exercice 9.1.10}\etiquette[9.1.10]{IV.exer9.1.10}
(\emph{Intersection de sous-espaces; nouveaux topos sans points; non
existence du complémentaire d'un sous-topos}).

Soient $X$ un espace topologique, $E=\Top(X)=\Ouv(X)^{\sim}$
(\hyperref[IV.subsec2.1]{2.1}).

\begin{enumerate}
\renewcommand{\theenumi}{\alph{enumi}}
\renewcommand{\labelenumi}{\theenumi)}
\item Soit $R$ la relation suivante sur $\Ouv(X):(U,U')\in
R\Longleftrightarrow U\cap U'$ est dense dans $U$ et dans $U'$.
Montrer que $R$ est une relation d'équivalence satisfaisant les
conditions Q0 1) et Q0 2) de \hyperref[IV.exer9.1.9]{9.1.9} d), et définit donc
un sous-topos $E'$ de $E$. Montrer que pour un point $x\in X$, le
point correspondant est dans l'image essentielle de $\Ppoint(E')$ si
et seulement si $x$ est \emph{épais} \ie appartient à l'adhérence de
l'intérieur de $\bar{x}$. Montrer que $E'$ est le topos vide
(\hyperref[IV.subsec2.2]{2.2}) si et seulement si $X$ est vide.

\item Si $X$ est sobre, montrer que la catégorie $\Ppoint(E')$ est
équivalente à la catégorie définie par le sous-ensemble ordonné de
$X$ (pour la relation de spécialisation) défini par les points épais
de $X$. Montrer que si les points de $X$ sont fermés (\resp si $X$
est noethérien) alors les points épais de $X$ sont les points isolés
de $X$ \sisi{\ie tel que}{\ie tels que} $\{x\}$ soit ouvert (\resp sont les points
% correction SGA4-CORR-0055 / SGA-SRC-000442: SGA4-EIV-04-L5847-L5853-851C52DCAA24.
maximaux de $X$). En conclure que si $X$ est\pageoriginale un espace
topologique séparé non vide sans points isolés, alors le sous-topos
$E'$ de $E=\Top(X)$ est un topos non « vide »  qui n'a pas de
points.

\item Soit $X'$ une partie de $X$, et $E_{X'}$ le sous-topos de $E$
associé à $X'$ par le procédé \hyperref[IV.exer9.1.8]{9.1.8}. a). Montrer que si
$X'$ est dense, alors $E_{X'}$ contient $E'$. Montrer que si $X$ est
sobre et si $X'$ et $X''$ sont deux parties sobres de $X$, leur
intersection $X'\cap X''$ est sobre et $E_{X'\cap X''}\subset
E_{X'}\cap E_{X''}$, et que l'inclusion peut être stricte
\footnote[1]{Voir cependant \hyperref[IV.exer9.1.11]{9.1.11} b) ci-dessous pour
le cas $X$ localement noethérien.}. (Prendre un espace séparé $X$
non vide admettant deux parties disjointes denses $X'$, $X''$).
Montrer que l'inclusion précédente est une égalité si et seulement
si le sous-topos $E_{X'}\cap E_{X''}$ a assez de points. Montrer
qu'il en est ainsi si $X'$ ou $X''$ est une partie \emph{localement
fermée} de~$X$.

\item Soit $X'$ une partie de $X$, et $X''$ le complémentaire de
$X'$. Un sous-topos $F$ de $E$ contient les $E_{\{x\}}$ pour $x\in
X''$ si et seulement si il contient $E_{X''}$ (\cf
\hyperref[IV.exer9.1.8]{9.1.8} b)). En conclure que si $X''$ est réunion de
parties localement fermées de $X$ (par exemple si les points de
$X''$ sont fermés), alors l'ensemble des sous-topos $F$ de $E$ dont
l'intersection avec $E_{X'}$ est le sous-topos vide $\Top(∅)$
\sisi{contient un plus grand élément}{ne contient un plus grand élément} (qui mérite alors le nom de
% correction SGA4-CORR-0056 / SGA-SRC-000443: SGA4-EIV-04-L5872-L5884-B8F065F7E378.
sous-topos de $E$ \emph{complémentaire faible de} $E_{X'}$ \cf
\hyperref[IV.exer9.1.13]{9.1.13}), \emph{que si} $E_{X'}\cap E_{X''}$ est le
sous-topos vide (condition qui n'est pas satisfaite si $X\neq ∅$
et si $X'$ et $X''$ sont tous deux denses dans $X$, $X$, $X'$, $X''$
étant sobres, \cf c)). Montrer que lorsque cette condition est en
défaut, alors dans l'ensemble ordonné des sous-topos de $E$, le Inf
de deux éléments n'est pas distributif par rapport aux $\Sup$
quelconques. (Le rédacteur ignore s'il est distributif par rapport
aux $\Sup$ finis; \cf cependant \hyperref[IV.exer9.1.11]{9.1.11} e) et
\hyperref[IV.exer9.1.12]{9.1.12} d)).
\end{enumerate}
\end{enonce*}

\begin{enonce*}[remark]{Exercice 9.1.11}\etiquette[9.1.11]{IV.exer9.1.11}
(\emph{Sous-topos des topos localement noethériens. Cas des espaces
noethériens}).
\begin{enumerate}
\renewcommand{\theenumi}{\alph{enumi}}
\renewcommand{\labelenumi}{\theenumi)}
\item Soit $f:E'\rightarrow E$ un plongement de topos. Montrer que si
$X$ est un objet prénoethérien de $E$ (\ie (\sisi{VI 1.}{\hyperrefext[VI][VI.def1.30]{1.30}}) toute suite
croissante de sous-objets\pageoriginale \sisi{de $X$ est stationnaire}{de $X$ est stationnaire})
% correction SGA4-CORR-0057 / SGA-SRC-000444: SGA4-EIV-04-L5899-L5902-FDA2DB226330.
alors $X'=f^{\ast}(X)$ est un objet prénoethérien de $E'$.
(Introduisant les topos induits ${E'}_{/X'}$ et $E_{/X}$, et
utilisant \hyperref[IV.exer9.1.6]{9.1.6} e), se ramener d'abord au cas où $X$
est l'objet final $e$ de $E$, donc $X'$ est l'objet final $e'$ de
$E'$. Noter ensuite que toute suite croissante de sous-objets de
$e'$ définit, par application de $f_{\ast}$, une suite croissante de
sous-objets de $f_{\ast}(e')=e$). En conclure que si $E$ admet une
famille génératrice formée d'objets prénoethériens (\resp si $E$ est
un topos noethérien, \resp localement noethérien (\sisi{VI 2.)}{\hyperrefext[VI][VI.def2.11]{2.11}}) alors $E'$
satisfait à la même condition. (Utiliser \hyperref[IV.exer9.1.9]{9.1.9} b)).

\item Soit $E$ un topos localement noethérien (\sisi{VI 2.}{\hyperrefext[VI][VI.def2.11]{2.11}}). Montrer que
tout sous-topos de $E$ est de la forme $E_Z$ (\hyperref[IV.exer9.1.8]{9.1.8}
a)), pour une sous-catégorie pleine convenable $Z$ de $\Ppoint(E)$.
(Utiliser a) et le théorème de \textsc{Deligne} que tout topos
localement noethérien admet suffisamment de points). En conclure que
si $E$ est de la forme $\Top(X)$, où $X$ est un espace topologique
sobre localement noethérien, alors $Z\mapsto E_Z$ est un
isomorphisme de l'ensemble ordonné des sous-espaces sobres de $X$,
avec l'ensemble ordonné \sisi{des sous-topos de $X$}{des sous-topos de $E$}. (Utiliser
% correction SGA4-CORR-0058 / SGA-SRC-000445: SGA4-EIV-04-L5918-L5922-2742A8C660C7.
\hyperref[IV.exer9.1.8]{9.1.8} e)). L'intersection de sous-espaces sobres $Z_i$
\sisi{(est nécessairement sobre et)}{est nécessairement sobre et} définit l'intersection des sous-topos
% correction SGA4-CORR-0059 / SGA-SRC-000446: SGA4-EIV-04-L5922-L5924-8C785513FF6A.
$E_{Z_i}$ (contrairement à ce qui peut se passer dans le cas général
(\hyperref[IV.exer9.1.10]{9.1.10} c)).

\item Soient $X$ un espace sobre localement noethérien, $X'$ une
partie de $X$, $X''$ son complémentaire. Montrer que les conditions
suivantes sont équivalentes:
\begin{enumerate}
\renewcommand{\theenumii}{\roman{enumii}}
\renewcommand{\labelenumii}{(\theenumii)}
\item $X'$ est constructible.

\item $X'$ et $X''$ sont sobres.

\item $X'$ est sobre, et le sous-topos $E_{X'}\approxeq \Top(X')$ de
$E=\Top(X)$ admet un sous-topos « complémentaire »  (\cf
\sisi{\hyperref[IV.exer9.1.13]{9.1.13} e)}{\hyperref[IV.exer9.1.13]{9.1.13}}.

\item L'intersection des sous-topos $E_{X'}\approx\Top(X')$ et
$E_{X''}\approx\Top(X'')$ de $E=\Top(X)$ est le sous-topos
« vide » de $E$.
\end{enumerate}

Montrer que si $E'$ et $E''$ sont deux sous-topos de $E$, ils sont
complémentaires faibles \sisi{si et seulement si ils}{si et seulement s'ils} sont complémentaires
% correction SGA4-CORR-0060 / SGA-SRC-000447: SGA4-EIV-04-L5946-L5948-64D4A5838458.
(\hyperref[IV.exer9.1.13]{9.1.13}).

\item Soit\pageoriginale $X$ un espace sobre noethérien. Montrer
l'équivalence des conditions suivantes:
\begin{itemize}
\item[(i)] Tout sous-espace sobre $X'$ de $X$ est constructible, \ie
(c)) son complémentaire est sobre.

\item[(i bis)] Pour tout point maximal $x$ de $X$, $X-\{x\}$ est
sobre, \ie $\{x\}$ est constructible.

\item[(ii)] $X$ est fini.

\item[(iii)] Tout sous-topos de $\Top(X)$ admet un sous-topos
complémentaire (\cf \hyperref[IV.exer9.1.13]{9.1.13}).

\item[(iv)] Dans l'ensemble ordonné des sous-topos de $\Top(X)$, le
Inf de deux éléments est distributif par rapport aux $\Sup$
quelconques.
\end{itemize}
(Utiliser l'argument de \hyperref[IV.exer9.1.10]{9.1.10} d), grâce à la formule
d'intersection prouvée en b)).

\item Soit $X$ un espace localement noethérien. Montrer que dans
l'ensemble ordonné des sous-topos de $\Top(X)$, le Inf de deux
éléments est distributif par rapport aux $\Sup$ \emph{finis}. (Se
ramener au cas $X$ sobre, et utiliser b)).

\item Pour tout sous-espace sobre $Z$ de l'espace sobre localement
noethérien $X$, soit $T_Z$ la topologie sur la catégorie $\Ouv(X)$
des ouverts de $X$ pour laquelle une famille d'inclusions
$U_i\rightarrow U$ est couvrante si et seulement si on a $U\cap
Z=\bigcup_i U_i \cap Z$. Montrer que $Z\mapsto T_Z$ établit une
bijection (renversant les relations d'inclusion) entre l'ensemble
des sous-espaces sobres $Z$ de $X$, et l'ensemble des topologies sur
$\Ouv(X)$ plus fines que la topologie canonique $T_X$. (Utiliser b)
et \hyperref[IV.rem9.1.2]{9.1.2} c)).

\item Soit $X$ un espace topologique sobre. Une partie $X'$ de $X$ est
sobre si et seulement si son complémentaire est réunion de parties
localement fermées; si $X$ est localement noethérien, cela signifie
aussi que $X'$ est « proconstructible »  (EGA IV 1.9.4). Montrer
qu'il existe sur $X$ une topologie $T_{\cons'}$, de telle façon
que les parties fermées de $X_{\cons'}=(X,T_{\cons'})$ soient les
parties sobres de $X$. Supposons désormais $X$ localement
noethérien, $X_{\cons'}$ n'est donc autre que l'espace
$X_{\cons}$\pageoriginale de EGA IV 9.1.13. Déduire alors de b) et c)
que l'ensemble ordonné $\Sigma$ des sous-topos de $\Top(X)$ est
isomorphe à l'ensemble ordonné des parties fermées de $X_{\cons}$,
et dans cette correspondance, les sous-topos admettant un
complémentaire correspondent aux parties à la fois ouvertes et
fermées de $X_{\cons}$. Montrer que \sisi{$X$}{$X_{cons}$} est localement compact, et
qu'il est compact si et seulement si $X$ est noethérien \ie le topos
\sisi{$\Top(E)$}{$\Top(X)$} est noethérien.

\item Soit $E$ un topos localement noethérien. Supposons qu'il ait la
propriété envisagée dans \hyperref[IV.subsubsec9.1.14]{9.1.14} b) (condition
peut-être toujours remplie\ldots). Montrer que les images
essentielles des foncteurs canoniques $\Ppoint(E')\rightarrow
\Ppoint(E)$, où $E'$ parcourt l'ensemble $\Sigma$ des sous-topos de
$E$, définissent dans l'ensemble $X=\Point(E)$ des classes
d'isomorphie de points de $E$ un ensemble de parties, qui est
l'ensemble des parties fermées pour une \emph{topologie} $T_{\cons}$
\emph{sur} $X$, et qu'on obtient de cette façon un isomorphisme
d'ensembles ordonnés entre l'ensemble $\Sigma$ des sous-topos de $E$
et l'ensemble des parties fermées de $X_{\cons}$. En conclure que
dans $\Sigma$, le $\Sup$ de deux éléments est distributif par
rapport aux Inf \emph{quelconques}. Montrer que $X_{\cons}$, qui
n'est pas nécessairement $𝒰$-petit, a un espace sobre
associé (\hyperref[IV.subsubsec4.2.1]{4.2.1}) qui est $𝒰$-petit.
\end{enumerate}
\end{enonce*}

\begin{enonce*}[remark]{Exercice 9.1.12}\etiquette[9.1.12]{IV.exer9.1.12}
(\emph{Topos finis}). Un topos $E$ est dit \emph{fini} s'il est
équivalent à un topos de la forme $\widehat{C}$, avec $C$ une
catégorie finie.
\begin{enumerate}
\renewcommand{\theenumi}{\alph{enumi}}
\renewcommand{\labelenumi}{\theenumi)}
\item (Dictionnaire). Soit $𝒱$ un univers tel que $𝒰\in 𝒱$. Soit
(Karfiness) la 2-catégorie définie comme 2-sous-catégorie pleine de
la catégorie $(\Cat)_{𝒱}$, formée des catégories éléments de
$𝒱$, qui sont \emph{Karoubiennes et équivalentes à une
catégorie finie} et (Topfin) la 2-catégorie définie comme la
2-sous-catégorie pleine de
$(𝒱\text{-}𝒰\text{-}\Top)$ (\hyperref[IV.subsubsec3.3.1]{3.3.1})
formée des $𝒰$-topos \emph{finis} qui sont $\in 𝒱$.
Montrer qu'on a des 2-équivalences quasi-inverses l'une de l'autre:
\begin{align*}
C & \longmapsto \widehat{C}: (\text{Karfiness})\longrightarrow
(\text{Topfin})\tag{9.1.12.1}\\
E & \longmapsto \Ppoint(E): (\text{Topfin})\longrightarrow
(\text{Karfiness}).\tag{9.1.12.2.}
\end{align*}\etiquette[9.1.12.2]{IV.eq9.1.12.2}
(Utiliser \hyperref[IV.exer7.6]{7.6} h), e)).

\item Montrer\pageoriginale qu'un topos fini est noethérien \sisi{(VI 2)}{\hyperrefext[VI][VI.def2.11]{2.11}}.

\item Montrer que tout sous-topos d'un topos fini $E$ est un topos
fini. Plus précisément, si $E$ est équivalent à
$\widehat{C}$, où
$C$ est une catégorie karoubienne finie (ou, plus généralement,
équivalente à une catégorie finie), montrer qu'on obtient un
isomorphisme ordonné entre l'ensemble des sous-topos $E'$ de $E$ et
l'ensemble des sous-catégories strictement pleines $C'$ de $C$ qui
sont karoubiennes (ou, ce qui revient au même, stables dans $C$ par
facteurs directs), en associant à toute $C'$ l'image essentielle de
$f_{\ast}:\widehat{C}'\rightarrow \widehat{C}$, où $f:C'\rightarrow C$ est
le foncteur d'inclusion. (Utiliser \hyperref[IV.subsec5.6]{5.6}. pour
s'assurer que $f_{\ast}$ est pleinement fidèle, et b) et
\hyperref[IV.exer9.1.11]{9.1.11} b) pour le fait que tout sous-topos $E'$ de $E$
s'obtient comme indiqué).

\item Soit $E$ un topos fini. Construire sur l'ensemble fini
$X=\Point(E)$ des classes d'isomorphie de points de $E$ une
topologie qui en fasse un espace topologique sobre, et telle que
l'ensemble ordonné des sous-topos de $E$ soit canoniquement
isomorphe à l'ensemble des parties fermées de $X$. (Prendre la « topologie constructible »  $\Cc_{\cons}$ de $X$, définie via la
relation d'ordre $(x\leq y)\Leftrightarrow$ ($x$ est un facteur
direct de $y$), pour laquelle l'adhérence d'un point $x$ est formée
des $y$ tels que $y\leq x$). En particulier l'ensemble $\Sigma$ des
sous-topos de $E$ est fini, et \sisi{le Inf $y$ est distributif}{le Inf est distributif} par
% correction SGA4-CORR-0061 / SGA-SRC-000448: SGA4-EIV-04-L6062-L6071-4FFA8904FC39.
rapport aux $\Sup$ quelconques.

\item Soit $C$ une catégorie équivalente à une catégorie
finie. Montrer que pour toute sous-catégorie pleine $C'$ de $C$, il
existe sur $C$ une topologie $T_{C'}$ pour laquelle une famille
$(X_i\rightarrow X)_{i\in I}$ est couvrante si et seulement si pour
tout $Y\in \Ob C'$, la famille des applications
$\Hom(Y,X_i)\rightarrow \Hom(Y,X)$ $(i\in I)$ est surjective.
Montrer que la catégorie des faisceaux sur $(C,T_{C'})$ est
équivalente à $\widehat{C}'$, et que $C'\mapsto T_{C'}$ est une
bijection (renversant les relations d'ordre) entre l'ensemble des
sous-catégories strictement pleines $C'$ de $C$ qui sont
karoubiennes (ou, ce qui revient au même, stables dans $C$ par
facteurs directs), et l'ensemble des topologies (\hyperrefext[II][II.def1.1]{1.1})
sur la catégorie $C$. (Utiliser $c$) et \hyperref[IV.rem9.1.2]{9.1.2}
c)).\pageoriginale En particulier, si $C$ est un site dont la
catégorie sous-jacente est équivalente à une catégorie finie, alors
le topos $C^{\sim}$ est un topos fini.

\item Soit $f:C'\rightarrow C$ un foncteur de $𝒰$-catégories, avec
$C$ équivalente à une catégorie finie. Montrer que le morphisme de
topos $\widehat{f}:\widehat{C}'\rightarrow \widehat{C}$
(\hyperref[IV.subsec4.6]{4.6}) est conservatif si et seulement si tout objet
de $C$ est isomorphe à un facteur direct d'un objet dans l'image de
$f$. (Se ramener au cas où $f$ est une inclusion
$f:C'\hookrightarrow C$, avec $C'$ comme dans c)).

\item Soit $C$ la catégorie ayant deux objets $e$ (l'objet final) et
$a$, et, en plus des flèches identiques, trois flèches
$f:a\rightarrow e$, $g:e\rightarrow a$, $p=gf$ (donc $p^2=p$).
Montrer que $E=\widehat{C}$ a, en plus des sous-topos $E$ et
$\Top(∅)$, exactement un sous-topos $E'$, savoir
$E'=\widehat{\{e\}}$; par suite, pour tout sous-topos $E''$ de $E$,
$E''\neq \Top(∅)$, on a $E'\cap E''\neq \Top(∅)$.
\end{enumerate}
\end{enonce*}

\begin{enonce*}[remark]{Exercice 9.1.13}\etiquette[9.1.13]{IV.exer9.1.13}
(\emph{Sous-topos complémentaires d'un topos}).

Soient $E$ un topos, $E'$ et $E''$ deux sous-topos. On dit que $E'$
et $E''$ sont \emph{complémentaires} l'un de l'autre si $E'\cap
E''=\Top(∅)$, $E'\vee E''=E$ (où le signe $\vee$ désigne le
$\Sup$ dans l'ensemble ordonné des sous-topos de $E$).

On supposera que dans l'ensemble $\Sigma$ des sous-topos de $E$, le
Inf de deux éléments est distributif par rapport au $\Sup$ fini.

\begin{enumerate}
\renewcommand{\theenumi}{\alph{enumi}}
\renewcommand{\labelenumi}{\theenumi)}
\item Si $E'$ et $E''$ sont complémentaires, $E''$ est le plus grand
parmi les \sisi{sous-topos $E''_1$ de tels que}{sous-topos $E''_1$ de $E$ tels que} $E'\cap E''_1=\Top(∅)$
% correction SGA4-CORR-0062 / SGA-SRC-000449: SGA4-EIV-04-L6122-L6125-DC316178D087.
(\ie $E''$ est un complémentaire faible de $E'$, dans la
terminologie de \hyperref[IV.exer9.1.10]{9.1.10} d)).

\item Montrer que les conditions suivantes sur $E$ sont équivalentes:
\begin{itemize}
\item[(i)] Si $E'$ et $E''$ sont deux sous-topos de $E$ tels que $E''$
soit un complémentaire faible de $E'$, alors $E''$ est un
complémentaire de $E'$.

\item[(ii)] Pour tout sous-topos $E'$ de $E$ tel que $E'\neq E$, il
existe un sous-topos $E''$ de $E$ tel que $E''\neq \Top(∅)$ et
$E'\cap E''=\Top(∅)$.

\item[(i bis)] Si $E'$ et $E''$ sont deux sous-topos de $E$ tels que
$E''$ soit maximal dans l'ensemble des sous-topos $E''_1$ tels que
$E'\cap E''_1=\Top(∅)$, alors $E''$ est \sisi{un complémentaire de $E$.}{un complémentaire de $E'$.}
% correction SGA4-CORR-0063 / SGA-SRC-000450: SGA4-EIV-04-L6137-L6140-2CB79A78F39A.
\end{itemize}

Montrer que même si $E$ est un topos fini (\hyperref[IV.exer9.1.12]{9.1.12}),
ces\pageoriginale conditions ne sont pas nécessairement vérifiées
(\hyperref[IV.exer9.1.12]{9.1.12} g)). Elles le sont cependant si $E$ est de la
forme $\Top(X)$, $X$ un espace localement noethérien. (Utiliser
\hyperref[IV.exer9.1.11]{9.1.11} b)).

\item Un sous-topos $E'$ de $E$ est appelé \emph{complémenté} s'il
admet un complémentaire (comparer \hyperref[IV.exer9.1.12]{9.1.12} c)). Prouver
que l'ensemble des sous-topos complémentés est stable par Inf et
$\Sup$ finis, et par complémentarité, et que cette dernière
transforme Inf en $\Sup$, $\Sup$ en Inf.

\item Énoncer les duales des observations a), b), c) (qui étaient en
fait des énoncés généraux triviaux sur un ensemble ordonné abstrait
$\Sigma$; noter que l'hypothèse faite sur $\Sigma$ est en fait
autoduale).
\end{enumerate}
\end{enonce*}

\setcounter{subsubsection}{13}
\subsubsection{}\etiquette[9.1.14]{IV.subsubsec9.1.14}
\emph{Questions ouvertes}\nde{Des réponses positives aux questions a) et d) ont été apportées 
dans [17].}. Soit $E$ un topos.

\begin{itemize}
\item[a)] Dans l'ensemble des sous-topos de $E$, le Inf de deux
éléments est-il distributif par rapport aux $\Sup$ finis ?

\item[b)] Soient $E'$, $E''$ deux sous-topos de $E$ tels que $E=E'\vee
E''$. Alors est-il vrai que tout point de $E$ est isomorphe à
l'image d'un point de $E'$ ou d'un point de $E''$ ?

On notera que si la réponse à b) est affirmative pour $E$ et tous
ses sous-topos, alors la réponse à a) est affirmative du moins pour
le cas de sous-topos ayant suffisamment de points (donc pour tous
les sous-topos, si $E$ est localement noethérien
(\hyperref[IV.exer9.1.11]{9.1.11} b)). La réponse à a) (et, à fortiori à b)),
n'est pas connue cependant pour tous les topos noethériens.
Cependant la réponse à a) et b) est affirmative si $E$ est de la
forme $\Top(X)$, $X$ est un espace localement noethérien
(\hyperref[IV.exer9.1.11]{9.1.11}), ou si $E$ est un topos fini
(\hyperref[IV.exer9.1.12]{9.1.12}).

On notera que a) peut se reformuler sous la forme suivante
(appliquée à tous les sous-topos $F$ de $E$):

\item[$\rm a'$)] Si $E'$ et $E''$ sont deux sous-topos d'un topos $F$,
tels que le morphisme canonique (\cf \hyperref[IV.exer8.7]{8.7} b))
\begin{equation*}
E'\amalg E''\longrightarrow F\tag{9.1.14.1}
\end{equation*}\etiquette[9.1.14.1]{IV.eq9.1.14.1}
soit conservatif (ce qui signifie aussi $E=E'\vee E''$
(\hyperref[IV.exer9.1.7.2]{9.1.7.2} e)), alors il en est de même pour le morphisme
déduit par 2-changement de base $F_1\hookrightarrow
F$,\pageoriginale où $F_1$ est un sous-topos de $F$.

On a une reformulation analogue $\rm b'$) de b), en prenant un
changement de base $F_1\rightarrow F$ par un topos \emph{ponctuel}
$F_1$ (\hyperref[IV.subsec2.2]{2.2}). On voit donc qu'une réponse affirmative
à a), b) résulterait d'une réponse affirmative à la question
suivante (où on prend pour $F$ n'importe quel sous-topos de $E$):

\item[c)] Si $E'$ et $E''$ sont deux sous-topos du topos $F$, dont le
$\Sup$ soit $F$ \ie tel que le morphisme de topos \eqhyperref[IV.eq9.1.14.1]{9.1.14.1} soit
conservatif, alors ce morphisme est-il même \emph{universellement
conservatif}, \ie reste-t-il conservatif après tout 2-changement de
base $F_1\rightarrow F$ ?

Cet énoncé a un sens, grâce au résultat de \textsc{Deligne} (voir [16] proposition 4.47)
affirmant l'existence des 2-produits-fibrés de topos. Voir cependant
l'exemple \hyperref[IV.exer9.1.7.2]{9.1.7.2} d).

D'autre part, on voit par \hyperref[IV.exer9.1.11]{9.1.11} h) qu'une réponse
affirmative à b) fournirait dans le cas localement noethérien une
réponse affirmative à la question suivante:

\item[d)] Dans l'ensemble ordonné $\Sigma$ des sous-topos de $E$, le
$\Sup$ de deux éléments est-il distributif par rapport aux Inf
\emph{quelconques} (pas nécessairement finis) ?

Cela signifie aussi que $\Sigma$ s'interprète comme l'ensemble
ordonné des « sous-topos fermés »  d'un topos convenable
$E_{\cons'}$ (savoir $\Sigma^{0^{\sim}}$, où l'ensemble ordonné
opposé $\Sigma^0$ de $\Sigma$, considéré comme une
catégorie, est
munie de sa topologie canonique, \cf \hyperref[IV.exer7.8]{7.8} b)), qui est
engendré par les sous-objets de l'objet final, donc est très proche
d'un espace topologique ordinaire. Il resterait à étudier ce topos
$E_{\cons'}$, en s'inspirant des exemples \hyperref[IV.exer9.1.11]{9.1.11} g) et
\hyperref[IV.exer9.1.12]{9.1.12} d), et notamment à déterminer s'il est
noethérien si $E$ l'est, ce qui revient à la question suivante pour
les sous-topos $F$ de $E$, qui a un sens indépendamment de d):

\item[e)] Soient $F$ un topos noethérien, $(E_i)_{i\in I}$ une famille
de sous-topos de $E$ dont l'intersection soit $\Top(∅)$. Alors
existe-t-il une sous-famille finie ayant la même propriété ?

Dans un ordre d'idées assez différent, rappelons également la
question\pageoriginale soulevée dans (\hyperref[IV.exer9.1.6]{9.1.6} e)), qui
mériterait d'être éclaircie.
\end{itemize}

\subsection{Cas des sous-topos ouverts}\etiquette[9.2]{IV.subsec9.2}

\subsubsection{}\etiquette[9.2.1]{IV.subsubsec9.2.1}

Soit $U$ un ouvert d'un topos $E$, \ie un \sisi{sous objet}{sous-objet} de l'objet
% correction SGA4-CORR-0239 / SGA-SRC-000919: SGA4-EIV-04-L6248-L6249-FCB4CF072111.
final $e$ de $E$ (\hyperref[IV.def8.3]{8.3}). Considérons le morphisme de
localisation (\hyperref[IV.subsec5.2]{5.2})
\begin{equation*}
j:E_{/U}\longrightarrow E{\quoi}.\tag{9.2.2}
\end{equation*}\etiquette[9.2.2]{IV.eq9.2.2}
Comme $U\rightarrow e$ est ici un monomorphisme, le foncteur $j_{!}$
(qui s'interprète comme le foncteur oubli) est pleinement fidèle,
donc son biadjoint $j_\ast$ l'est également (\sisi{1.5.7  a)}{\hyperrefext[I][I.rem5.7]{5.7}}). En d'autres
termes, \emph{le morphisme de localisation $j$ $(9.2.2)$ associé à
un ouvert $U$ du topos $E$ est un plongement de topos}
(\hyperref[IV.def9.1.1]{9.1.1} a)).

Un plongement de topos $F\rightarrow E$ est appelé un
\emph{plongement ouvert} s'il est isomorphe au plongement de topos
défini par un ouvert $U$ de $E$. Un sous-topos $E'$
(\hyperref[IV.def9.1.1]{9.1.1} b)) de $E$ est appelé un \emph{sous-topos ouvert}
s'il est défini à l'aide d'un ouvert de $E$, \ie si le morphisme
d'inclusion canonique $E'\rightarrow E$ est un plongement ouvert. On
notera que l'ouvert $U$ de $E$ associé à un plongement ouvert de
topos $j:E'\rightarrow E$ est déterminé de façon unique comme
égal au sous-objet $j_{!}(e')$ de $e$, où $e'$ désigne l'objet final
de $E'$. Par suite, on trouve une correspondance biunivoque entre
ouverts $U$ de $E$ et sous-topos ouverts de $E$.

\begin{enonce*}[remark]{Remarque 9.2.3}\etiquette[9.2.3]{IV.rem9.2.3}
Conformément aux définitions précédentes, le
sous-topos ouvert de
$E$ associé à l'ouvert $U$ de $E$ est la sous-catégorie strictement
pleine de $E$ image essentielle du foncteur image directe
$$
j_{\ast}:E_{/U}\longrightarrow E
$$
On se gardera de confondre cette sous-catégorie de $E$ avec la
sous-catégorie strictement pleine, image essentielle de $j_{!}$,
formée des objets $X$ de $E$ tels que le morphisme structural
$X\rightarrow e$ se factorise par $U$. Il est clair cependant que
ces deux sous-catégories se déterminent mutuellement, et qu'elles
sont canoniquement équivalentes. C'est pourquoi certains auteurs ont
pu être tentés d'appeler (voire, ont pu appeler) « sous-topos
ouvert défini par $U$ »  la sous-catégorie strictement pleine
image essentielle de $j_{!}$, dont la description en termes de $U$
est plus simple que celle\pageoriginale de l'image essentielle de
$j_{\ast}$. Cette terminologie ne présente pas d'inconvénient tant
qu'on ne se propose pas d'étudier d'autres sortes de sous-topos que
des sous-topos ouverts. Comme ce n'est pas notre cas, nous ne
suivrons pas ici les auteurs sus-mentionnés.

Pour la commodité du lecteur, nous allons résumer les propriétés
spéciales les plus importantes des plongements ouverts:
\end{enonce*}

\begin{enonce*}{Proposition 9.2.4}\etiquette[9.2.4]{IV.prop9.2.4}
Soit
$$
j:\mathcal{U}\longrightarrow E
$$
un plongement ouvert de topos $(\hyperref[IV.subsubsec9.2.1]{9.2.1})$. Alors le
foncteur $j_{!}$ adjoint à gauche de $j^{\ast}$ existe, de sorte
qu'on a une suite de trois foncteurs adjoints
\[
\xymatrix@R=.35cm@C=.5cm{
           & j_!,\quad j^{\ast}, \quad j_{\ast} : &\\
&&\\
\mathcal{U} \ar @<-16pt>[rr]^{j_{\ast}} \ar
           @<14pt>[rr]^{j_{!}}&            &\ar @<1pt>[ll]_{j^{\ast}} E
}
\]

De plus:
\begin{enumerate}
\renewcommand{\theenumi}{\alph{enumi}}
\renewcommand{\labelenumi}{\theenumi)}
\item Le foncteur $j_{!}$ et le foncteur $j_{\ast}$ sont pleinement
fidèles.

\item Le foncteur $j_{!}$ commute aux produits fibrés, aux produits
indexés par des petits ensembles $I\neq ∅$, et aux limites
projectives relatives à de petites catégories d'indices
co-filtrantes (\hyperrefext[I][I.def2.7]{2.7}).

\item Pour tout objet $X$ de $E$, le morphisme d'adjonction
$$
j_{!}j^{\ast}(X)\longrightarrow X
$$
est un monomorphisme.
\end{enumerate}
\end{enonce*}

\noindent \emph{Démonstration}. On peut supposer que $j$ est le
morphisme de localisation défini par un ouvert $U$ de $E$. Alors a)
est mis pour mémoire, b) provient de l'interprétation de $j_{!}$
comme un foncteur oubli et du fait que $U\rightarrow e$ est un
monomorphisme. Enfin c) est immédiat en notant que le morphisme
envisagé n'est autre que le morphisme déduit de l'inclusion
% correction SGA4-CORR-0240 / SGA-SRC-000920: SGA4-EIV-04-L6341-L6343-AF315B890934.
$U\rightarrow e$ par le changement de base $X\rightarrow e$, compte
tenu que par changement de base un monomorphisme est transformé en
monomorphisme.

\begin{enonce*}[remark]{Remarques 9.2.4.1}\etiquette[9.2.4.1]{IV.rem9.2.4.1}
On peut trouver des plongements de topos $j:E'\rightarrow E$ tels
que $j_{!}$ existe, mais que $j_{!}$ ne commute pas aux produits de
deux objets, ni aux limites projectives cofiltrantes (exercice
\hyperref[IV.exer9.5.9]{9.5.9} c)).
\end{enonce*}

\setcounter{subsubsection}{4}
\subsubsection{}\etiquette[9.2.5]{IV.subsubsec9.2.5}
Soit\pageoriginale maintenant
$$
g:E'\longrightarrow E
$$
un morphisme de topos quelconque, et soit $U'=g^{\ast}(U)$ l'image
inverse de l'ouvert $U$ de $E$. On a vu dans \hyperref[IV.prop5.11]{5.11} que
le diagramme de topos
\[
\xymatrix@C=2cm@R=1.5cm{
E'/_{U'} \ar[d]_{g_U} \ar[r]^{j'}& E' \ar[d]^{g}\\
E/_{U}  \ar[r]^{j} & E }
\]
est 2-cartésien. On en conclut en particulier que l'image inverse
par $g$ d'un sous-topos ouvert de $E$ (\cf \hyperref[IV.exer9.1.6]{9.1.6} d))
est un sous-topos ouvert de $E'$, ou ce qui revient au même, que la
notion de plongement ouvert de topos (\hyperref[IV.subsubsec9.2.1]{9.2.1}) est
stable par 2-changements de base dans la 2-catégorie des
$𝒰$-topos (éléments d'un univers donné~$𝒱$).

Notons aussi que pour que le morphisme de topos $g:E'\rightarrow E$
se factorise, à isomorphisme près, par $j:E_{/U}\rightarrow E$, il
faut et il suffit évidemment que $j':E'_{/U'}\rightarrow E'$ soit
une équivalence de catégories, ce qui signifie encore que
l'inclusion $U'\hookrightarrow e'$ (= objet final de $E'$) est un
isomorphisme. Ceci précise donc \hyperref[IV.prop9.1.4]{9.1.4}, en caractérisant
de façon simple l'image essentielle du foncteur envisagé dans
\loccit

\subsubsection{}\etiquette[9.2.6]{IV.subsubsec9.2.6}
Appliquons \hyperref[IV.subsubsec9.2.5]{9.2.5} au cas où $E'$ est de la forme
$E_{/V}$, où $V$ est un deuxième ouvert de $E$, $g$ étant le
morphisme de localisation. Alors $U'=U\cap V$, et on trouve que le
produit 2-cartésien de $E_{/U}$ et $E_{/V}$ sur $E$ s'identifie à
$E_{/U\cap V}$. On en conclut \emph{qu'une intersection finie de
sous-topos ouverts de $E$ est un sous-topos ouvert de $E$}, plus
précisément l'application
\begin{equation*}
U\longmapsto \text{ sous-topos de $E$ défini par $U$}
\tag{9.2.6.1}
\end{equation*}\etiquette[9.2.6.1]{IV.eq9.2.6.1}
% correction SGA4-CORR-0241 / SGA-SRC-000921: SGA4-EIV-04-L6391-L6395-F108BF35513B.
commute aux intersections finies.

\begin{enonce*}[remark]{Exercice 9.2.7}\etiquette[9.2.7]{IV.exer9.2.7}
Prouver que l'application \eqhyperref[IV.eq9.2.6.1]{9.2.6.1} commute également aux $\Sup$
\emph{quelconques}. (Utiliser \hyperref[IV.exer9.1.7.2]{9.1.7.2} e)).
\end{enonce*}

\subsection{Construction du sous-topos fermé complémentaire d'un
sous-topos ouvert}\etiquette[9.3]{IV.subsec9.3}\pageoriginale

\subsubsection{}\etiquette[9.3.1]{IV.subsubsec9.3.1}
Soient $\mathsf{T}$ un espace topologique, et $U$ un ouvert de
$\mathsf{T}$, de sorte que $\Top(U)$ s'identifie à un sous-topos
ouvert de $\Top(\mathsf{T})$ (notations de \hyperref[IV.subsec2.1]{2.1}). Soit
$Y$ le sous-espace topologique fermé de $\mathsf{T}$ complémentaire
de $U$. On peut alors, (à équivalence près) considérer $\Top(Y)$
comme un sous-topos de $\Top(\mathsf{T})$, savoir le sous-topos
formé des objets $F$ de $\Top(\mathsf{T})$ dont la restriction à $U$
est \sisi{le topos final de $U$}{l'objet final de $U$}. Cette description d'une sous-catégorie de
$E=\Top(\mathsf{T})$ garde un sens chaque fois qu'on a un topos et
un ouvert $U$ de celui-ci, et on verra qu'elle fournit toujours un
sous-topos de $E$. De plus, il est immédiat dans le cas particulier
envisagé d'abord (et avec la terminologie introduite dans l'exercice
\hyperref[IV.exer9.1.13]{9.1.13}) que $\Top(Y)$ et $\Top(U)$ sont des sous-topos
complémentaires l'un de l'autre (utiliser \hyperref[IV.subsubsec9.2.5]{9.2.5} et
\hyperref[IV.exer9.1.7.2]{9.1.7.2} e)). Il en sera encore de même dans le cas
général, et nous verrons que cette propriété caractérise de façon
unique le sous-topos envisagé. Il mérite donc à tous points de vue,
le nom de sous-topos fermé complémentaire de l'ouvert envisagé $U$,
ou du sous-topos ouvert défini par $U$. Le détail de la construction
de ce topos $\mathcal{F}$, et du foncteur image inverse
$i^{\ast}:E\rightarrow \mathcal{F}$, sera donné dans la présente
section, tandis que la section \hyperref[IV.subsec9.4]{9.4} en développera les
premières propriétés.

\subsubsection{}\etiquette[9.3.2]{IV.subsubsec9.3.2}
Soient donc, comme dans \hyperref[IV.subsubsec9.2.1]{9.2.1}, $E$ un topos, $U$ un
ouvert de $E$, et considérons le morphisme de localisation
\begin{equation*}
j:\mathcal{U}=E_{/U}\longrightarrow E{\quoi},\tag{9.3.2.1}
\end{equation*}\etiquette[9.3.2.1]{IV.eq9.3.2.1}
qui est un plongement ouvert.

Pour tout objet $X$ de $E$, posons
\begin{equation*}
X_{\complement U}=U\coprod_{U\times X}
X{\quoi},\tag{9.3.2.2}
\end{equation*}\etiquette[9.3.2.2]{IV.eq9.3.2.2}
où l'amalgamation est faite pour les projections canoniques
$$
\pr_1:U\times X\longrightarrow U,\quad \pr_2:U\times
X\longrightarrow X
$$
% correction SGA4-CORR-0242 / SGA-SRC-000922: SGA4-EIV-04-L6443-L6447-B99FBEB3E93E.
On a donc un diagramme \emph{cocartésien} (\hyperrefext[I][I.sec10]{10}) dans $E$,
dépendant fonctoriellement de $X$:\pageoriginale
\[\vcenter{
\xymatrix@C=1.75cm@R=1.5cm{
U\times X \ar[r]^{\pr_2} \ar[d]_{\pr_1}& X \ar[d]\\
U \ar[r]        & X_{\complement U} }}\leqno{(9.3.2.3)}
\]\etiquette[9.3.2.3]{IV.eq9.3.2.3}
On notera que dans l'exemple envisagé dans \hyperref[IV.subsubsec9.3.1]{9.3.1},
$X_{\complement U}$ est canoniquement isomorphe à
$i_{\ast}i^{\ast}(X)$ (où $i:Y\rightarrow T$ est l'inclusion), et
$X\rightarrow X_{\complement U}$ s'identifie au morphisme
d'adjonction.

\begin{enonce*}{Proposition 9.3.3}\etiquette[9.3.3]{IV.prop9.3.3}
Avec les notations précédentes, on a ce qui suit:
\begin{enumerate}
\renewcommand{\theenumi}{\alph{enumi}}
\renewcommand{\labelenumi}{\theenumi)}
\item Pour tout objet $X$ de $E$, les conditions suivantes sont
équivalentes:
\begin{itemize}
\item[(i)] Il existe un objet $Y$ de $E$ et un isomorphisme $X\simeq
Y_{\complement U}$.

\item[(ii)] Le morphisme canonique $X\xrightarrow{m}X_{\complement U}$ est
un isomorphisme.

\item[(iii)] Le morphisme canonique $\pr_2:X\times U\rightarrow U$ est
un isomorphisme (\ie $j^{\ast}(X)$ est un objet final de $E_{/U}$).
\end{itemize}

\item Pour tout morphisme
$$
m:\sisi{E}{X}\longrightarrow X'
$$
% correction SGA4-CORR-0243 / SGA-SRC-000923: SGA4-EIV-04-L6479-L6483-CF236A2D6F6D.
dans $E$, les conditions suivantes sont équivalentes:
\begin{itemize}
\item[${\rm (i')}$] Le diagramme
\[
\xymatrix@C=2cm@R=1.5cm{
X \ar[r]^{m} \ar[d]& X'\ar[d]\\
X_{\complement U} \ar[r]& X'_{\complement U} }
\]
est cartésien.

\item[${\rm (ii')}$] Le morphisme $m\times \id_U:X\times U\rightarrow
X'\times U$ est un isomorphisme, \ie $j^{\ast}(m)$ est un
isomorphisme.
\end{itemize}

\item le foncteur
$$
X\longrightarrow X_{\complement U}
$$
de $E$ dans $E$ est exact à gauche, transforme les familles
épimorphiques en familles épimorphiques, commute aux limites
inductives filtrantes et\pageoriginale aux sommes amalgamées.
\end{enumerate}
\end{enonce*}

\noindent \emph{Démonstration}. On démontre d'abord la proposition
dans le cas où $E$ est de la forme $\widehat{C}$, où $C$ est une
petite catégorie, et pour cela on est ramené par I.3.1 au cas où $E$
est le topos « ponctuel »  $(\Ens)$, où la proposition est
évidente. On passe de là au cas général par les
procédés standard de
\hyperrefext[II][II.sec4]{4}, par utilisation du foncteur « faisceau associé ».

\begin{enonce*}{Proposition 9.3.4}\etiquette[9.3.4]{IV.prop9.3.4}
Avec les notations de $\hyperref[IV.subsubsec9.3.2]{9.3.2}$, la sous-catégorie
strictement pleine $\mathcal{F}$ de $E$ définie par les \sisi{sous-objets}{objets}
% correction SGA4-CORR-0244 / SGA-SRC-000924: SGA4-EIV-04-L6517-L6520-E0E4165BBC45.
$X$ de $E$ tels que $j^{\ast}(X)$ soit un objet final de
$\mathcal{U}$ {\rm (\cf \hyperref[IV.prop9.3.3]{9.3.3} a))} est un sous-topos de
$E$, \ie $(\hyperref[IV.def9.1.1]{9.1.1}\quad )$) c'est un topos, et le foncteur
d'inclusion $i_{\ast}:\mathcal{F}\rightarrow E$ est le foncteur
image directe par un morphisme de topos $i:\mathcal{F}\rightarrow
E$, dont le foncteur image réciproque est le foncteur $X\mapsto
i^{\ast}X=X_{\complement U}$. Le morphisme d'adjonction
$X\rightarrow \sisi{i^{\ast}i_{\ast}X}{i_{\ast}i^{\ast}X}$ est le morphisme canonique
% correction SGA4-CORR-0245 / SGA-SRC-000925: SGA4-EIV-04-L6525-L6527-3E7AB5039BCF.
$X\rightarrow X_{\complement U}$ de $\eqhyperref[IV.eq9.3.2.3]{9.3.2.3}$
\end{enonce*}

\noindent \emph{Démonstration}: Pour tout objet $X$ de $E$,
désignons par $u(X)$ le morphisme canonique $X\rightarrow
X_{\complement U}$. Le morphisme $u(X)$ est fonctoriel en $X$ et
pour tout objet $X$, $u(X_{\complement U})$ est un isomorphisme. Il
résulte alors formellement et trivialement de ces deux propriétés
que le foncteur $i^{\ast}$ est adjoint à gauche au foncteur
$i_{\ast}$ et que le morphisme d'adjonction est $u(X)$. Comme le
foncteur $X\rightarrow X_{\complement U}$ est exact à gauche, le
foncteur $i^{\ast}$ est exact à gauche. Il résulte alors de
\hyperrefext[II][II.theo5.5]{5.5} que $\mathcal{F}$ est un topos, et de la
définition \hyperref[IV.def3.1]{3.1} que le couple $(i^{\ast}, i_{\ast})$ est
un morphisme de topos.

\setcounter{subsubsection}{4}
\subsubsection{}\etiquette[9.3.5]{IV.subsubsec9.3.5}
Le sous-topos $\mathcal{F}$ de $E$ décrit dans \hyperref[IV.prop9.3.4]{9.3.4}
est appelé le \emph{sous-topos fermé de $E$ complémentaire de
l'ouvert $U$ de $E$,} ou (au choix) du sous-topos ouvert de $E$
correspondant à $U$. Conformément à \hyperref[IV.def9.1.1]{9.1.1} b), le
morphisme $i:\mathcal{F}\rightarrow E$ construit dans
\hyperref[IV.prop9.3.4]{9.3.4} est appelé le \emph{morphisme d'inclusion.} Une
sous-catégorie pleine $\mathcal{F}$ de $E$ est \sisi{appelé}{appelée} un
% correction SGA4-CORR-0246 / SGA-SRC-000926: SGA4-EIV-04-L6550-L6553-681102120B89.
\emph{sous-topos fermé} de $E$ s'il existe un ouvert $U$ de $E$ tel
que $\mathcal{F}$ soit le sous-topos fermé complémentaire de $U$. On
notera que cet $U$ est déterminé de façon unique comme
$i_{\ast}(\phi_{\mathcal{F}})$, où $\phi_{\mathcal{F}}$ est l'objet
initial de $\mathcal{F}$, égal à $i^{\ast}(\phi_E)$; on l'appelle
aussi \emph{l'ouvert de $E$ complémentaire du sous-topos fermé
$\mathcal{F}$} et\pageoriginale le sous-topos de $E$ défini par $U$
s'appelle aussi le \emph{sous-topos ouvert de $E$ complémentaire de
$\mathcal{F}$}. Enfin, un plongement de topos (\hyperref[IV.def9.1.1]{9.1.1} a))
$i:F\rightarrow E$ est dit \emph{plongement fermé} si le sous-topos
correspondant de $E$ est fermé.

Si $G$ est un Groupe sur $E$, $s\in \Hom(e_E,G)$, on appelle
\emph{support} de $G$ (\resp de $s$) le sous-topos fermé de $E$
complémentaire de l'ouvert cosupport de $G$ (\resp $s$)
(\hyperref[IV.subsubsec8.5.1]{8.5.1}, \hyperref[IV.subsubsec8.5.2]{8.5.2}).
\subsubsection{}\etiquette[9.3.6]{IV.subsubsec9.3.6}
Avec les notations de \hyperref[IV.subsubsec9.3.2]{9.3.2} et \hyperref[IV.prop9.3.4]{9.3.4}
on trouve donc deux morphismes de topos
\begin{equation*}
\mathcal{U}\xrightarrow{j}
E\xleftarrow{i}\mathcal{F}{\quoi},\tag{9.3.6.1}
\end{equation*}\etiquette[9.3.6.1]{IV.eq9.3.6.1}
définissant cinq foncteurs formant deux suites de foncteurs adjoints
$(j_{!}, j^{\ast}, j_{\ast})$ et $(i^{\ast}, i_{\ast})$:
\[
\xymatrix@C=2cm@R=1.5cm@M=.2cm{ \Uu \ar@<6pt>[r]^{j_!}
\ar@<-13pt>[r]|{j_{\ast}} & E \ar@<3pt>[r]^{i^{\ast}}
\ar@<4pt>[l]|{j^{\ast}}& \Ff
\ar@<13pt>[l]|{i_{\ast}}.}\leqno{(9.3.6.2)}
\]\etiquette[9.3.6.2]{IV.eq9.3.6.2}

Le foncteur
\begin{equation*}
\rho=\sisi{i_{\ast}j^{\ast}}{i^{\ast}j_{\ast}}:\mathcal{U}\longrightarrow
\mathcal{F}\tag{9.3.6.3}
\end{equation*}\etiquette[9.3.6.3]{IV.eq9.3.6.3}
% correction SGA4-CORR-0247 / SGA-SRC-000927: SGA4-EIV-04-L6584-L6588-FE73B7C7219F.
est appelé le \emph{foncteur de recollement} relatif à l'ouvert $U$
de $E$; plus généralement si $\mathcal{U}$ et $\mathcal{F}$ sont
deux topos quelconques, un foncteur $\rho:\mathcal{U}\rightarrow
\mathcal{F}$ est appelé \emph{foncteur de recollement} s'il est
exact à gauche et accessible (I.9.2), ou ce qui revient au même
(\sisi{\hyperref[IV.rem1.8]{1.8}\quad}{\hyperrefext[I][I.sec8]{8}}), s'il admet un pro-adjoint. Les raisons de
cette terminologie vont apparaître dans \hyperref[IV.prop9.4.1]{9.4.1} d) et
\hyperref[IV.theo9.5.4]{9.5.4} ci-dessous.

\subsection{Premières propriétés du sous-topos fermé
$\mathcal{F}$ et de $i:\mathcal{F}\rightarrow
E$}\etiquette[9.4]{IV.subsec9.4}

\begin{enonce*}{Proposition 9.4.1}\etiquette[9.4.1]{IV.prop9.4.1}
Les notations sont celles de $\hyperref[IV.subsubsec9.3.2]{9.3.2}$ et
$\hyperref[IV.prop9.3.4]{9.3.4}$
\begin{enumerate}
\renewcommand{\theenumi}{\alph{enumi}}
\renewcommand{\labelenumi}{\theenumi)}
\item Le foncteur $i_{\ast}:\mathcal{F}\rightarrow E$ transforme les
familles épimorphiques en familles épimorphiques. Il commute à la
% correction SGA4-CORR-0248 / SGA-SRC-000928: SGA4-EIV-04-L6608-L6611-CF99EBFBDAE0.
formation des sommes amalgamées et des limites inductives
filtrantes.

\item Pour tout objet $X$ de $E$ tel que le morphisme de projection
$U\times X\rightarrow U$ soit un épimorphisme, le morphisme
d'adjonction $X\rightarrow i_{\ast} i^{\ast}X$ est un épimorphisme.

\item Le\pageoriginale couple de foncteurs $(i^{\ast}, j^{\ast})$ de
$E$ dans $\mathcal{F}$ et $\mathcal{U}$ respectivement est
conservatif {\rm (I 6.1)} \ie un morphisme $m$ de $E$ est un
isomorphisme si et seulement si $i^{\ast}(m)$ et $j^{\ast}(m)$ sont
des isomorphismes (ceci équivaut à dire que le couple de foncteurs
$(i^{\ast}, j^{\ast})$ est fidèle $(\hyperref[IV.subsubsec6.4.0]{6.4.0})$).

\item Le foncteur de recollement
$\rho=i^{\ast}j_{\ast}:\mathcal{U}\rightarrow \mathcal{F}$ est un
foncteur exact à gauche et accessible {\rm (\hyperrefext[I][I.def9.2]{9.2})}, ou ce
qui revient au même {\rm(\hyperrefext[I][I.sec8]{8}.)} il admet un pro-adjoint
$\sigma:\Pro\mathcal{F}\rightarrow \Pro \mathcal{U}$.
\end{enumerate}
\end{enonce*}

\noindent \emph{Démonstration}: a) résulte \sisi{le}{de} 9.3.3 c). Pour
% correction SGA4-CORR-0249 / SGA-SRC-000929: SGA4-EIV-04-L6632-L6634-B6B507C7699A.
démontrer b) il suffit de se reporter à la définition
\eqhyperref[IV.eq9.3.2.2]{9.3.2.2} de $X_{\complement U}$ compte-tenu de ce que
$i_{\ast}i^{\ast}X=X_{\complement U}$ (\hyperref[IV.prop9.3.4]{9.3.4}).
L'assertion c) résulte de \hyperref[IV.prop9.3.3]{9.3.3} b). Les foncteurs
$i_{\ast}$ et $j_{\ast}$ sont exacts à gauche et par suite $\rho$
est exact à gauche. Le foncteur $j_{\ast}$ est un foncteur image
directe par un morphisme de topos et par suite il est accessible
(\hyperrefext[I][I.prop9.5]{9.5}); on note qu'un topos est une catégorie accessible
(\hyperrefext[I][I.rem9.11.3]{9.11.3})). Comme le foncteur $i^{\ast}$ commute aux
limites inductives (c'est un foncteur image réciproque), le foncteur
$\rho$ composé de deux foncteurs accessibles est accessible.

\begin{enonce*}{Proposition 9.4.2}\etiquette[9.4.2]{IV.prop9.4.2}
Soit $E'$ un topos. Alors le foncteur $h\mapsto i\circ h$:
$$
\sisi{\Homtop(E',\mathcal{F})}{\Hhomtop(E',\mathcal{F})}\longrightarrow
\sisi{\Homtop(E',E)}{\Hhomtop(E',E)}
$$
% correction SGA4-CORR-0250 / SGA-SRC-000930: SGA4-EIV-04-L6645-L6652-55F35A597EDD.
est pleinement fidèle, et son image essentielle est formée des
morphismes de topos $g:E'\rightarrow E$ tels que $g^{\ast}(U)$ soit
un objet initial $\phi_{E'}$ de $E'$.
\end{enonce*}

Par \hyperref[IV.prop9.1.4]{9.1.4}, on sait que le foncteur envisagé est
pleinement fidèle et que $g$ est dans son image essentielle si et
seulement si pour tout objet $X'$ de $E'$, l'objet $g_{\ast}(X')$ de
$E$ appartient à $\mathcal{F}$, \ie qu'on a
\begin{equation*}
j^{\ast}(g_{\ast}(X'))\simeq e_{\mathcal{U}}\text{ pour tout } X'\in
\Ob E'{\quoi}.\tag{*}
\end{equation*}
Or, posons $U'=g^{\ast}(U)$, et considérons le diagramme de topos
\[
\xymatrix@C=1.5cm{
E'_{/U'}\ar[r]^{j'} \ar[d]_{g_U}& E'\ar[d]^{g}\\
E_{/U}   \ar[r]^{j}             & E\quad , }
\]
il\pageoriginale est clair qu'on a un isomorphisme canonique
$$
j^{\ast}(g_{\ast}(X'))\simeq g_{U_{\ast}}({j'}^{\ast}(X')){\quoi},
$$
d'autre part il est immédiat aussi que le foncteur ${j'}^{\ast}$ est
essentiellement surjectif, donc la condition (*) équivaut à la
% correction SGA4-CORR-0251 / SGA-SRC-000931: SGA4-EIV-04-L6671-L6678-2CB3AC8F47FB.
condition
\begin{equation*}
g_{U\ast}(Y')\simeq e_{\mathcal{U}}\text{ pour tout }Y'\in \Ob
E'_{/U'}{\quoi}.\tag{**}
\end{equation*}
Utilisant la propriété d'adjonction de $g_{U\ast}$ et $g_{U}^{\ast}$,
on voit aussitôt que cela signifie que $g^{\ast}_U(Y)$ est un objet
initial de $E'_{/U'}$ pour tout objet $Y$ de $E_{/U}$, ou encore
(utilisant que l'objet initial de $E'_{/U'}$ est strict
(\hyperrefext[II][II.def4.5]{4.5}) qu'il en est ainsi pour $Y=e_{\mathcal{U}}$, objet
final de $\mathcal{U}$. Mais cela signifie aussi que l'objet final
$U'$ de $E'_{/U'}$ est initial, ou, ce qui revient manifestement au
même, si et seulement si $U'$ est un objet initial $\phi_{E'}$ de
$E'$. \cqfd

\begin{enonce*}{Corollaire 9.4.3}\etiquette[9.4.3]{IV.cor9.4.3}
Soient $g:E'\rightarrow E$ un morphisme de topos $U'=g^{\ast}(U)$,
$\mathcal{U}'=E'_{/U'}$ et $\mathcal{F}'$ le sous-topos fermé de
$E'$ complémentaire de l'ouvert $U'$, enfin
$i':\mathcal{F}'\rightarrow E'$ le morphisme d'inclusion. Alors le
morphisme de topos $g\circ i':\mathcal{F}'\rightarrow E$ se factorise
% correction SGA4-CORR-0252 / SGA-SRC-000932: SGA4-EIV-04-L6690-L6697-5250A2976FF1.
à isomorphisme près en un morphisme de topos
$g_{\mathcal{F}}:\mathcal{F}'\rightarrow \mathcal{F}$, et le
diagramme de topos correspondant
\[\vcenter{
\xymatrix@=.5cm{ \Ff' \ar[rr]^{i'}\ar[dd]_{g_{\Ff}}  &&
E'\ar[dd]^{g}
\\
                          &            &  \\
\Ff  \ar[rr]^(.4){i}  &  \ar@/^.5pc/[ru]  &   E }}\leqno{(9.4.3.1)}
\]\etiquette[9.4.3.1]{IV.eq9.4.3.1}
est $2$-cartésien {\rm (\cf \hyperref[IV.prop5.11]{5.11}).}
\end{enonce*}

Résulte formellement de \hyperref[IV.prop9.4.2]{9.4.2} et des définitions.

Avec la terminologie introduite dans \hyperref[IV.exer9.1.6]{9.1.6} d), on peut
donc dire, en particulier, que l'image inverse par un morphisme de
topos $g:E'\rightarrow E$ d'un sous-topos \emph{fermé} $\mathcal{F}$
de $E$ est un sous-topos \emph{fermé} $\mathcal{F}'$ de $E'$ (savoir
le sous-topos fermé complémentaire de l'ouvert $U'=g^{\ast}(U)$, où
$U$ est l'ouvert de $E$ complémentaire de $\mathcal{F}$).

\begin{enonce*}{Corollaire 9.4.4}\etiquette[9.4.4]{IV.cor9.4.4}
Soient $E$ un topos, $\mathcal{U}$ un sous-topos ouvert de $E$,
$\mathcal{F}$ le sous-topos fermé complémentaire.
\begin{enumerate}
\renewcommand{\theenumi}{\alph{enumi}}
\renewcommand{\labelenumi}{\theenumi)}
\item On a\pageoriginale
$$
\mathcal{U} \cap \mathcal{F} \approxeq \Top(∅)\quav
\mathcal{U}\vee \mathcal{F}=E{\quoi},
$$
où le signe $\vee$ désigne le $\Sup$ dans l'ensemble de tous les
sous-topos de $E$ {\rm (\hyperref[IV.rem9.1.2]{9.1.2} c))}. (Avec la
terminologie introduite dans $\hyperref[IV.exer9.1.13]{9.1.13}$, $\mathcal{U}$ et
$\mathcal{F}$ sont des sous-topos de $E$ complémentaires l'un de
l'autre).

\item $\mathcal{U}$ (\resp $\mathcal{F}$) est le plus-grand parmi les sous-topos
$E'$ de $E$ tels que l'on ait $E'\cap \mathcal{F}\approxeq
\Top(∅)$ (\resp $\mathcal{U} \cap E'\approxeq \Top(∅)$).
\end{enumerate}
\end{enonce*}

\noindent \emph{Démonstration.}
\begin{enumerate}
\renewcommand{\theenumi}{\alph{enumi}}
\renewcommand{\labelenumi}{\theenumi)}
\item la première relation revient à dire que si $E'$ est un topos et
$g:E'\rightarrow E$ est un morphisme de topos qui se factorise à
isomorphisme près par $\mathcal{U}$ et par $\mathcal{F}$, alors
$E'\approxeq \Top(∅)$; or en vertu de \hyperref[IV.subsubsec9.2.5]{9.2.5} et
\hyperref[IV.prop9.4.2]{9.4.2}, l'hypothèse signifie que $g^{\ast}(U)$ est à la
fois un objet initial et un objet final de $E'$, d'où la conclusion.
Pour la deuxième relation, on note que par \hyperref[IV.exer9.1.7.2]{9.1.7.2} e)
elle équivaut à \hyperref[IV.prop9.4.1]{9.4.1} c).

\item Supposons que $E'\cap \mathcal{F}\approxeq \Top(∅)$; en vertu de
9.4.3 $E'\cap \mathcal{F}$ est équivalent au sous-topos de $E'$
    complémentaire de $g^{\ast}(U)$, où $g:E'\rightarrow E$ est le
    morphisme d'inclusion, donc la condition envisagée signifie que
    $U'=g^{\ast}(U)$ est un objet final de $E'$, \ie (\hyperref[IV.subsubsec9.2.5]{9.2.5}) que
    $E'\subset \mathcal{U}$. Supposons que $\mathcal{U}\cap E'\approxeq
    \Top(∅)$; en vertu de \hyperref[IV.prop5.11]{5.11} $\mathcal{U}\cap E'$ est équivalent à
    $E'_{/U'}$, donc la condition envisagée signifie que $U'$ est
    objet initial de $E'$, \ie (\hyperref[IV.prop9.4.2]{9.4.2}) que $E'\subset \mathcal{F}$.
\end{enumerate}

\begin{enonce*}{Corollaire 9.4.5}\etiquette[9.4.5]{IV.cor9.4.5}
Soit $(\mathcal{F}_i)_{i\in I}$ une famille de sous-topos fermés du
topos $E$, et pour tout $i\in I$, soit $U_i$ l'ouvert de $E$
% correction SGA4-CORR-0253 / SGA-SRC-000933: SGA4-EIV-04-L6765-L6770-387E3971CC7C.
complémentaire de $\Ff_i$. Alors le sous-topos de $E$ intersection
de $\Ff_i$ $(\hyperref[IV.prop9.1.3]{9.1.3})$ est un sous-topos fermé, dont
l'ouvert complémentaire est $U=\Sup_i U_i$.
\end{enonce*}

Cela résulte formellement de la propriété universelle de
l'intersection comme 2-produit et de \hyperref[IV.prop9.4.2]{9.4.2}, compte tenu
que pour un morphisme de topos $g:E'\rightarrow E$,
$g^{\ast}(U)=\Sup_i g^{\ast}(U_i)$ est un objet initial si et
seulement si les $g^{\ast}(U_i)$ le sont.

On prouve de façon analogue:

\begin{enonce*}{Corollaire 9.4.6}\etiquette[9.4.6]{IV.cor9.4.6}
Avec les notations de  $9.4.5$ supposons $I$ fini. Alors
le\pageoriginale sous-topos de $E$ borne supérieure de $\mathcal{F}_{i}$ {\rm
(\hyperref[IV.rem9.1.2]{9.1.2} c))} est un sous-topos fermé de $E$, dont
l'ouvert complémentaire est égal à $\bigcap_i U_i$.
% correction SGA4-CORR-0254 / SGA-SRC-000934: SGA4-EIV-04-L6781-L6785-EBA4EEBABC6F.
\end{enonce*}

\begin{enonce*}[remark]{Exercice 9.4.7}\etiquette[9.4.7]{IV.exer9.4.7}
(\emph{Recollement de sous-topos}). Soit $E$ un topos.
\begin{enumerate}
\renewcommand{\theenumi}{\alph{enumi}}
\renewcommand{\labelenumi}{\theenumi)}
\item Soient $E'$ un sous-topos de $E$, $(S_i)_{i\in I}$ une famille
d'objets de $E$ recouvrant l'objet final, et pour tout $i\in I$,
soit $S'_i$ l'image inverse de $S_i$ dans $E'$, et $E'_i\approxeq
E'_{/S'_i}$ le sous-topos de $E_i=E_{/S_i}$ image inverse par le
morphisme de localisation $E_i\rightarrow E$ du sous-topos $E'$ de
$E$ (\hyperref[IV.exer9.1.6]{9.1.6} d) et \hyperref[IV.prop5.11]{5.11}). Montrer que pour
un objet $X$ de $E$, $X$ appartient à $E'$ si et seulement si pour
tout $i\in I$, l'image inverse $X| S_i$ de $X$ dans $E_i$ appartient
à $E'_i$.

\item Supposons que $E$ soit de la forme $C^{\sim}$, où $C$ est un
$𝒰$-site. Montrer que le préfaisceau sur $C$
\begin{center}
$S\mapsto$ ensemble des sous-topos du topos
$(C_{/S})^{\sim}\approxeq \sisi{E^{\sim}_{/\varepsilon(S)}}{E_{/\varepsilon(S)}}$
% correction SGA4-CORR-0255 / SGA-SRC-000935: SGA4-EIV-04-L6803-L6808-8CC8FB69BC6A.
\end{center}
est un faisceau. $E$ étant de nouveau quelconque, en conclure qu'il
existe un élément de $E$ qui représente le foncteur
$$
\text{$S\mapsto$ ensemble des sous-topos du topos $E_{/S}$;}
$$

\item Avec les notations de a), montrer que pour que le sous-topos
$E'$ de $E$ soit un sous-topos ouvert (\resp un sous-topos fermé),
il faut et il suffit que pour tout $i\in I$, il en soit de même du
sous-topos induit $E'_i$ de $E_i$.
\end{enumerate}
\end{enonce*}

\begin{enonce*}[remark]{Exercice 9.4.8}\etiquette[9.4.8]{IV.exer9.4.8}
(\emph{Intérieur, extérieur, adhérence et frontière d'un
sous-topos.}) Soient $E$ un topos, $E'$ un sous-topos,
$i:E'\rightarrow E$ le morphisme d'inclusion.
\begin{enumerate}
\renewcommand{\theenumi}{\alph{enumi}}
\renewcommand{\labelenumi}{\theenumi)}
\item Montrer que $i_{\ast}(\phi_{E'})$ est le plus grand ouvert $U$
de $E$ tel que l'on ait $E'| \mathcal{U} \approxeq \Top(∅)$, où
$\mathcal{U}$ est le sous-topos ouvert de $E$ défini par $U$. On
appellera cet $U$, ou $\mathcal{U}$, l'\emph{extérieur} du
sous-topos $E'$ de $E$, et le sous-topos fermé de $E$ complémentaire
de $U$ l'\emph{adhérence} de $E'$.

\item Montrer qu'il existe un plus grand ouvert $V$ de $E$ tel que le
sous-topos ouvert correspondant $\mathcal{V}$ de $E$ soit contenu
dans $E'$. (Utiliser \hyperref[IV.exer9.2.7]{9.2.7}) On appellera ce $V$, ou
$\mathcal{V}$, l'\emph{intérieur} du sous-topos $E'$ de $E$. Les
ouverts $U$ et $V$ de $E$ sont disjoints $(U\cap V=\phi_E)$; le
sous-topos fermé de $E$ complémentaire\pageoriginale de
$\Sup(U,V)=\mathcal{U}\amalg \mathcal{V}$ s'appelle la
\emph{frontière} du sous-topos $E'$ de $E$.

\item Pour que $E'$ soit un sous-topos ouvert (\resp fermé) de $E$, il
faut et il suffit qu'il soit égal à son intérieur (\resp qu'il soit
égal à son adhérence). Pour que la frontière de $E'$ soit vide, il
faut et il suffit que $E'$ soit un sous-topos ouvert et fermé de
$E$, ou encore, qu'il soit le sous-topos ouvert de $E$ défini par un
ouvert $V$ de $E$ qui soit \emph{\sisi{sommande direct}{sommande directe}} de l'objet final
% correction SGA4-CORR-0256 / SGA-SRC-000936: SGA4-EIV-04-L6849-L6851-3947D4E19478.
$e$, \ie tel qu'il existe un sous-objet \sisi{V}{U} de $e$ avec $e\simeq
% correction SGA4-CORR-0257 / SGA-SRC-000937: SGA4-EIV-04-L6850-L6852-2AA67ED0B6B3.
U\amalg V$. Pour que l'extérieur de $E'$ soit « vide »  \ie pour
que l'adhérence soit égale à $E$, il faut et il suffit que
$i:E'\rightarrow E$ soit dominant (\hyperref[IV.exer8.8]{8.8}).

\item Soient $S$ un objet de $E$, $F=E_{/S}$ le topos induit, $F'\simeq
F_{/S'}$ le sous-topos de $F$ induit par le sous-topos $E'$ de $E$.
Montrer que l'extérieur (\resp l'adhérence, \resp
l'intérieur, \resp
la frontière) de $F'$ est induit par l'extérieur (\resp\ldots) de
$E'$.
\end{enumerate}
\end{enonce*}

\begin{enonce*}[remark]{Exercice 9.4.9}\etiquette[9.4.9]{IV.exer9.4.9}
(\emph{Sous-topos localement fermés d'un topos.}) Un sous-topos $F$
d'un topos $E$ est dit \emph{sous-topos localement fermé} de $E$
s'il est une intersection d'un sous-topos ouvert et d'un sous-topos
fermé.
\begin{enumerate}
\renewcommand{\theenumi}{\alph{enumi}}
\renewcommand{\labelenumi}{\theenumi)}
\item Prouver que l'intersection d'une famille finie de sous-topos
localement fermés de $E$ est localement fermé. Si $g:E'\rightarrow
E$ est un morphisme de topos, prouver que l'image inverse
(\hyperref[IV.exer9.1.6]{9.1.6} d)) d'un sous-topos localement fermé de $E$ est
un sous-topos localement fermé de $E'$.

\item Soient $X$ un espace topologique, et $E=\Top(X)$. Pour toute
partie $X'$ de $X$, considérons le sous-topos $T(X')$ de $E$ image
essentielle de $\Top(X')$ dans $\Top(X)$. Montrer que si $X'$ est
une partie localement fermée, $T(X')$ est un sous-topos localement
fermé de $E$, la réciproque étant vraie si on suppose de plus $X$ et
$X'$ sobres. Prouver que $X'\mapsto T(X')$ établit un isomorphisme
d'ensembles ordonnés entre l'ensemble des parties localement fermées
$X'$ de $X$ et l'ensemble des sous-topos localement fermés de~$E$.

\item Soit $F$ un sous-topos de $E$. Montrer que tout sous-topos ouvert
% correction SGA4-CORR-0258 / SGA-SRC-000938: SGA4-EIV-04-L6888-L6890-5B98819A3952.
de $F$\pageoriginale est induit par un sous-topos ouvert de $E$.
(Utiliser \hyperref[IV.exer9.1.8]{9.1.8} c)). Montrer que $F$ est un sous-topos
localement fermé de $E$ si et seulement si c'est un sous-topos
ouvert de son adhérence (\hyperref[IV.exer9.4.8]{9.4.8} a)). En conclure, en
utilisant \hyperref[IV.exer9.4.8]{9.4.8} d), que la propriété pour $F$
d'être un
sous-topos localement fermé de $E$ est une propriété
locale sur $E$.

\item Soit $F$ un sous-topos localement fermé de $E$. Montrer que
parmi toutes les façons d'écrire $F$ comme une intersection d'un
% correction SGA4-CORR-0259 / SGA-SRC-000939: SGA4-EIV-04-L6898-L6901-939255FEBA5D.
sous-topos ouvert $\mathcal{U}$ et d'un sous-topos fermé
$\mathcal{F}$ de $E$, il en est une avec $\mathcal{U}$ le plus grand
possible, et $\mathcal{F}$ le plus petit possible. (Prendre
$\mathcal{F} =$ adhérence de $F$, et $\mathcal{U} =$ plus grand
sous-topos ouvert induisant sur $\mathcal{F}$ le sous-topos $F$.)
Compatibilité de la formation de $\mathcal{U}$, $\mathcal{F}$ avec
la localisation sur $E$.
\end{enumerate}
\end{enonce*}

\begin{enonce*}[remark]{Exercice 9.4.10}\etiquette[9.4.10]{IV.exer9.4.10}
Développer la notion de \emph{sous-topos constructible, localement
constructible, quasi-constructible, localement quasi-constructible},
sur le modèle des notions familières dans le cas des espaces
topologiques (EGA ${\rm 0_{III}}$ \hyperref[IV.subsec9.1]{9.1}, EGA IV 10.1).
Dans le cas d'un topos de la forme $\Top(X)$, on retrouvera une
bijection entre l'ensemble des parties constructibles (\resp\ldots)
de $X$, et l'ensemble des sous-topos constructibles (\resp\ldots) de
$\Top(X)$.
\end{enonce*}

\subsection{Le théorème de recollement}\etiquette[9.5]{IV.subsec9.5}

\subsubsection{}\etiquette[9.5.1]{IV.subsubsec9.5.1}
Soit $f:A\rightarrow B$ un foncteur entre deux catégories. Désignons
par $(B,A,f)$ la catégorie suivante: les objets de $(B,A,f)$ sont
les triples
$$
(X,Y,u)
$$
où $X$ est un objet de $B$, $Y$ un objet de $A$, et $u$ un morphisme
de $X$ dans $f(Y)$; les morphismes entre deux objets $(X,Y,u)$ et
$(X',Y',u')$ sont les couples $(m,m')$, où $m$ est un morphisme de
$X$ dans $X'$ et $m'$ un morphisme de $Y$ dans $Y'$, tels que le
diagramme ci-après soit commutatif:
\[
\xymatrix@C=1.5cm@R=1cm{
X \ar[r]^{u} \ar[d]_{m} & **[r]f(Y) \ar[d]^{f(m')}\\
X' \ar[r]^{u'}& **[r]f(Y') }
\]
la\pageoriginale composition des morphismes est définie de la
manière évidente.

\subsubsection{}\etiquette[9.5.2]{IV.subsubsec9.5.2}
Remarquons que la catégorie $(B,A,f)$ dépend de manière fonctorielle
du système des données $A$, $B$, $f$, en un sens que nous laissons
au lecteur le soin de préciser. En particulier, si $f'$ est un
deuxième foncteur de $A$ dans $B$, alors tout morphisme de foncteurs
$$
f\xrightarrow{u} f'
$$
définit un foncteur \sisi{$(A,B,f)\rightarrow (A,B,f')$}{$(B,A,f)\rightarrow (B,A,f')$}, qui est un
% correction SGA4-CORR-0260 / SGA-SRC-000941: SGA4-EIV-04-L6944-L6952-49D5BB793FFD.
isomorphisme si $u:f\rightarrow f'$ l'est.

\subsubsection{}\etiquette[9.5.3]{IV.subsubsec9.5.3}
Revenons à la situation de \hyperref[IV.subsec9.3]{9.3}, et considérons la
catégorie $(\mathcal{F},\mathcal{U}, \rho)$ définie dans
\hyperref[IV.subsubsec9.5.1]{9.5.1}. À chaque objet $X$ de $E$ correspond un
objet $(i^{\ast}X, j^{\ast}X, h(X):i^{\ast}X\rightarrow \rho
j^{\ast}X)$ de $(\mathcal{F}, \mathcal{U},\rho)$ où $h(X)$ s'obtient
en transformant par le foncteur $i^{\ast}$ le morphisme d'adjonction
$X\rightarrow j_{\ast}j^{\ast}(X)$ (on a $\rho=i^{\ast}j_{\ast}$
\sisi{(9.10.)}{\hyperref[IV.eq9.3.6.3]{9.3.6.3}}). On a donc défini ainsi un foncteur
\begin{equation*}
\Phi:E\longrightarrow (\mathcal{F}, \mathcal{U},
\rho){\quoi}.\tag{9.5.3.1}
\end{equation*}\etiquette[9.5.3.1]{IV.eq9.5.3.1}

\begin{enonce*}{Théorème 9.5.4}\etiquette[9.5.4]{IV.theo9.5.4}
\begin{enumerate}
\renewcommand{\theenumi}{\alph{enumi}}
\renewcommand{\labelenumi}{\theenumi)}
\item Soient $E$ un topos, $\mathcal{U}$ un sous-topos ouvert de $E$,
$\mathcal{F}$ le sous-topos fermé complémentaire. Le foncteur
$$
\rho=i^{\ast}j_{\ast}:\mathcal{U} \longrightarrow \mathcal{F}
$$
est exact à gauche et accessible et le foncteur $\Phi$ $(9.5.3.1)$
est une équivalence de catégories.

\item Réciproquement, soient $\mathcal{U}$ et $\mathcal{F}$ deux topos et
$\rho:\mathcal{U}\rightarrow \mathcal{F}$ un foncteur de recollement
$(\hyperref[IV.subsubsec9.3.6]{9.3.6})$ (\ie accessible {\rm (\hyperrefext[I][I.def9.2]{9.2})}
et exact à gauche). La catégorie $E=(\mathcal{F},\mathcal{U},\rho)$
est un topos. Soient $\phi_{\mathcal{F}}$ un objet initial de
$\mathcal{F}$, $e_{\mathcal{U}}$ un objet final de $\mathcal{U}$, et
soit $U=(\phi_{\mathcal{F}},e_{\mathcal{U}},
\phi_{\mathcal{F}}\rightarrow \rho(e_{\mathcal{U}}))\in \Ob E$. Le
foncteur $X\mapsto (\phi_{\mathcal{F}},
X,\phi_{\mathcal{F}}\rightarrow \rho(X))$ induit une équivalence de
$\mathcal{U}$ sur $E_{/U}$, donc avec le sous-topos ouvert
$\mathcal{U}'$ de $E$ défini par $U$. Le foncteur $Y\mapsto
(Y,e,Y\rightarrow \rho(e))$ est un plongement fermé
$(\hyperref[IV.subsubsec9.3.5]{9.3.5})$, \ie induit une équivalence de
$\mathcal{F}$ avec un sous-topos fermé $\mathcal{F}'$ de $E$. Les
sous-topos $\mathcal{U}'$ et $\mathcal{F}'$\pageoriginale sont
complémentaires. Soit $\rho':\mathcal{U}'\rightarrow \mathcal{F}'$
le foncteur de recollement. Les foncteurs $\rho$ et $\rho'$ sont
compatibles avec les équivalences explicitées ci-dessus.
\end{enumerate}
\end{enonce*}

\begin{enonce*}{Définition 9.5.5}\etiquette[9.5.5]{IV.def9.5.5}
Soient $\mathcal{U}$ et $\mathcal{F}$ deux topos et
$\rho:\mathcal{U}\rightarrow \mathcal{F}$ un foncteur de recollement
$(\hyperref[IV.subsubsec9.3.6]{9.3.6})$. Le topos $(\mathcal{F},
\mathcal{U},\rho)$ est appelé le topos construit à partir de
$\mathcal{U}$ et de $\mathcal{F}$ par recollement à l'aide du
foncteur $\rho$.
\end{enonce*}

\setcounter{subsubsection}{5}
\subsubsection{}\etiquette[9.5.6]{IV.subsubsec9.5.6}
\emph{Démonstration de} \hyperref[IV.theo9.5.4]{9.5.4} a) \footnote[1]{Pour un
énoncé de recollement plus général que \hyperref[IV.def9.5.5]{9.5.5}, \cf
exercice\nde{Cet exercice ne semble pas exister.}.} La première assertion n'est qu'un rappel de
\hyperref[IV.prop9.4.1]{9.4.1} d). Le couple de foncteurs $(i^{\ast},j^{\ast})$
est conservatif, ou encore, fidèle (\hyperref[IV.prop9.4.1]{9.4.1} c)). A
fortiori le foncteur $\Phi$ est fidèle. Montrons qu'il est
pleinement fidèle. Il résulte de \hyperref[IV.prop9.3.3]{9.3.3} b) que, pour
tout objet $X$ de $E$, le diagramme ci-après est cartésien:
\[\vcenter{
\xymatrix@C=2.5cm{
X  \ar[r] \ar[d]               & j_{\ast}j^{\ast}X \ar[d]\\
i_{*}i^{*}X   \ar[r]     & i_{\ast}i^{\ast}j_{\ast}j^{\ast}X ;
}}\leqno{(*)}
\]
Soient alors $X$ et $Y$ deux objets de $E$, et $(m,m')$:
\[\vcenter{
\xymatrix@C=2.75cm@R=1cm{
 i^{\ast}X  \ar[r] \ar[d]_{m}               & i^{\ast}j_{\ast}j^{\ast}X \ar[d]_{i^{\ast}j_{\ast}m'}\\
i^{*}Y   \ar[r]     & i^{\ast}j_{\ast}j^{\ast}Y ; }}
\]
un morphisme entre $\Phi(X)$ et $\Phi(Y)$. En utilisant la
fonctorialité du produit cartésien et les diagrammes $(\ast)$, on
obtient un morphisme $w:X\rightarrow Y$ qui rend \sisi{commutatif}{commutatifs} les
% correction SGA4-CORR-0261 / SGA-SRC-000960: SGA4-T1-IV7035-01-9926B6E7.
diagrammes ci-après:
\begin{displaymath}
\xymatrix@C=.5cm@R=1.25cm{
X \ar@{-}[rr]^{w} \ar[d]  &             &Y \ar[d]\\
 j_{\ast}j^{\ast}X \ar[rr]^{j_{\ast}m'}       &   &
j_{\ast}j^{\ast}Y , }\qquad\qquad \xymatrix@C=.5cm@R=1.25cm{
X \ar@{-}[rr]^{w} \ar[d]& &Y \ar[d]\\
 i_{\ast}i^{\ast}X \ar[rr]^{i_{\ast}m}&   & i_{\ast}i^{\ast}Y
}
\end{displaymath}
En\pageoriginale appliquant respectivement les foncteurs $j^{\ast}$
et $i^{\ast}$ on obtient: $j^{\ast}w=m'$, $i^{\ast}w=m$. Montrons
maintenant que le foncteur $\Phi$ est essentiellement surjectif.
Soit $(Y,X,u:Y\rightarrow i^{\ast}j_{\ast}(X))$ un objet de
$(\mathcal{F},\mathcal{U},i^{\ast} j_{\ast})$. On en déduit un
diagramme:
\[
\xymatrix@C=2cm{
                             & j_{\ast}X \ar[d]\\
i_{\ast}Y \ar[r]^{i_{\ast}u} & i_{\ast}i^{\ast}j_{\ast}X }
\]
D'où, en prenant le produit fibré, un diagramme cartésien:
\[\vcenter{
  \xymatrix@C=2.75cm@R=1.25cm{
W \ar[r] \ar[d]& j_{\ast}X \ar[d]\\
i_{\ast}Y \ar[r]^{i_{\ast}u}& i_{\ast}i^{\ast}j_{\ast} X\\
}}\leqno{(**)}
\]

Nous laissons au lecteur, le soin de montrer que $\Phi(W)$ est
isomorphe à $(Y,X,u:Y\rightarrow i^{\ast}j_{\ast}(X))$ et que le
diagramme (${\ast}{\ast}$) est isomorphe au diagramme ($\ast$)
relatif à $W$.

\subsubsection{}
\emph{Démonstration} de \hyperref[IV.theo9.5.4]{9.5.4} \sisi{1)}{b)}. Nous nous bornerons à
montrer que $(\mathcal{F},\mathcal{U},\rho)$ est un topos. La
démonstration des autres assertions est laissée au lecteur.
Vérifions les propriétés a), b), c), d) de \hyperref[IV.subsubsec1.1.2]{1.1.2}.
\begin{enumerate}
\renewcommand{\theenumi}{\alph{enumi}}
\renewcommand{\labelenumi}{\theenumi)}
\item \emph{Les limites projectives finies sont représentables}: Clair
car les catégories $\mathcal{U}$ et $\mathcal{F}$ sont des topos et
$\rho$ commute aux limites projectives finies.

\item \emph{Les sommes directes sont représentables, disjointes et
universelles}: Soit $(Y_i\xrightarrow{u_i}\rho(X_i))$, $i\in I$, une
famille d'objets de $(\mathcal{F},\mathcal{U},\rho)$. Par définition
des limites inductives, on a un morphisme canonique
$$
\coprod_{i\in I}\rho(X_i)\longrightarrow \rho(\coprod_{i\in I} X_i)
$$
d'où,\pageoriginale en composant avec le morphisme $\coprod{u_i}$,
un morphisme
$$
n:\coprod_{i\in \sisi{1}{I}} Y_i\longrightarrow
\sisi{(\coprod_{i\in I}X_i)}{\rho(\coprod_{i\in I}X_i)}{\quoi}.
% correction SGA4-CORR-0262 / SGA-SRC-000961: SGA4-T1-IV7092-02-47D253B6.
$$
On vérifie immédiatement que ce dernier objet est la somme directe
de la famille considérée, et que cette somme directe est disjointe
et universelle.

\item \emph{Les relations d'équivalence sont effectives et
universelles}:

Soit
\[
\xymatrix@C=2cm@R=1.5cm{
R_1 \ar@<3pt>[r] \ar@<-3pt>[r] \ar[d]& Y \ar[d]^{w}\\
\rho(R_2) \ar@<3pt>[r]^{f(u)} \ar@<-3pt>[r]_{f(v)} & \rho(X) }
\]
une relation d'équivalence sur l'objet $Y\xrightarrow{w}\rho(X)$. La
relation $\xymatrix@C=.6cm{R_1 \ar@<-2pt>[r]\ar@<2pt>[r]& Y}$ est
alors une relation d'équivalence sur $Y$ dans $\mathcal{F}$, et la
relation $\xymatrix{
R_2 \ar@<2pt>[r]^{u} \ar@<-2pt>[r]_{v}& X\\
}$ est \sisi{un}{une} relation d'équivalence sur $X$ dans $\mathcal{U}$. Les
% correction SGA4-CORR-0263 / SGA-SRC-000962: SGA4-T1-IV7112-03-CE696FD7.
quotients $Y/R_1$, $\rho(X)/\rho(R_2)$, $X/R_2$ sont représentables
dans $\mathcal{F}$ et $\mathcal{U}$ car ces catégories sont des
topos. De plus, par définition des limites inductives, le morphisme
canonique $\rho(X)\rightarrow \rho(X/R_2)$ se factorise par
$\rho(X)/\rho(R_2)$. On en déduit, par la fonctorialité des limites
inductives, un morphisme canonique
$Y/R_1\xrightarrow{t}\rho(X/R_2)$. On vérifie que ce dernier objet
est le quotient de $Y\xrightarrow{w}\rho(X)$ par la relation
d'équivalence considérée, que ce quotient est effectif (I~10) et que
toutes ces propriétés sont conservées par changement de base.

\item $(\mathcal{F},\mathcal{U},\rho)$ \emph{admet une petite famille
génératrice.} Cette propriété résulte de \sisi{(1.9.2 5)}{\hyperrefext[I][I.cor9.25]{9.25}} (compte tenu de
ce que $\rho$ est accessible), en prenant
$B=\Delta_1=(0\xrightarrow{f}1)$, $E_0=\mathcal{F}$,
$E_{1}=\mathcal{U}$, $f^{\ast}=\rho$: $E_1\rightarrow E_0$, d'où
$\Hhom_B(B,E)=(\mathcal{F},\mathcal{U},\rho)$.

On peut aussi éviter le recours à \sisi{(1.9.2 5)}{\hyperrefext[I][I.cor9.25]{9.25}} (dont la démonstration
donnée était assez pénible !) et utiliser l'hypothèse sur $\rho$
sous la forme que $\rho$ admet un pro-adjoint
$$
\sigma:\Pro(\mathcal{F})\longrightarrow \Pro(\mathcal{U}){\quoi},
$$
(cette\pageoriginale interprétation (\hyperrefext[I][I.sec8]{8}) reposant sur
(\hyperrefext[I][I.sec9]{9}), dont la démonstration est nettement plus
compréhensible). Il suffit alors de noter que si $C'$ (\resp $C''$)
est une sous-catégorie génératrice de $\mathcal{U}$ (\resp
$\mathcal{F}$), et si pour tout $X''\in \Ob C''$, on représente
$\sigma(X'')$ sous forme d'un système projectif
$(\sigma(X'')_i)_{i\in \sisi{I(X')}{I(X'')}}$, où
$\sisi{I(X')}{I(X'')}$ est un petit ensemble
% correction SGA4-CORR-0264 / SGA-SRC-000963: SGA4-T1-IV7143-04-28086F30.
ordonné filtrant, alors la famille des objets de $E$ qui sont, soit
de la forme $(\phi_{\mathcal{F}}, X',\phi_{\mathcal{F}}\rightarrow
\sisi{\phi(X')}{\rho(X')})$ avec $X'\in \Ob C'$, soit de la forme $(X'',\sigma
% correction SGA4-CORR-0265 / SGA-SRC-000964: SGA4-T1-IV7145-05-28086F30.
(X'')_i, u_{X''}:X''\rightarrow \rho(\sigma (X'')_i)$, où $X''\in \Ob
{C''}$, $i\in I (X'')$, et $u_{X''}$ correspond par adjonction au
morphisme canonique $\sigma(X'')\rightarrow \sigma (X'')_i$, --est
une famille génératrice (évidemment petite) si $C'$, $C''$ sont
choisis petits. La vérification de ce fait est effectivement
immédiate, et laissée au lecteur.
\end{enumerate}

\setcounter{subsubsection}{4} \setcounter{paragraph}{7}
\paragraph{}\etiquette[9.5.4.8]{IV.par9.5.4.8}
Utilisons les notations de 9.5.7 b), et montrons comment on peut
expliciter, à isomorphismes canoniques près de foncteurs, le système
de foncteurs \eqhyperref[IV.eq9.3.6.2]{9.3.6.2}, où on fait
$E=(\mathcal{F},\mathcal{U},\rho)$. Le détail des vérifications des
assertions ci-dessous (essentiellement mécaniques à l'aide de
\hyperref[IV.theo9.5.4]{9.5.4}) est laissé au lecteur.

\setcounter{subsubsection}{8} \setcounter{paragraph}{0}
\paragraph{}\etiquette[9.5.8.1]{IV.par9.5.8.1}
Le foncteur
$$
i^{\ast}:(\mathcal{F},\mathcal{U},\rho)\longrightarrow \mathcal{F}
$$
est donné par
$$
i^{\ast}(X,Y,u:X\longrightarrow \rho(Y))=X{\quoi}.
% correction SGA4-CORR-0266 / SGA-SRC-000966: SGA4-T1-IV7172-07-7A61D30E.
$$

\paragraph{}\etiquette[9.5.8.2]{IV.par9.5.8.2}
Le foncteur
$$
i_{\ast}:\mathcal{F}\longrightarrow (\mathcal{F},\mathcal{U},\rho)
$$
est donné par $i_{\ast}(X)=(X,e,X\rightarrow \rho(e))$ ($e=$ objet
final de $\mathcal{U}$).

\paragraph{}\etiquette[9.5.8.3]{IV.par9.5.8.3}
Le morphisme d'adjonction
$$
\id\longrightarrow i_{\ast}i^{\ast}
$$
est isomorphe au morphisme fonctoriel
\[
\xymatrix@C=2cm{
X \ar[r]^{u} \ar[d]_{id}& \rho(Y) \ar[d]\\
X \ar[r]                & \rho(e) }
\]

\paragraph{}\etiquette[9.5.8.4]{IV.par9.5.8.4}
Le\pageoriginale foncteur
$$
j^{\ast}:(\mathcal{F},\mathcal{U},\rho)\longrightarrow \mathcal{U}
$$
est donné par\qquad $j^{\ast}(X,Y,u:X\longrightarrow \rho(Y))=Y$.

\paragraph{}\etiquette[9.5.8.5]{IV.par9.5.8.5}
Le foncteur
$$
j_{!}:\mathcal{U}\longrightarrow (\mathcal{F}, \mathcal{U}, \rho)
$$
est donné par\qquad
$j_{!}(Y)=(\phi_{\mathcal{F}},Y,\phi_{\mathcal{F}}\longrightarrow
\rho(Y))$

\noindent \sisi{($\phi$ objet initial de)}{($\phi$ objet initial de $\mathcal F$)}
% correction SGA4-CORR-0267 / SGA-SRC-000967: SGA4-T1-IV7211-08-7A09DEB3.

\paragraph{}\etiquette[9.5.8.6]{IV.par9.5.8.6}
Le foncteur
$$
j_{\ast}:\mathcal{U}\longrightarrow (\mathcal{F}, \mathcal{U},\rho)
$$
est donné par\qquad $j_{\ast}(Y)=(\rho(Y),Y,\id
:\rho(Y)\longrightarrow \rho(Y))$.

\paragraph{}\etiquette[9.5.8.7]{IV.par9.5.8.7}
Le morphisme d'adjonction
$$
j_{!}j^{\ast}\longrightarrow \id
$$
est isomorphe au morphisme fonctoriel
\[
\xymatrix{
\phi_{\Ff} \ar[r] \ar[d] & \rho(Y) \ar[d]^{\rho(\id)}\\
X \ar[r]^{u}                 & \rho(Y)     &. }
\]

\paragraph{}\etiquette[9.5.8.8]{IV.par9.5.8.8}
Le morphisme d'adjonction
$$
\id\longrightarrow j_\ast j^{\ast}
$$
est isomorphe \sisi{aux}{au} morphisme fonctoriel
% correction SGA4-CORR-0268 / SGA-SRC-000968: SGA4-T1-IV7238-09-5083DEA5.
\[
\xymatrix{
X  \ar[r]^{u} \ar[d]_{u}& \rho(Y) \ar[d]^{\rho(\id)}\\
\rho(Y) \ar[r]^{\id}& \rho(Y) }
\]
\begin{enonce*}[remark]{Exercice 9.5.9}\etiquette[9.5.9]{IV.exer9.5.9}
(\emph{Propriétés d'exactitude d'un topos recollé.}) Soient
$\mathcal{U}$, $\mathcal{F}$ deux topos,
$\rho:\mathcal{U}\rightarrow \mathcal{F}$ un foncteur de
recollement, $E=(\mathcal{F}, \mathcal{U},\rho)$ le topos recollé,
$j:\mathcal{U}\rightarrow E$ et $i:\mathcal{F}\rightarrow E$ les
morphismes de topos canoniques, de sorte que
$\rho=i^{\ast}j_{\ast}$.
\begin{enumerate}
\renewcommand{\theenumi}{\alph{enumi}}
\renewcommand{\labelenumi}{\theenumi)}
\item Prouver que $i^{\ast}$ commute aux (petits) produits (\ie (\hyperref[IV.rem1.8]{1.8}
ou \sisi{\hyperref[IV.rem1.8]{1.8}}{\hyperrefext[I][I.sec8]{8}}) le foncteur $i^{\ast}$ admet un adjoint à gauche
$i_{!}$) si et seulement si il en est de même de $\rho$, \ie
(\loccit) si et seulement si $\rho$ admet un adjoint à gauche
$\sigma:\mathcal{F}\rightarrow \mathcal{U}$ (ou encore si et
seulement si le proadjoint de $\rho$ restreint\pageoriginale à
$\mathcal{F}$ \eqhyperref[IV.eq9.6.3]{9.6.3} se factorise à isomorphisme près par
$\mathcal{U}$): on a alors $\sigma=j^{\ast}i_{!}$. Donc les
foncteurs de recollement $\rho$ ayant cette propriété correspondent
à isomorphisme près aux foncteurs $\sigma:\mathcal{F}\rightarrow
\mathcal{U}$ qui ont un adjoint à droite, \ie (\hyperref[IV.cor1.6]{1.6}) qui
commutent aux $𝒰$-limites inductives, ou encore qui sont
continus. Lorsque $\mathcal{F}$ est équivalent à un topos
$C^{\sim}$, où $C$ est un $𝒰$-site, ces foncteurs
correspondent donc à isomorphisme près aux foncteurs continus
$C\rightarrow \mathcal{U}$ (\hyperrefext[III][III.prop1.2]{1.2} et \hyperrefext[III][III.prop1.7]{1.7}).

\item Supposons que le foncteur $i_{!}$ soit défini. Montrer que pour
que $i_{!}$ commute à un type déterminé de petites limites
projectives, il faut et il suffit que $\sigma:\mathcal{F}\rightarrow
\mathcal{U}$ y commute. (Pour la nécessité, utiliser que $j^{\ast}$
commute aux petites limites projectives; pour la suffisance,
utiliser la description équivalente de $E$ en termes de $\sigma$
comme formé des triples $(Y,X,u)$ avec $Y\in \Ob \mathcal{F}$, $X\in
\Ob \mathcal{U}$, $u:\sigma (Y)\rightarrow X$.)

\item Déduire de a) et b) un exemple de plongement fermé de topos
$i:\mathcal{F}\rightarrow E$ tel que $i_{!}$ existe, mais ne commute
ni aux produits fibrés, ni aux limites projectives filtrantes. (Il
suffit de trouver un foncteur continu entre deux topos qui ne
possède aucune de ces deux propriétés d'exactitude.)
\end{enumerate}
\end{enonce*}

\begin{enonce*}[remark]{Exercice 9.5.10}\etiquette[9.5.10]{IV.exer9.5.10}
(\emph{Recollement de topos de la forme $\widehat{C}$.}) Soient $C$
une petite catégorie, et $E=\widehat{C}$, qui est un
$𝒰$-topos.
\begin{enumerate}
\renewcommand{\theenumi}{\alph{enumi}}
\renewcommand{\labelenumi}{\theenumi)}
\item Se rappeler que les ouverts $U$ de $E$ correspondent aux
\emph{cribles} $C'$ de $C$ \sisi{(1.4.1, \hyperref[IV.subsec4.2]{4.2})}{(\hyperrefext[I][I.def4.1]{4.1}, \hyperref[IV.subsubsec8.4.4]{8.4.4}}, $E_{/U}$
étant alors équivalent à $\widehat{C'}$, le foncteur
$j^{\ast}:E\rightarrow E_{/U}$ s'identifiant au foncteur restriction
\sisi{(4.5.11)}{\ref{VI.lem5.4}}, \ie le morphisme d'inclusion $j:E_{/U}\rightarrow E$ étant
$\widehat{f}:\widehat{C'}\rightarrow \widehat{C}$
(\hyperref[IV.subsubsec4.6.1]{4.6.1}), où $f:C'\rightarrow C$ est l'inclusion.
Soit $C''$ la sous-catégorie pleine de $C$ complémentaire de $C'$
(\ie $\Ob C''$ est le complémentaire de $\Ob C'$ dans $\Ob C$),
$g:C''\rightarrow C$ l'inclusion. Alors le sous-topos fermé
$\mathcal{F}$ de $E=\widehat{C}$, complémentaire du sous-topos
ouvert défini par $C'$, est canoniquement équivalent à
$\widehat{C''}$, le morphisme d'inclusion $i:\mathcal{F}\rightarrow
E$ s'identifiant à $\widehat{g}:\widehat{C''}\rightarrow
\widehat{C}$.

\item Considérer\pageoriginale le foncteur
\begin{equation*}
h:{C'}^0\times C''\longrightarrow (\Ens),\quad
h(X',X'')=\Hom_C(X',X''){\quoi},\tag{9.5.10.1}
\end{equation*}\etiquette[9.5.10.1]{IV.eq9.5.10.1}
et noter que $C$ se reconstitue à isomorphisme près, avec sa
sous-catégorie pleine $C'$, par la connaissance des catégories $C'$,
$C''$ et du bifoncteur $h$ (ce dernier pouvant être choisi
arbitrairement). (Comparer \hyperref[IV.subsec9.7.1]{9.7.1} plus bas.) La donnée
de $h$ équivaut à celle d'un foncteur $C''\rightarrow
\widehat{C'}=\Hhom({C'}^0,(\Ens))$, ou encore (\hyperrefext[III][III.prop1.2]{1.2} et
\hyperrefext[III][III.prop1.7]{1.7}) à celle d'un foncteur
\begin{equation*}
\sigma:\sisi{\widehat{C}''}{\widehat{C''}}\longrightarrow
\sisi{\widehat{C}'}{\widehat{C'}}\tag{9.5.10.2}
% correction SGA4-CORR-0269 / SGA-SRC-000969: SGA4-T1-IV7325-10-3ABBEA8D.
\end{equation*}\etiquette[9.5.10.2]{IV.eq9.5.10.2}
continu, \ie (\hyperref[IV.cor1.6]{1.6}) commutant aux petites limites
inductives, ou encore, admettant un adjoint à droite.
\begin{equation*}
\rho:\widehat{C}'\longrightarrow \widehat{C''}{\quoi}.\tag{9.5.10.3}
\end{equation*}\etiquette[9.5.10.3]{IV.eq9.5.10.3}
Montrer que $\rho$ n'est autre que le foncteur de recollement
associé au sous-topos ouvert $\widehat{C}'$ de $E=\widehat{C}$.

\item Déduire de a) et b) que si on se donne un foncteur de
recollement \eqhyperref[IV.eq9.5.10.3]{9.5.10.3} entre deux catégories de
préfaisceaux,
d'où un topos recollé $E$, les conditions suivantes sont
équivalentes:
\begin{itemize}
\item[(i)] $E$ est équivalent à un topos de la forme $\widehat{C}$.

\item[(ii)] Le foncteur $\rho$ commute aux petits produits (ou encore, il
admet un adjoint à gauche \eqhyperref[IV.eq9.5.10.2]{9.5.10.2}\sisi{}{)}.
% correction SGA4-CORR-0270 / SGA-SRC-000970: SGA4-T1-IV7343-11-9708B68A.

\item[(iii)] Le morphisme d'inclusion $i:\widehat{C''}\rightarrow E$ est « essentiel »  \ie $i^{\ast}$ commute aux produits, ou encore
$i^{\ast}$ admet un adjoint à gauche $i_{!}$.
\end{itemize}
Lorsqu'il en est ainsi, expliciter une catégorie $C$ telle que
$E\approx \widehat{C}$ à l'aide du bifoncteur \eqhyperref[IV.eq9.5.10.1]{9.5.10.1}
défini par la restriction de \eqhyperref[IV.eq9.5.10.2]{9.5.10.2} à $C''$.

\item Montrer qu'un topos obtenu en recollant deux topos ponctuels
n'est pas nécessairement équivalent à un topos de la forme
$\widehat{C}$.
\end{enumerate}
\end{enonce*}

\begin{enonce*}[remark]{Exercice 9.5.11}\etiquette[9.5.11]{IV.exer9.5.11}
(Morphisme d'un topos recollé dans un topos.)
\begin{enumerate}
\renewcommand{\theenumi}{\alph{enumi}}
\renewcommand{\labelenumi}{\theenumi)}
\item Soient: $\rho \sisi{}{: \mathcal{U}} \rightarrow \mathcal{F}$ un foncteur entre deux
catégories, $E=(\mathcal{F},\mathcal{U},\rho)$\pageoriginale la
catégorie « recollée »  de \hyperref[IV.subsubsec9.5.1]{9.5.1}, et $F$ une
catégorie quelconque. Considérer le foncteur $f\mapsto \rho\circ f$:
$$
\bar{\rho}:\bar{\mathcal{U}}=\Hhom(F,\mathcal{U})\longrightarrow
\bar{\mathcal{F}}=\Hhom(F,\mathcal{F}){\quoi},
$$
et définir une équivalence de catégories
$$
\Hhom(F,E)\xrightarrow{\approx}
\sisi{(\bar{\mathcal{U}},\bar{\mathcal{F}},\bar{\rho})}
{(\bar{\mathcal{F}},\bar{\mathcal{U}},\bar{\rho})}{\quoi}.
% correction SGA4-CORR-0271 / SGA-SRC-000971: SGA4-T1-IV7369-12-23D06E4B.
$$

\item Soit $f:F\rightarrow E$ un foncteur, et considérons ses composés
avec les foncteurs canoniques $j^{\ast}:E\rightarrow \mathcal{U}$ et
$i^{\ast}:E\rightarrow \mathcal{F}$. Montrer que $f$ commute à un
type déterminé de limites inductives (supposé représentable dans
$\mathcal{U}$ et $\mathcal{F}$, donc dans $E$) si et seulement si
$j^{\ast}f$ et $i^{\ast}f$ y commutent. Montrer que si $\rho$
commute à un type déterminé de limites projectives (qui est donc
représentable dans $E$), alors $f$ y commute si et seulement si il
en est de même de $j^{\ast}f$ et $i^{\ast}f$.

\item Supposons maintenant que $\mathcal{U}$ et $\mathcal{F}$ soient
  des topos, et
$\rho$ un foncteur de recollement. Soit $F$ un topos,
$f:F\rightarrow E$ un foncteur. Conclure de b) que $f$ est un
foncteur image inverse associé à un morphisme de topos $E\rightarrow
F$ si et seulement si il en est de même des foncteurs
$j^{\ast}f:F\rightarrow \mathcal{U}$ et $i^{\ast}f:F\rightarrow
\mathcal{F}$. Conclure de ceci et de a) que l'on a une équivalence
entre la catégorie $\Hhomtop(E,F)$, et la catégorie des triples
$(f',f'',u)$, où $f'$ (\resp $f''$) est un objet de
$\Hhomtop(\mathcal{U}, F)$ (\resp $\Hhomtop(\mathcal{F}, F)$), et où
$u:{f''}^{\ast}\rightarrow \rho\circ {f'}^{\ast}$ est un
homomorphisme de foncteurs $F\rightarrow \mathcal{F}$: si $f$ est un
objet de $\Hhomtop(E,F)$, $(f',f'',u)$ le triple correspondant,
alors $f'=f\circ j$, $f''=f\circ i$.

\item Préciser en quel sens on peut dire que c) fournit une
caractérisation à équivalence près (dans une 2-catégorie de topos)
du topos $E$ en termes du triple $(\mathcal{U},\mathcal{F},\rho)$.
\end{enumerate}
\end{enonce*}

\begin{enonce*}[remark]{Exercice 9.5.12}\etiquette[9.5.12]{IV.exer9.5.12}
(\emph{Recollement de deux espaces topologiques.})
\begin{enumerate}
\renewcommand{\theenumi}{\alph{enumi}}
\renewcommand{\labelenumi}{\theenumi)}
\item Soient $X$ un espace topologique, $\rho:\Top(X)\rightarrow
(\Ens)=\Top(pt)$ un foncteur de recollement, \ie (\hyperrefext[I][I.sec8]{8}) un
foncteur proreprésentable, $A\in \Ob \Pro(\Top(X))$ l'objet qui le
pro-représente, $E$ le topos déduit de $\rho$ par recollement.
Montrer que $E$ est équivalent à un topos de la forme $\Top(Y)$ si
et seulement si $A$ est isomorphe à un objet de $\Pro(\Ouv(X))$.
En\pageoriginale conclure que si $X$ n'est pas vide, on peut trouver
$\rho$ tel que $E$ ne soit pas équivalent à un topos de la forme
$\Top(Y)$.

\item Plus généralement, considérons un foncteur recollement
$\rho:\Top(X)\rightarrow \Top(Y)$, avec $X$, $Y$ deux espaces
topologiques. En déduire une application
$$
\phi:\sisi{X}{Y}\longrightarrow \Ob \Pro(\Top(\sisi{Y}{X}))
$$
(composée de $\sisi{X}{Y}\rightarrow
\Ob \Ppoint(\Top(\sisi{X}{Y}))\rightarrow
\Ob\Pro(\Top(\sisi{X}{Y}))$ (\hyperref[IV.subsubsec6.8.5]{6.8.5}) et du pro-adjoint
$\sigma$ de $\rho$. (N.B. pour $\sisi{x}{y}\in \sisi{X}{Y}$,
$\phi(\sisi{x}{y})$ s'identifie à la
\emph{fibre en $\sisi{x}{y}$ du faisceau de recollement} \eqhyperref[IV.eq9.6.4]{9.6.4}
\emph{défini par} $\rho$.) Montrer que si $E$ est équivalent à un
topos de la forme $\Top(Z)$, alors $\phi$ prend ses valeurs dans
l'image essentielle de $\Pro(\Ouv(\sisi{Y}{X}))$.
% correction SGA4-CORR-0272 / SGA-SRC-000972: SGA4-T1-IV7414-13-EBD2ED5E.
\end{enumerate}
\end{enonce*}

\begin{enonce*}[remark]{Problème}
$A$-t-on une réciproque ?
\end{enonce*}

\subsection{Faisceau de recollement}\etiquette[9.6]{IV.subsec9.6}

Nous avons vu dans \hyperref[IV.theo9.5.4]{9.5.4} que la donnée d'un topos $E$
muni d'un ouvert $U$ revient essentiellement à celle de deux topos
$\mathcal{U}$ et $\mathcal{F}$ et d'un foncteur de recollement
\begin{equation*}
\rho:\mathcal{U}\longrightarrow \mathcal{F}\tag{9.6.1}
\end{equation*}\etiquette[9.6.1]{IV.eq9.6.1}
\ie d'un foncteur admettant un pro-adjoint
\begin{equation*}
\sigma:\Pro(\mathcal{F})\longrightarrow
\Pro(\mathcal{U}){\quoi}.\tag{9.6.2}
\end{equation*}\etiquette[9.6.2]{IV.eq9.6.2}
D'après le sorite sur les foncteurs proadjoints (\hyperrefext[I][I.sec8]{8}), le
foncteur $\rho$ est connu à isomorphisme unique près (et par suite
le topos recollé $E=(\mathcal{F},\mathcal{U},\rho)$ est connu à
équivalence de topos près) quand on connaît le foncteur $\sigma$, ou
encore, le foncteur induit par $\sigma$
\begin{equation*}
\sigma_{\circ}:\mathcal{F}\longrightarrow
\Pro(\mathcal{U}){\quoi}.\tag{9.6.3}
\end{equation*}\etiquette[9.6.3]{IV.eq9.6.3}
Le foncteur $\sigma$ est déterminé, à isomorphisme unique près, par
la propriété de commuter aux petites limites projectives filtrantes
et de prolonger le foncteur $\sigma_\circ$. D'autre part, pour qu'un
foncteur donné \eqhyperref[IV.eq9.6.3]{9.6.3} soit isomorphe à un pro-adjoint, il
faut et il suffit évidemment que pour tout objet $X$ de
$\mathcal{U}$, le foncteur contravariant
$$
Y\longmapsto \Hom_{\Pro(\mathcal{U})}(\sigma_{\circ}(Y),X)
$$
sur\pageoriginale $\mathcal{F}$ soit représentable, ou encore, que
ce soit un faisceau sur $\mathcal{F}$ pour la topologie canonique
(\sisi{\hyperref[IV.subsubsec1.1.2]{1.1.2} (iii)}{\hyperref[IV.theo1.2]{1.2} iii)}). On voit tout de suite que ceci
signifie aussi que le foncteur opposé
\begin{equation*}
\sigma_{\circ}^{\circ}:\mathcal{F}^{\circ}\longrightarrow
\Pro(\mathcal{U})^{\circ}\subset
\check{\mathcal{U}}=\Hhom(\mathcal{U},(\Ens))\tag{9.6.4}
\end{equation*}\etiquette[9.6.4]{IV.eq9.6.4}
est un \emph{faisceau} sur $\mathcal{F}$ à valeurs dans
$\Pro(\mathcal{U})^{\circ}$ (\hyperrefext[II][II.def6.1]{6.1}). Ainsi, \emph{la
donnée d'un foncteur de recollement} \eqhyperref[IV.eq9.6.1]{9.6.1} \emph{équivaut
aussi} (à isomorphisme près) \emph{à celle d'un faisceau sur
$\mathcal{F}$ à valeurs dans} $\Pro(\mathcal{U})^{\circ}$. Le
faisceau ainsi associé à un foncteur de recollement (ou à une
situation $(E,U,\mathcal{F})$ comme dans (\hyperref[IV.subsec9.3]{9.3})
s'appelle le \emph{faisceau de recollement} associé au foncteur de
recollement envisagé $\rho$ (\resp à l'ouvert $U$ du topos $E$).

\begin{enonce*}[remark]{Exercice 9.6.5}\etiquette[9.6.5]{IV.exer9.6.5}
Préciser \hyperref[IV.theo9.5.4]{9.5.4} en un énoncé de 2-équivalence de
2-catégories, entre la 2-catégorie formée des couples $(E,U)$ d'un
topos $E$ et d'un ouvert $U$ dudit, et la 2-catégorie formée des
triples $(\mathcal{U},\mathcal{F},\rho)$ de topos $\mathcal{U}$,
$\mathcal{F}$ et d'un foncteur de recollement
$\rho:\mathcal{U}\rightarrow \mathcal{F}$ (la description explicite
des 2-catégories en question étant à la charge du lecteur). Donner
une variante de cet énoncé faisant intervenir la 2-catégorie des
triples $(\mathcal{U},\mathcal{F}, \phi)$, où $\phi$ est maintenant
un faisceau sur $\mathcal{F}$ à valeurs dans
$\Pro(\mathcal{U})^{\circ}$.
\end{enonce*}

\subsection{Points d'un topos recollé}\etiquette[9.7]{IV.subsec9.7}

\subsubsection{}\etiquette[9.7.1]{IV.subsec9.7.1}
Soient $C$ une catégorie, et
$$
u:D'\longrightarrow C\quav v:D''\longrightarrow C
$$
deux foncteurs pleinement fidèles, tels que les images essentielles
$C'$, $C''$ de $u$ et de $v$ soient telles que
$$
\Ob C'\cap \Ob C''=\phi\quav \Ob C'\cup \Ob C''=\Ob C{\quoi}.
$$
Supposons de plus que
$$
X'\in \Ob C'\quav X''\in \Ob C''\Rightarrow
\Hom(X'',X')=\phi{\quoi},
$$
ou ce qui revient au même, qu'on ait
$$
\Hom(v(A''), u(A'))=\phi\mbox{ pour } A'\in \Ob D', A''\in \Ob
D''{\quoi}.
$$
Il est alors immédiat que la catégorie $C$ se reconstitue, à
% correction SGA4-CORR-0273 / SGA-SRC-000973: SGA4-T1-IV7530-14-B2CBAC9A.
isomorphisme canonique\pageoriginale près, par la connaissance des
catégories $C'$ et $C''$ et du foncteur
$$
(X',X'')\longmapsto \Hom(X',X''):{C'}^{\circ}\times
C''\longrightarrow (\Ens){\quoi};
$$
de même, $C$ se reconstitue à équivalence près (celle-ci définie à
isomorphisme unique près) par la connaissance de $D'$, $D''$ et du
foncteur
$$
(A',A'')\longmapsto \Hom(u(A'), v(A'')):{D'}^{\circ}\times
D''\longrightarrow (\Ens){\quoi}.
$$
Cette remarque s'applique en particulier à la situation décrite dans
la proposition suivante, et nous montre que la catégorie
$\Ppoint(E)$ est déterminée, à équivalence près, par la connaissance
des catégories $\Ppoint(\mathcal{U})$ et $\Ppoint(\mathcal{F})$ et
d'un certain foncteur
$$
\Ppoint(\mathcal{U})^{\circ}\times
\Ppoint(\mathcal{F})\longrightarrow (\Ens)
$$
déduit du foncteur de recollement $\rho:\mathcal{U}\rightarrow
\mathcal{F}$.

\begin{enonce*}{Proposition 9.7.2}\etiquette[9.7.2]{IV.prop9.7.2}
Soient $E$ un topos, $U$ un ouvert de $E$, $\mathcal{U}=E_{/U}$ et
$\mathcal{F}$ le sous-topos fermé de $E$ complémentaire de $U$.
Considérons les foncteurs $p\mapsto p\circ j$ et $q\mapsto q\circ i$
induits par les morphismes d'inclusion $j:\mathcal{U}\rightarrow E$
et $i:\mathcal{F} \rightarrow E$:
\begin{equation*}
u:\Ppoint(\mathcal{U})\longrightarrow \Ppoint(E),\quad
v:\Ppoint(\mathcal{F})\longrightarrow
\Ppoint(E){\quoi}.\tag{9.7.2.1}
\end{equation*}\etiquette[9.7.2.1]{IV.eq9.7.2.1}
Alors on a ce qui suit:
\begin{enumerate}
\renewcommand{\theenumi}{\alph{enumi}}
\renewcommand{\labelenumi}{\theenumi)}
\item Les foncteurs précédents $u$ et $v$ sont pleinement fidèles, et
tout point de $E$ appartient à l'image essentielle de l'un ou
l'autre des foncteurs $u$ et $v$ exclusivement.

\item Soient $p$ et $q$ deux points de $E$, tels que $p$ (\resp $q$)
appartienne à l'image essentielle de $u$ (\resp de $v$). Alors on a
\begin{equation*}
\Hom(q,p)=\phi{\quoi}.\tag{9.7.2.2}
\end{equation*}\etiquette[9.7.2.2]{IV.eq9.7.2.2}

\item Soient $p$ (\resp $q$) un point de $\mathcal{U}$ (\resp $\mathcal{F}$). On a
alors un isomorphisme, fonctoriel en $(q,p)$:
\begin{equation*}
\Hom(u(p),v(q))\simeq \Hom(\Pro(\rho)(p),q)\simeq
\Hom(p,\sigma(q)){\quoi},\tag{9.7.2.3}
\end{equation*}\etiquette[9.7.2.3]{IV.eq9.7.2.3}
où
$$
\sigma:\Pro(\mathcal{F})\longrightarrow \Pro(\mathcal{U})
$$
est le foncteur $\eqhyperref[IV.eq9.6.2]{9.6.2}$ pro-adjoint du foncteur de
recollement $\rho:\mathcal{U}\rightarrow \mathcal{F}$, et où
$$
\rho:\mathcal{U}\longrightarrow \mathcal{F}{\quoi},
$$
on\pageoriginale a identifié (à équivalence près) les catégories
$\Ppoint(\mathcal{U})$, $\Ppoint(\mathcal{F})$ à des sous-catégories
pleines de $\Pro(\mathcal{U})$, $\Pro(\mathcal{F})$ respectivement
(\hyperref[IV.subsubsec6.8.5]{6.8.5}).
\end{enumerate}
\end{enonce*}

\noindent \emph{Démonstration}
\begin{enumerate}
\renewcommand{\theenumi}{\alph{enumi}}
\renewcommand{\labelenumi}{\theenumi)}
\item La pleine fidélité de $u$ et $v$ est connue (\hyperref[IV.prop9.1.4]{9.1.4}). L'assertion
sur les images essentielles de $u$, $v$ résulte des critères
\hyperref[IV.subsubsec9.2.5]{9.2.5} et \hyperref[IV.prop9.2.4]{9.2.4} pour qu'un point
$P\rightarrow E$ se factorise à isomorphisme près par $\mathcal{U}$
\resp $\mathcal{F}$, compte tenu du fait que le topos ponctuel $P$
n'a que deux ouverts, savoir l'objet initial et l'objet final de
$P$.

\item Compte tenu de a), l'assertion signifie que si le point
$p:P\rightarrow E$ se factorise par $\mathcal{U}$, il en est de même
de tout point $p':P\rightarrow E$ tel que $\Hom(p',p)\neq \phi$. Or
si $\alpha :p'\rightarrow p$ il induit
$p^{\ast}(\sisi{\mathcal{U}}{U})\rightarrow
{p'}^{\ast}(\sisi{\mathcal{U}}{U})$, et
alors $p^{\ast}(\sisi{\mathcal{U}}{U})=e_p$ implique
% correction SGA4-CORR-0274 / SGA-SRC-000974: SGA4-T1-IV7619-15-555C7D1D.
${p'}^{\ast}(\sisi{\mathcal{U}}{U})=e_p$, \cqfd

\item Grâce au lemme \hyperref[IV.lem9.7.2.4]{9.7.2.4} ci-dessous, appliqué aux inclusions
$j:\mathcal{U}\rightarrow E$ et $i:\mathcal{F}\rightarrow E$, on a
un diagramme de foncteurs commutatif à des isomorphismes canoniques
près
\[
\xymatrix@C=2cm@R=.55cm@M=.18cm{ \Point(\Uu) \ar@{^{(}->}[d]
\ar[r]^{u}& \Point(E) \ar@{^{(}->}[d]&
\Point(\Gg) \ar@{^{(}->}[d] \ar[l]_{v}\\
\Pro(\Uu) \ar[r]^{\Pro(j_{\ast})}  & \Pro(E)   & \Pro(\Gg)
\ar[l]_{\Pro(i_{\ast})} , }
\]
où les flèches verticales désignent les foncteurs pleinement fidèles
canoniques (\hyperref[IV.subsubsec6.8.5]{6.8.5}). D'autre part, comme $i_{\ast}$
admet un adjoint à gauche $i^{\ast}$, $\Pro(i_{\ast})$ admet un
adjoint à gauche $\Pro(i^{\ast})$ \sisi{\hyperref[IV.rem1.8]{1.8}}{\ref{I.sec8.11.5}}, et on obtient
finalement des isomorphismes fonctoriels
\begin{align*}
\Hom(u(p),v(q)) \simeq & \Hom_{\Pro(E)}(\Pro(j_{\ast})(p),
\Pro(i_{\ast})(q))\\
& \sisi{}{\simeq}\Hom_{\Pro(\mathcal{F})}(\Pro(i^{\ast})(\Pro(j_{\ast})(p)),q)\\
& \sisi{}{\simeq}\Hom_{\Pro(\mathcal{F})}(\Pro(i^{\ast}j_{\ast})(p),q),
% correction SGA4-CORR-0275 / SGA-SRC-000975: SGA4-T1-IV7639-16-32EF81CE.
\end{align*}
ce qui n'est autre que la première formule (9.7.2.3), la deuxième
étant alors triviale par définition de $\sigma$ comme pro-adjoint de
$\rho$. Cela prouve donc \hyperref[IV.prop9.7.2]{9.7.2}, modulo le
\end{enumerate}

\begin{enonce*}{Lemme 9.7.2.4}\etiquette[9.7.2.4]{IV.lem9.7.2.4}
Soit $f:F\rightarrow E$ un morphisme de topos. Alors le diagramme de
foncteurs
\[\vcenter{
\xymatrix@C=2cm@R=.5cm{
  \Point(F)\ar[r]^{p \mapsto f \circ p} \ar[d] & \Point(E)\ar[d]\\
\Pro(F) \ar[r]^{\Pro(f_{\ast})}     & \Pro(E) }}\leqno{(9.7.2.4.1)}
\]\etiquette[9.7.2.4.1]{IV.eq9.7.2.4.1}
est\pageoriginale commutatif à isomorphisme canonique près (les
flèches verticales désignant les foncteurs pleinement fidèles
$\hyperref[IV.subsubsec6.8.5]{6.8.5}$).
\end{enonce*}

La vérification est immédiate à partir des définitions, et laissée
au lecteur.

\begin{enonce*}{Corollaire 9.7.3}\etiquette[9.7.3]{IV.cor9.7.3}
Soit $E$ un topos admettant assez de points, alors il en est de même
pour tout sous-topos fermé $\mathcal{F}$ de $E$. Plus précisément,
si $(p_i)_{i\in I}$ est une famille conservative de points de $E$,
alors la famille des points de $\mathcal{F}$, définie par la
sous-famille $(p_i)_{i\in J}$ formée des $p_i$ qui proviennent de
points de $\mathcal{F}$, est conservative. (Comparer
$\hyperref[IV.cor6.7.3]{6.7.3}$, $\hyperref[IV.rem6.7.4]{6.7.4}$.)
\end{enonce*}

En effet, soit $f:X\rightarrow Y$ une flèche de $\mathcal{F}$
transformée en isomorphisme par les foncteurs fibre associés aux
points envisagés de $\mathcal{F}$. Comme $i^{\ast}i_{\ast}\simeq
\id_{\mathcal{F}}$, on voit que $i_{\ast}(f):i_{\ast}(X)\rightarrow
i_{\ast}(Y)$ est transformé en isomorphisme par les foncteurs fibre
associés aux $p_i$ avec $i\in J$; il en est de même pour les $p_i$
avec $i\in I-J$, \ie provenant de points de $U$, puisque
$j^{\ast}i_{\ast}$ est le foncteur constant de valeur l'objet final.
Donc $i_{\ast}(f)$ est un isomorphisme, donc aussi $f\simeq
i^{\ast}i_{\ast}(f)$, \cqfd

\begin{enonce*}[remark]{Remarque 9.7.4}\etiquette[9.7.4]{IV.rem9.7.4}
Si $E$ est un topos ayant assez de points, il est clair qu'un ouvert
$U$ de $E$ est déterminé quand on connaît la sous-catégorie pleine
de $\Ppoint(E)$ image essentielle de $\Ppoint(E_{/U})$; en fait, si
$C$ est une sous-catégorie pleine de $\Ppoint(E)$ définissant une
famille conservative de points de $E$, il suffit même de connaître
la sous-catégorie $C_{/U}$ des éléments de $C$ appartenant à l'image
essentielle de $\Ppoint(U)$. De ceci et de \hyperref[IV.prop9.7.2]{9.7.2} a) il
résulte qu'un sous-topos fermé $\mathcal{F}$ de $E$ est déterminé
quand on connaît l'image essentielle de $\Ppoint(\mathcal{F})$ dans
$\Ppoint(E)$ ou même seulement l'intersection de celle-ci avec $C$.
\end{enonce*}

\begin{enonce*}[remark]{Exercice 9.7.5}\etiquette[9.7.5]{IV.exer9.7.5}
Soit $E$ un topos. Pour tout sous-topos $F$ de $E$, soit $P(F)$ le
sous-ensemble de $\Point(E)$ formé des classes d'isomorphie de
points de $E$ qui se factorisent par \sisi{$E$}{$F$}. Ainsi, $P(F)$ est
homéomorphe à $\Point(F)$ (\hyperref[IV.exer9.1.8]{9.1.8} c)). Montrer que si
$F$ est un sous-topos ouvert (\resp fermé, \resp localement fermé
(\hyperref[IV.exer9.4.9]{9.4.9})) de $E$, alors $P(F)$ est une partie
ouverte\pageoriginale (\resp fermée, \resp. localement fermée) de
\sisi{$E$}{$\Point(E)$}. Montrer que lorsque $E$ a suffisamment de points, il en est de
même de tout sous-topos localement fermé de $E$, et l'application
$F\mapsto P(F)$ induit une bijection entre l'ensemble des sous-topos
ouverts (\resp fermés, \resp localement fermés) de $E$, et
l'ensemble des parties ouvertes (\resp fermées, \resp localement
fermées) de $\Point(E)$. Généraliser les
considérations précédentes
au cas où on remplace $\Point(E)$ par n'importe quel sous-espace $X$
de $\Point(E)$ correspondant à une famille conservative de points de
$E$.
\end{enonce*}

\subsection{Compléments sur certains topos liés aux espaces topologiques}\etiquette[9.8]{IV.subsec9.8}

Nous indiquons dans ce numéro quelques compléments, qui seront
donnés sous forme d'exercices. Il est conseillé au lecteur de les
parcourir, pour s'habituer au point de vue des topos dans diverses
situations de type classique.

\begin{enonce*}[remark]{Exercice 9.8.1}\etiquette[9.8.1]{IV.exer9.8.1}
(\emph{Descente de faisceaux sur un espace topologique.})

Soit $f:X\rightarrow Y$ une application continue d'espaces
topologiques $\in 𝒰$, d'où un foncteur
$$
f^{\ast}:\Top(Y)\longrightarrow \Top(X){\quoi}.
$$
\begin{enumerate}
\renewcommand{\theenumi}{\alph{enumi}}
\renewcommand{\labelenumi}{\theenumi)}
\item Les conditions suivantes sont équivalentes:
\begin{enumerate}
\renewcommand{\theenumii}{\roman{enumii}}
\renewcommand{\labelenumii}{(\theenumii)}
\item $f^{\ast}$ est fidèle.

\item $f^{\ast}$ est conservatif.

\item $f(X)$ est une partie très dense (EGA IV 10.1.3) de $Y$.
\end{enumerate}
Il suffit pour ceci que $f$ ou $f_{\sob}$ soit surjectif. Donner un
exemple où $f^{\ast}$ est fidèle et où $f=f_{\sob}$ n'est pas
surjectif (prendre $X$ discret, $f$ injectif, $f(X)=$ ensemble des
points fermés de $Y$, $Y$ un espace topologique noethérien non
discret). Donner un exemple où $f_{\sob}$ est surjectif mais non
$f$, un autre où $f$ est surjectif mais non $f_{\sob}$.

\item Supposons $f(X)$ très dense dans $Y$, et sa topologie quotient
de celle de $X$. Prouver que la flèche $f$ de $(\Esp)$ est un
morphisme de descente pour la catégorie fibrée des faisceaux
d'ensembles sur des espaces variables, \ie que le foncteur
$f^{\ast}$ induit un foncteur \emph{pleinement fidèle}\pageoriginale
de la catégorie $\Top(Y)$ dans la catégorie des objets de $\Top(X)$
munis d'une donnée de descente relativement à $f:X\rightarrow Y$. Se
ramener \sisi{un}{au} cas $f$ surjectif, et calquer le raisonnement de VIII 9.1
% correction SGA4-CORR-0276 / SGA-SRC-000976: SGA4-T1-IV7762-17-792CEB3F.
donné dans le contexte des topologies étales de schémas.) Donner une
réciproque.

\item Supposons que $f(X)$ soit très dense dans $Y$, que sa topologie
soit quotient de celle de $X$, et que les fibres de $f$ soient
connexes. Prouver que le foncteur $f^{\ast}$ est pleinement fidèle.
(Calquer le raisonnement de \SGA 1 IX 3.4.) Donner une réciproque,
tout au moins si les points de $Y$ sont fermés.

\item Supposons que $Y$ soit réunion d'ouverts $Y_i$ au-dessus
desquels $X$ admette des sections. Prouver qu'alors $f$ est un
\emph{morphisme de descente effective} pour la catégorie fibrée des
faisceaux d'ensembles sur des espaces topologiques variables, \ie
que le foncteur envisagé dans \sisi{$c$}{$b$}) est même une équivalence de
% correction SGA4-CORR-0277 / SGA-SRC-000977: SGA4-T1-IV7776-18-58B9E6C0.
catégories.
\end{enumerate}
\end{enonce*}

\begin{enonce*}[remark]{Exercice 9.8.2}\etiquette[9.8.2]{IV.exer9.8.2}
(\emph{Topos quotients d'espaces topologiques. Étendues
topologiques.})

\begin{enumerate}
\renewcommand{\theenumi}{\alph{enumi}}
\renewcommand{\labelenumi}{\theenumi)}
\item Soit
\[
\xymatrix@C=1.5cm@R=1.5cm@M=.2cm{
X_2 \ar@<6pt>[r] \ar@<-1.5pt>[r] \ar@<-8pt>[r]& X_1 \ar@<2pt>[r]^{p_1} \ar@<-4pt>[r]_{p_2}& X_0\\
}\leqno{\sisi{(9.8.5.1)}{(9.8.2.1)}}
% correction SGA4-CORR-0278 / SGA-SRC-000978: SGA4-T1-IV7792-19-1C8DEEAC.
\]\etiquette[9.8.2.1]{IV.eq9.8.2.1}
un objet semi-simplicial tronqué à l'ordre 2 de la catégorie
$(\Esp)$ des espaces topologiques $\in 𝒰$, d'où par les
foncteurs images inverses un diagramme de catégories de faisceaux
\[
\xymatrix@M=.2cm@C=1.25cm{
  \Top(X_0) \ar@<2pt>[r]^{p_1^{\ast}} \ar@<-2pt>[r]_{p_2^{\ast}}&
  \Top(X_1) \ar@<-5pt>[r] \ar[r] \ar@<5pt>[r]& \Top(X_2),
}
\]
(de façon plus précise, on a une catégorie cofibrée sur la
catégorie des simplexes-types $\Delta_{n}$, $0\leq n\leq 2$).
Montrer que la catégorie $\varprojlim$ de ce diagramme (\ie la
catégorie des faisceaux d'ensembles sur $X_0$, \emph{munis} d'une
donnée de descente relativement au diagramme \sisi{(9.8.6.1)}{\hyperref[IV.eq9.8.2.1]{9.8.2.1}}, \ie d'un
isomorphisme $p^{\ast}_1(F)\simeq p^{\ast}_2(F)$ tel que\ldots), est
un $𝒰$-topos $T$. (Pour un énoncé plus général de stabilité
des topos par opérations $\varprojlim$, \cf
\sisi{VI)}{\hyperrefext[VI][VI.sec8]{8})} Définir un morphisme
% correction SGA4-CORR-0279 / SGA-SRC-000979: SGA4-T1-IV7810-20-7A4139DC.
de topos $q:\Top(X_0)\rightarrow T$, et prouver que pour tout
$𝒰$-topos\pageoriginale $E$, $\Hhomtop(T,E)$ est
équivalent via $q^{\ast}$ à la catégorie $\varprojlim$ du diagramme
suivant de catégories, déduit de (9.8.2.1):
\[
\xymatrix@M=.2cm{ \Hhomtop(\Top(\sisi{X_o}{X_0}),E)\ar@<2pt>[r] \ar@<-4pt>[r]&
% correction SGA4-CORR-0280 / SGA-SRC-000980: SGA4-T1-IV7816-21-290E737A.
\Hhomtop(\Top(X_1),E)\ar[r]\ar@<6pt>[r]\ar@<-6pt>[r]&
\Hhomtop(\Top(X_2),E). }
\]
\item En particulier, soit $R$ une relation d'équivalence dans un
objet $X$ de $(\Esp)$, \ie (\sisi{\hyperrefext[I][I.sec10]{10}}{\hyperrefext[I][I.sec10.9]{10.9}}) un sous-objet $R$ de
$X\times X$ (N.B. l'inclusion $R\hookrightarrow X\times X$ est
continue, mais la topologie de $R$ n'est pas nécessairement induite
par celle de $X\times X$), tel que pour tout objet $Y$ de $(\Esp)$,
$\Hom(Y,R)$ soit le graphe d'une relation d'équivalence dans
$\Hom(Y,X)$. On en déduit un diagramme de la forme (9.8.2.1), avec
$X_0=X$, $X_1=R$, $X_2=\fprod{(R,p_2)}{R,p_1}{X}$, où $p_1$, $p_2$
sont les deux projections de $R$ dans $X$, d'où un topos, qu'on
notera $\Top(X)/R$, ou même $\Top(X/R)$, voire $X/R$, par abus de
notations, et qui joue le rôle d'un quotient de $X$ par $R$, en un
sens précisé par a). Supposons que la topologie de $R$ soit induite
par celle de $X\times X$, et que, désignant par $X/R$ l'espace
topologique quotient ordinaire, $X$ admette localement des sections
sur $X/R$, prouver qu'alors $T$ est équivalent à $\Top(X/R)$.
(Utiliser \hyperref[IV.exer9.8.1]{9.8.1} d)).

\item Soit $H\rightarrow X$ un morphisme injectif de groupes
topologiques, \ie un monomorphisme d'objets groupes dans $(\Esp)$,
et soit $R$ la relation d'équivalence qu'il définit dans $X$, \ie
$R=H\times X$, avec $p_1=\sisi{\pr_1}{\pr_2}$ et $p_2$ défini par l'action de $H$
% correction SGA4-CORR-0281 / SGA-SRC-000981: SGA4-T1-IV7839-22-3A49E2E2.
sur $X$ via translations à gauche. Montrer que pour que la topologie
de $R$ soit induite par celle de $X\times X$, il faut et il suffit
que la topologie de $H$ soit induite par celle de $X$. Donner des
exemples où cette condition n'est pas remplie, et où la topologie
quotient ordinaire de $X/H$ est la topologie grossière avec a)
$H=\mathbb{Z}\times \mathbb{Z}$, $X=\mathbb{R}$, et b) $H=\mathbb{R}$,
$X=\mathbb{T}\times \mathbb{T}$, avec
$\mathbb{T}=\mathbb{R}/\mathbb{Z}$, (« géodésiques du tore »). Prouver que les deux topos obtenus sont équivalents et ne sont
pas équivalents à $\Top(X/H)$ (qui est un topos final
(\hyperref[IV.subsec2.2]{2.2}), ni à aucun topos de la forme $\Top(Y)$).

\item Soit $G$ un groupe topologique opérant sur un espace topologique
$X$, d'où de façon bien connue un objet semi-simplicial
\[
\xymatrix@M=.2cm{ ... G\times G \times X
\ar@<6pt>[r]\ar@<-1.5pt>[r]\ar@<-8pt>[r] &
  G\times X \ar@<3pt>[r] \ar@<-4pt>[r] & X.
}
\]
Montrer\pageoriginale que le topos $T$ qui s'en déduit en vertu de
a) s'identifie à la catégorie des espaces $X'$ à groupe
d'opérateurs
$G$ au-dessus de $X$, tels que $X'\rightarrow X$ soit un étalement
(compatible à l'action de $G$). En particulier, lorsque $G$ est
discret, on retrouve le topos $\Top(X,G)$ de \hyperref[IV.subsec2.3]{2.3}.

Les considérations qui précèdent justifient la notation $\Top(X)/G$,
voire $\Top(X/G)$ ou même simplement $X/G$, pour le topos précédent
$T$. On fera attention cependant que lorsque $G$ est un groupe
discret (pour fixer les idées) opérant \emph{proprement} sur $X$, de
sorte que l'espace quotient ordinaire $X/G$ possède des propriétés
assez raisonnables, le morphisme de topos naturel $T\rightarrow
\Top(X/G)$ déduit de la caractérisation universelle a) de $T$ n'est
une équivalence de topos que si $G$ opère librement \ie sans points
fixes, \ie lorsque $G\times X\rightarrow X\times X$ est un
monomorphisme; donc dans le cas d'une « pré-relation
d'équivalence » (ou « groupoïde » au sens de (\SGA 3 V 1) qui n'est pas une
relation d'équivalence, la notion de passage au quotient au sens des
topos (ou « passage au quotient fin ») ne correspond pas en
général (via la correspondance $X\rightarrow \Top(Y)$) au passage au
quotient topologique habituel.

\item Soit $T$ un $𝒰$-topos. Montrer que les conditions suivantes
sont équivalentes:
\begin{itemize}
\item[(i)] Il existe une famille $(S_i)_{i\in I}$ d'objets de $T$
couvrant l'objet final de $\sisi{E}{T}$, telle que pour tout $i\in I$, le
topos induit $\sisi{T_{/S_1}}{T_{/S_i}}$ soit équivalent à un topos de la forme
$\Top(X)$, avec $X\in(\Esp)$.

\item[${\rm (i')}$] Il existe un $S\in \Ob \sisi{E}{T}$ couvrant l'objet final
% correction SGA4-CORR-0282 / SGA-SRC-000982: SGA4-T1-IV7886-23-2E418EDB.
tel que le topos induit $T_{/S}$ soit équivalent à un topos de la
forme $\Top(X)$.

\item[(ii)] Il existe un espace topologique $X$, et une pré-relation
d'équivalence $R$ dans $X$ (au sens de la catégorie $(\Esp)$) qui
soit \emph{étale} (\ie telle que $p_1:R\rightarrow X$ soit un
étalement), tels que $T$ soit équivalent à $\Top(X/R)$, où les
notations sont celles de b).
\end{itemize}

On\pageoriginale dira alors que le topos $T$ est \emph{localement un
espace topologique,} ou encore que $T$ est une \emph{étendue
topologique} (ou simplement une \emph{étendue,} si aucune confusion
n'est à craindre).

\item Montrer que le topos $T$ construit dans a) a suffisamment de
points, et plus précisément, qu'il admet une famille conservative de
points paramétrée par le (petit) ensemble $X_0$. En particulier,
toute étendue \sisi{à}{a} une petite famille conservative de points, donc a
% correction SGA4-CORR-0283 / SGA-SRC-000983: SGA4-T1-IV7910-24-D1B5455A.
suffisamment de points. Lorsqu'une étendue $T$ est réalisée sous la
forme $\Top(X/R)$ comme dans e) ii), prouver que tout point de $T$
est isomorphe à l'image d'un point de $\Top(X)$, donc est défini par
un point ordinaire de $X_{\sob}$ (ou encore de $X$, lorsque $X$ est
sobre).

\item Montrer que le topos $\Top(X/G)$ de \hyperref[IV.subsec2.3]{2.3} est une
étendue. (Utiliser sa description d) ou \hyperref[IV.exer7.1.10]{7.1.10} c).)
Montrer que pour tout $x\in X$, le monoïde des endomorphismes du
point de $\Top(X,G)=\Top(X)/G$ défini par $x$ est canoniquement
isomorphe au groupe de stabilité $G_x$ de $x$. En particulier, les
points d'une étendue peuvent avoir des groupes d'automorphismes non
triviaux. Déterminer $\Ppoint(\sisi{J}{T})$ pour l'étendue $T$ définie par une
% correction SGA4-CORR-0284 / SGA-SRC-000984: SGA4-T1-IV7923-25-D744DC52.
prérelation d'équivalence étale dans un espace $X$, et montrer que
tout endomorphisme d'un point de $T$ est un automorphisme.

\item Soit $T$ un topos. Montrer que les conditions suivantes sont
équivalentes:
\begin{itemize}
\item[i)] Pour tout point $p$ de $T$, tout endomorphisme de $p$
(\resp tout automorphisme de $p$) est l'identité.

\item[ii)] Pour tout topos $S$ ayant suffisamment de \sisi{point}{points}, et tout
% correction SGA4-CORR-0285 / SGA-SRC-000985: SGA4-T1-IV7933-26-7FCE3F30.
morphisme de topos $f:S\rightarrow T$, tout endomorphisme de $f$
(\resp tout automorphisme de $f$) est l'identité.
\end{itemize}
Montrer de même que les conditions suivantes sont équivalentes:
\begin{itemize}
\item[${\rm i')}$] Pour deux points $p$, $q$ de $T$, il existe au
plus un morphisme de $p$ dans $q$.

\item[${\rm ii')}$] Pour deux morphismes $f$, $g:S\rightrightarrows
T$ d'un topos $S$ ayant suffisamment de\pageoriginale points dans
$T$, il existe au plus un morphisme de $f$ dans $g$.
\end{itemize}

Lorsque ces dernières conditions sont vérifiées, on dit que le topos
$T$ est \emph{ponctuellement rigide,} ou simplement \emph{rigide}.
Montrer que si $X$ est un espace topologique muni d'une pré-relation
d'équivalence étale $R$, l'étendue quotient $T=\Top(X)/R$ est rigide
si et seulement si $R$ est une relation d'équivalence.

\item Soient $(X,G)$ un espace topologique à groupe discret d'opérateurs,
$H$ un sous-groupe distingué de $G$ tel que l'opération induite de
$H$ sur $X$ soit propre et libre, et soient $X'=X/H$, $G'=G/H$, de
sorte que $G'$ opère sur $X'$ de façon évidente. Prouver que le
morphisme de topos
$$
\Top(X,G)\longrightarrow \Top(X',G')
$$
induit par le morphisme canonique $(X,G)\rightarrow (X',G')$
d'espaces à opérateurs (\hyperref[IV.subsubsec4.1.2]{4.1.2}) est une équivalence
de topos.

Montrer que l'ensemble des ouverts du topos $\Top(X/G)$ s'identifie
à l'ensemble des ouverts de $X$ stables par l'action de $G$, ou
encore à l'ensemble des ouverts de l'espace topologique quotient
$X/G$. En conclure que $\Top(X/G)$ est connexe si et seulement si
toute partie de $X$ à la fois ouverte et fermée et stable sous
l'action de $G$ est vide ou égale à $X$. Lorsque les composantes
connexes de $X$ sont ouvertes, $\Top(X,G)$ est connexe si et
seulement si $G$ opère transitivement sur l'ensemble $\pi_0(X)$ des
composantes connexes de $X$; plus généralement $\pi_0(\Top(X/G))$
(\hyperref[IV.exer8.7]{8.7} g)) est un pro-ensemble essentiellement constant
isomorphe à l'ensemble quotient $\pi_0(X)/G$.

Montrer que lorsque $X$ est localement connexe et localement
simplement connexe, et $T=\Top(X,G)$ connexe \ie $G$ transitif sur
$\pi_0(X)$, alors parmi toutes les façons de réaliser $T$ à
l'aide d'un espace à opérateurs $(X',G')$ (\cf ci-dessus pour
certaines telles façons), il y en a une, unique à isomorphisme
non unique près, pour laquelle $X'$ est connexe et simplement
connexe. Montrer que le choix d'un tel $(X',G')$ revient au choix
d'un « revêtement universel »  de l'objet final de $T$, et que
$G'$ est isomorphe au groupe $\pi_1(T)$ relatif à ce revêtement
universel (\cf \hyperref[IV.exer2.7.5]{2.7.5}).

\noindent (Hint:\pageoriginale montrer que la donnée d'une
équivalence de $T$ avec un $\Top(X,G)$ revient à la donnée d'un
$G_T$-torseur $P$ dans $T$, et que si $(X,G)$ et $P$ se
correspondent, on a une équivalence canonique
$\sisi{\Top(X)-T_{/p}}{\Top(X)\simeq T_{/P}}$).
% correction SGA4-CORR-0286 / SGA-SRC-000986: SGA4-T1-IV7988-27-9AABB68D.

Conclure de ceci une démonstration triviale de la dernière assertion
de c). Caractériser les topos $T$ équivalents à des $\Top(X,G)$, où
$G$ est un groupe discret opérant sur un espace topologique
localement connexe et localement simplement connexe en permutant
transitivement les composantes connexes, comme étant les étendues
connexes, localement connexes et localement simplement connexes dont
le revêtement universel $P$ de l'objet final $e_T$ de $T$ soit tel
que le topos induit $T_{/P}$ soit équivalent à un topos de la forme
$\Top(X)$ (ou, comme on dira simplement par abus de langage $T_{/P}$
« est »  un espace topologique).

Donner un exemple d'une étendue, localement isomorphe à $\mathbb{R}$,
qui est connexe et simplement connexe mais qui n'est pas un espace
topologique. (Prendre le revêtement universel de la droite
$\mathbb{R}$ avec origine dédoublée.)

\item Un topos annelé $(T,{\Oo}_T)$ (\cf \S 11) est appelé une
\emph{étendue différentiable} (\resp une \emph{étendue analytique
réelle.} \resp une \emph{étendue analytique complexe}) s'il existe
des objets $S_i$ de $T$ couvrant l'objet final, tels que le topos
annelé induit $(T_{/\sisi{E}{S_i}}, {\Oo}_T/\sisi{E}{S_i})$ soit équivalent au topos
% correction SGA4-CORR-0287 / SGA-SRC-000987: SGA4-T1-IV8012-28-C3CE1667.
annelé défini par une variété différentiable avec son faisceau de
fonctions réelles $C^{\infty}$ (\resp au topos annelé défini par un
espace analytique réel, \resp au topos annelé défini par un espace
analytique complexe). Donner une description constructive de ces
topos annelés, du type de e) ii) (où on prendra pour $X$
respectivement une variété différentiable, un espace analytique
\sisi{complexe}{réel} ou un espace analytique \sisi{réel}{complexe}). Donner des exemples de tels
% correction SGA4-CORR-0288 / SGA-SRC-000988: SGA4-T1-IV8018-29-8786D002.
topos annelés qui ne soient pas équivalents à des topos annelés
associés à des espaces topologiques, en reprenant les exemples
envisagés dans c).
\end{enumerate}
\end{enonce*}

\begin{enonce*}[remark]{Exercice 9.8.3}\etiquette[9.8.3]{IV.exer9.8.3}
(\emph{Topos\pageoriginale associés aux relations d'équivalence
locales et aux feuilletages}) \footnote[1]{Cet exercice est donné
sous toutes réserves, ayant été rédigé hâtivement et insuffisamment
vérifié. Monsieur N. \textsc{Saavedra} a vérifié les parties a) à
e). Prière au lecteur de nous communiquer ses observations
éventuelles.}.

Soit $X$ un espace topologique.
\begin{enumerate}
\renewcommand{\theenumi}{\alph{enumi}}
\renewcommand{\labelenumi}{\theenumi)}
\item Pour tout ouvert $U$ de $X$, soit $\Quot(U)$ l'ensemble des
relations d'équivalence dans $U$, et pour un ouvert $V\subset U$,
considérons l'application naturelle $\Quot(U)\rightarrow \Quot(V)$,
d'où un préfaisceau $\Quot_X$ sur $X$. Soit $Q_X$ ou simplement $Q$
le faisceau associé, et soit $r$ une section de $Q$ (« \emph{relation d'équivalence locale sur $X$} »). De la relation
d'ordre sur l'ensemble $\Quot(U)$ des relations d'équivalence sur un
ouvert $U$ de $X$, déduire sur le préfaisceau $\Quot$, et sur le
faisceau associé $Q$, une structure d'ordre \ie une structure de
préfaisceau \resp de faisceau à valeurs dans la catégorie des
ensembles ordonnés.

\item Considérer, pour toute relation d'équivalence $R$ sur $X$, la
section $\loc(R)$ de $Q$ qu'elle définit. Soit $E(r)$ l'ensemble des
relations d'équivalence sur $X$ telles que $\loc(R)$ soit moins fine
que $r$, et soit $\glob(r)$ la relation d'équivalence borne
supérieure de $E(r)$, (dont le graphe est l'intersection des graphes
des $R\in E(r)$). Soit $(U_x)_{x\in X}$ une famille de voisinage
ouverts des $x\in X$, et pour tout $x\in X$ soit $R_x$ une relation
d'équivalence dans $U_x$ dont le germe en $x$ soit $r_x$. Pour toute
famille $𝒱$ d'ouverts $V_x$ $(x\in X, x\in V_x\subset U_x)$,
soit $R_{𝒱}$ la relation d'équivalence dans $X$ engendrée
par la famille des relations $R_x| V_x$. Montrer que si
$𝒱'\leq 𝒱$ (dans un sens évident) on a
$R_{𝒱'}\geq R_{𝒱}$, et que $\glob(r)$ est la borne
supérieure de la famille filtrante croissante de relations
d'équivalence $R_{𝒱}$. Donner un exemple où $\glob(r)$ n'est
pas dans $E(r)$, \ie où il existe un $a\in X$ tel que,
pour\pageoriginale tout voisinage ouvert $W\subset U_a$ de $a$, il
existe un $𝒱=(V_x)$ et deux points $y$, $z$ de $W$ qui sont
équivalents pour $R_a$, mais qui ne sont pas équivalents pour
$R_{𝒱}$. Montrer que pour toute $R\in \Quot(X)$ on a
$\glob\loc(R)\geq R$.

\item On dit que $r$ est une \emph{relation d'équivalence locale
précohérente (\resp cohérente)} si $\glob(r)\in E(r)$ \ie
$\loc(\glob(r))\leq r$ (\resp si pour tout ouvert $U$ de $X$, la
restriction de $r$ à $U$ est précohérente). On dit que $r$ est
globalement cohérente si elle est cohérente, et si, de plus, $r=\loc
\glob(r)$. On dit qu'une relation d'équivalence $R$ sur $X$ est
\emph{localement cohérente} si $r=\loc(R)$ est cohérente,
\emph{cohérente} si de plus $R=\glob(r)$ \ie $R=\glob(\loc(R))$.
Montrer que les relations d'équivalence \sisi{\emph{globalement}}{\emph{localement}}
cohérentes \sisi{$R$}{$r$} sur $X$ correspondent biunivoquement aux relations
d'équivalence cohérentes \sisi{}{$R$} sur $X$, par $r\mapsto \glob(r)$ et
$R\mapsto \loc(R)$;

\item Montrer que pour que $r$ soit cohérente il suffit que pour tout
$a\in X$ et tout voisinage ouvert $W'\subset U_a$ de $a$, il existe
un voisinage ouvert $W\subset W'$ de $a$, tel que toute classe
d'équivalence de $R_a| W$ soit contenue dans une composante connexe
d'une classe d'équivalence de $R_a| W'$ (à fortiori, il suffit qu'il
existe une famille fondamentale de voisinages ouverts $W_a\subset
U_a$ de $a$ tels que les fibres de la \sisi{relations}{relation} d'équivalence
% correction SGA4-CORR-0289 / SGA-SRC-000989: SGA4-T1-IV8091-30-724EFD92.
induite $R_a| W_a$ soient connexes). (Hint: se ramener à établir la
pré-cohérence dans le cas où $r$ est définie par une $R\in
\Quot(X)$, et noter dans ce cas que les fibres des relations
d'équivalence $R_{𝒱}$ sont des parties relativement ouvertes
- donc aussi relativement fermées- des fibres de $R$). Montrer que
pour qu'une relation d'équivalence globale $R$ soit cohérente, il
suffit qu'elle soit localement cohérente et que ses fibres soient
connexes; prouver que cette condition est nécessaire si les fibres
sont fermées. Donner un exemple d'une relation d'équivalence à
fibres connexes et localement connexes, et qui n'est pas localement
cohérente (?).

\item Soit\pageoriginale $f:X'\rightarrow X$ une application
continue. Définir un homomorphisme de préfaisceaux naturel
$f_{\ast}(\Quot_{X'})\leftarrow \Quot_{X}$, induisant un
homomorphisme de faisceaux $f_{\ast}(Q_{X'})\leftarrow Q_X$ d'où
$f^{\ast}(Q_X)\rightarrow Q_{X'}$. Si $r$ est une relation
d'équivalence locale sur $X$, on dit que $f$ est une « \emph{application fibre} »  pour la relation d'équivalence locale
$r$ sur $X$, si l'image inverse de $r$ est la relation d'équivalence
locale grossière sur $X'$. Soit, pour tout espace topologique $X'$,
$\Homfib_r(X',X)$ l'ensemble des applications continues de $X'$ dans
$X$ qui sont des applications fibres relativement à $r$. Montrer que
pour $X'$ variable, on obtient un contrafoncteur en $X'$.

\item Supposons $r$ cohérente. Montrer que le foncteur précédent est
représentable par un espace topologique $X^r$ au-dessus de $X$.
Montrer que l'application fibre universelle $\Phi:X^r\rightarrow X$
est bijective, et que tout $x\in X^r$ admet un voisinage ouvert $U$
tel que l'application $U\rightarrow X$ induite par $\Phi$ soit un
homéomorphisme de $U$ sur son image. Montrer que les composantes
connexes \sisi{$X^r$}{de $X^r$} sont ouvertes et correspondent par la bijection
% correction SGA4-CORR-0290 / SGA-SRC-000990: SGA4-T1-IV8122-31-A3882CB6.
$\Phi$ aux fibres de la relation d'équivalence $R=\glob(r)$. Donner
un exemple où la restriction de $\Phi$ aux composantes connexes de
$X^r$ n'induit pas des homéomorphismes de ces espaces avec leurs
images dans $X$.

\item La relation d'équivalence locale $r$ est dite \emph{ouverte} si
elle est une section du sous-faisceau $\Qouv_X$ de $Q$ provenant du
sous-préfaisceau $\Quotouv_X$ de $\Quot_X$ dont la valeur en tout
ouvert $U$ de $X$ est l'ensemble des relations d'équivalence
\emph{ouvertes} de $U$. On dira que la relation d'équivalence locale
$r$ est \emph{strictement ouverte} si elle est cohérente et ouverte
ou ce qui revient au même, si elle est cohérente et si pour tout
ouvert $U$ de $X$, $\glob(r| U)$ est une relation d'équivalence
ouverte dans $U$. On dit que la relation d'équivalence $R$ dans $X$
est \emph{strictement ouverte} si elle est cohérente et si $\loc(R)$
est une relation d'équivalence locale strictement ouverte. Prouver
que pour que $r$ supposée ouverte soit strictement ouverte, il
suffit qu'elle satisfasse à\pageoriginale la condition envisagée
dans d); si $r$ est localement définie par une $R$ à fibres fermées,
alors cette condition est aussi nécessaire.

\item Soit $F$ un faisceau d'ensembles sur $X$. Pour tout ouvert $U$
de $X$, soit $\Quot(U,F)$ l'ensemble des couples formés d'une
relation d'équivalence $R_U$ dans $U$ et d'une relation
d'équivalence $R_{F| U}$ dans $F| U$ ($F$ étant interprété comme
espace étalé sur $X$) tels que le morphisme structural $p:F|
U\rightarrow U$ soit compatible avec les relations d'équivalence
$R_U$, $R_{F| U}$, et que le diagramme correspondant d'espaces
topologiques
\[
\xymatrix@C=2cm@R=1cm{
F|U \ar[r]\ar[d] \ar[d]_{p}& (F|U)/R_{F|U} \ar[d]_{q}\\
U   \ar[r]&  U/R_{U} }
\]
soit cartésien avec $q$ un étalement (de sorte que $F| U$
s'identifie à l'image inverse du faisceau $(F| U)/R_{F| U}$ sur
$U/R_U$). Montrer que les $\Quot(U,F)$ pour $U$ variable définissent
un préfaisceau sur $X$, dont le faisceau associé sera noté $Q(F/X)$.
Définir un homomorphisme de faisceaux $Q(F/X)\rightarrow Q_X$. Pour
une section donnée $r$ de $Q_X$, on appelle $r$-\emph{structure} sur
le faisceau $F$ toute section de $Q(F/X)$ au-dessus de la section
donnée $r$ de $Q_X$. Définir la catégorie $\Top(X/r)$ des
$r$-faisceaux sur $X$, et définir un foncteur conservatif et fidèle:
% correction SGA4-CORR-0291 / SGA-SRC-000991: SGA4-T1-IV8165-32-0A76E8E7.
« oubli de la $r$-structure » $\Top(X/r)\rightarrow \Top(X)$.
Prouver que dans $\Top(X/r)$ les limites projectives finies et les
limites inductives finies sont représentables, et que le foncteur
précédent commute aux dites limites. En conclure que dans
$\Top(X/r)$ les sommes finies sont disjointes et universelles et les
relations d'équivalence sont effectives \sisi{}{et }universelles. Prouver que
% correction SGA4-CORR-0292 / SGA-SRC-000992: SGA4-T1-IV8170-33-F3FDAA7D.
$\Top(X/r)$ admet une petite famille génératrice. Donner un exemple
avec $r$ cohérente, où $\Top(X/r)$ n'admet pas des sommes directes
infinies (?), donc n'est pas un topos.

\item Soit $f:X\rightarrow Y$ une application continue compatible avec
$R=\glob(r)$. Définir un foncteur
$$
f^{\ast}_r:\Top(Y)\longrightarrow \Top(X/r)
$$
dont\pageoriginale le composé avec le foncteur d'inclusion
$\Top(X/r)\rightarrow \Top(X)$ soit $f^{\ast}$. Montrer que si $r$
est globalement cohérente et strictement ouverte (\ie si $R$ est
strictement ouverte et $r=\loc(R)$), et si $f$ est l'application
canonique $X\rightarrow Y=X/R$, alors $f^{\ast}_r$ est une
équivalence de catégories, donc $\Top(X/r)$ est un
$𝒰$-topos, et $\Top(X/r)\rightarrow \Top(X)$ est le foncteur
image inverse d'un morphisme de topos.

\item Conclure de i) que si $r$ est strictement ouverte, alors
$\Top(X/r)$ est un topos, et le foncteur d'inclusion
$\Top(X/r)\rightarrow \Top(X)$ est le foncteur image inverse associé
à un morphisme de topos
$$
p:\Top(X)\longrightarrow \Top(X/r){\quoi}.
$$

\item Définir une loi fonctorielle du topos $\Top(X/r)$ et du
morphisme de topos précédent $p$, par rapport au couple $(X,r)$ d'un
espace topologique $X$ muni d'une relation d'équivalence locale
strictement ouverte. Montrer que si $X'\rightarrow X$ est un
étalement et si $r'$ est l'image inverse de $r$ dans $X'$, alors le
morphisme induit $\Top(X'/r')\rightarrow \Top(X/r)$ est équivalent à
un morphisme de localisation $\Top(X/r)_{/E}\rightarrow \Top(X/r)$,
où $E$ est un objet de $\Top(X/r)$ (déterminé à isomorphisme unique
près). Pour que $E$ couvre l'objet final de $\Top(X/r)$, il faut et
il suffit que le saturé sous $R=\glob(r)$ de l'image de $X'$ dans
$X$ soit égal à $X$.

\item Déduire de k) que $\Top(X/r)$ est une \emph{étendue rigide}
(\hyperref[IV.exer9.8.2]{9.8.2} b)). Montrer que si $X$ est sobre l'ensemble des
classes
     d'isomorphie de points de $\Top(X/r)$ est homéomorphe, pour sa
     topologie canonique (\hyperref[IV.exer7.8]{7.8}) à l'espace topologique quotient
     $X/\glob(r)$, et que pour deux points de $\Top(X/r)$, provenant
     de points $x$ et $x'$ de $X$, il existe au plus un morphisme de
     l'un dans l'autre; il y en a effectivement un si et seulement si
     $x'$ et $x$ sont équivalents mod $\glob(r)$ à des points $x'_1$
     et $x_1$ tels que $x_1$ générise $x'_1$.

\item Étudier le morphisme de topos canonique $\Top(X)\rightarrow
     \Top(X/r)$, en notant\pageoriginale que pour tout objet $U$ de
     $\Top(X/r)$ tel que le topos induit sur $U$ « soit »  un
     espace topologique ordinaire, le morphisme induit
     $\sisi{\Top(X)_{/p}\ast_{(U)}\rightarrow \Top(X/r)/U}
     {\Top(X)_{/p^*(U)}\rightarrow \Top(X/r)_{/U}}$ « est »  une
     % correction SGA4-CORR-0293 / SGA-SRC-000993: SGA4-T1-IV8221-34-B09B7EC8.
     application continue d'espaces topologiques ordinaires
     $X_U\rightarrow U$. Montrer que les fibres de l'application
     $X_U\rightarrow U$ sont homéomorphes à des composantes connexes
     de l'espace $X^r$ de $f\sisi{}{)}$.
     % correction SGA4-CORR-0294 / SGA-SRC-000994: SGA4-T1-IV8229-35-43C85C16.

\item Soit $X$ une variété différentiable munie d'un feuilletage, \ie
     d'un sous-faisceau $F$ localement facteur direct du faisceau
     tangent de $X$, stable par crochets. Définir sur $X$ une relation
     d'équivalence locale strictement ouverte associée au feuilletage,
     et définir sur le topos quotient $T=\Top(X/r)$ un \sisi{Anneau}{anneau} qui en
     % correction SGA4-CORR-0295 / SGA-SRC-000995: SGA4-T1-IV8235-36-DC3DE9FA.
     fasse une \emph{étendue différentiable} (\hyperref[IV.exer9.8.2]{9.8.2} j)). Définir une
     équivalence \sisi{de}{entre} la catégorie des Modules localement libres sur
     l'étendue différentiable $T$ (appelés aussi fibrés vectoriels
     différentiables sur $T$), et la catégorie des Modules localement
     % correction SGA4-CORR-0296 / SGA-SRC-000996: SGA4-T1-IV8237-37-5465DBD0.
     libres sur $X$ munis d'une connexion relativement au sous-fibré
     $F$ donné du fibré tangent. Donner des variantes dans le cas
     analytique réel et analytique complexe.
\end{enumerate}
\end{enonce*}

\section{Faisceaux de morphismes}\etiquette[10]{IV.sec10}\pageoriginale

\begin{enonce*}{Proposition 10.1}\etiquette[10.1]{IV.prop10.1}
Soient $E$ un topos, $X$ et $Y$ deux objets de $E$; le foncteur
$Z\mapsto \Hom_E(Z\times X,Y)\simeq \Hom_{E_{/Z}}(X_Z, Y_Z)$ est
représentable.
\end{enonce*}

En effet, les limites inductives sont universelles dans $E$ (II
4.3). Le foncteur $Z\mapsto Z\times X$ commute donc aux limites
inductives. Par suite, le foncteur $Z\mapsto \Hom_E(Z\times X,Y)$
transforme les limites inductives de l'argument $Z$ en limites
projectives. Il est donc représentable \sisi{(IV \hyperref[IV.prop1.4]{1.4} et
\hyperref[IV.theo1.2]{1.2})}{(\hyperref[IV.cor1.2.1]{1.2.1})}.

\setcounter{subsection}{1}
\subsection{}\etiquette[10.2]{IV.subsec10.2}
L'objet représentant le foncteur $Z\mapsto \Hom_{E}(Z\times X, Y)$
est noté $\scrHom_{E}(X,Y)$ (ou plus simplement $\scrHom(X,Y)$) et
est appelé le \emph{faisceau des morphismes de $X$ dans $Y$}. C'est
un bifoncteur en $X$ et $Y$. On a donc un isomorphisme
trifonctoriel.
\begin{equation*}
\Hom_E(Z,\scrHom_E(X,Y))\simeq \sisi{\scrHom_E}{\Hom_E}(Z\times X,
Y){\quoi}.\tag{10.2.1}
% correction SGA4-CORR-0297 / SGA-SRC-000997: SGA4-T1-IV8268-38-9201C125.
\end{equation*}\etiquette[10.2.1]{IV.eq10.2.1}

Il résulte alors de la formule \eqhyperref[IV.eq10.2.1]{10.2.1} que le bifoncteur
$(X,Y)\mapsto \scrHom(X,Y)$ transforme les limites inductives de
l'argument $X$ (\resp les limites projectives de l'argument $Y$) en
limites projectives dans $E$.

\begin{enonce*}{Proposition 10.3}\etiquette[10.3]{IV.prop10.3}
Soit $v:E\rightarrow E'$ un morphisme de topos. Pour un objet
variable $X$ de $E$ et un objet variable $Y$ de $E'$, on a un
isomorphisme bifonctoriel.
\begin{equation*}
v_{\ast}\scrHom_E(v^{\ast}Y,X)\simeq
\scrHom_{E'}(Y,v_{\ast}X){\quoi}.\tag{10.3.1}
\end{equation*}\etiquette[10.3.1]{IV.eq10.3.1}
\end{enonce*}

\noindent \emph{Démonstration}. Pour tout objet $Z$ de $E'$, on a
une suite d'isomorphismes, fonctoriels en tous les arguments:
% correction SGA4-CORR-0298 / SGA-SRC-000998: SGA4-T1-IV8289-39-2F01DCA0.
\begin{alignat*}{4}
&\Hom_{E'}(Z,v_{\ast}\scrHom_E(v^{\ast}Y,X)) &&\simeq
&& \Hom_E(v^{\ast}Z,\scrHom_E(v^{\ast}Y,X)) &&\quad \text{(adjonction)}\\
& \Hom_E(v^{\ast}Z,\scrHom_E(v^{\ast}Y,X)) && \simeq
&& \Hom_E(v^{\ast}Z\times v^{\ast}Y,X) &&\quad \eqhyperref[IV.eq10.2.1]{10.2.1}\\
&\Hom_E(v^{\ast}Z\times v^{\ast}Y,X) &&\simeq &&
\Hom_E(v^{\ast}(Z\times Y),X) &&\quad \text{($v^{\ast}$ exact à
gauche)}\\
&\Hom_E(v^{\ast}(Z\times Y),X) &&\simeq && \Hom_{E'}(Z\times Y,
v_{\ast}X) &&\quad \text{(adjonction)}\\
& \Hom_{E'}(Z\times Y, v_{\ast}X) &&\simeq &&
\Hom_{E'}(Z,\scrHom_{E'}(Y,v_{\ast}X)) &&\quad \eqhyperref[IV.eq10.2.1]{10.2.1}
\end{alignat*}
Les deux membres de \eqhyperref[IV.eq10.3.1]{10.3.1} représentent donc des foncteurs
isomorphes, \cqfd

\begin{enonce*}{Corollaire 10.4}\etiquette[10.4]{IV.cor10.4}
Soient\pageoriginale $C$ un $𝒰$-site, $X$ un
$𝒰$-préfaisceau sur $C$, $Y$ un $𝒰$-faisceau sur
$C$, $𝒱$ un univers contenant $𝒰$ tel que $C$ soit
$𝒱$-petit, $C^{\widehat{\phantom{,}}}_{𝒱}$ le topos
des $𝒱$-préfaisceaux sur $C$, $C^{\sim}_{𝒰}$ le
topos des $𝒰$-faisceaux sur $C$. Le préfaisceau
$\scrHom_{C^{\widehat{\phantom{,}}}_{𝒱}}(X,Y)$ est un
$𝒰$-faisceau. Soit $X\rightarrow \underline{a}X$ le morphisme
canonique de $X$ dans son faisceau associé. Le morphisme
correspondant
$\scrHom_{C^{\widehat{\phantom{,}}}_{𝒱}}(\underline{a}X,Y)\rightarrow
\scrHom_{C^{\widehat{\phantom{,}}}_{𝒱}}(X,Y)$ est un
isomorphisme. On a un isomorphisme canonique
$\scrHom_{C^{\widehat{\phantom{,}}}_{𝒱}}(\underline{a}X,Y)\simeq
\scrHom_{C^{\sim}_{𝒰}}(\underline{a}X,Y)$.
\end{enonce*}

\setcounter{subsection}{4} \setcounter{subsubsection}{0}
\subsubsection{}\etiquette[10.4.1]{IV.subsubsec10.4.1}
Supposons d'abord que $C$ soit un petit site et que
$𝒱=𝒰$. On a alors un morphisme de topos:
$C^{\sim}_{𝒰}\rightarrow
C^{\widehat{\phantom{,}}}_{𝒰}$ (\sisi{IV }{}\hyperref[IV.subsec4.8]{4.8}), d'où
\eqhyperref[IV.eq10.3.1]{10.3.1} les assertions dans ce cas. Pour passer de là au cas
général, on remarque tout d'abord que le foncteur « faisceau
associé »  ne dépend pas de l'univers (\hyperrefext[II][II.prop3.6]{3.6}) et que le
topos des $𝒱$-faisceaux sur $C$ est équivalent au topos des
$𝒱$-faisceaux sur $C^{\sim}_{𝒰}$ (\hyperrefext[III][III.sec4]{4} et
\hyperref[IV.sec1]{1}). Il suffit donc de démontrer le lemme suivant:

\begin{enonce*}{Lemme 10.4.2}\etiquette[10.4.2]{IV.lem10.4.2}
Soient $E$ un $𝒰$-topos, $𝒱$ un univers contenant
$𝒰$, $E^{\sim}_{𝒱}$ le topos des
$𝒱$-faisceaux sur $E$, $X$ et $Y$ deux objets de $E$. Il
existe un isomorphisme canonique $\scrHom_E(X,Y)\simeq
\scrHom_{E^{\sim}_{𝒱}}(X,Y)$.
% correction SGA4-CORR-0299 / SGA-SRC-000999: SGA4-T1-IV8341-40-4F812E70.
\end{enonce*}

\setcounter{subsubsection}{2}
\subsubsection{}\etiquette[10.4.3]{IV.subsubsec10.4.3}
Le foncteur $\epsilon_{E}:E\rightarrow E_{\widetilde{𝒱}}$
est pleinement fidèle et exact à gauche. Par suite on a, pour tout
objet $Z$ de $E$ un isomorphisme
$$
\Hom_{E_{\widetilde{𝒱}}}(\epsilon_E Z, \epsilon_E
\scrHom_E(X,Y))\simeq \Hom_{E_{\widetilde{𝒱}}}(\epsilon_E
Z\times \epsilon_E X, \epsilon_E Y)
$$
Mais tout objet de $E_{\widetilde{𝒱}}$ est limite inductive
d'objets provenant de $E$, d'où \eqhyperref[IV.eq10.2.1]{10.2.1} l'assertion.

\subsection{}\etiquette[10.5]{IV.subsec10.5}
Soit $v:E\rightarrow E'$ un morphisme de topos. On se propose de
définir quatre morphismes bifonctoriels.
\begin{alignat*}{4}
\Phi_v & : v^{\ast}\scrHom(X,Y) &&\longrightarrow
\scrHom(v^{\ast}X,v^{\ast}Y) &\quav
& X \text{ et } Y \text{ objets de } E'{\quoi},\\
\Psi_v & : v_{\ast}\scrHom(X,Y) &&\longrightarrow \scrHom(v_{\ast}X,
v_{\ast}Y) &\quav & X \text{ et } Y \text{ objets de } E{\quoi},\\
\Xi_v & : \scrHom(v_{!}X,Y) &&\longrightarrow
v_{\ast}\scrHom(X,v^{\ast}Y) &\quav & X\in \ob E, Y\in \ob
E'{\quoi},\\
\Lambda_v & : v_{!}(X\times v^{\ast}Y) &&\longrightarrow
v_{!}(X)\times Y & & X\in \ob E, Y\in \ob E'{\quoi},
\end{alignat*}
où,\pageoriginale pour définir les morphismes $\Xi_v$ et $\Lambda_v$
on suppose que le foncteur $v^{\ast}$ image réciproque par $v$ admet
un adjoint à gauche $v_{!}$.

\subsubsection{}\etiquette[10.5.1]{IV.subsubsec10.5.1}
\emph{Définition de} $\Psi_v$. Soit $X$ un objet de $E$. On a un
morphisme d'adjonction $v^{\ast}v_{\ast}X\rightarrow X$, d'où, pour
tout objet $Z$ de $E'$, un morphisme bifonctoriel $v^{\ast}(Z\times
v_{\ast}X)\rightarrow (v^{\ast}Z)\times X$. On a donc, pour tout
objet $Y$ de $E$, un morphisme trifonctoriel
$$
\Hom_E((v^{\ast}Z)\times X, Y)\longrightarrow
\Hom_E(v^{\ast}(Z\times v_{\ast}X),Y){\quoi}.
$$
On en déduit, par adjonction, un morphisme trifonctoriel
$$
\Hom_{E'}(Z,v_{\ast}\scrHom(X,Y))\longrightarrow \Hom_{E'}(Z,
\scrHom(v_{\ast}X, v_{\ast}Y)){\quoi};
$$
d'où le morphisme $\Psi_v$.

\subsubsection{}\etiquette[10.5.2]{IV.subsubsec10.5.2}
\emph{Définition de} $\Lambda_v$. Pour tout objet $X$ de $E$, on a
un morphisme d'adjonction $X\rightarrow v^{\ast}v_{!}X$, d'où, pour
tout objet $Y$ de $\sisi{E}{E'}$, un morphisme bifonctoriel $X\times
% correction SGA4-CORR-0300 / SGA-SRC-001000: SGA4-T1-IV8395-41-A8786C19.
v^{\ast}Y\rightarrow v^{\ast}v_{!}(X)\times v^{\ast}Y\simeq
v^{\ast}(v_{!}(X)\times Y)$; d'où, par adjonction, le morphisme
$\Lambda_v$.

\subsubsection{}\etiquette[10.5.3]{IV.subsubsec10.5.3}
\emph{Définition de} $\Xi_v$. Soient $X$ un objet de $E$, $Z$ et $Y$
deux objets de $E'$. Le morphisme $\Lambda_v:v_{!}(X\times
v^{\ast}Z)\rightarrow v_{!}(X)\times Z$ fournit un morphisme
trifonctoriel $\Hom_{E'}(v_{!}(X)\times Z,Y)\rightarrow
\Hom_{E'}(v_{!}(X\times v^{\ast}Z),Y)$; d'où, par adjonction, un
% correction SGA4-CORR-0301 / SGA-SRC-001001: SGA4-T1-IV8406-42-161129D9.
morphisme trifonctoriel
$$
\Hom_{E'}(Z,\scrHom(v_{!}X,Y))\longrightarrow
\Hom_{E'}(Z,v_{\ast}\scrHom(X,v^{\ast}Y)){\quoi};
$$
d'où le morphisme $\Xi_v$.

\subsubsection{}\etiquette[10.5.4]{IV.subsubsec10.5.4}
\emph{Définition de} $\Phi_v$. Le morphisme trifonctoriel de
(\hyperref[IV.subsubsec10.5.3]{10.5.3})
$$
\Hom_{E'}(v_{!}(Z)\times X,Y)\longrightarrow \Hom_{E'}(v_{!}(Z\times
v^{\ast}X),Y){\quoi};
$$
où $Z$ est un objet de $E$ et $X$, $Y$ sont des objets de $E'$,
permet d'obtenir, par adjonction, un morphisme trifonctoriel,
$$
\Hom_E(Z,v^{\ast}\scrHom(X,Y))\longrightarrow
\Hom_E(Z,\scrHom(v^{\ast}X, v^{\ast}Y)){\quoi};
$$
d'où une définition du morphisme $\Phi_v$. Il existe une deuxième
manière de définir $\Phi_v$. Soient $X$, $Y$, $Z$ trois objets de
$E'$. Le foncteur $v^{\ast}:E'\rightarrow E$ fournit un morphisme
trifonctoriel
$$
\Hom_{E'}(Z\times X,Y)\longrightarrow \Hom_E (v^{\ast}(Z)\times
v^{\ast}(X), v^{\ast}Y){\quoi};
$$
d'où, par adjonction, un morphisme trifonctoriel
$$
\Hom_{E'}(Z,\scrHom(X,Y))\longrightarrow
\Hom_{E'}(Z,v_{\ast}\scrHom(v^{\ast}X, v^{\ast}Y)){\quoi}.
$$
On\pageoriginale a donc un morphisme bifonctoriel
$$
\scrHom(X,Y)\longrightarrow v_{\ast}\scrHom(v^{\ast}X,
v^{\ast}Y){\quoi};
$$
d'où, par adjonction, un morphisme dont on laisse au lecteur le soin
de vérifier que c'est bien le morphisme $\Phi_v$ défini ci-dessus.

\begin{enonce*}{Proposition 10.6}\etiquette[10.6]{IV.prop10.6}
Soit $v:E\rightarrow E'$ un morphisme de topos.
\begin{enumerate}
\renewcommand{\labelenumi}{\theenumi)}
\item Le morphisme $\Psi_v$ $(\hyperref[IV.subsubsec10.5.1]{10.5.1})$ est un
  isomorphisme pour tout
objet $Y$ de $E$ si et seulement si le morphisme d'adjonction
$v^{\ast}v_{\ast}X\rightarrow X$ est un isomorphisme. En
particulier, le morphisme $\Psi_v$ est un isomorphisme pour tout
couple d'objets $(X,Y)$ si et seulement si le foncteur
$v_{\ast}:E\rightarrow E'$ est pleinement fidèle, \ie si $v$ est un
plongement de topos.

\item Les conditions suivantes sont équivalentes:
\begin{enumerate}
\renewcommand{\theenumii}{\roman{enumii}}
\renewcommand{\labelenumii}{(\theenumii)}
\item Pour tout couple $(X,Y)$ d'objets de $E'$, le morphisme $\Phi_v$
$(\hyperref[IV.subsubsec10.5.4]{10.5.4})$ est un isomorphisme.

\item Le foncteur $v^{\ast}$ admet un adjoint à gauche $v_{!}$ et pour
tout objet $X$ de $E$ et tout objet $Y$ de $E'$, le morphisme
$\Xi_v$ $(\hyperref[IV.subsubsec10.5.3]{10.5.3})$ est un isomorphisme.

\item Le foncteur $v^{\ast}$ admet un adjoint à gauche $v_{!}$ et pour
tout objet $X$ de $E$ et pour tout objet $Y$ de $E'$, le morphisme
$\Lambda_v$ $(\hyperref[IV.subsubsec10.5.2]{10.5.2})$ est un isomorphisme.
\end{enumerate}
\end{enumerate}
\end{enonce*}

\setcounter{subsection}{6} \setcounter{subsubsection}{0}
\subsubsection{}\etiquette[10.6.1]{IV.subsubsec10.6.1}
La première assertion résulte immédiatement de
(\hyperref[IV.subsubsec10.5.1]{10.5.1}). Il résulte aussi immédiatement de
(\hyperref[IV.subsubsec10.5.2]{10.5.2}), (\hyperref[IV.subsubsec10.5.3]{10.5.3}),
(\hyperref[IV.subsubsec10.5.4]{10.5.4}) que (iii) $\Leftrightarrow$ (ii)
$\Leftrightarrow$ ((i) et le foncteur $v^{\ast}$ admet un adjoint à
gauche $v_{!}$). Il reste donc à montrer que (i) implique que
$v^{\ast}$ admet un adjoint à gauche, \ie (\hyperref[IV.rem1.8]{1.8}) que (i)
implique que $v^{\ast}$ commute aux petits produits. Soient $e'$
l'objet final de $E'$ et $I$ un petit ensemble. Comme $v^{\ast}$ est
exact à gauche, $v^{\ast}(e')$ est un objet final de $E$ et comme
$v^{\ast}$ commute aux limites inductives $v^{\ast}(\amalg_I
e')\simeq \amalg_I v^{\ast}(e')$. Pour tout objet $Y$ de $E'$, on a
$\scrHom(\amalg_I e', Y)\simeq \Pi_I\scrHom(e',Y)\simeq \Pi_I Y$ et
$\sisi{\scrHom(\amalg_I v^{\ast}(e')$, $v^{\ast}Y)}
{\scrHom(\amalg_I v^{\ast}(e'),v^{\ast}Y)}\simeq
\sisi{\Pi_i}{\Pi_I}\scrHom(v^{\ast}(e'), v^{\ast}Y)\simeq \Pi_I
% correction SGA4-CORR-0302 / SGA-SRC-001002: SGA4-T1-IV8495-43-EED03651.
v^{\ast}Y$.\pageoriginale Le morphisme fonctoriel $\Phi_v$ induit
donc un isomorphisme $v^{\ast}\Pi_I Y\simeq \Pi_I v^{\ast}Y$ dont on
vérifie que c'est l'isomorphisme canonique. \cqfd

\begin{enonce*}{Corollaire 10.7}\etiquette[10.7]{IV.cor10.7}
Soient $E$ un topos, $Z$ un objet de $E$, $j_Z:E_{/Z}\rightarrow E$
le morphisme de localisation (\hyperref[IV.subsec5.2]{5.2}), $X$, $Y$
deux objets de $E$.
\begin{enumerate}
\renewcommand{\labelenumi}{\theenumi)}
\item Il existe un isomorphisme bifonctoriel
$$
j^{\ast}_Z\scrHom(X,Y)\simeq \scrHom(j^{\ast}_ZX,
j^{\ast}_ZY){\quoi}.
$$

\item Soit $X'$ un objet de $E_{/Z}$. Il existe un isomorphisme
bifonctoriel
$$
\scrHom_E(j_{Z!}X',Y)\simeq j_{Z\ast}\scrHom_{E/Z}(X',j^{\ast}_Z
Y){\quoi}.
$$

\item Soit $Y'$ un objet de $E/Z$. Lorsque $Z$ est un ouvert de $E$,
il existe un isomorphisme bifonctoriel
$$
j_{Z\ast}\scrHom(X',Y')\simeq \scrHom(j_{Z\ast}X',
j_{Z\ast}Y'){\quoi}.
$$
\end{enumerate}
\end{enonce*}

Les assertions 1) et 2) se démontrent en remarquant que le morphisme
$\Lambda_{j_Z}$ est un isomorphisme. Pour l'assertion 3), il suffit
de remarquer que $j_{Z\ast}$ est pleinement fidèle, car $j_{Z!}$ est
pleinement fidèle (\hyperrefext[I][I.rem5.7]{5.7}. a)).

\begin{enonce*}{Corollaire 10.8}\etiquette[10.8]{IV.cor10.8}
Soient $E$ un topos, $X$, $Y$ et $Z$ trois objets de $E$,
$j_Z:E_{/Z}\rightarrow E$ le morphisme de localisation.
\begin{enumerate}
\renewcommand{\labelenumi}{\theenumi)}
\item On a un isomorphisme canonique
$$
j_{Z\ast}j^{\ast}_ZY\simeq \scrHom(Z,Y){\quoi}.
$$

\item On a un isomorphisme canonique
$$
\Hom_E(Z,\scrHom(X,Y))\simeq \Hom_{E_{/Z}}(j^{\ast}_ZX,
j^{\ast}_ZY){\quoi}.
$$
\end{enumerate}
\end{enonce*}

Soit $e_Z$ l'objet final de $E_{/Z}$ (\ie l'objet
$\id_Z:Z\rightarrow Z$)). Il résulte de \eqhyperref[IV.eq10.2.1]{10.2.1} qu'on a un
isomorphisme canonique
$$
\scrHom_{E/Z}(e_Z, j^{\ast}_ZY)\simeq j^{\ast}_ZY{\quoi};
$$
d'où, d'après (\hyperref[IV.cor10.7]{10.7} 2), un isomorphisme
$$
j_{Z\ast}j^{\ast}_Z Y\simeq \scrHom(j_{Z!}e_Z, Y){\quoi.}
$$
Comme $j_{Z!}e_Z=Z$, on a démontré 1). Démontrons 2). De
(\hyperref[IV.cor10.7]{10.7}), on tire un\pageoriginale isomorphisme canonique
$$
\Hom_{E_{/Z}}(e_Z, j^{\ast}_Z\scrHom(X,Y))\simeq
\Hom_{E_{/Z}}(j^{\ast}_Z X, j^{\ast}_ZY){\quoi};
$$
d'où, par adjonction sur le premier membre,
$$
\Hom_E(Z,\scrHom(X,Y))\simeq \Hom_{E_{/Z}}(j^{\ast}_ZX,
j^{\ast}_ZY){\quoi}.
$$

\section{Topos annelés, localisation dans les topos annelés}\etiquette[11]{IV.sec11}

\setcounter{subsection}{1} \setcounter{subsubsection}{0}
\subsubsection{}\etiquette[11.1.1]{IV.subsubsec11.1.1}
Soit $𝒰$ un univers. On appelle $𝒰$-\emph{topos
annelé} un couple $(E,A)$ où $E$ est un $𝒰$-topos et $A$ un
objet muni d'une structure d'anneau. On appelle $𝒰$-site
annelé un couple $(C,A)$ où $C$ est un $𝒰$-site et $A$ un
$𝒰$-\sisi{faisceaux}{faisceau} d'anneaux sur $C$. On ne mentionne pas
% correction SGA4-CORR-0303 / SGA-SRC-001003: SGA4-T1-IV8581-44-E8F45FAC.
l'univers lorsque le contexte ne prête pas à confusion. À un site
annelé $(C,A)$ est associé le topos annelé $(C^{\sim},A)$. À un
topos annelé $(E,A)$ est associé le site annelé
constitué par le
site $E$ et le faisceau d'anneaux représenté par $A$.

\subsubsection{}\etiquette[11.1.2]{IV.subsubsec11.1.2}
Soit $(E,A)$ un topos annelé. On note ${}_AE$ (\resp $E_A$) la
catégorie des faisceaux de $A$-modules (nous écrirons aussi
$A$-Modules) à gauche (\resp à droite) unitaires. La catégorie
${}_AE$ (\resp $E_A$) est une catégorie abélienne
(\hyperrefext[II][II.prop6.7]{6.7}). Soient $M$ et $N$ deux faisceaux de $A$-modules
à gauche (\resp à droite). Le groupe commutatif des morphismes de
$A$-Modules de $M$ dans $N$ est noté $\Hom_A(M,N)$.

\subsubsection{}\etiquette[11.1.3]{IV.subsubsec11.1.3}
Soient $A$, $B$ et $C$ trois anneaux d'un topos $E$, $M$ un faisceau
de $A-B$ bimodules, $N$ un faisceau de $A-C$ bimodules à gauche.
Soit $e$ un objet final de $E$ et posons $\Hom(e,B)=\Gamma(B)$,
$\Hom(e,C)=\Gamma(C)$. La structure de $A-B$ bimodule de $M$ fournit
un homomorphisme d'anneaux de $\Gamma(B)^{\circ}$ dans l'anneau des
endomorphismes du $A$-module $M$ et de même, la structure de $A-C$
bimodule de $N$ fournit un homomorphisme d'anneaux de
$\Gamma(C)^{\circ}$ dans l'anneau des endomorphismes du $A$-module
$N$. On en déduit, par fonctorialité, une structure de
$\Gamma(B)-\Gamma(C)$ bimodule sur le groupe commutatif
$\Hom_A(M,N)$. En particulier, lorsque $A$ est un faisceau d'anneaux
commutatifs, le groupe $\Hom_A(M,N)$ est muni canoniquement d'une
structure de $\Gamma(A)$-module.

\subsubsection{}\etiquette[11.1.4]{IV.subsubsec11.1.4}
Soient\pageoriginale $E$ un topos, $A$ et $B$ deux anneaux de $E$,
$u:A\rightarrow B$ un morphisme de faisceaux d'anneaux, $M$ un
$B$-module (à gauche pour fixer les idées). Le faisceau $M$
\sisi{peut-être}{peut être} considéré comme un faisceau de $A$-modules par
% correction SGA4-CORR-0304 / SGA-SRC-001004: SGA4-T1-IV8616-45-FFEBCCE6.
l'intermédiaire de $u$: pour tout objet $X$ de $E$, $M(X)$ est muni
de la structure de $A(X)$-module déduite de sa structure de
$B(X)$-module et de l'homomorphisme $u(X):A(X)\rightarrow B(X)$ par
restriction des scalaires. On obtient ainsi un foncteur « \emph{restriction des scalaires par} $u$ »
$$
\Res(u):{}_BE\longrightarrow {}_AE{\quoi}.
$$

\subsubsection{}\etiquette[11.1.5]{IV.subsubsec11.1.5}
En particulier lorsque $A=\ZZ$ (faisceau constant $\sisi{\underline{\underline{\ZZ}}}{\underline{\ZZ}}$ \ie faisceau
associé au préfaisceau constant $\ZZ$) et lorsque $u:\sisi{\underline{\underline{\ZZ}}}{\underline{\ZZ}}\rightarrow B$
est l'unique morphisme canonique, on obtient un foncteur de ${}_BE$
% correction SGA4-CORR-0305 / SGA-SRC-001005: SGA4-T1-IV8629-46-59AD3CE6.
dans la catégorie ${}_{\sisi{\underline{\underline{\ZZ}}}{\underline{\ZZ}}}E$ qui n'est autre que la catégorie des
faisceaux abéliens notée $E_{Ab}$. Ce foncteur est appelé le
\emph{foncteur faisceau abélien sous-jacent}.

\begin{enonce*}{Proposition 11.1.6}\etiquette[11.1.6]{IV.prop11.1.6}
Le foncteur restriction des scalaires par $u$ commute aux limites
inductives et projectives. Il est conservatif.
\end{enonce*}

La proposition est vraie pour le foncteur restriction des scalaires
pour les modules ordinaires, \ie lorsque $E$ est le topos ponctuel.
Elle est donc vraie lorsque $E$ est le topos des préfaisceaux sur un
petit site. Il résulte alors de la détermination des limites
inductives et projectives à l'aide du foncteur faisceau associé
(\hyperrefext[II][II.prop6.4]{6.4}) que la proposition est vraie dans le cas général.

\setcounter{subsection}{2} \setcounter{subsubsection}{0}
\subsubsection{}\etiquette[11.2.1]{IV.subsubsec11.2.1}
Soient $(E,A)$ un topos annelé, $X$ un objet de $E$,
$j_X:E_{/X}\rightarrow E$ le morphisme de localisation
\sisi{(\hyperref[IV.sec8]{8})}{\hyperref[IV.subsec5.2]{5.2}}. D'après \hyperrefext[III][III.prop1.7]{1.7} ou
\hyperref[IV.subsubsec3.1.2]{3.1.2}, le faisceau $j^{\ast}_XA$ est muni
canoniquement d'une structure d'anneau. Le faisceau d'anneaux
$j^{\ast}_XA$ est noté le plus souvent $A| X$ ou bien encore,
abusivement, $A$. Sauf mention du contraire, le topos $E_{/X}$ sera
annelé par $A| X$. Le faisceau $j_{X\ast}j^{\ast}_XA$ est muni
canoniquement d'une structure d'anneau. Le morphisme d'adjonction
$A\rightarrow j_{X\ast}j^{\ast}_XA$ est un morphisme de faisceaux
d'anneaux.

\subsubsection{}\etiquette[11.2.2]{IV.subsubsec11.2.2}
Soit\pageoriginale de plus $M$ un $A$-module (à gauche pour fixer
les idées). Le faisceau $j^{\ast}_XM$ est muni d'une structure de
$A| X$-module (\hyperrefext[III][III.prop1.7]{1.7}; ou \hyperref[IV.subsubsec3.1.2]{3.1.2}); d'où
un foncteur:
$$
j^{\ast}_X:{}_AE\longrightarrow {}_{A| X}E_{/X}
$$
appelé foncteur de restriction à $E_{/X}$, ou encore foncteur de
restriction à $X$. Le foncteur $j^{\ast}_X$ commute aux limites
inductives et projectives (\loccit). Il est en particulier exact.
Soit maintenant $N$ un $A| X$-Module. Le faisceau $j_{X\ast}N$ est
un faisceau de $j_{X\ast}A| X$-Modules (\hyperrefext[III][III.prop1.7]{1.7} ou 
\hyperref[IV.subsubsec3.1.2]{3.1.2}); d'où, par restriction des scalaires par le
morphisme d'adjonction $A\rightarrow j_{X\ast}(A| X)$
(\hyperref[IV.subsubsec11.1.4]{11.1.4}), un $A$-Module encore noté
$j_{X\ast}N$. On a donc défini un foncteur
$$
j_{X\ast}:{}_{A| X}E_{/X}\longrightarrow {}_AE
$$
qui commute aux limites projectives (\hyperref[IV.prop11.1.6]{11.1.6} et
\hyperrefext[III][III.prop1.7]{1.7}). Pour tout $A$-Module $M$, le morphisme
d'adjonction $M\rightarrow j_{X\ast}j^{\ast}_XM$ est un morphisme de
$A$-Modules. Pour tout $A| X$-Module $N$ et tout $A$-Module $M$, le
morphisme d'adjonction $M\rightarrow j_{X\ast}j^{\ast}_XM$ définit
un morphisme bifonctoriel en $M$ et $N$
\begin{equation*}
\Hom_{A| X}(j^{\ast}_XM,N)\longrightarrow
\Hom_A(M,j_{X\ast}N){\quoi}.\tag{11.2.2.1}
\end{equation*}\etiquette[11.2.2.1]{IV.eq11.2.2.1}
Ce dernier morphisme est un isomorphisme \ie les foncteurs
$j^{\ast}_X$ et $j_{X\ast}$ pour les $A$-Modules, sont adjoints (III
1.7).

\subsubsection{}\etiquette[11.2.3]{IV.subsubsec11.2.3}
Les foncteurs $j^{\ast}_X$ et $j_{X\ast}$ pour les Modules,
commutent au foncteur « ensemble sous-jacent »
(\hyperrefext[III][III.prop1.7]{1.7} 4)). Ils commutent donc aux foncteurs restriction
des scalaires et en particulier au foncteur faisceau abélien
sous-jacent.

\setcounter{subsection}{3} \setcounter{subsubsection}{0}
\begin{enonce*}{Proposition 11.3.1}\etiquette[11.3.1]{IV.prop11.3.1}
Soient $(E,A)$ un topos annelé, $X$ un objet de $E$. Le foncteur
$$
j^{\ast}_X:{}_{A}E\longrightarrow {}_{A| X}E_{/X}
$$
admet un adjoint à gauche noté $j_{X!}:{}_{A| X}E_{/X}\rightarrow
{}_{A}E$ et appelé le prolongement par zéro. Le foncteur
prolongement par zéro est exact et fidèle et commute aux limites
inductives. Les foncteurs $j_{X!}$ pour les Modules\pageoriginale
commutent aux foncteurs restriction des scalaires et, en
particulier, ils commutent au foncteur faisceau abélien sous-jacent.
\end{enonce*}

L'existence du foncteur $j_{X!}$ résulte de \hyperrefext[III][III.prop1.7]{1.7}. Comme
$j_{X!}$ est un adjoint à gauche, il commute aux limites inductives
(\hyperrefext[I][I.sec2]{2}). Pour démontrer les autres assertions, supposons
d'abord que $E$ soit le topos des préfaisceaux d'ensembles sur une
petite catégorie $C$ contenant l'objet $X$. On sait alors que
$E_{/X}$ est équivalent à la catégorie des préfaisceaux sur $C_{/X}$
et que, modulo cette équivalence, le foncteur
$j^{\ast}_X:E\rightarrow E_{/X}$ n'est autre que la composition avec
le foncteur d'oubli $C_{/X}\rightarrow C$ (\hyperrefext[I][I.prop5.11]{5.11}). Il
résulte alors immédiatement de la construction explicite de $j_{X!}$
(\hyperrefext[I][I.prop5.1]{5.1}) que pour tout $A| X$-Module $N$ et pour tout objet
$Y$ de $C$, on a
$$
j_{X!}N(Y)=\bigoplus_{u\in \Hom_C(Y,X)}N(u){\quoi};
$$
d'où l'exactitude de $j_{X!}$ et le fait que le prolongement par
zéro commute aux foncteurs restriction des scalaires dans ce cas.
Dans le cas général, on peut supposer que $E$ est le topos des
faisceaux sur un petit site $C$ et que $X$ provient d'un objet de
$C$ (\hyperref[IV.cor1.2.1]{1.2.1}). Le topos $E_{/X}$ est alors équivalent au
topos des faisceaux sur $C_{/X}$ (muni de la topologie induite)
(\hyperrefext[III][III.prop5.4]{5.4}) et le morphisme de topos $j_X:E_{/X}\rightarrow
E$ provient du foncteur d'oubli $C_{/X}\rightarrow C$ qui est
continu et cocontinu (\hyperrefext[III][III.prop5.2]{5.2}). Il résulte alors de
\hyperrefext[III][III.prop1.7]{1.7} que le prolongement par zéro pour les
% correction SGA4-CORR-0306 / SGA-SRC-001006: SGA4-T1-IV8738-47-18F0F667.
faisceaux s'obtient en composant le prolongement par zéro pour les
préfaisceaux avec le foncteur faisceau associé; d'où l'assertion
d'exactitude et la commutation aux restrictions des scalaires.

Pour démontrer la fidélité, il revient au même de
montrer que pour
$A| X$-Module $N$, le morphisme d'adjonction
$$
\ad_N:N\longrightarrow j^{\ast}_X j_{X!}N
$$
est un monomorphisme. Or, en notant $j_{\widehat{X}!}$ le foncteur
prolongement par zéro pour les préfaisceaux et $\underline{a}$ le
foncteur « faisceau associé », on a un diagramme commutatif
\[
\xymatrix@C=2.75cm@R=1.5cm{ N \ar[r]^{\ad_{N}}
\ar[rd]_{\ad_{\hat{N}}} & j^{\ast}_{X} \underline{a}
j_{\widehat{X}!} N \\
                   & \underline{a} j^{\ast}_X j_{\widehat X!} N,\ar[u]
% correction SGA4-CORR-0307 / SGA-SRC-001007: SGA4-T1-IV8757-48-587E80B1.
}
\]
où\pageoriginale $\ad_{\widehat{N}}$ est le morphisme d'adjonction
pour les préfaisceaux. Le morphisme $\ad_{\widehat{N}}$ est un
monomorphisme d'après ce qui précède, donc
$\underline{a}(\ad_{\widehat{N}})$ est un monomorphisme. De plus, comme
le foncteur d'oubli $C_{/X}\rightarrow C$ est continu et cocontinu,
le morphisme canonique $\underline{a}j^{\ast}_X\rightarrow
j^{\ast}_X\underline{a}$ est un isomorphisme (\hyperrefext[III][III.prop2.3]{2.3}). Par
suite $\ad_N$ est un monomorphisme.

\begin{enonce*}[remark]{Remarque 11.3.2}\etiquette[11.3.2]{IV.rem11.3.2}
Soit $N$ un $A| X$-Module. Le faisceau d'ensemble sous-jacent à
$j_{X!}N$ n'est pas, en général, isomorphe au faisceau obtenu en
prolongeant par le vide (\ref{III.subsec5.3} et \hyperref[IV.subsec5.2]{5.2}) le faisceau
d'ensemble sous-jacent à $N$. On prendra donc garde de ne pas
confondre le foncteur \emph{prolongement par zéro} noté
$j_{X!}:{}_{A|X}E_{/X}\rightarrow {}_{A}E$ dans \hyperref[IV.prop11.3.1]{11.3.1},
et le foncteur prolongement par le vide noté encore
$j_{X!}:E_{/X}\rightarrow E$ dans III 5.3 et IV 5.2. Dans la plupart
des cas rencontrés dans la pratique, l'abus de notation signalé
ci-dessus n'amène pas de confusions. Lorsqu'une confusion est
néanmoins possible, nous utiliserons les notations $j^{\ab}_{X!}$ et
$j^{\ens}_{X!}$ pour désigner respectivement les foncteurs
prolongement par zéro et prolongement par le vide.
\end{enonce*}

\begin{enonce*}{Proposition 11.3.3}\etiquette[11.3.3]{IV.prop11.3.3}
Soient $(E,A)$ un topos annelé et $X$ un objet de $E$. Le $A$-Module
$j_{X!}(A| X)$, noté le plus souvent $A_X$ ou $A_{X,E}$, est le
$A$-Module libre engendré par $X$ {\rm (\hyperrefext[II][II.prop6.5]{6.5})} \ie pour
tout $A$-Module $M$, on a un isomorphisme canonique, fonctoriel en
$M$:
$$
\Hom_E(X,M)\xleftarrow{\sim} \Hom_A (A_X, M){\quoi}.
$$
Soit $(X_i)_{i\in I}$ une famille topologiquement génératrice de $E$
{\rm (\hyperrefext[II][II.def3.0.1]{3.0.1}).} La famille $(A_{X_i})_{i\in I}$ est une
famille génératrice de la catégorie des $A$-Modules.
\end{enonce*}

Soit $e_X$ l'objet final du topos $E_{/X}$
% correction SGA4-CORR-0308 / SGA-SRC-001008: SGA4-T1-IV8799-49-5776A488.
$(e_X=X\xrightarrow{\id}X)$. On a un isomorphisme
$$
\Hom_{A| X}(A| X, j^{\ast}_X M)\xrightarrow{\sim} \Hom_{E_{/X}}(e_X,
j^{\ast}_X M){\quoi};
$$
déduit\pageoriginale de la section unité $e_X\rightarrow \sisi{A/X}{A|X}$; d'où,
% correction SGA4-CORR-0309 / SGA-SRC-001009: SGA4-T1-IV8805-50-37406B1D.
par adjonction, un isomorphisme
$$
\Hom_A(j^{\ab}_{X!}(A| X),
M)\xrightarrow{\sim}\Hom_E(j^{\ens}_{X!}(e_X),M){\quoi}.
$$
Comme $j^{\ens}_{X!}(e_X)=X$, on obtient l'isomorphisme annoncé. La
dernière assertion résulte de \hyperrefext[II][II.cor6.6]{6.6}.

\begin{enonce*}[remark]{Remarque 11.3.4}\etiquette[11.3.4]{IV.rem11.3.4}
Soient $A$ et $B$ deux faisceaux d'anneaux sur un topos $E$, $X$ un
objet de $E$ et $N$ un $A|X$-$B| X$-bi Module. Le faisceau abélien
obtenu en prolongeant $N$ par zéro est muni canoniquement d'une
structure de $A$-$B$-biModule ainsi qu'il résulte immédiatement de
sa description explicite (\hyperref[IV.prop11.3.1]{11.3.1}). En particulier le
$A$-module libre engendré $A_X$ est un $A$-biModule.
\end{enonce*}

\begin{enonce*}[remark]{Exercice 11.3.5}\etiquette[11.3.5]{IV.exer11.3.5}
Soient $E$ un topos, $X$ un objet de $E$, $x:P\rightarrow E$ un
point de $E$ (\hyperref[IV.def6.1]{6.1}), $(x_i)_{i\in I}$ la famille des
points de $E_{/X}$ au-dessus de $x$ (en correspondance biunivoque
avec la fibre $X_x$ (\hyperref[IV.eq6.7.2]{6.7.2}). Montrer que pour tout
% correction SGA4-CORR-0310 / SGA-SRC-001010: SGA4-T1-IV8825-51-5821ED8B.
faisceau abélien $M$ sur $E_{/X}$, la fibre $(j_{X!}M)_x$ est
canoniquement isomorphe à $\oplus_i M_{x_i}$.
\end{enonce*}

\section{Opération sur les modules}\etiquette[12]{IV.sec12}

\begin{enonce*}{Proposition 12.1}\etiquette[12.1]{IV.prop12.1}
Soient $(E,A)$ un topos annelé, $M$ et $N$ deux $A$-Modules à gauche
(\resp à droite). Le foncteur sur $E$ qui a tout objet $X$ de $E$
associe le groupe commutatif $\Hom_A(M,\scrHom_E(X,N))$ est
représentable par un faisceau abélien noté $\scrHom_A(M,N)$ (ou
parfois $\scrHom(M,N)$ lorsqu'aucune confusion n'en résulte). Pour
tout objet $X$ de $E$ on a un isomorphisme canonique
\begin{equation*}
\scrHom_A(M,N)(X)\simeq \Hom_{A| X}(j^{\ast}_XM,
j^{\ast}_XN){\quoi}.\tag{12.1.1}
\end{equation*}\etiquette[12.1.1]{IV.eq12.1.1}
\end{enonce*}

On a un isomorphisme canonique $\scrHom_E(X,N)=j_{X\ast}j^{\ast}_XN$
(\hyperref[IV.cor10.8]{10.8}) et par suite $\scrHom_E(X,N)$ est muni
fonctoriellement en $X$ d'une structure de $A$-Module à gauche
(\resp à droite). Le foncteur $X\rightarrow \scrHom_E(X,N)$
transforme les limites inductives de l'argument $X$ en limites
projectives de $A$-Modules (\hyperref[IV.subsec10.2]{10.2} et
\hyperref[IV.subsubsec11.2.3]{11.2.3}) et par suite le foncteur
$$
X\longrightarrow \Hom_A(M,\scrHom_E(X,N))
% correction SGA4-CORR-0311 / SGA-SRC-001011: SGA4-T1-IV8854-52-06331E21.
$$
transforme les limites inductives\pageoriginale de l'argument $X$ en
limites projectives de groupes commutatifs, donc en limites
projectives des ensembles sous-jacents. Il est donc représentable
(\hyperref[IV.prop1.4]{1.4} et \hyperref[IV.theo1.2]{1.2}) et l'objet qui le représente
est muni d'une structure de faisceau abélien. On a, par définition,
pour tout objet $X$ de $E$, un isomorphisme canonique
\begin{equation*}
\scrHom_A(M,N)(X)=\Hom(X,\scrHom_A(M,N))\simeq
\Hom_A(M,\scrHom(X,N))\tag{12.1.2}
\end{equation*}\etiquette[12.1.2]{IV.eq12.1.2}
et l'isomorphisme (\hyperref[IV.eq12.1.1]{12.1.1}) résulte de \eqhyperref[IV.eq12.1.2]{12.1.2}, de
l'isomorphisme $\scrHom(X,N)\approx j_{X\ast}j^{\ast}_XN$
(\hyperref[IV.cor10.8]{10.8}) et des formules d'adjonction de 10.

\setcounter{subsection}{1}
\subsection{}\etiquette[12.2]{IV.subsec12.2}
Le faisceau abélien $\scrHom_A(M,N)$ est appelé le \emph{faisceau
% correction SGA4-CORR-0312 / SGA-SRC-001012: SGA4-T1-IV8873-53-B501F382.
des morphismes de $A$-Modules de $M$ dans $N$}. C'est un bifoncteur
en $M$ et $N$. Il résulte de sa définition et de
(\hyperref[IV.subsec10.2]{10.2}) qu'il transforme les limites projectives de
l'argument $N$ (\resp les limites inductives de $M$) en limites
projectives de faisceaux abéliens. En particulier il est exact à
gauche en ses deux arguments.

\begin{enonce*}{Proposition 12.3}\etiquette[12.3]{IV.prop12.3}
Soient $(E,A)$ un topos annelé, $M$ et $N$ deux $A$-Modules à droite
(\resp à gauche), $X$ un objet de $E$:
\begin{enumerate}
\renewcommand{\theenumi}{\alph{enumi}}
\renewcommand{\labelenumi}{\theenumi)}
\item On a des isomorphismes canoniques:
$$
\scrHom_A(M,N)(X)\approx \Hom_A(M, j_{X\ast}j^{\ast}_XN)\simeq
\Hom_{A| X}(j^{\ast}_XM, j^{\ast}_XN)\simeq
\Hom_A(j_{X!}j^{\ast}_XM, N)
$$

\item On a un isomorphisme canonique:
$$
\Phi_X:j^{\ast}_X\scrHom_A (M,N)\simeq \scrHom_{A| X}(j^{\ast}_X M,
j^{\ast}_XN)
$$
Soit de plus $P$ un $A| X$-Module à droite (\resp à gauche):

\item On a un isomorphisme canonique:
$$
j_{X\ast}\scrHom_{A| X}(j^{\ast}_X M,P)\simeq
\scrHom_A(M,j_{X\ast}P)
$$

\item On a un isomorphisme canonique:
$$
\Xi_X:\scrHom_A (j_{X!}P,N)\simeq j_{X\ast}\scrHom_{A|
X}(P,j^{\ast}_XN)
$$
\end{enumerate}
\end{enonce*}

\noindent Les isomorphismes de a) résultent de \eqhyperref[IV.eq12.1.2]{12.1.2}, de
l'isomorphisme $\scrHom_E(X,N)\simeq j_{X\ast}j^{\ast}_XN$
(\hyperref[IV.cor10.8]{10.8}) et des formules d'adjonction de 10.

\begin{itemize}
\item[Démontrons \sisi{d}{b}).] Pour tout objet $Y$ de $E_{/X}$ on a la suite
% correction SGA4-CORR-0313 / SGA-SRC-001013: SGA4-T1-IV8920-54-6BF2430C.
d'isomorphismes:

$\Hom(Y,j^{\ast}_X\scrHom_A(M,N))\simeq \Hom(j_{X!}Y,\scrHom_A
(M,N))$ (adjonction pour les faisceaux d'ensembles)
\begin{alignat*}{2}
& \Hom(j_{X!}Y,\scrHom_A(M,N))\simeq \Hom_A (M,\scrHom_E(j_{X!}Y,
N))
&\quad& (\sisi{12.2.1}{12.1.2})\\
& \sisi{\Hom}{\Hom_A}(M, \scrHom_E(j_{X!}Y,N))\simeq \Hom_A
% correction SGA4-CORR-0315 / SGA-SRC-001015: SGA4-T1-IV8929-56-1766266A.
(M,j_{X\ast}\scrHom_{E_{/X}}(Y,j^{\ast}_XN)) && (\hyperref[IV.cor10.7]{10.7})\\
& \Hom_{\sisi{\Lambda}{A}}(M, j_{X\ast}\scrHom_{E_{/X}}(Y, j^{\ast}_XN))\simeq
% correction SGA4-CORR-0316 / SGA-SRC-001016: SGA4-T1-IV8931-57-47001EA9.
\Hom_{\sisi{A/X}{A|X}}(j^{\ast}_XM, \scrHom_E (Y,j^{\ast}_XN)) && \eqhyperref[IV.eq11.2.2.1]{11.2.2.1}\\
 & \Hom_{A| X}(j^{\ast}_XM, \scrHom_E (Y, j^{\ast}_XN))\simeq
\Hom(Y,\scrHom_{A| X}(j^{\ast}_X M, j^{\ast}_X\sisi{M}{N})) && (\sisi{12.2.1}{12.1.2})
% correction SGA4-CORR-0318 / SGA-SRC-001018: SGA4-T1-IV8934-59-7D10B7E1.
\end{alignat*}\pageoriginale

\item[Démontrons c).] Pour tout objet $Y$ de $E$, on a la suite
d'isomorphismes:
\begin{alignat*}{2}
& \Hom(Y,j_{X\ast}\scrHom_{A| X}(j^{\ast}_XM,P))\simeq
\Hom(j^{\ast}_XY, \scrHom_{A| X}(j^{\ast}_XM; P)) &\quad&
\text{(adjonction)}\\
&\Hom (j^{\ast}_XY, \scrHom_{A| X}(j^{\ast}_XM, P))\simeq
\Hom_{\sisi{A/X}{A|X}}(j^{\ast}_XM, \scrHom_{E_{/X}}(j^{\ast}_XY,P)) && (\sisi{12.2.1}{12.1.2})\\
&\Hom_{A| X}(j^{\ast}_XM, \scrHom_{E_{/X}}(j^{\ast}_XY,P))\simeq
\Hom_A(M, j_{X\ast}\scrHom_{E_{/X}}(j^{\ast}_XY,P)) && \eqhyperref[IV.eq11.2.2.1]{11.2.2.1}\\
&\Hom_A(M,j_{X\ast}\scrHom_{E_{/X}}(j^{\ast}_XY,P))\simeq
\Hom_A(M,\scrHom_E(Y, j_{X\ast}P)) && (\hyperref[IV.prop10.3]{10.3})\\
&\Hom_A(M,\scrHom_E(Y, j_{X\ast}P))\simeq \Hom(Y,
\scrHom_A(M,j_{X\ast}P)) && (\sisi{12.2.1}{12.1.2})
\end{alignat*}

\item[Démontrons d).] Pour tout objet $Y$ de $E$, on a la suite
% correction SGA4-CORR-0319 / SGA-SRC-001019: SGA4-T1-IV8953-60-0B3DDC2A.
d'isomorphismes:
\begin{alignat*}{2}
& \Hom(Y,\scrHom_A(j_{X!}P,N))\simeq \Hom_A(j_{X!}P,\scrHom_E(Y,N))
&\quad& (\sisi{12.2.1}{12.1.2})\\
&\Hom_A(j_{X!}P,\scrHom_E(Y,N))\simeq \Hom_{A|
X}(P,j^{\ast}_X\scrHom_E(Y,N)) && (\hyperref[IV.prop11.3.1]{11.3.1})\\
&\Hom_{\sisi{A/X}{A|X}}(P,j^{\ast}_X\scrHom_E(Y,N))\simeq \Hom_{A|
X}(P,\scrHom_{E_{/X}}(j^{\ast}_XY, j^{\ast}_XN)) && (\hyperref[IV.cor10.7]{10.7})\\
&\Hom_{\sisi{A/X}{A|X}}(P,\scrHom_{E_{/X}}(j^{\ast}_XY, j^{\ast}_XN))\simeq
\Hom(j^{\ast}_XY, \scrHom_{\sisi{A/X}{A|X}}(P,j^{\ast}_XN)) && (\sisi{12.2.1}{12.1.2})\\
% correction SGA4-CORR-0314 / SGA-SRC-001014: six defining citations.
% correction SGA4-CORR-0317 / SGA-SRC-001017: all five restricted-ring occurrences.
&\Hom(j^{\ast}_XY, \scrHom_{A| X}(P,j^{\ast}_XN))\simeq \Hom(Y,
% correction SGA4-CORR-0320 / SGA-SRC-001020: SGA4-T1-IV8964-61-508DD0E6.
j_{X\ast}\scrHom_{A| X}(P, j^{\ast}_XN)) && \text{(adjonction)}
\end{alignat*}
\end{itemize}

\begin{enonce*}{Corollaire 12.4}\etiquette[12.4]{IV.cor12.4}
Soient $C$ un $𝒰$-site, $A'$ un $𝒰$-préfaisceau
d'anneaux sur $C$, $M$ un $𝒰$-préfaisceau de $A'$-modules (à
gauche pour fixer les idées), $N$ un $𝒰$-faisceau de
$A'$-modules (à gauche). Notons $A$ le faisceau associé à $A'$, de
sorte que les faisceaux $N$ et $\underline{a}M$ (faisceau associé à $M$)
sont des $A$-Modules. Le préfaisceau $X\mapsto \Hom_{A'|
X}(j^{\ast}_XM, j^{\ast}_XN)$ est un $𝒰$-faisceau abélien
canoniquement isomorphe à $\scrHom_A(\underline{a}M,N)$.
\end{enonce*}

En effet il résulte de \hyperrefext[III][III.prop5.5]{5.5} et \hyperrefext[III][III.prop2.3]{2.3} que
le foncteur de localisation à $C_{/X}$ commute avec les foncteurs
faisceaux associés (pour $C$ et $C_{/X}$). Par suite, on a des
isomorphismes fonctoriels en l'objet variable $X$ de $C$:
\begin{align*}
\Hom_{A'| X}(j^{\ast}_XM, j^{\ast}_XN) &\simeq \Hom_{A|
X}(\underline{a}j^{\ast}_XM, j^{\ast}_XN)\\
&\simeq \Hom_{A| X}(j^{\ast}_X \underline{a}M, j^{\ast}_XN)\simeq
\scrHom_A (\underline{a}M, N)(X){\quoi}.
\end{align*}

\begin{enonce*}{Corollaire 12.5}\etiquette[12.5]{IV.cor12.5}
Soient\pageoriginale $E$ un topos, $A$, $B$, $C$ trois faisceaux
d'anneaux sur $E$, $M$ un $A-B$ biModule, $N$ un $A-C$ biModule. Le
faisceau abélien $\scrHom_A(M,N)$ est muni canoniquement d'une
structure de $B-C$ biModule. En particulier lorsque $A$ est un
faisceau d'anneaux commutatifs et lorsque $M$ et $N$ sont des
$A$-Modules, $\scrHom_A(M,N)$ est muni canoniquement d'une structure
de $A$-Module. Les isomorphismes canoniques {\rm b), c), d)} de
$\hyperref[IV.prop12.3]{12.3}$ sont des isomorphismes de biModules.
\end{enonce*}

Nous nous bornerons à donner des indications. Pour tout objet $X$ de
$E$, on a $\scrHom_A(M,N)(X)\simeq \Hom_{A| X}(j^{\ast}_XM,
j^{\ast}_XN)$ (\hyperref[IV.prop12.3]{12.3}). Par suite, le groupe commutatif\\
$\scrHom_A(M,N)(X)$ est muni canoniquement d'une structure de
$B(X)-C(X)$ bimodule (\hyperref[IV.subsubsec11.1.3]{11.1.3}) dont on vérifie
qu'elle est fonctorielle en $X$. Les autres assertions sont laissées
au lecteur.

\begin{enonce*}{Corollaire 12.6}\etiquette[12.6]{IV.cor12.6}
Soient $(E,A)$ un topos annelé, $X$ un objet de $\sisi{A}{E}$, $N$ un
% correction SGA4-CORR-0321 / SGA-SRC-001021: SGA4-T1-IV9011-62-C998766D.
$A$-Module (à gauche pour fixer les idées). On a des isomorphismes
canoniques de $A$-Modules:
$$
\scrHom_E(X,N)\simeq j_{X\ast}j^{\ast}_XN\simeq \scrHom_A(A_X,
N){\quoi};
$$
la structure de $A$-Module sur $\scrHom_A(A_X,N)$ provenant de la
structure de biModule sur $A_X$ $(\hyperref[IV.rem11.3.4]{11.3.4})$.
\end{enonce*}

Le premier isomorphisme résulte de \hyperref[IV.cor10.7]{10.7}, le deuxième de
l'isomorphisme de $A$-Modules $N\simeq \scrHom_A(A,N)$ et de
\hyperref[IV.prop12.3]{12.3}.

\begin{enonce*}{Proposition 12.7}\etiquette[12.7]{IV.prop12.7}
Soient $(E,A)$ un topos annelé, $M$ un $A$-Module à droite et $N$ un
$A$-Module à gauche. Le foncteur qui à tout faisceau abélien $P$
associe le groupe commutatif
$\Hom_A(M,\scrHom_{\sisi{\underline{\underline{\ZZ}}}{\underline{\ZZ}}}(N,P))$ est représentable par
un faisceau abélien noté $M\otimes_A N$ et appelé le produit
tensoriel sur $A$ de $M$ et de $N$.
\end{enonce*}

On a une injection canonique de $\Hom_A(M,\scrHom_{\sisi{\underline{\underline{\ZZ}}}{\underline{\ZZ}}}(N,P))$ dans
l'ensemble \\ $\Hom(M,\scrHom(N,P))\simeq \Hom(M\times N, P)$; On
constate aussitôt, en revenant aux définitions, qu'un morphisme $f$
de faisceaux d'ensembles de $M\times N$ dans $P$ provient \sisi{d'une}{d'un}
élément de $\Hom_A(M,\scrHom_Z(N,P))$ si et seulement si pour tout
% correction SGA4-CORR-0322 / SGA-SRC-001022: SGA4-T1-IV9039-63-6264364A.
objet $X$ de $E$, $f(X):M(X)\times N(X)\rightarrow P(X)$ est une
application\pageoriginale $A(X)$-bilinéaire de $M(X)\times N(X)$
dans $P(X)$. Appelons morphisme $A$-bilinéaire les morphismes de
faisceaux d'ensembles de $M\times N$ dans $P$ qui possèdent cette
propriété. Il s'agit donc de représenter le foncteur des morphismes
% correction SGA4-CORR-0323 / SGA-SRC-001023: SGA4-T1-IV9044-64-0731D6C2.
$A$-bilinéaires de $M\times N$ dans $P$. Pour cela on procède comme
dans le cas ordinaire \ie comme dans le cas où $E$ est le topos
ponctuel. Considérons les huit morphismes:
\begin{itemize}
\item[(i):] $M\times M\times N \longrightarrow M\times N$\qquad $1\leq
i\leq 3$

\item[(i):] $M\times N\times N\longrightarrow M\times N$\qquad $4\leq
i\leq 6$

\item[(i):] $M\times A\times N\longrightarrow M\times N$\qquad $7\leq
i\leq 8$
\end{itemize}
définis par les formules:
\begin{tabbing}
\quad\= (1) \= $(m_1, m_2, n)$ \= $\longmapsto$ \= $(m_1+m_2, n)$\\
\> (2) \> $(m_1, m_2, n)$ \> $\longmapsto$ \> $(m_1, n)$\\
\> (3) \> $(m_1, m_2, n)$ \> $\longmapsto$ \> $(m_2, n)$\\
\> (4) \> $(m, n_1, n_2)$ \> $\longmapsto$ \> $(m, n_1+n_2)$\\
\> (5) \> $(m, n_1, n_2)$ \> $\longmapsto$ \> $(m, n_1)$\\
\> (6) \> $(m, n_1, n_2)$ \> $\longmapsto$ \> $(m, n_2)$\\
\> (7) \> $(m, a, n)$ \> $\longmapsto$ \> $(ma, n)$\\
\> (8) \> $(m, a, n)$ \> $\longmapsto$ \> $(m, an)$;
\end{tabbing}
où la formule (i) décrit l'application (i)(X) pour les objets
variables $X$ de $E$. Notons $L$ le foncteur « faisceau abélien
libre engendré ». Il est clair que le plus grand quotient (au sens
des faisceaux abéliens) de $L(M\times N)$ qui égalise le morphisme
$L(1)$ à $L(2)+L(3)$, $L(4)$ à $L(5)+L(6)$ et $L(7)$ à $L(8)$
représente le foncteur des morphismes $A$-bilinéaires de $M\times N$
dans $P$.

\setcounter{subsection}{7}





\subsection{}\etiquette[12.8]{IV.subsec12.8}
On constate que le foncteur $P\mapsto \Hom_A(N,\scrHom_{\sisi{\underline{\underline{\ZZ}}}{\underline{\ZZ}}}(M,P))$ est
aussi canoniquement isomorphe au foncteur des applications
$A$-bilinéaires de $M\times N$ dans $P$. Par suite le produit
tensoriel $M\otimes_A N$ représente aussi le foncteur
$\Hom_A(N,\scrHom_{\sisi{\underline{\underline{\ZZ}}}{\underline{\ZZ}}}(M,P))$. On a donc des isomorphismes, \sisi{fonctoriel}{fonctoriels}
en \sisi{tout}{tous} les arguments:
% correction SGA4-CORR-0324 / SGA-SRC-001024: SGA4-T1-IV9088-65-D768D282.
\begin{align*}%\etiquette[12.8.1]{IV.eq12.8.1}
\Hom_Z(M\otimes_A N,P) &\simeq \Hom_A(M,\scrHom_{\sisi{\underline{\underline{\ZZ}}}{\underline{\ZZ}}}(N,P)){\quoi},\tag{12.8.1}\\
\Hom_Z(M\otimes_A N,P) &\simeq \Hom_A (N,\scrHom_{\sisi{\underline{\underline{\ZZ}}}{\underline{\ZZ}}}(M,P)){\quoi}.\tag{12.8.2}
\end{align*}

\subsection{}\etiquette[12.9]{IV.subsec12.9}
Il\pageoriginale résulte de (12.8.1), ou bien de
(12.8.2) et de (\hyperref[IV.subsec12.2]{12.2}) que le foncteur
$\otimes_A$ commute aux limites inductives en ses deux arguments.

\begin{enonce*}{Proposition 12.10}\etiquette[12.10]{IV.prop12.10}
Soient $C$ un $𝒰$-site, $A'$ un préfaisceau d'anneaux sur
$C$, $M'$ (\resp $N'$) un $A'$-module à droite (\resp à gauche),
$A$, $M$, $N$ les faisceaux associés à $A'$, $M'$ et $N'$
respectivement. Le faisceau associé au préfaisceau $X\mapsto
M'(X)\otimes_{\sisi{A(X)}{A'(X)}}N'(X)$ $(X\in \ob C)$ est canoniquement isomorphe
% correction SGA4-CORR-0325 / SGA-SRC-001025: SGA4-T1-IV9104-66-341D342C.
au faisceau $M\otimes_A N$.
\end{enonce*}

Le préfaisceau $X\mapsto M'(X)\otimes_{A'(X)}N'(X)$ peut se
construire à partir des préfaisceaux $M'$, $N'$ et $A'$ par les
opérations indiquées dans la démonstration de \hyperref[IV.prop12.7]{12.7}.
Comme le foncteur « faisceau associé » commute aux limites
projectives et inductives finies et aux foncteurs « objet abélien
libre engendré », la formation du produit tensoriel commute au
foncteur « faisceau associé ».

\begin{enonce*}{Proposition 12.11}\etiquette[12.11]{IV.prop12.11}
Soient $(E,A)$ un topos annelé, $M$ un $A$-Module à droite, $N$ un
$A$-Module à gauche, $X$ un objet de $E$.
\begin{enumerate}
\renewcommand{\theenumi}{\alph{enumi}}
\renewcommand{\labelenumi}{\theenumi)}
\item On a un isomorphisme canonique
$$
j^{\ast}_X(M\otimes_A N)\simeq (j^{\ast}_XM)\otimes_{A|X}
% correction SGA4-CORR-0326 / SGA-SRC-001026: SGA4-T1-IV9125-67-CA45475C.
(j^{\ast}_XN){\quoi}.
$$
Soient de plus $P$ un $\sisi{A/X}{A|X}$-Module à droite et $Q$ un $\sisi{A/X}{A|X}$ module à
% correction SGA4-CORR-0327 / SGA-SRC-001027: SGA4-T1-IV9128-68-FF49CB55.
gauche.

\item On a des isomorphismes canoniques (formules de projection):
\begin{align*}
j_{X!}(P\otimes_{A| X} j^{\ast}_XN) &\simeq (j_{X!}P)\otimes_A N\\
j_{X!}(j^{\ast}_XM\otimes_{A| X}Q) &\simeq M\otimes_A
(j_{X!}Q){\quoi}.
\end{align*}
\end{enumerate}
\end{enonce*}

\begin{itemize}
\item[Démontrons a).] Pour tout faisceau abélien $R$ sur $E_{/X}$, on
a la suite d'isomorphismes fonctoriels en $R$:
\begin{alignat*}{2}
&\Hom_{\sisi{\underline{\underline{\ZZ}}}{\underline{\ZZ}}}(j^{\ast}_X(M\otimes_A N),R)\simeq \Hom_{\sisi{\underline{\underline{\ZZ}}}{\underline{\ZZ}}} (M\otimes_A N,
j_{X\ast}R) &\quad& \eqhyperref[IV.eq11.2.2.1]{11.2.2.1}\\
&\Hom_{\sisi{\underline{\underline{\ZZ}}}{\underline{\ZZ}}}(M\otimes_A N, j_{X\ast}R)\simeq \Hom_A(M, \scrHom_Z(N,
j_{X\ast}R)) && (12.8.1)\\
&\Hom_A (M, \scrHom_Z(N, j_{X\ast}R))\simeq
\Hom_A(M,j_{X\ast}\scrHom_{\sisi{\underline{\underline{\ZZ}}}{\underline{\ZZ}}}(j^{\ast}_XN,R)) && (\hyperref[IV.prop12.3]{12.3})\\
& \Hom_A(M, j_{X\ast}\scrHom_{\sisi{\underline{\underline{\ZZ}}}{\underline{\ZZ}}}(j^{\ast}_XN,R))\simeq \Hom_{A|
X}(j^{\ast}_XM, \scrHom_{\sisi{\underline{\underline{\ZZ}}}{\underline{\ZZ}}} (j^{\ast}_XN,R)) && \eqhyperref[IV.eq11.2.2.1]{11.2.2.1}\\
&\Hom_{\sisi{A/X}{A|X}}(j^{\ast}_XM, \scrHom_{\sisi{\underline{\underline{\ZZ}}}{\underline{\ZZ}}} (j^{\ast}_XN,R)\sisi{}{)}\simeq \Hom_{\sisi{\underline{\underline{\ZZ}}}{\underline{\ZZ}}}
(j^{\ast}_XM\sisi{\otimes_A}{\otimes_{A|X}} j^{\ast}_X N,R) && (12.8.1)
% correction SGA4-CORR-0328 / SGA-SRC-001028: SGA4-T1-IV9152-69-53936BE2.
\end{alignat*}

Exhibons\pageoriginale le premier isomorphisme de b). Pour tout
faisceau abélien $R$, on a une suite d'isomorphismes fonctoriels en
$R$:
\begin{alignat*}{2}
& \Hom_Z(j_{X!}(P\otimes_{A| X}j^{\ast}_XN),R)\simeq
\Hom_Z(P\otimes_{A| X} j^{\ast}_X N, j^{\ast}_X R) &\quad&
(\hyperref[IV.prop11.3.1]{11.3.1})\\
& \Hom_Z(P\otimes_{A| X}j^{\ast}_XN, j^{\ast}_XR)\simeq
\Hom_{A/X}(P,\scrHom_Z(j^{\ast}_XN, j^{\ast}_XR)) && (12.8.1)\\
&\Hom_{A| X}(P,\scrHom_Z (j^{\ast}_XN, j^{\ast}_XR))\simeq
\Hom_{A| X}(P, j^{\ast}_X\scrHom_Z (N,R)) && (\hyperref[IV.prop12.3]{12.3})\\
&\Hom_{A| X}(P, j^{\ast}_X\scrHom_Z(N,R))\simeq \Hom_A (j_{X!}P,
\scrHom_Z (N,R)) && (\hyperref[IV.prop11.3.1]{11.3.1})\\
&\Hom_A (j_{X!}P, \scrHom_Z(N,R))\simeq \Hom_Z (j_{X!}P\otimes_A
N,R) && (12.8.1)
\end{alignat*}

Le deuxième isomorphisme s'obtient de manière analogue à l'aide de
(12.8.2).
\end{itemize}

\begin{enonce*}{Corollaire 12.12}\etiquette[12.12]{IV.cor12.12}
Soient $E$ un topos, $A$, $B$, $C$ trois faisceaux d'anneaux sur
$E$, $M$ un $B-A$ biModule, $N$ un $A-C$ biModule. Le faisceau
abélien $M\otimes_A N$ est muni canoniquement d'une structure de
$B-C$ biModule. En particulier, lorsque $A$ est un faisceau
d'anneaux commutatifs et lorsque $M$ et $N$ sont des $A$-Modules,
$M\otimes_AN$ est muni canoniquement d'une structure de $A$-Module.
Les isomorphismes de $\sisi{\hyperref[IV.subsec12.9]{12.9}}{\hyperref[IV.prop12.11]{12.11}}$ sont des isomorphismes de
% correction SGA4-CORR-0329 / SGA-SRC-001029: SGA4-T1-IV9184-70-0DDD9B84.
biModules.
\end{enonce*}

Pour tout objet $X$ de $E$, $j^{\ast}_XM$ est un $A| X$-Module à
droite sur lequel opère, à gauche, l'anneau $B(X)$ et de même,
$j^{\ast}_XN$ est un $A/X$-Module à gauche sur lequel opère, à
droite, $C(X)$. Par fonctorialité, le faisceau abélien
$j^{\ast}_XM\otimes_{A/X} j^{\ast}_XN\simeq j^{\ast}_X (M\otimes_A
N)$ est muni d'une structure de $B(X)-C(X)$ « \sisi{objet}{bimodule} », structure
qui varie fonctoriellement en $X$. Par suite $M\otimes_A N$ est muni
d'une structure de $B-C$ biModule. Cette structure de $B(X)-C(X)$
\sisi{objet}{bimodule} sur $j^{\ast}_X(M\otimes_A N)$ se reflète de façon évidente
% correction SGA4-CORR-0330 / SGA-SRC-001030: SGA4-T1-IV9193-71-95BB0047.
sur le foncteur « morphisme $A$-bilinéaire ». On constate alors
que, lorsque $A$ est commutatif et lorsque $M$ et $N$ sont des
$A$-Modules, les structures de $A$-Modules à droite et à gauche
qu'on obtient sont « égales »  et par suite $M\otimes_A N$ est
dans ce cas un $A$-Module. La dernière assertion est laissée au
lecteur.

\begin{enonce*}{Corollaire 12.13}\etiquette[12.13]{IV.cor12.13}
Soient $(E,A)$ un topos annelé, $M$ un $A$-Module à
gauche\pageoriginale (\resp à droite), $X$ un objet de $E$. On a
(avec la notation $A_X$ de $\hyperref[IV.prop11.3.3]{11.3.3}$) un isomorphisme
canonique
$$
A_X\otimes_AM\simeq j_{X!} j^{\ast}_XM\text{ (\resp } M\otimes_A
% correction SGA4-CORR-0331 / SGA-SRC-001031: SGA4-T1-IV9210-72-A5DC58C6.
A_X\simeq j_{X!}j^{\ast}_XM)){\quoi}.
$$
\end{enonce*}

\noindent \sisi{}{Ceci }résulte des formules de projections (\hyperref[IV.prop12.11]{12.11}).
% correction SGA4-CORR-0332 / SGA-SRC-001032: SGA4-T1-IV9215-73-92DB5EF5.

\begin{enonce*}{Proposition 12.14}\etiquette[12.14]{IV.prop12.14}
Soient $E$ un topos, $A$ et $B$ deux faisceaux d'anneaux sur $E$,
$M$ un $A$-Module à droite, $N$ un $A-B$ biModule, $P$ un $B$-Module
à droite. On a un isomorphisme canonique
$$
\Hom_B(M\otimes_A N,P)\simeq \Hom_A (M,\scrHom_B (N,P)){\quoi}.
$$
\end{enonce*}

\noindent De (12.8.1) on tire un morphisme $A$-bilinéaire
canonique $\scrHom_Z(N,P)\times N\rightarrow P$; d'où, en se
restreignant à $\scrHom_B(N,P)\times N$ (qui est un sous-faisceau),
un morphisme $A$-bilinéaire $\scrHom_B(N,P)\times N\rightarrow P$
dont on vérifie immédiatement qu'il est $B$-linéaire sur
le deuxième
facteur. On a donc un morphisme canonique de $B$-Modules
$\scrHom_B(N,P)\otimes_A N\rightarrow P$; d'où une application
canonique, fonctorielle en $M$:
\begin{equation*}
\Hom_A(M,\scrHom_B(N,P))\longrightarrow \Hom_B (M\otimes_A
N,P){\quoi}.\tag{12.14.1}
\end{equation*}\etiquette[12.14.1]{IV.eq12.14.1}
Montrons que ce morphisme de foncteurs est un isomorphisme. Comme
les deux membres transforment les limites inductives de $M$ en
limites projectives, il suffit de montrer que \eqhyperref[IV.eq12.14.1]{12.14.1} est un
isomorphisme, lorsque $M$ parcourt une famille génératrice de la
catégorie $E_A$. Il suffit donc de montrer que \eqhyperref[IV.eq12.14.1]{12.14.1} est
un isomorphisme lorsque $M=A_X$ où $X$ est un objet de $E$. La
vérification est alors immédiate.


\section{Morphisme de topos annelés}\etiquette[13]{IV.sec13}

\begin{enonce*}{Définition 13.1}\etiquette[13.1]{IV.def13.1}
Soient $(E,A)$ et $(E',A')$ deux topos annelés. Un morphisme de
topos annelé $u:(E,A)\rightarrow (E',A')$ est un couple $(m,\theta)$
où $m:E\rightarrow E'$ est un morphisme de topos {\rm
(\hyperref[IV.def3.1]{3.1})} et $\theta:m^{\ast}A'\rightarrow A$ est un
morphisme d'Anneaux.
\end{enonce*}

\setcounter{subsection}{1}
\subsubsection{}\etiquette[13.1.1]{IV.subsubsec13.1.1}
Comme le foncteur $m^{\ast}:E'\rightarrow E$ est adjoint à gauche au
foncteur $m_{\ast}:E\rightarrow E'$, se donner un morphisme de topos
annelés $(m,\theta):(E,A)\rightarrow (E',A')$ revient à se donner un
morphisme de topos $m:E\rightarrow E'$ et un morphisme de faisceaux
d'anneaux $\Theta:A'\rightarrow m_{\ast}A$.

\subsection{}\etiquette[13.2]{IV.subsec13.2}
A\pageoriginale un morphisme de topos annelés
$u=(m,\theta):(E,A)\rightarrow (E',A')$, on associe deux
% correction SGA4-CORR-0333 / SGA-SRC-001033: SGA4-T1-IV9268-74-5079AFCA.
foncteurs remarquables entre les catégories de Modules:

\subsubsection{Le foncteur image directe pour les
  Modules:}\etiquette[13.2.1]{IV.subsubsec13.2.1}
Soit $M$ un $A$-Module à gauche (\resp à droite). L'objet
$m_{\ast}M$ est muni canoniquement d'une structure de
$m_{\ast}A$-Module; d'où par restriction des scalaires par le
morphisme canonique $\Theta':A'\rightarrow m_{\ast}A$, un
$A'$-Module noté $u_{\ast}(M)$ et appelé \emph{l'image directe de
$M$ par le morphisme $u$.}

\subsubsection{Le foncteur image réciproque pour les
  Modules:}\etiquette[13.2.2]{IV.subsubsec13.2.2}
Soit $N$ un $A'$-Module à gauche (\resp à droite). L'objet
$m^{\ast}N$ de $E$ est muni canoniquement d'une structure de
$m^{\ast}A'$-Module à gauche (\resp à droite). Le $A$-Module à
gauche $A\otimes_{m^{\ast}A'}m^{\ast}N$ (\resp à droite
$m^{\ast}N\otimes_{m^{\ast}A'}A$) où $A$ est muni de la structure de
% correction SGA4-CORR-0334 / SGA-SRC-001034: SGA4-T1-IV9285-75-39E4B20F.
$m^{\ast}A'$-Module définie par $\theta:m^{\ast}A'\rightarrow A$,
est noté $u^{\ast}N$ et est appelé l'\emph{image réciproque du
Module $N$ par le morphisme de topos annelé} $u$.

\subsubsection{}\etiquette[13.2.3]{IV.subsubsec13.2.3}
On notera que pour tout $A$-Module $M$, le faisceau d'ensembles et
le faisceau abélien sous-jacent à $u_{\ast}M$ est le faisceau
$m_{\ast}M$. Aussi emploie-t-on le plus souvent la notation
$u_{\ast}$ pour désigner le foncteur $m_{\ast}$: image directe pour
les faisceaux d'ensembles. En revanche pour un $A'$-Module $N$, le
faisceau d'ensembles ou le faisceau abélien sous-jacent à
$u^{\ast}N$ n'est ni égal ni isomorphe en général à $m^{\ast}N$
(sauf toutefois lorsque $\theta$ est un isomorphisme). Il y a donc
lieu de distinguer entre l'image réciproque pour les Modules et
l'image réciproque pour les faisceaux d'ensembles ou les faisceaux
abéliens. On utilise le plus souvent la notation
$u^{-1}:E'\rightarrow E$ pour désigner le foncteur $m^{\ast}$ image
inverse pour les faisceaux d'ensembles. Le foncteur $u^{-1}$ est
appelé le foncteur \emph{image réciproque « ensembliste »} par
le morphisme de topos annelé $u$, par opposition avec l'image
réciproque « au sens modules »  $u^{\ast}$.

\begin{enonce*}{Définition 13.3}\etiquette[13.3]{IV.def13.3}
Soient $(C,A)$ et $(C',A')$ deux $𝒰$-sites annelés. Un
morphisme de sites annelés $u:(C,A)\rightarrow (C',A')$ est un
couple $(m,\theta)$ où $m$ est un morphisme du site $C$ dans le site
$C'$ (\hyperref[IV.subsec4.9]{4.9}) et $\theta:m^{\ast}A'\rightarrow
A$\pageoriginale \sisi{}{est }un morphisme d'Anneaux.
% correction SGA4-CORR-0335 / SGA-SRC-001035: SGA4-T1-IV9313-76-94BA9B47.
\end{enonce*}

\setcounter{subsection}{3} \setcounter{subsubsection}{0}
\subsubsection{}\etiquette[13.3.1]{IV.subsubsec13.3.1}
Un morphisme de topos annelé est un morphisme de $𝒰$-sites
annelés. Un morphisme de sites annelés donne naissance à un
morphisme entre les topos annelés correspondants
(\hyperref[IV.subsubsec11.1.1]{11.1.1} et \hyperref[IV.subsubsec4.9.1]{4.9.1}).

\begin{enonce*}{Proposition 13.4}\etiquette[13.4]{IV.prop13.4}
Soit $u:(E,A)\rightarrow (E',A')$ un morphisme de topos annelés.
\begin{enumerate}
\renewcommand{\theenumi}{\alph{enumi}}
\renewcommand{\labelenumi}{\theenumi)}
\item Pour un $A$-Module à gauche (\resp à droite) variable $M$, et
pour un $A'$-Module à gauche (\resp à droite) variable $N$, on a des
isomorphismes bifonctoriels canoniques (dits isomorphismes
d'adjonction):
\begin{align*}
\Hom_A(u^{\ast}N,M) &\simeq
\Hom_{A'}(N,u_{\ast}M){\quoi},\tag{13.4.1}\\
u_{\ast}\scrHom_A(u^{\ast}N,M) &\simeq
\scrHom_{A'}(N,u_{\ast}M){\quoi}.\tag{13.4.2}
\end{align*}

\item Pour un objet variable $X$ de $E'$, on a un isomorphisme
canonique, fonctoriel en $X$:
\begin{equation*}
u^{\ast}A'_X\sisi{\longrightarrow}{\simeq} A_{u^{-1}(X)}{\quoi},\tag{13.4.3}
% correction SGA4-CORR-0336 / SGA-SRC-001036: SGA4-T1-IV9342-77-E6F725C7.
\end{equation*}\etiquette[13.4.3]{IV.eq13.4.3}
où $A'_X$ et $A_{u^{-1}(X)}$ désignent les Modules libres engendrés
$(\hyperref[IV.prop11.3.3]{11.3.3})$.

\item Lorsque $A'$ est commutatif et lorsque le morphisme canonique
$u^{-1}A'\rightarrow A$ est central (\resp lorsque le morphisme
canonique $u^{-1}A'\rightarrow A$ est un isomorphisme) on a, pour un
$A'$-Module à droite $M$ variable et un $A'$-Module à gauche $N$
variable, un isomorphisme bifonctoriel canonique:
\begin{align*}
u^{\ast}M\otimes_A u^{\ast}N &\simeq
u^{\ast}(M\otimes_{A'}N)\tag{13.4.4}\\
u^{\ast}M\otimes_A u^{\ast}N &\simeq
u^{-1}(M\otimes_{A'}N)\sisi{)}{}{\quoi}.\tag*{(\resp (13.4.5)\sisi{}{)}}
% correction SGA4-CORR-0337 / SGA-SRC-001037: SGA4-T1-IV9356-78-58E54788.
\end{align*}
\end{enumerate}
\end{enonce*}

On a tout d'abord un isomorphisme canonique
$\Hom_A(u^{\ast}N,M)\simeq \Hom_{u^{-1}A'}(u^{-1}N,M)$
(\sisi{\hyperref[IV.cor12.12]{12.12}}{\hyperref[IV.prop12.14]{12.14}}), puis un isomorphisme
$\Hom_{u^{-1}A'}(u^{-1}N,M)\simeq \Hom_{A'}(N,u_{\ast}M)$
(\hyperrefext[III][III.prop1.7]{1.7}); d'où (13.4.1). Exhibons l'isomorphisme
(13.4.2). Pour tout objet $X$ de $E'$, on a la suite
d'isomorphismes fonctoriels:
\begin{alignat*}{3}
& \Hom_{E'}(X,u_{\ast}\scrHom_A (u^{\ast}N,M))\simeq && \Hom(
u^{-1}X,
 \scrHom_A (u^{\ast}N, M)) &\quad& \text{(adjonction){\quoi},}\\
&\Hom_A(u^{\ast}N,\scrHom_E (u^{-1}X, M)) \simeq &&\qquad ''\qquad ''\qquad
''\qquad '' && (\hyperref[IV.prop12.1]{12.1}){\quoi},\\
& \Hom_A (u^{\ast}N,\scrHom_E(u^{-1}X, M))\simeq &&\Hom_{A'}(N,
u_{\ast} \scrHom_E (u^{-1} X, M)) && (13.4.1){\quoi},\\
& \Hom_{A'}(N,\scrHom_{E'}(X,u_{\ast}M))\simeq &&\qquad ''\qquad  ''\qquad
 ''\qquad '' && \eqhyperref[IV.eq10.3.1]{10.3.1}{\quoi},\\
& \qquad ''\qquad ''\qquad ''\qquad '' \simeq &&
 \Hom_{\sisi{A'}{E'}}(X,\scrHom_{A'}(N,u_{\ast}M)) && (\hyperref[IV.prop12.1]{12.1}){\quoi}.
% correction SGA4-CORR-0338 / SGA-SRC-001038: SGA4-T1-IV9379-79-6766F9E2.
\end{alignat*}
La\pageoriginale formule (13.4.5) s'obtient alors à partir de
la formule \sisi{13.3.2}{13.4.2}. par adjonction (définition du produit tensoriel
% correction SGA4-CORR-0339 / SGA-SRC-001039: SGA4-T1-IV9382-80-F9AB528E.
(\hyperref[IV.prop12.7]{12.7})). La formule (13.4.4) se déduit de
(13.4.5) en utilisant la commutativité du produit tensoriel
lorsque l'anneau de base est commutatif (\hyperref[IV.subsec12.8]{12.8}).
Enfin, pour démontrer \eqhyperref[IV.eq13.4.3]{13.4.3}, on considère la suite
d'isomorphismes
\begin{alignat*}{3}
& \Hom_A(u^{\ast}A'_X, M) && \simeq\Hom_{A'}(A'_X, u_{\ast}M)
&\quad&
(13.4.1){\quoi},\\
& \Hom_{E'}(X, u_{\ast}M) && \simeq\quad ''\qquad ''\qquad ''&&
(\hyperref[IV.prop11.3.3]{11.3.3}){\quoi},\\
&\quad ''\qquad ''\qquad ''&& \simeq\Hom_E(u^{-1}X, M) &&
\text{(adjonction)}{\quoi},\\
&\Hom_A(A_{u^{-1}X},M) &&\simeq\quad ''\qquad ''\qquad ''&&
(\hyperref[IV.prop11.3.3]{11.3.3}){\quoi}.
\end{alignat*}

\begin{enonce*}{Corollaire 13.5}\etiquette[13.5]{IV.cor13.5}
Soient $(E,A)$ un topos annelé, $x:P\rightarrow E$ un point de $E$,
$F\mapsto F_x$ le foncteur fibre associé (\hyperref[IV.def6.1]{6.1}). Le
foncteur fibre en $x$ transforme le produit tensoriel des
$A$-Modules en produit tensoriel (ordinaire) des $A_x$-modules.
\end{enonce*}

Le produit tensoriel dans $P$ est le produit tensoriel ordinaire des
modules ($P$ le topos ponctuel est la catégorie des ensembles).
L'assertion résulte donc de (13.4.4).

\begin{enonce*}{Corollaire 13.6}\etiquette[13.6]{IV.cor13.6}
Soit $u:(E,A)\rightarrow (E',A')$ un morphisme de topos annelés. Le
foncteur $u_{\ast}$ image directe pour les Modules à droite ou à
gauche commute aux limites projectives et en particulier est exact à
gauche. Le foncteur $u^{\ast}$ image réciproque pour les Modules à
droite ou à gauche commute aux limites inductives et en particulier
est exact à droite.
\end{enonce*}

Ceci résulte de la formule (13.4.1) (\hyperrefext[I][I.prop2.11]{2.11}).

\setcounter{subsection}{6}
\subsection{}\etiquette[13.7]{IV.subsec13.7}
On notera que le foncteur $u^{\ast}$, image réciproque pour les
Modules, n'est pas, en général, exact, alors que le foncteur
$u^{-1}$, image réciproque ensembliste, est exact. On peut cependant
affirmer l'exactitude de $u^{\ast}$ lorsque le morphisme canonique
$u^{-1}A'\rightarrow A$ est plat (à droite\pageoriginale ou à
gauche) (V 1.8). C'est en particulier le cas lorsque le morphisme
canonique $u^{-1}A'\rightarrow A$ est un isomorphisme.

\subsection{}\etiquette[13.8]{IV.subsec13.8}
Soient $C$ un $𝒰$-site, $A'$ un préfaisceau d'anneaux sur
$C$, et notons $A$ le faisceau associé à $A'$. La catégorie des
préfaisceaux de $A'$-modules qui sont des faisceaux est équivalente
à la catégorie des $A$-Modules. On utilise parfois la notation
$\Hom_{A'}(M,N)$ pour désigner le groupe $\Hom_A(M,N)$
(\hyperref[IV.subsubsec11.1.2]{11.1.2}). De même, et abusivement, on utilise les
notations $\scrHom_{A'}(M,N)$ et $M\otimes_{A'}N$ pour désigner les
faisceaux $\scrHom_A(M,N)$ et $M\otimes_A N$. Lorsque $A'=k_E$ est
le préfaisceau constant, défini par un anneau ordinaire $k$, on
écrit aussi $\scrHom_k$, $\otimes_k$ au lieu de $\scrHom_{k_E}$,
$\otimes_{k_E}$

\begin{enonce*}[remark]{Exercice 13.9}\etiquette[13.9]{IV.exer13.9}
\emph{Topos localement annelés} (\cf [\hyperref[IV.sec9]{9}] pour plus de
renseignements dans l'ordre d'idées qui suit).

Soit $(E,{\Oo}_E)$ un topos commutativement annelé. Pour $X\in
\ob E$ et $f\in {\Oo}_E(X)$, soit $X_f$ le plus grand sous-objet
de $X$ sur lequel $f$ soit inversible.
\begin{enumerate}
\renewcommand{\theenumi}{\alph{enumi}}
\renewcommand{\labelenumi}{\theenumi)}
\item Montrer que pour $X'\in \ob E_{/X}$, $f_{X'}$ est inversible si
et seulement si le morphisme structural $X'\rightarrow X$ se
factorise par $X_f$, et que pour $f$, $g\in {\Oo}_E(X)$, on a
$$
X_{fg}=X_f\cap X_g{\quoi}.
$$

\item Montrer que les conditions suivantes (i) à (iii) sont
équivalentes:
\begin{enumerate}
\renewcommand{\theenumii}{\roman{enumii}}
\renewcommand{\labelenumii}{(\theenumii)}
\item Pour $X\in \ob E$ et $f$, $g\in {\Oo}_E(X)$, on a
$$
X_{f+g}\subset \Sup(X_f,X_g)
$$

\item Pour $X\in \ob E$ et $f$, $g\in {\Oo}_E(X)$ tels que $f+g$ soit
inversible, on a $X=\Sup(X_f,X_g)$.

\item Pour $X\in \ob E$ et $f\in {\Oo}_E(X)$, on a
$$
X=\Sup(X_f, X_{1-f}){\quoi}.
$$

Montrer que ces conditions impliquent la condition suivante (iv),
et\pageoriginale sont équivalentes à cette dernière si $E$ a
suffisamment de points:

\item Pour tout point $p$ de $E$, l'anneau fibre ${\Oo}_{E,p}$ est un
\emph{anneau local}.
\end{enumerate}

Lorsque les conditions équivalentes (i) à (iii) sont satisfaites, on
dira que $(E,{\Oo}_E)$ est un \emph{topos localement annelé}, et
on dit de même qu'un site annelé est un \emph{site localement
annelé} si le topos annelé correspondant est un topos localement
annelé.

\item Soient $(E,{\Oo}_E)$ et $(E',{\Oo}_{E'})$ deux topos localement
annelés, et $f$ un morphisme de topos annelés du premier dans le
second. Montrer que la condition (i) suivante implique la condition
(ii) et lui est équivalente si $E$ a suffisamment de points:
\begin{enumerate}
\renewcommand{\theenumii}{\roman{enumii}}
\renewcommand{\labelenumii}{(\theenumii)}
\item Pour tout $\sisi{X}{X'}\in \ob E'$ et $s'\in {\Oo}_{E'}(X')$, posant
$X=f^{-1}(\sisi{X}{X'})$, $s=f^{\ast}(s')$, on~a
$$
\sisi{X'_{s'}=f^{-1}(X_s)}{X_s=f^{-1}(X'_{s'})}{\quoi}.
% correction SGA4-CORR-0340 / SGA-SRC-001040: SGA4-T1-IV9502-81-B0D29043.
$$

\item Pour tout point $p$ de $E$, posant $p'=f(p)$, l'homomorphisme
naturel sur les fibres
$$
{\Oo}_{E',p'}\longrightarrow {\Oo}_{E,p}
$$
est un homomorphisme \emph{local} d'anneaux locaux.
\end{enumerate}

Lorsque la condition (i) est satisfaite, on dira que $f$ est un
\emph{morphisme de topos localement annelés}, ou encore un
\emph{morphisme admissible de topos localement annelés} si une
confusion est à craindre. On désigne par $\Hhomtoplocan(E,E')$ la
sous-catégorie pleine de la catégorie $\Hhomtopan(E,E')$ de tous les
morphismes de topos annelés de $E$ dans $E'$ définie par les
morphismes admissibles de $E$ dans $E'$.

\item Supposons que $(E,{\Oo}_E)$ soit le topos annelé
  associé à un
espace topologique annelé $(X,{\Oo}_X)$ (\hyperref[IV.subsec2.1]{2.1}).
Montrer que pour que $(E,{\Oo}_E)$ soit localement annelé, il
faut et il suffit que $(X,{\Oo}_X)$ soit localement annelé, \ie
que pour tout $x\in X$, la fibre ${\Oo}_{X,x}$ soit un anneau
local. (Noter que dans le critère (iv) de b), il suffit de prendre
le point $p$ dans une famille \emph{conservatrice} de points de
$E$). Soit\pageoriginale $f:(X,{\Oo}_X)\rightarrow (X',
{\Oo}_{X'})$ un morphisme d'espaces annelés, où ${\Oo}_X$ et
${\Oo}_{X'}$ sont des faisceaux d'anneaux locaux. Montrer que le
morphisme de topos annelés correspondant
$\Top(X,{\Oo}_X)\rightarrow \Top(X', {\Oo}_{X'})$ est
admissible si et seulement si il en est de même pour le morphisme
$f$, \ie; si et seulement si pour tout $x\in X$,
${\Oo}_{X',f(x)}\rightarrow {\Oo}_{X,x}$ est un
homomorphisme \emph{local} d'anneaux locaux.

\item Soient $(E,{\Oo}_E)$ et $(E',{\Oo}_{E'})$ deux topos localement
annelés, tels que $E$ ait suffisamment de points et que
$(E',{\Oo}_{E'})$ soit équivalent au topos annelé défini par un
\emph{schéma} $(X', {\Oo}_{X'})$; prouver que la catégorie
$\Hhomtoplocan(E,E')$ est équivalente à une catégorie
\emph{discrète}, \ie que c'est un groupoïde (tout morphisme est un
isomorphisme) rigide (les groupes d'automorphismes des objets sont
les groupes unité).

\noindent (Hint: se ramener au cas où $(E,{\Oo}_E)$ est le topos
ponctuel annelé par un corps). On se rappellera que, par contre,
$\Hhomtop(E,E')$ n'est pas en général équivalent à une
catégorie
discrète, même si $E$ et $E'$ sont des topos définis par des
schémas
(et même si $E$ est le topos ponctuel), \cf
\hyperref[IV.subsubsec4.2.3]{4.2.3}.

\item Soient $E$ un topos localement annelé, $P$ le topos annelé
défini par l'espace annelé $\Spec k$, où $k$ est un corps. Montrer
que les morphismes admissibles de topos localement annelés
$P\rightarrow E$ correspondent aux couples $(p,u)$ d'un point $p$ de
$E$, et d'une injection de $k(p)$ dans $k$, où $k(p)$ est le corps
% correction SGA4-CORR-0341 / SGA-SRC-001041: SGA4-T1-IV9564-82-64CA36C8.
résiduel de l'anneau local ${\Oo}_{E,p}$. Généraliser en un
énoncé exhibant les morphismes admissibles d'un topos localement
annelé ponctuel dans $E$ (généralisant EGA I 2.4.4).

On appelle \emph{point géométrique d'un topos localement
  annelé} $E$
tout morphisme admissible dans $E$ du topos localement annelé défini
par un espace annelé de la forme $\Spec(k)$, où $k$ est un corps
\emph{algébriquement clos}. Définir la \emph{catégorie des points
géométriques de} $E$ (sous-entendu: correspondants à des corps $k\in
𝒰$), notée $\Pptgeom(E)$, et un foncteur canonique
$\Pptgeom(E)\rightarrow \Ppoint(E)$. Montrer que ce foncteur est
fidèle si\pageoriginale et seulement si $E$ n'a pas de point \ie
$\Ppoint(E)$ est la catégorie vide.

\item Soient $E$ un $𝒰$-topos, $𝒱$ un univers tel que $𝒰\in
𝒱$. Définir une 2-catégorie
$(𝒱\text{-}𝒰\text{-}\Top)_{/E}$ en termes de l'objet
$E$ de la 2-catégorie $(𝒱\text{-}𝒰\text{-}\Top)$
(\hyperref[IV.subsubsec3.3.1]{3.3.1}), en s'inspirant de \hyperref[IV.exer5.14]{5.14} a).
Lorsque $E$ est annelé par un Anneau $A$, définir de même une
2-catégorie $(𝒱\text{-}𝒰\text{-}\text{Topan})_{/E}$, et
% correction SGA4-CORR-0342 / SGA-SRC-001042: SGA4-T1-IV9580-83-B5805DC4.
lorsque $E$ est localement annelé, utilisant la notion de morphisme
admissible (\cf c)), définir une 2-catégorie
$(𝒱\text{-}𝒰\text{-}\text{Topanloc})_{/E}$,
admettant $\Pptgeom(E)$ (\cf f)) comme sous-catégorie pleine.
\end{enumerate}
\end{enonce*}

\section{Modules sur un topos défini par recollement}\etiquette[14]{IV.sec14}

\subsection{}\etiquette[14.1]{IV.subsec14.1}
Soient $(E,A)$ un topos annelé, $U$ un sous-topos ouvert de $E$, $F$
le sous-topos fermé complémentaire, $j:U\rightarrow E$ et
$i:F\rightarrow E$ les morphismes canoniques (\hyperref[IV.subsec9.3]{9.3}).
Le topos $U$ est annelé par $j^{\ast}A$ (\sisi{\hyperref[IV.subsubsec11.2.2]{11.2.2}}{\hyperref[IV.subsubsec11.2.1]{11.2.1}})
et le morphisme de topos $\sisi{i}{j}$, complété de manière évidente, devient
un morphisme de topos annelés. Dans ce numéro, le topos $F$ sera
annelé par le faisceau $i^{\ast}A$, noté $A| F$, et le morphisme
$\sisi{j}{i}$, complété de manière évidente, est alors un morphisme de topos
% correction SGA4-CORR-0343 / SGA-SRC-001043: SGA4-T1-IV9600-84-F5A47F94.
annelés.


\subsection{}\etiquette[14.2]{IV.subsec14.2}
On a donc deux morphismes de topos annelés:
$$
j:(U,A| U)\longrightarrow (E,A)\text{ et } i:(F,A| F)\longrightarrow
(E,A){\quoi};
$$
d'où cinq foncteurs entre les catégories de Modules (à gauche pour
fixer les idées) correspondantes:
\[\vcenter{
\xymatrix@C=2cm@R=1.5cm@M=.25cm{ ({}_{A|U}U) \ar@<14pt>[r]^{j_!}
\ar@<-14pt>[r]^{j_{\ast}}&
  ({}_{A}E) \ar[l]_{j^{\ast}} \ar@<7pt>[r]^{i^{\ast}} &
  ({}_{A|F}F)\quad. \ar@<7pt>[l]_{i_{\ast}}
}}\leqno(14.2.1)
\]\etiquette[14.2.1]{IV.eq14.2.1}
Chaque foncteur du diagramme \eqhyperref[IV.eq14.2.1]{14.2.1} est adjoint à gauche à
celui qui se trouve au-dessous de lui.

\subsection{}\etiquette[14.3]{IV.subsec14.3}
On\pageoriginale sait que le foncteur $X\mapsto (j^{\ast}X,
i^{\ast}X; i^{\ast}X\rightarrow i^{\ast}j_{\ast}j^{\ast}X)$ est une
équivalence de $E$ dans le topos $(U,F, i^{\ast}j_{\ast})$
(\hyperref[IV.theo9.5.4]{9.5.4}). Cette équivalence induit une équivalence entre
les catégories de Modules correspondantes et par suite le foncteur
$P\mapsto (j^{\ast}P, i^{\ast}P; i^{\ast}P\rightarrow
i^{\ast}j_{\ast}j^{\ast}P)$ est une équivalence, notée $\Phi$, de la
catégorie ${}_AE$ dans la catégorie $({}_{A| U}U, {}_{A| F}F,
i^{\ast}j_{\ast})$. Les foncteurs de (14.2.1) composés avec $\Phi$
ou $\Phi^{-1}$ sont alors les foncteurs:
\begin{align*}
\Phi\circ j_{\ast} &: M\mapsto (M,i^{\ast} j_{\ast}M;
i^{\ast}j_{\ast}M\xrightarrow{\id} i^{\ast}j_{\ast}M){\quoi},\\
j^{\ast}\circ \Phi^{-1} &: (M,N; N\rightarrow
i^{\ast}j_{\ast}M)\mapsto M{\quoi},\\
\Phi\circ j_{!} &: M\mapsto (M,0;0\rightarrow i^{\ast}
j_{\ast}M){\quoi},\\
\Phi\circ i_{\ast} &: N\mapsto (0,N;N\rightarrow 0){\quoi},\\
i^{\ast}\circ \Phi^{-1} &: (M,N;N\rightarrow
i^{\ast}j_{\ast}M)\mapsto N{\quoi}.
\end{align*}
Le lecteur pourra, à titre d'exercice, expliciter les différents
morphismes d'adjonction.

\subsection{}\etiquette[14.4]{IV.subsec14.4}
Notons
\begin{equation*}
i^{!}:{}_AE\longrightarrow {}_{A| F}F\tag{14.4.1}
\end{equation*}\etiquette[14.4.1]{IV.eq14.4.1}
le foncteur défini par la formule:
\begin{equation*}
i^{!}\circ
\Phi^{-1}(M,N;N\xrightarrow{u}i^{\ast}j_{\ast}M)=\Ker(u){\quoi}.\tag{14.4.2}
\end{equation*}\etiquette[14.4.2]{IV.eq14.4.2}

\begin{enonce*}{Proposition 14.5}\etiquette[14.5]{IV.prop14.5}
Le foncteur $i^{!}$ est adjoint à droite au foncteur $i_{\ast}$. Le
morphisme d'adjonction $i_{\ast}i^{!}\rightarrow \id$ est un
monomorphisme. Les foncteurs $i^{!}$ (pour des anneaux variables)
commutent aux foncteurs restriction des scalaires.
\end{enonce*}

Il est clair que tout morphisme
$$
(0,N;0)\longrightarrow (M,N;N\xrightarrow{u}i^{\ast}j_{\ast}M)
$$
se factorise d'une manière unique par $(0,\ker(u);0)$; ce qui
démontre la propriété d'adjonction. Les autres assertions sont
triviales.

\begin{enonce*}{Proposition 14.6}\etiquette[14.6]{IV.prop14.6}
Pour\pageoriginale tout objet $P$ de ${}_AE$, on a les suites
exactes fonctorielles en $P$:
\begin{gather}
0\longrightarrow j_{!}j^{\ast}P\longrightarrow P\longrightarrow
i_{\ast}i^{\ast}P\longrightarrow 0{\quoi};\tag{14.6.1}\\
0\longrightarrow i_{\ast}i^{!}P\longrightarrow P\longrightarrow
% correction SGA4-CORR-0344 / SGA-SRC-001044: SGA4-T1-IV9683-85-F97749AC.
j_{\ast}j^{\ast}P{\quoi},\tag{14.6.2}
\end{gather}
où les flèches non triviales sont les flèches d'adjonction.
\end{enonce*}

Remarquons d'abord qu'une suite
$$
(M', N'; u')\longrightarrow (M,N;u)\longrightarrow (M'',N'';u'')
$$
de $({}_{\sisi{A/U}{A|U}}U, {}_{\sisi{A/F}{A|F}}F; i^{\ast}j_{\ast})$ est exacte si et
% correction SGA4-CORR-0345 / SGA-SRC-001045: SGA4-T1-IV9693-86-3B314E04.
seulement si les suites correspondantes
\begin{align*}
& M'\longrightarrow M\longrightarrow M''\\
& N'\longrightarrow N\longrightarrow N''
\end{align*}
sont exactes.

Les suites (14.6.1) et (14.6.2) sont transformées par
l'équivalence $\Phi$ en des suites du type \eqhyperref[IV.subsec14.3]{14.3}:
\begin{align*}
& 0\longrightarrow (M,0;0)\longrightarrow (M,N;u)\longrightarrow
(0,N;0)\longrightarrow 0\\
& 0\longrightarrow (0,\ker(u);0)\longrightarrow
(M,N,u)\longrightarrow (M,i^{\ast}j_{\ast}M; \id){\quoi}.
\end{align*}
La vérification de l'exactitude de ces suites est triviale.

\setcounter{subsection}{6}
\subsection{}\etiquette[14.7]{IV.subsec14.7}
La suite exacte (14.6.2) permet d'obtenir une nouvelle
interprétation du foncteur $i_{\ast}i^{!}$. En effet soit $X$ un
% correction SGA4-CORR-0346 / SGA-SRC-001046: SGA4-T1-IV9713-87-7650C6FD.
objet de $E$. De (14.6.2) on tire la suite exacte de groupes
commutatifs:
$$
0\longrightarrow \Hom_E(X,i_{\ast}i^{!}P)\longrightarrow
\Hom_E(X,P)\longrightarrow \Hom_E(X,j_{\ast}j^{\ast}P){\quoi}.
$$
Notons encore $U$ l'ouvert de $E$, objet final du sous-topos ouvert
$U$. Il résulte des propriétés d'adjonction des foncteurs
$j_{\ast}$, $j^{\ast}$ et $j_{!}^{\ens}$ que le groupe
$\Hom_E(X,j_{\ast}j^{\ast}P)$ est canoniquement isomorphe à
$\Hom_E(X\times U,P)$ et que le morphisme $\Hom_E(X,P)\rightarrow
\Hom_E(X,j_{\ast}j^{\ast}P)$ n'est autre que le morphisme défini par
le monomorphisme $X\times U\rightarrow X$ (\hyperref[IV.sec5]{5}). En
d'autres termes (\hyperref[IV.subsubsec8.5.2]{8.5.2}):

\begin{enonce*}{Proposition 14.8}\etiquette[14.8]{IV.prop14.8}
Les\pageoriginale sections de $i_{\ast}i^{!}P$ sur $E_{/X}$ sont les
sections de $P$ sur $E_{/X}$ dont le support
$(\hyperref[IV.subsubsec9.3.5]{9.3.5})$ \sisi{contient}{est contenu dans} $F$ \sisi{(notation de $5.9.1$)}{}. En
% correction SGA4-CORR-0347 / SGA-SRC-001047: SGA4-T1-IV9731-88-BEEF2A95.
d'autres termes, $i_{\ast}i^{!}P$ est le plus grand sous-faisceau de
$P$ à support contenu dans $F$.
\end{enonce*}

\setcounter{subsection}{8}
\subsection{}\etiquette[14.9]{IV.subsec14.9}
Par abus de langage, on dit parfois que $i_{\ast}i^{!}P$ est le
sous-Module de $P$ défini par les sections de $P$ à support dans
$F$. Il résulte de \hyperref[IV.rem8.5.3]{8.5.3} que cette terminologie ne fait
qu'étendre aux topos généraux une terminologie
utilisée pour les
topos de faisceaux sur des espaces topologiques.

\begin{enonce*}{Proposition 14.10}\etiquette[14.10]{IV.prop14.10}
\begin{enumerate}
\renewcommand{\labelenumi}{\theenumi)}
\item Le foncteur $P\mapsto i_{\ast}i^{\ast}P$ est isomorphe au
foncteur $P\mapsto i_{\ast}i^{\ast}A\otimes_A P$.

\item Le foncteur $P\mapsto i_{\ast}i^{!}P$ est isomorphe au foncteur
$P\mapsto \scrHom_A(i_{\ast}i^{\ast}A,P)$.
% correction SGA4-CORR-0348 / SGA-SRC-001048: SGA4-T1-IV9753-89-79D4291F.
\end{enumerate}
\end{enonce*}

La suite exacte (14.6.1) s'écrit, dans le cas particulier où
$P$ est le $A$-Module $A$:
\begin{equation*}
0\longrightarrow A_U\longrightarrow A\longrightarrow
i_{\ast}i^{\ast}A\longrightarrow 0\tag{14.10.1}
\end{equation*}\etiquette[14.10.1]{IV.eq14.10.1}
où $A_U$ est le $A$-Module libre engendré par l'ouvert $U$
correspondant au sous-topos ouvert $U$ (9 et 11.3.3). Pour tout
$A$-Module $P$, les $A$-Modules $j_{!}j^{\ast}P$ et
$j_{\ast}j^{\ast}P$ sont canoniquement isomorphes respectivement à
$A_U\otimes_A P$ et $\scrHom_A(A_{U}, P)$ (\hyperref[IV.prop12.3]{12.3}. et
\hyperref[IV.cor12.6]{12.6}). De plus, les morphismes canoniques
$j_{!}j^{\ast}P\rightarrow P$ et $P\rightarrow j_{\ast}j^{\ast}P$
proviennent, modulo ces isomorphismes, du monomorphisme
$A_U\rightarrow A$. On tire donc de la suite exacte \sisi{15.10.1}{\ref{IV.eq14.10.1}}, deux
% correction SGA4-CORR-0349 / SGA-SRC-001049: SGA4-T1-IV9772-90-29D2631A.
suites exactes (\hyperref[IV.subsec12.2]{12.2} et \hyperref[IV.prop12.11]{12.11}):
\begin{align*}
& j_{!}j^{\ast}P\longrightarrow P\longrightarrow
i_{\ast}i^{\ast}A\otimes_AP\longrightarrow 0{\quoi},\\
& 0\longrightarrow \scrHom_A(i_{\ast}i^{\ast}A, P)\longrightarrow
P\longrightarrow j_{\ast} j^{\ast}P{\quoi};
\end{align*}
d'où les isomorphismes annoncés par comparaison avec (14.6.1)
et (14.6.2).

\begin{thebibliography}{99}\pageoriginale
\bibitem{SGA4-04:1} M. Artin et B. Mazur, Homotopy of varieties in the etale
topology, in: Proceedings of a Conference on Local Fields,
Driebergen 1966, Springer.
\bibitem{SGA4-04:2} P. Gabriel et G. Zisman, Calculus of fractions and
homotopy theory, Ergebnisse der Mathematik, Bd 35.

\bibitem{SGA4-04:3} J. Giraud, Algèbre homologique non commutative:
Grundlehren Springer 1971.

\bibitem{SGA4-04:4} A. Grothendieck, Sur quelques points d'Algèbre
homologique, Tohoku Math. Journal, (cité [Toh]).

\bibitem{SGA4-04:5} A. Grothendieck, Fondements de la Géométrie
Algébrique, (Recueil d'exposés Bourbaki 1957/62), Secrétariat
mathématique, 11 rue P. Curie Paris.
% correction SGA4-CORR-0350 / SGA-SRC-001050: SGA4-T1-IV9797-91-55A4F43F.

\bibitem{SGA4-04:6} A. Grothendieck, Crystals and the De Rham cohomology
of schemes, (notes by I. Coates et O. Jussila), in: Dix exposés sur
la cohomologie des schémas, North Holland Pub, Cie, 1969.

\bibitem{SGA4-04:7} A. Grothendieck, Classes de Chern des
représentations linéaires des groupes discrets, in: Dix exposés sur
la cohomologie des schémas, North Holland Pub. Cie, 1969.

\bibitem{SGA4-04:8} R. Godement, Théorie des faisceaux,
Act. Scient. Ind. ${\rm n^{\circ}}$ 1 252, (1958), Hermann (Paris)
(cité [TF]).

\bibitem{SGA4-04:9} M. Hakim, Topos annelés et schémas relatifs, Thèse
multigraphiée, Orsay 1967.

\bibitem{SGA4-04:10} D. Mumford, Picard groups of moduli problems, in
Arithmetic Algebraic Geometry, Harper's Series in Modern
Mathematics.

\bibitem{SGA4-04:11} Nguyen Dinh Ngoc, Thèse Sciences Mathématiques,
Paris 1963, ${\rm n^{\circ}}$ 4 995.

\bibitem{SGA4-04:12} J.E. Roos, Distributivité des $\varinjlim$ par
\sisi{rapports}{rapport} aux $\varprojlim$ des topos
% correction SGA4-CORR-0351 / SGA-SRC-001051: SGA4-T1-IV9822-92-13772179.

\begin{tabular}{@{}r@{\;}l@{\hspace{1cm}}r@{\;}l}
a) & CR t.259 p. 969-972 & b) & CR T. 259 p. 1605-1608\\
c) & CR t.259 p. 1801-1804 & & (août et septembre 1964).
\end{tabular}

\bibitem{SGA4-04:13} P. Deligne - D. Mumford, The \sisi{irreductibility}{irreducibility} of
% correction SGA4-CORR-0352 / SGA-SRC-001052: SGA4-T1-IV9830-93-20198B70.
the space of curves of given genus, Pub. math. ${\rm n^{\circ}}$ 36
(1969).

\bibitem{SGA4-04:14} B. Mitchell, Theory of categories, Academic Press
(1965).

\bibitem{SGA4-04:15} J. Lurie, Higher Topos Theory, Annals of Mathematics Studies, 170. Princeton University Press, 
Princeton, NJ, 2009. xviii+925 pp.

\bibitem{SGA4-04:16}  P. Johnstone, Topos theory, 
London Mathematical Society Monographs, Vol. 10. 
Academic Press [Harcourt Brace Jovanovich, Publishers], London-New York, 1977. xxiii+367 pp.

\bibitem{SGA4-04:17} F. Borceux, G. Van den Bossche, Structure des topologies
d'un topos, Cahiers de topologie et géométrie différentielle catégoriques,
t. 25, n°1 (1984), p 37-39.

\end{thebibliography}

\ifx\danslelivre\undefined
\end{document}
\fi
